% Generated mechanically from the frozen perfect.tex target.
% Do not edit this generated file; edit the bound locale source instead.

% OK, start here.
%

\setcounter{chapter}{35}
\chapter{概形的导出范畴}



\phantomsection
\label{perfect-section-phantom}


\section{引言}
\label{perfect-section-introduction}

\noindent
本章讨论概形上模的导出范畴。这里讨论的大部分材料可见于
\cite{TT}、\cite{Bokstedt-Neeman}、\cite{BvdB} 和 \cite{LN}。
当然还有许多其他参考文献。


\section{约定}
\label{perfect-section-conventions}

\noindent
若 $\mathcal{A}$ 为 Abel 范畴，且 $M$ 为 $\mathcal{A}$ 的对象，则也用
$M$ 表示 $K(\mathcal{A})$ 和/或 $D(\mathcal{A})$ 中对应于如下复形的
对象：它在次数 $0$ 上为 $M$，在所有其他次数上为零。

\medskip\noindent
若有环 $A$，则 $K(A)$ 表示 $A$-模复形的同伦范畴，$D(A)$ 表示相应的
导出范畴。类似地，若有赋环空间 $(X, \mathcal{O}_X)$，则
$K(\mathcal{O}_X)$ 表示 $\mathcal{O}_X$-模复形的同伦范畴，
$D(\mathcal{O}_X)$ 表示相应的导出范畴。










\section{拟凝聚模的导出范畴}
\label{perfect-section-derived-quasi-coherent}

\noindent
本节讨论概形 $X$ 上拟凝聚模与所有模之间的关系。可参考
\cite[Appendix B]{TT}。由 Schemes, Section
\ref{schemes-section-quasi-coherent} 中的讨论，嵌入
$\QCoh(\mathcal{O}_X) \subset \textit{Mod}(\mathcal{O}_X)$
把 $\QCoh(\mathcal{O}_X)$ 展示为 $\mathcal{O}_X$-模范畴的弱 Serre
子范畴。用
$$
D_\QCoh(\mathcal{O}_X) \subset D(\mathcal{O}_X)
$$
表示上同调层为拟凝聚层的复形所成的子范畴，见 Derived Categories,
Section \ref{derived-section-triangulated-sub}。于是得到典范函子
\begin{equation}
\label{perfect-equation-compare}
D(\QCoh(\mathcal{O}_X))
\longrightarrow
D_\QCoh(\mathcal{O}_X)
\end{equation}
，见 Derived Categories, Equation (\ref{derived-equation-compare})。

\begin{lemma}
\label{perfect-lemma-quasi-coherence-direct-sums}
设 $X$ 为概形。那么 $D_\QCoh(\mathcal{O}_X)$ 有直和。
\end{lemma}

\begin{proof}
由 Injectives, Lemma \ref{injectives-lemma-derived-products}，导出范畴
$D(\mathcal{O}_X)$ 有直和，并且可通过对任意代表逐项取直和来计算。由于
取直和（在任意 Grothendieck Abel 范畴中）是正合函子，直和的上同调层
显然是各上同调层的直和。拟凝聚层的直和仍为拟凝聚层，故结论成立，见
Schemes, Section \ref{schemes-section-quasi-coherent}。
\end{proof}

\noindent
我们需要一些关于导出极限的知识。提醒读者：在下面的引理中，导出极限通常
不是 $D_\QCoh$ 的对象。

\begin{lemma}
\label{perfect-lemma-Rlim-quasi-coherent}
设 $X$ 为概形。设 $(K_n)$ 为 $D_\QCoh(\mathcal{O}_X)$ 中的逆系统，其在
$D(\mathcal{O}_X)$ 中的导出极限为 $K = R\lim K_n$。假设对所有
$q \in \mathbf{Z}$ 及 $n \geq 1$，态射
$H^q(K_{n + 1}) \to H^q(K_n)$ 都是满射。那么：
\begin{enumerate}
\item $H^q(K) = \lim H^q(K_n)$；
\item $R\lim H^q(K_n) = \lim H^q(K_n)$；并且
\item 对每个仿射开集 $U \subset X$，当 $p > 0$ 时有
$H^p(U, \lim H^q(K_n)) = 0$。
\end{enumerate}
\end{lemma}

\begin{proof}
令 $\mathcal{B}$ 为 $X$ 的仿射开集全体。因为 $H^q(K_n)$ 拟凝聚，由
Cohomology of Schemes, Lemma
\ref{coherent-lemma-quasi-coherent-affine-cohomology-zero}，对
$U \in \mathcal{B}$ 有 $H^p(U, H^q(K_n)) = 0$。此外，由 Schemes,
Lemma \ref{schemes-lemma-equivalence-quasi-coherent}，对
$U \in \mathcal{B}$，映射
$H^0(U, H^q(K_{n + 1})) \to H^0(U, H^q(K_n))$ 是满射。我们刚验证了
Cohomology, Lemma \ref{cohomology-lemma-derived-limit-suitable-system} 的
条件，故 (1) 由其推出。(2) 与 (3) 由 Cohomology, Lemma
\ref{cohomology-lemma-inverse-limit-is-derived-limit} 推出。
\end{proof}

\noindent
下面的引理将帮助我们在 $D_\QCoh(\mathcal{O}_X)$ 的对象上“计算”右导出
函子。

\begin{lemma}
\label{perfect-lemma-nice-K-injective}
设 $X$ 为概形，设 $E$ 为 $D_\QCoh(\mathcal{O}_X)$ 的对象。那么 Derived
Categories, Remark
\ref{derived-remark-map-into-derived-limit-truncations}
中的态射 $E \to R\lim \tau_{\geq -n}E$ 是同构\footnote{特别地，由
Derived Categories, Lemma \ref{derived-lemma-difficulty-K-injectives}，
$E$ 有 K-内射代表。}。
\end{lemma}

\begin{proof}
用 $\mathcal{H}^i = H^i(E)$ 表示 $E$ 的第 $i$ 个上同调层。令
$\mathcal{B}$ 为 $X$ 的仿射开子集全体。由 Cohomology of Schemes,
Lemma \ref{coherent-lemma-quasi-coherent-affine-cohomology-zero}，对所有
$p > 0$、所有 $i \in \mathbf{Z}$ 以及所有 $U \in \mathcal{B}$，有
$H^p(U, \mathcal{H}^i) = 0$。因此结论由 Cohomology, Lemma
\ref{cohomology-lemma-is-limit-dimension} 推出。
\end{proof}

\begin{lemma}
\label{perfect-lemma-application-nice-K-injective}
设 $X$ 为概形。设 $F : \textit{Mod}(\mathcal{O}_X) \to \textit{Ab}$ 为
加性函子，$N \geq 0$ 为整数。假设：
\begin{enumerate}
\item $F$ 与可数直积交换；
\item 对所有 $p \geq N$ 及拟凝聚 $\mathcal{F}$，有
$R^pF(\mathcal{F}) = 0$。
\end{enumerate}
那么对 $E \in D_\QCoh(\mathcal{O}_X)$：
\begin{enumerate}
\item 当 $i \leq a$ 时，
$H^i(RF(\tau_{\leq a}E)) \to H^i(RF(E))$ 是同构；
\item 当 $i \geq b$ 时，
$H^i(RF(E)) \to H^i(RF(\tau_{\geq b - N + 1}E))$ 是同构；
\item 若对某个 $-\infty \leq a \leq b \leq \infty$，当
$i \not \in [a, b]$ 时有 $H^i(E) = 0$，则当
$i \not \in [a, b + N - 1]$ 时有 $H^i(RF(E)) = 0$。
\end{enumerate}
\end{lemma}

\begin{proof}
陈述 (1) 即 Derived Categories, Lemma
\ref{derived-lemma-negative-vanishing}。

\medskip\noindent
证明陈述 (2)。记 $E_n = \tau_{\geq -n}E$。由 Lemma
\ref{perfect-lemma-nice-K-injective}，有 $E = R\lim E_n$。于是由 Injectives,
Lemma \ref{injectives-lemma-RF-commutes-with-Rlim}，在
$D(\textit{Ab})$ 中有 $RF(E) = R\lim RF(E_n)$。因此对每个
$i \in \mathbf{Z}$，有短正合列
$$
0 \to R^1\lim H^{i - 1}(RF(E_n)) \to H^i(RF(E)) \to \lim H^i(RF(E_n)) \to 0
$$
，见 More on Algebra, Remark
\ref{more-algebra-remark-compare-derived-limit}。为证明 (2)，我们将说明左端
项为零，并且对任意满足 $i \geq b$ 的 $b$，右端项等于
$H^i(RF(E_{-b + N - 1}))$。

\medskip\noindent
对每个 $n$，有可区别三角
$$
H^{-n}(E)[n] \to E_n \to E_{n - 1} \to H^{-n}(E)[n + 1]
$$
，它位于 $D(\mathcal{O}_X)$ 中（Derived Categories, Remark
\ref{derived-remark-truncation-distinguished-triangle}）。因为
$H^{-n}(E)$ 拟凝聚，有
$$
H^i(RF(H^{-n}(E)[n])) = R^{i + n}F(H^{-n}(E)) = 0
$$
（当 $i + n \geq N$ 时），以及
$$
H^i(RF(H^{-n}(E)[n + 1])) = R^{i + n + 1}F(H^{-n}(E)) = 0
$$
（当 $i + n + 1 \geq N$ 时）。由此可知
$$
H^i(RF(E_n)) \to H^i(RF(E_{n - 1}))
$$
在 $n \geq N - i$ 时是同构。因此系统 $H^i(RF(E_n))$ 都满足 ML 条件，
短正合列中的 $R^1\lim$ 项为零（见 More on Algebra, Section
\ref{more-algebra-section-Rlim} 中的讨论）。此外，系统
$H^i(RF(E_n))$ 从 $n = N - i - 1$ 起为常值，正合所需。

\medskip\noindent
证明 (3)。在关于 $E$ 的假设下，有 $\tau_{\leq a - 1}E = 0$，故由 (1)
得到：当 $i \leq a - 1$ 时，$H^i(RF(E))$ 为零。类似地，
$\tau_{\geq b + 1}E = 0$，故由 (2) 得到：当 $i \geq b + N$ 时，
$H^i(RF(E))$ 为零。
\end{proof}

\noindent
下面的引理是本章许多结果的关键要素。

\begin{lemma}
\label{perfect-lemma-affine-compare-bounded}
设 $X = \Spec(A)$ 为仿射概形。图中的所有函子
$$
\xymatrix{
D(\QCoh(\mathcal{O}_X)) \ar[rr]_{(\ref{perfect-equation-compare})}
& &
D_\QCoh(\mathcal{O}_X) \ar[ld]^{R\Gamma(X, -)} \\
& D(A) \ar[lu]^{\widetilde{\ \ }}
}
$$
都是三角范畴的等价。此外，对 $D_\QCoh(\mathcal{O}_X)$ 中的 $E$，有
$H^0(X, E) = H^0(X, H^0(E))$。
\end{lemma}

\begin{proof}
函子 $R\Gamma(X, -)$ 给出函子 $D(\mathcal{O}_X) \to D(A)$，因而限制后
给出函子
\begin{equation}
\label{perfect-equation-back}
R\Gamma(X, -) : D_\QCoh(\mathcal{O}_X) \longrightarrow D(A).
\end{equation}
利用 $X$ 上拟凝聚模与 $A$-模范畴之间的等价，我们将说明此函子是
(\ref{perfect-equation-compare}) 的拟逆。

\medskip\noindent
说明。用 $(Y, \mathcal{O}_Y)$ 表示其环层由 $A$ 给出的单点空间，用
$\pi : (X, \mathcal{O}_X) \to (Y, \mathcal{O}_Y)$
表示显然的赋环空间态射。于是 $R\Gamma(X, -)$ 可等同于 $R\pi_*$，并且
通过等价
$\textit{Mod}(\mathcal{O}_Y) = \text{Mod}_A = \QCoh(\mathcal{O}_X)$
，函子 (\ref{perfect-equation-compare}) 可等同于
$L\pi^* = \pi^* = \widetilde{\ }$（见 Modules, Lemma
\ref{modules-lemma-construct-quasi-coherent-sheaves} 以及 Schemes, Lemmas
\ref{schemes-lemma-compare-constructions} 和
\ref{schemes-lemma-equivalence-quasi-coherent}）。因此函子
$$
\xymatrix{
D(A) \ar@<1ex>[r] & D(\mathcal{O}_X) \ar@<1ex>[l]
}
$$
互为伴随（由 Cohomology, Lemma \ref{cohomology-lemma-adjoint}）。特别地，
得到典范伴随映射
$$
a : \widetilde{R\Gamma(X, E)} \longrightarrow E
$$
（对 $D(\mathcal{O}_X)$ 中的 $E$），以及
$$
b : M^\bullet \longrightarrow R\Gamma(X, \widetilde{M^\bullet})
$$
（对 $A$-模复形 $M^\bullet$）。

\medskip\noindent
设 $E$ 为 $D_\QCoh(\mathcal{O}_X)$ 的对象。可以把 Lemma
\ref{perfect-lemma-application-nice-K-injective} 应用于函子
$F(-) = \Gamma(X, -)$；由 Cohomology of Schemes, Lemma
\ref{coherent-lemma-quasi-coherent-affine-cohomology-zero}，可取
$N = 1$。因此
$$
H^0(R\Gamma(X, E)) = H^0(R\Gamma(X, \tau_{\geq 0}E)) = \Gamma(X, H^0(E))
$$
（最后一个等号来自典范截断的定义）。利用这一点，我们将说明伴随映射
$a$ 与 $b$ 诱导同构 $H^0(a)$ 与 $H^0(b)$。于是 $a$ 与 $b$ 都是拟同构
（因为该陈述在平移下不变），从而引理得证。

\medskip\noindent
在两种情形下，都使用 $\widetilde{\ }$ 为正合函子这一事实（Schemes,
Lemma \ref{schemes-lemma-spec-sheaves}）。也就是说，这说明
$$
H^0\left(\widetilde{R\Gamma(X, E)}\right) =
\widetilde{H^0(R\Gamma(X, E))} =
\widetilde{\Gamma(X, H^0(E))}
$$
等于 $H^0(E)$，因为 $H^0(E)$ 拟凝聚。因此 $H^0(a)$ 是同构。另一方面，
有
$$
H^0(R\Gamma(X, \widetilde{M^\bullet})) =
\Gamma(X, H^0(\widetilde{M^\bullet})) =
\Gamma(X, \widetilde{H^0(M^\bullet)}) =
H^0(M^\bullet)
$$
，这证明 $H^0(b)$ 是同构。
\end{proof}

\begin{lemma}
\label{perfect-lemma-affine-K-flat}
设 $X = \Spec(A)$ 为仿射概形。若 $K^\bullet$ 为 $A$-模的 K-平坦复形，
则 $\widetilde{K^\bullet}$ 为 $\mathcal{O}_X$-模的 K-平坦复形。
\end{lemma}

\begin{proof}
由 More on Algebra, Lemma \ref{more-algebra-lemma-base-change-K-flat}，对每个
$\mathfrak p \in \Spec(A)$，$K^\bullet \otimes_A A_\mathfrak p$ 是
$A_\mathfrak p$-模的 K-平坦复形。因此由 Cohomology, Lemma
\ref{cohomology-lemma-check-K-flat-stalks}（以及 Schemes, Lemma
\ref{schemes-lemma-spec-sheaves}），可知 $\widetilde{K^\bullet}$ K-平坦。
\end{proof}

\begin{lemma}
\label{perfect-lemma-quasi-coherence-pushforward}
若 $f : X \to Y$ 为由环映射 $A \to B$ 给出的仿射概形态射，则图
$$
\xymatrix{
D(B) \ar[d] \ar[r] & D_\QCoh(\mathcal{O}_X) \ar[d]^{Rf_*} \\
D(A) \ar[r] & D_\QCoh(\mathcal{O}_Y)
}
$$
交换。
\end{lemma}

\begin{proof}
由 Lemma \ref{perfect-lemma-affine-compare-bounded} 以及
$R\Gamma(Y, Rf_*K) = R\Gamma(X, K)$ 得证；后一个等式见 Cohomology,
Lemma \ref{cohomology-lemma-Leray-unbounded}。
\end{proof}

\begin{lemma}
\label{perfect-lemma-quasi-coherence-pullback}
设 $f : Y \to X$ 为概形态射。
\begin{enumerate}
\item 函子 $Lf^*$ 把 $D_\QCoh(\mathcal{O}_X)$ 映入
$D_\QCoh(\mathcal{O}_Y)$。
\item 若 $X$ 与 $Y$ 都仿射，且 $f$ 由环映射 $A \to B$ 给出，则图
$$
\xymatrix{
D(B) \ar[r] & D_\QCoh(\mathcal{O}_Y) \\
D(A) \ar[r] \ar[u]^{- \otimes_A^\mathbf{L} B} &
D_\QCoh(\mathcal{O}_X) \ar[u]_{Lf^*}
}
$$
交换。
\end{enumerate}
\end{lemma}

\begin{proof}
先证明图
$$
\xymatrix{
D(B) \ar[r] & D(\mathcal{O}_Y) \\
D(A) \ar[r] \ar[u]^{- \otimes_A^\mathbf{L} B} &
D(\mathcal{O}_X) \ar[u]_{Lf^*}
}
$$
交换。这由 Lemma \ref{perfect-lemma-affine-K-flat} 及有关函子的构造立即可知。为证
(1)，设 $E$ 为 $D_\QCoh(\mathcal{O}_X)$ 的对象。要说明 $Lf^*E$ 的
上同调层拟凝聚，可以在 $X$ 上局部地工作。注意 $Lf^*$ 与限制到开子概形
相容。因此可假设 $f$ 是 (2) 中那样的仿射概形态射。由 Lemma
\ref{perfect-lemma-affine-compare-bounded}，$E$ 来自一个 $A$-模复形。由证明中
第一个图的交换性，$Lf^*E$ 也有同样性质，从而 (1) 成立。
\end{proof}

\begin{lemma}
\label{perfect-lemma-quasi-coherence-tensor-product}
设 $X$ 为概形。
\begin{enumerate}
\item 对 $D_\QCoh(\mathcal{O}_X)$ 的对象 $K, L$，导出张量积
$K \otimes^\mathbf{L}_{\mathcal{O}_X} L$ 属于
$D_\QCoh(\mathcal{O}_X)$。
\item 若 $X = \Spec(A)$ 仿射，则
$$
\widetilde{M^\bullet} \otimes_{\mathcal{O}_X}^\mathbf{L} \widetilde{K^\bullet}
=
\widetilde{M^\bullet \otimes_A^\mathbf{L} K^\bullet}
$$
对任意一对 $A$-模复形 $K^\bullet$、$M^\bullet$ 成立。
\end{enumerate}
\end{lemma}

\begin{proof}
(2) 中的等式由 Lemma \ref{perfect-lemma-affine-K-flat} 和导出张量积的构造立即
推出。为证 (1)，设 $K, L$ 为 $D_\QCoh(\mathcal{O}_X)$ 的对象。要检验
$K \otimes^\mathbf{L} L$ 属于 $D_\QCoh(\mathcal{O}_X)$，可以在 $X$ 上
局部地工作，故可假设 $X = \Spec(A)$ 仿射。由 Lemma
\ref{perfect-lemma-affine-compare-bounded}，可用 $A$-模复形表示 $K$ 与 $L$。
于是结论由 (2) 推出。
\end{proof}





\section{全直像}
\label{perfect-section-total-direct-image}

\noindent
下面的引理是 Cohomology of Schemes, Lemma
\ref{coherent-lemma-quasi-coherence-higher-direct-images} 的类似物。

\begin{lemma}
\label{perfect-lemma-quasi-coherence-direct-image}
设 $f : X \to S$ 为概形态射。假设 $f$ 拟分离且拟紧。
\begin{enumerate}
\item 函子 $Rf_*$ 把 $D_\QCoh(\mathcal{O}_X)$ 映入
$D_\QCoh(\mathcal{O}_S)$。
\item 若 $S$ 拟紧，则存在整数 $N = N(X, S, f)$，使得对
$D_\QCoh(\mathcal{O}_X)$ 的对象 $E$，只要对 $m > 0$ 有
$H^m(E) = 0$，就对 $m \geq N$ 有 $H^m(Rf_*E) = 0$。
\item 事实上，若 $S$ 拟紧，则可以找到 $N = N(X, S, f)$，使得对每个
概形态射 $S' \to S$，函子 $R(f')_*$ 都有同样结论，其中
$f' : X' \to S'$ 是 $f$ 的基变换。
\end{enumerate}
\end{lemma}

\begin{proof}
设 $E$ 为 $D_\QCoh(\mathcal{O}_X)$ 的对象。为证 (1)，必须说明 $Rf_*E$
有拟凝聚上同调层。该问题在 $S$ 上为局部问题，故可假设 $S$ 拟紧。按照
Cohomology of Schemes, Lemma
\ref{coherent-lemma-quasi-coherence-higher-direct-images} 选取
$N = N(X, S, f)$。于是对所有拟凝聚 $\mathcal{O}_X$-模 $\mathcal{F}$
及所有 $p \geq N$，有 $R^pf_*\mathcal{F} = 0$，并且基变换后仍然如此。

\medskip\noindent
先假设 $E$ 下有界。我们将说明，对这样的 $E$，以上选取的 $N$ 使 (1)、
(2) 与 (3) 成立。在这种情形下，例如可使用谱序列
$$
R^pf_*H^q(E) \Rightarrow R^{p + q}f_*E
$$
（Derived Categories, Lemma
\ref{derived-lemma-two-ss-complex-functor}），结合
$R^pf_*H^q(E)$ 的拟凝聚性以及 $R^pf_*H^q(E)$ 在 $p \geq N$ 时的消没，
得到此情形下
(1)、(2) 与 (3) 成立。

\medskip\noindent
下面证明 (2) 与 (3)。假设当 $m > 0$ 时 $H^m(E) = 0$。设
$U \subset S$ 为仿射开集。由 Cohomology of Schemes, Lemma
\ref{coherent-lemma-quasi-coherence-higher-direct-images-application}
以及 $N$ 的选取，对 $p \geq N$ 及任意拟凝聚 $\mathcal{O}_X$-模
$\mathcal{F}$，有 $H^p(f^{-1}(U), \mathcal{F}) = 0$。因此可把 Lemma
\ref{perfect-lemma-application-nice-K-injective} 应用于函子
$\Gamma(f^{-1}(U), -)$，得到
$$
R\Gamma(U, Rf_*E) = R\Gamma(f^{-1}(U), E)
$$
在次数 $\geq N$ 上的上同调消没。因为这对所有仿射开集 $U \subset S$
都成立，所以当 $m \geq N$ 时 $H^m(Rf_*E) = 0$。

\medskip\noindent
最后在一般情形下证明 (1)。回忆存在可区别三角
$$
\tau_{\leq -n - 1}E \to E \to \tau_{\geq -n}E \to
(\tau_{\leq -n - 1}E)[1]
$$
，它位于 $D(\mathcal{O}_X)$ 中，见 Derived Categories, Remark
\ref{derived-remark-truncation-distinguished-triangle}。由 (2)，
$Rf_*\tau_{\leq -n - 1}E$ 的上同调层在次数 $\geq -n + N$ 上消没。因此，
给定整数 $q$，可知对某个 $n$，$R^qf_*E$ 等于
$R^qf_*\tau_{\geq -n}E$，于是可应用上面的结果。
\end{proof}

\begin{lemma}
\label{perfect-lemma-acyclicity-lemma}
设 $f : X \to S$ 为拟分离且拟紧的概形态射。设
$\mathcal{F}^\bullet$ 为拟凝聚 $\mathcal{O}_X$-模复形，其每一项关于
$f_*$ 都右无环。那么 $f_*\mathcal{F}^\bullet$ 在 $D(\mathcal{O}_S)$ 中
表示 $Rf_*\mathcal{F}^\bullet$。
\end{lemma}

\begin{proof}
总有典范映射 $f_*\mathcal{F}^\bullet \to Rf_*\mathcal{F}^\bullet$。我们的
任务是说明它在上同调层上为同构。因为该陈述在平移下不变，只需证明
$H^0(f_*(\mathcal{F}^\bullet)) \to R^0f_*\mathcal{F}^\bullet$
是同构。该陈述在 $S$ 上为局部问题，故可假设 $S$ 仿射。由 Lemma
\ref{perfect-lemma-quasi-coherence-direct-image}，对所有充分大的 $n$，有
$R^0f_*\mathcal{F}^\bullet = R^0f_*\tau_{\geq -n}\mathcal{F}^\bullet$。
因此可假设 $\mathcal{F}^\bullet$ 下有界。按假设，每个
$\mathcal{F}^n$ 关于 $f_*$ 都右无环，故由 Leray 无环性引理（Derived
Categories, Lemma \ref{derived-lemma-leray-acyclicity}），
$f_*\mathcal{F}^\bullet \to Rf_*\mathcal{F}^\bullet$ 是拟同构。
\end{proof}

\begin{lemma}
\label{perfect-lemma-acyclicity-lemma-global}
设 $X$ 为拟分离拟紧概形。设 $\mathcal{F}^\bullet$ 为拟凝聚
$\mathcal{O}_X$-模复形，其每一项关于 $\Gamma(X, -)$ 都右无环。那么
$\Gamma(X, \mathcal{F}^\bullet)$ 在 $D(\Gamma(X, \mathcal{O}_X)$ 中表示
$R\Gamma(X, \mathcal{F}^\bullet)$。
\end{lemma}

\begin{proof}
把 Lemma \ref{perfect-lemma-acyclicity-lemma} 应用于典范态射
$X \to \Spec(\Gamma(X, \mathcal{O}_X))$。略去一些细节。
\end{proof}

\begin{lemma}
\label{perfect-lemma-spectral-sequence}
设 $X$ 为拟分离拟紧概形。对 $D_\QCoh(\mathcal{O}_X)$ 的任意对象 $K$，
谱序列
$$
E_2^{i, j} = H^i(X, H^j(K)) \Rightarrow H^{i + j}(X, K)
$$
（Cohomology, Example \ref{cohomology-example-spectral-sequence}）有界且收敛。
\end{lemma}

\begin{proof}
由 Cohomology, Lemma
\ref{cohomology-lemma-spectral-sequence-filtered-object} 中使用
$\tau_{\leq -p}K$ 所给滤过构造谱序列的方法可知，只需说明：给定
$n \in \mathbf{Z}$，有
$$
H^n(X, \tau_{\leq -p}K) = 0 \text{ 对 } p \gg 0
$$
以及
$$
H^n(X, K) = H^n(X, \tau_{\leq -p}K) \text{ 对 } p \ll 0
$$
。第一点来自 Lemma \ref{perfect-lemma-application-nice-K-injective}（取
$F = \Gamma(X, -)$）及 Cohomology of Schemes, Lemma
\ref{coherent-lemma-quasi-coherence-higher-direct-images} 中的界。对任意
赋环空间 $(X, \mathcal{O}_X)$ 及任意 $K \in D(\mathcal{O}_X)$，只要
$-p \leq n$，第二点就成立。
\end{proof}

\begin{lemma}
\label{perfect-lemma-quasi-coherence-pushforward-direct-sums}
设 $f : X \to S$ 为拟分离且拟紧的概形态射。那么
$Rf_* : D_\QCoh(\mathcal{O}_X) \to D_\QCoh(\mathcal{O}_S)$
与直和交换。
\end{lemma}

\begin{proof}
设 $E_i$ 为 $D_\QCoh(\mathcal{O}_X)$ 的一族对象，并令
$E = \bigoplus E_i$。要证明映射
$$
\bigoplus Rf_*E_i \longrightarrow Rf_*E
$$
是同构。我们将说明它在次数 $0$ 的上同调层上诱导同构，这就推出本引理。
为此可以在 $S$ 上局部地工作，故可以并且确实假设 $S$ 拟紧。按照 Lemma
\ref{perfect-lemma-quasi-coherence-direct-image} 选取整数 $N$。由所引引理，
$R^0f_*E = R^0f_*\tau_{\geq -N}E$ 且
$R^0f_*E_i = R^0f_*\tau_{\geq -N}E_i$。注意
$\tau_{\geq -N}E = \bigoplus \tau_{\geq -N}E_i$。因此可假设所有 $E_i$ 的
上同调层在次数 $< -N$ 上消没。接着使用谱序列
$$
R^pf_*H^q(E) \Rightarrow R^{p + q}f_*E
\quad\text{且}\quad
R^pf_*H^q(E_i) \Rightarrow R^{p + q}f_*E_i
$$
（Derived Categories, Lemma
\ref{derived-lemma-two-ss-complex-functor}），把问题约化到拟凝聚层直和的
情形。该情形由 Cohomology of Schemes, Lemma
\ref{coherent-lemma-colimit-cohomology} 处理。
\end{proof}









\section{仿射态射}
\label{perfect-section-affine-morphisms}

\noindent
本节汇集沿概形的仿射态射取直像的一些信息。

\begin{lemma}
\label{perfect-lemma-pushforward-affine-morphism}
设 $f : X \to S$ 为概形的仿射态射。设 $\mathcal{F}^\bullet$ 为拟凝聚
$\mathcal{O}_X$-模复形。那么
$f_*\mathcal{F}^\bullet = Rf_*\mathcal{F}^\bullet$.
\end{lemma}

\begin{proof}
合并 Lemma \ref{perfect-lemma-acyclicity-lemma} 与 Cohomology of Schemes, Lemma
\ref{coherent-lemma-relative-affine-vanishing}。另一种证明是在 $S$ 上仿射
局部地工作，并使用 Lemma \ref{perfect-lemma-quasi-coherence-pushforward}。
\end{proof}

\begin{lemma}
\label{perfect-lemma-affine-morphism}
设 $f : X \to S$ 为概形的仿射态射。那么
$Rf_* : D_\QCoh(\mathcal{O}_X) \to D_\QCoh(\mathcal{O}_S)$
反映同构。
\end{lemma}

\begin{proof}
该陈述是指：若 $Rf_*\alpha$ 是同构，则 $D_\QCoh(\mathcal{O}_X)$ 的态射
$\alpha : E \to F$ 是同构。可以在上同调层上检验这一点。特别地，该问题
在 $S$ 上为局部问题。因此可假设 $S$ 仿射，进而 $X$ 也仿射。在这种
情形下，由 Lemma \ref{perfect-lemma-affine-compare-bounded} 给出的导出范畴
$D_\QCoh(\mathcal{O}_X)$ 与 $D_\QCoh(\mathcal{O}_S)$ 的描述，结论显然。
略去一些细节。
\end{proof}

\begin{lemma}
\label{perfect-lemma-affine-morphism-pull-push}
设 $f : X \to S$ 为概形的仿射态射。对 $D_\QCoh(\mathcal{O}_S)$ 中的
$E$，有
$Rf_* Lf^* E = E \otimes^\mathbf{L}_{\mathcal{O}_S} f_*\mathcal{O}_X$.
\end{lemma}

\begin{proof}
因为 $f$ 仿射，映射 $f_*\mathcal{O}_X \to Rf_*\mathcal{O}_X$ 是同构
（Cohomology of Schemes, Lemma
\ref{coherent-lemma-relative-affine-vanishing}）。存在典范映射
$E \otimes^\mathbf{L} f_*\mathcal{O}_X =
E \otimes^\mathbf{L} Rf_*\mathcal{O}_X \to Rf_* Lf^* E$，它伴随于映射
$$
Lf^*(E \otimes^\mathbf{L} Rf_*\mathcal{O}_X) =
Lf^*E \otimes^\mathbf{L} Lf^*Rf_*\mathcal{O}_X \longrightarrow
Lf^* E \otimes^\mathbf{L} \mathcal{O}_X = Lf^* E
$$
；后者来自 $1 : Lf^*E \to Lf^*E$ 以及典范映射
$Lf^*Rf_*\mathcal{O}_X \to \mathcal{O}_X$。要检验这样构造的映射为同构，
可以在 $S$ 上局部地工作。因此可假设 $S$ 仿射，进而 $X$ 也仿射。在这种
情形下，由 Lemmas \ref{perfect-lemma-affine-compare-bounded} 和
\ref{perfect-lemma-quasi-coherence-pullback} 给出的导出范畴
$D_\QCoh(\mathcal{O}_X)$、$D_\QCoh(\mathcal{O}_S)$ 及函子 $Lf^*$ 的
描述，结论显然。略去一些细节。
\end{proof}

\noindent
设 $Y$ 为概形。设 $\mathcal{A}$ 为 $\mathcal{O}_Y$-代数层。用
$D_\QCoh(\mathcal{A})$ 表示 $D_\QCoh(\mathcal{O}_X)$ 在限制函子
$D(\mathcal{A}) \to D(\mathcal{O}_X)$ 下的逆像。换言之，
$K \in D(\mathcal{A})$ 属于 $D_\QCoh(\mathcal{A})$，当且仅当它的上同调层
作为 $\mathcal{O}_X$-模是拟凝聚的。若 $\mathcal{A}$ 本身拟凝聚，则这
等价于要求上同调层作为 $\mathcal{A}$-模拟凝聚，见 Morphisms, Lemma
\ref{morphisms-lemma-affine-equivalence-modules}。

\begin{lemma}
\label{perfect-lemma-affine-morphism-equivalence}
设 $f : X \to Y$ 为概形的仿射态射。那么 $f_*$ 诱导等价
$$
\Phi : D_\QCoh(\mathcal{O}_X) \longrightarrow D_\QCoh(f_*\mathcal{O}_X)
$$
，它与 $D_\QCoh(f_*\mathcal{O}_X) \to D_\QCoh(\mathcal{O}_Y)$ 的复合为
$Rf_* : D_\QCoh(\mathcal{O}_X) \to D_\QCoh(\mathcal{O}_Y)$。
\end{lemma}

\begin{proof}
回忆：对对象 $K \in D_\QCoh(\mathcal{O}_X)$，计算 $Rf_*$ 的方法是选取表示
$K$ 的 $\mathcal{O}_X$-模 K-内射复形 $\mathcal{I}^\bullet$，再取
$f_*\mathcal{I}^\bullet$。因此令 $\Phi(K)$ 为复形
$f_*\mathcal{I}^\bullet$，把它视作 $f_*\mathcal{O}_X$-模复形。用
$g : (X, \mathcal{O}_X) \to (Y, f_*\mathcal{O}_X)$ 表示显然的赋环空间
态射。那么 $g$ 是赋环空间的平坦态射（茎的描述见下文），且 $\Phi$ 是
$Rg_*$ 在 $D_\QCoh(\mathcal{O}_X)$ 上的限制。我们断言 $Lg^*$ 是拟逆。
首先注意，$Lg^*$ 把 $D_\QCoh(f_*\mathcal{O}_X)$ 映入
$D_\QCoh(\mathcal{O}_X)$，因为 $g^*$ 把拟凝聚模变成拟凝聚模（Modules,
Lemma \ref{modules-lemma-pullback-quasi-coherent}）。为完成证明，只需说明
伴随映射
$$
Lg^*\Phi(K) = Lg^*Rg_*K \to K
\quad\text{且}\quad
M \to Rg_*Lg^*M = \Phi(Lg^*M)
$$
对 $K \in D_\QCoh(\mathcal{O}_X)$ 及
$M \in D_\QCoh(f_*\mathcal{O}_X)$ 都是同构。这是局部问题，故可假设
$Y$ 仿射，进而 $X$ 也仿射。

\medskip\noindent
假设 $Y = \Spec(B)$ 且 $X = \Spec(A)$。设
$\mathfrak p = x \in \Spec(A) = X$ 为映到
$\mathfrak q = y \in \Spec(B) = Y$ 的点。那么
$(f_*\mathcal{O}_X)_y = A_\mathfrak q$ 且
$\mathcal{O}_{X, x} = A_\mathfrak p$，故 $g$ 平坦。因此 $g^*$ 正合，并且
对 $D(f_*\mathcal{O}_X)$ 中任意 $M$，有
$H^i(Lg^*M) = g^*H^i(M)$。对 $K \in D_\QCoh(\mathcal{O}_X)$，由高阶直像
的消没（Cohomology of Schemes, Lemma
\ref{coherent-lemma-relative-affine-vanishing}）以及 Lemma
\ref{perfect-lemma-application-nice-K-injective}（略去一个小细节），可知
$$
H^i(\Phi(K)) = H^i(Rf_*K) = f_*H^i(K)
$$
。因此只需说明
$$
g^*g_*\mathcal{F} \to \mathcal{F}
\quad\text{且}\quad
\mathcal{G} \to g_*g^*\mathcal{G}
$$
都是同构，其中 $\mathcal{F}$ 是拟凝聚 $\mathcal{O}_X$-模，而
$\mathcal{G}$ 是拟凝聚 $f_*\mathcal{O}_X$-模。这由 Morphisms, Lemma
\ref{morphisms-lemma-affine-equivalence-modules} 推出。
\end{proof}





\section{以闭子集为支集的上同调}
\label{perfect-section-cohomology-support}

\noindent
我们针对概形和拟凝聚模，详细阐述 Cohomology, Sections
\ref{cohomology-section-cohomology-support} 与
\ref{cohomology-section-cohomology-support-bis} 中的材料。

\begin{definition}
\label{perfect-definition-supported-on}
设 $X$ 为概形。设 $E$ 为 $D(\mathcal{O}_X)$ 的对象。设
$T \subset X$ 为闭子集。如果上同调层 $H^i(E)$ 的支集都在 $T$ 上，则称
$E$ {\it 支集在 $T$ 上}。
\end{definition}

\noindent
在该定义的情形下，我们复述 Cohomology, Section
\ref{cohomology-section-cohomology-support-bis} 中的部分讨论。设 $X$ 为
概形，设 $T \subset X$ 为闭子集。支集包含在 $T$ 中的
$\mathcal{O}_X$-模所成范畴，是所有 $\mathcal{O}_X$-模所成范畴的 Serre
子范畴，见 Homology, Definition
\ref{homology-definition-serre-subcategory} 和 Modules, Lemma
\ref{modules-lemma-support-section-closed}。因此可应用 Derived Categories,
Section \ref{derived-section-triangulated-sub} 中的讨论，得到
$D(\mathcal{O}_X)$ 的一个严格满、饱和三角子范畴，其对象都支集在 $T$
上；下文把该范畴记作 $D_T(\mathcal{O}_X)$。

\medskip\noindent
在 Definition \ref{perfect-definition-supported-on} 的情形下，用 $i : T \to X$
表示包含映射。回忆 Cohomology, Section
\ref{cohomology-section-cohomology-support-bis}：在这种情形下有函子
$R\mathcal{H}_T : D(\mathcal{O}_X) \to D(i^{-1}\mathcal{O}_X)$
，它是 $i_* : D(i^{-1}\mathcal{O}_X) \to D(\mathcal{O}_X)$ 的右伴随。

\begin{lemma}
\label{perfect-lemma-support-quasi-coherent}
设 $X$ 为概形。设 $T \subset X$ 为闭子集，使 $X \setminus T$ 是 $X$ 的
逆紧开子集。设 $i : T \to X$ 为包含映射。
\begin{enumerate}
\item 对 $D_\QCoh(\mathcal{O}_X)$ 中的 $E$，有
$i_*R\mathcal{H}_T(E)$ 属于 $D_{\QCoh, T}(\mathcal{O}_X)$。
\item 函子
$i_* \circ R\mathcal{H}_T : D_\QCoh(\mathcal{O}_X) \to
D_{\QCoh, T}(\mathcal{O}_X)$ 是包含函子
$D_{\QCoh, T}(\mathcal{O}_X) \to D_\QCoh(\mathcal{O}_X)$ 的右伴随。
\end{enumerate}
\end{lemma}

\begin{proof}
令 $U = X \setminus T$，并用 $j : U \to X$ 表示包含映射。由 Cohomology,
Lemma \ref{cohomology-lemma-triangle-sections-with-support-sheaves}，存在
可区别三角
$$
i_*R\mathcal{H}_T(E) \to E \to Rj_*(E|_U) \to i_*R\mathcal{H}_Z(E)[1]
$$
，它位于 $D(\mathcal{O}_X)$ 中。由 Lemma
\ref{perfect-lemma-quasi-coherence-direct-image}，复形 $Rj_*(E|_U)$ 有拟凝聚
上同调层（此处使用 $U$ 在 $X$ 中逆紧）。因此 (1) 成立。(2) 由此以及
Cohomology, Lemma
\ref{cohomology-lemma-complexes-with-support-on-closed} 中函子的伴随性推出。
\end{proof}

\begin{lemma}
\label{perfect-lemma-support-direct-sums}
设 $X$ 为概形。设 $T \subset X$ 为闭子集，使 $X \setminus T$ 是 $X$ 的
逆紧开子集。那么对 $D_\QCoh(\mathcal{O}_X)$ 的一族对象 $E_i$（
$i \in I$），有
$R\mathcal{H}_T(\bigoplus E_i) = \bigoplus R\mathcal{H}_T(E_i)$.
\end{lemma}

\begin{proof}
令 $U = X \setminus T$，并用 $j : U \to X$ 表示包含映射。由 Cohomology,
Lemma \ref{cohomology-lemma-triangle-sections-with-support-sheaves}，存在
可区别三角
$$
i_*R\mathcal{H}_T(E) \to E \to Rj_*(E|_U) \to i_*R\mathcal{H}_Z(E)[1]
$$
，它对 $D(\mathcal{O}_X)$ 中任意 $E$ 都位于 $D(\mathcal{O}_X)$ 中。由
Lemma \ref{perfect-lemma-quasi-coherence-pushforward-direct-sums}，函子
$E \mapsto Rj_*(E|_U)$ 在 $D_\QCoh(\mathcal{O}_X)$ 上与直和交换。由此，
函子 $i_* \circ R\mathcal{H}_T$ 也如此（略去细节）。由于
$i_* : D(i^{-1}\mathcal{O}_X) \to D_T(\mathcal{O}_X)$ 是等价
（Cohomology, Lemma
\ref{cohomology-lemma-complexes-with-support-on-closed}），结论成立。
\end{proof}

\begin{remark}
\label{perfect-remark-support-c-equations}
设 $X$ 为概形。设 $f_1, \ldots, f_c \in \Gamma(X, \mathcal{O}_X)$。用
$Z \subset X$ 表示由 $f_1, \ldots, f_c$ 截出的闭子概形。对
$0 \leq p < c$ 及 $1 \leq i_0 < \ldots < i_p \leq c$，用
$U_{i_0 \ldots i_p} \subset X$ 表示 $f_{i_0} \ldots f_{i_p}$ 可逆处的
开子概形。对任意 $\mathcal{O}_X$-模 $\mathcal{F}$，令
$$
\mathcal{F}_{i_0 \ldots i_p} =
(U_{i_0 \ldots i_p} \to X)_*(\mathcal{F}|_{U_{i_0 \ldots i_p}})
$$
在这种情形下，{\it 扩充交错 {\v C}ech 复形}是 $\mathcal{O}_X$-模复形
\begin{equation}
\label{perfect-equation-extended-alternating}
0 \to \mathcal{F} \to
\bigoplus\nolimits_{i_0} \mathcal{F}_{i_0} \to
\ldots \to
\bigoplus\nolimits_{i_0 < \ldots < i_p} \mathcal{F}_{i_0 \ldots i_p} \to
\ldots \to \mathcal{F}_{1 \ldots c} \to 0
\end{equation}
，其中 $\mathcal{F}$ 放在次数 $0$ 上。映射构造如下。给定
$1 \leq i_0 < \ldots < i_{p + 1} \leq c$ 及 $0 \leq j \leq p + 1$，
有典范映射
$$
\mathcal{F}_{i_0 \ldots \hat i_j \ldots i_{p + 1}} \to
\mathcal{F}_{i_0 \ldots i_p}
$$
，它来自包含
$U_{i_0 \ldots i_p} \subset U_{i_0 \ldots \hat i_j \ldots i_{p + 1}}$.
扩充交错复形的微分使用这些典范映射，并带符号 $(-1)^j$。
\end{remark}

\begin{lemma}
\label{perfect-lemma-extended-alternating-zero}
设 $X$、$f_1, \ldots, f_c \in \Gamma(X, \mathcal{O}_X)$ 及
$\mathcal{F}$ 如 Remark \ref{perfect-remark-support-c-equations}。复形
(\ref{perfect-equation-extended-alternating}) 在 $X \setminus Z$ 上的限制无环。
\end{lemma}

\noindent
指出：更一般地，对从有限开覆盖
$X \setminus Z = U_1 \cup \ldots \cup U_c$ 出发、按 Remark
\ref{perfect-remark-support-c-equations} 定义的任意扩充交错 {\v C}ech 复形，本
引理都成立。

\begin{proof}
设 $W \subset X \setminus Z$ 为开子集。在 $W$ 上取层复形
(\ref{perfect-equation-extended-alternating}) 的截面，得到复形
$$
\mathcal{F}(W) \to \bigoplus\nolimits_{i_0} \mathcal{F}(U_{i_0} \cap W) \to
\bigoplus\nolimits_{i_0 < i_1} \mathcal{F}(U_{i_0i_1} \cap W) \to \ldots
$$
换言之，得到覆盖 $W = \bigcup U_i \cap W$ 及 $\{1, \ldots, c\}$ 上标准
次序的扩充有序 {\v C}ech 复形，见 Cohomology, Section
\ref{cohomology-section-alternating-cech}。由 Cohomology, Lemma
\ref{cohomology-lemma-alternating-cech-trivial}，只要对某个
$1 \leq i \leq c$，$W$ 包含在 $V(f_i)$ 中，该复形就同伦于零。证毕。
\end{proof}

\begin{remark}
\label{perfect-remark-extended-alternating-map-to-support}
设 $X$、$f_1, \ldots, f_c \in \Gamma(X, \mathcal{O}_X)$ 及
$\mathcal{F}$ 如 Remark \ref{perfect-remark-support-c-equations}。用
$\mathcal{F}^\bullet$ 表示复形
(\ref{perfect-equation-extended-alternating})。由 Lemma
\ref{perfect-lemma-extended-alternating-zero}，$\mathcal{F}^\bullet$ 的上同调层
支集在 $Z$ 上，因而 $\mathcal{F}^\bullet$ 是 $D_Z(\mathcal{O}_X)$ 的
对象。另一方面，等式 $\mathcal{F}^0 = \mathcal{F}$ 决定
$D(\mathcal{O}_X)$ 中的典范映射
$\mathcal{F}^\bullet \to \mathcal{F}$。由于
$i_* \circ R\mathcal{H}_Z$ 是包含函子
$D_Z(\mathcal{O}_X) \to D(\mathcal{O}_X)$ 的右伴随（见 Cohomology,
Lemma \ref{cohomology-lemma-complexes-with-support-on-closed}），得到典范
交换图
$$
\xymatrix{
\mathcal{F}^\bullet \ar[rd] \ar[rr] & & \mathcal{F} \\
& i_*R\mathcal{H}_Z(\mathcal{F}) \ar[ru]
}
$$
，它位于 $D(\mathcal{O}_X)$ 中，并且关于 $\mathcal{O}_X$-模
$\mathcal{F}$ 函子性成立。
\end{remark}

\begin{lemma}
\label{perfect-lemma-extended-alternating-represented}
设 $X$、$f_1, \ldots, f_c \in \Gamma(X, \mathcal{O}_X)$ 及
$\mathcal{F}$ 如 Remark \ref{perfect-remark-support-c-equations}。若
$\mathcal{F}$ 拟凝聚，则复形 (\ref{perfect-equation-extended-alternating}) 在
$D_Z(\mathcal{O}_X)$ 中表示 $i_* R\mathcal{H}_Z(\mathcal{F})$。
\end{lemma}

\begin{proof}
用 $\mathcal{F}^\bullet$ 表示复形
(\ref{perfect-equation-extended-alternating})。本引理的陈述是指 Remark
\ref{perfect-remark-extended-alternating-map-to-support} 中的映射
$\mathcal{F}^\bullet \to i_*R\mathcal{H}_Z(\mathcal{F})$ 为同构。因为
$\mathcal{F}^\bullet$ 属于 $D_Z(\mathcal{O}_X)$（见所引评注），由
Cohomology, Lemma
\ref{cohomology-lemma-complexes-with-support-on-closed}，有
$i_*R\mathcal{H}_Z(\mathcal{F}^\bullet) = \mathcal{F}^\bullet$。态射
$U_{i_0 \ldots i_p} \to X$ 仿射，因为在 $X$ 的仿射开集上，它由把函数
$f_{i_0} \ldots f_{i_p}$ 取逆给出。因此
$$
\mathcal{F}_{i_0 \ldots i_p} =
(U_{i_0 \ldots i_p} \to X)_*\mathcal{F}|_{U_{i_0 \ldots i_p}} =
R(U_{i_0 \ldots i_p} \to X)_*\mathcal{F}|_{U_{i_0 \ldots i_p}}
$$
，这由 Cohomology of Schemes, Lemma
\ref{coherent-lemma-relative-affine-vanishing} 及 $\mathcal{F}$ 拟凝聚的假设
可知。由 Cohomology, Lemma
\ref{cohomology-lemma-sections-support-in-closed-disjoint-open}，可知
$R\mathcal{H}_Z(\mathcal{F}_{i_0 \ldots i_p}) = 0$。因此对 $p > 0$，有
$i_*R\mathcal{H}_Z(\mathcal{F}^p) = 0$。综合以上各点，得到
$$
\mathcal{F}^\bullet = i_*R\mathcal{H}_Z(\mathcal{F}^\bullet) =
i_*R\mathcal{H}_Z(\mathcal{F})
$$
，正合所需。
\end{proof}

\begin{lemma}
\label{perfect-lemma-supported-trivial-vanishing}
设 $X$ 为概形。设 $T \subset X$ 为闭子集，它局部可由结构层的至多 $c$
个元素截出。那么对 $i > c$ 及任意拟凝聚 $\mathcal{O}_X$-模
$\mathcal{F}$，有 $\mathcal{H}^i_Z(\mathcal{F}) = 0$。
\end{lemma}

\begin{proof}
这立即来自 Lemma \ref{perfect-lemma-extended-alternating-represented} 给出的
$R\mathcal{H}_T(\mathcal{F})$ 的局部描述。
\end{proof}

\begin{lemma}
\label{perfect-lemma-supported-vanishing}
设 $X$ 为概形。设 $Z \subset X$ 为闭子集，它局部可由含 $c$ 个元素的
Koszul 正则序列截出。那么对每个平坦拟凝聚 $\mathcal{O}_X$-模
$\mathcal{F}$，当 $i \not = c$ 时有
$\mathcal{H}^i_Z(\mathcal{F}) = 0$。
\end{lemma}

\begin{proof}
由 Lemma \ref{perfect-lemma-extended-alternating-represented} 给出的
$R\mathcal{H}_Z(\mathcal{F})$ 的描述，问题归结为如下代数陈述：给定环
$R$、Koszul 正则序列 $f_1, \ldots, f_c \in R$ 以及平坦 $R$-模 $M$，
扩充交错 {\v C}ech 复形
$M \to \bigoplus\nolimits_{i_0} M_{f_{i_0}} \to
\bigoplus\nolimits_{i_0 < i_1} M_{f_{i_0}f_{i_1}} \to
\ldots \to M_{f_1 \ldots f_c}$
（来自 More on Algebra, Section
\ref{more-algebra-section-alternating-cech}）仅在次数 $c$ 上有上同调。由
More on Algebra, Lemma
\ref{more-algebra-lemma-vanishing-extended-alternating-koszul}，得到 $R$ 的
扩充交错 {\v C}ech 复形所需的消没。由于 $M$ 的复形是把它与平坦
$R$-模 $M$ 作张量所得（More on Algebra, Lemma
\ref{more-algebra-lemma-extended-alternating-form-module}），结论成立。
\end{proof}

\begin{remark}
\label{perfect-remark-supported-map-c-equations}
设 $X$、$f_1, \ldots, f_c \in \Gamma(X, \mathcal{O}_X)$ 及
$\mathcal{F}$ 如 Remark \ref{perfect-remark-support-c-equations}。存在典范的
$\mathcal{O}_X|_Z$-线性映射
$$
c_{f_1, \ldots, f_c} :
i^*\mathcal{F}
\longrightarrow
\mathcal{H}^c_Z(\mathcal{F})
$$
，它关于 $\mathcal{F}$ 函子性成立。事实上，用 $\mathcal{F}^\bullet$
表示扩充交错 {\v C}ech 复形 (\ref{perfect-equation-extended-alternating})，则有
Remark \ref{perfect-remark-extended-alternating-map-to-support} 中的典范映射
$\mathcal{F}^\bullet \to i_*R\mathcal{H}_Z(\mathcal{F})$
。这在次数 $c$ 的上同调层上决定典范映射
$$
\Coker\left(\bigoplus \mathcal{F}_{1 \ldots \hat i \ldots c} \to
\mathcal{F}_{1 \ldots c}\right)
\longrightarrow
i_*\mathcal{H}^c_Z(\mathcal{F})
$$
。给定 $\mathcal{F}$ 的局部截面 $s$，可以考虑
$\mathcal{F}_{1 \ldots c}$ 的局部截面
$$
\frac{s}{f_1 \ldots f_c}
$$
。该截面在上述余核中的类，只依赖于 $s$ 模去映射
$(f_1, \ldots, f_c) : \mathcal{F}^{\oplus c} \to \mathcal{F}$.
的像。由于 $i_*i^*\mathcal{F}$ 等于
$(f_1, \ldots, f_c) : \mathcal{F}^{\oplus c} \to \mathcal{F}$
的余核，故得到 $\mathcal{O}_X$-模映射
$i_*i^*\mathcal{F} \to i_*\mathcal{H}_Z^c(\mathcal{F})$.
因为 $i_*$ 全忠实，得到映射 $c_{f_1, \ldots, f_c}$。
\end{remark}

\begin{example}
\label{perfect-example-affine-supported-map-c-equations}
设 $X = \Spec(A)$ 仿射，$f_1, \ldots, f_c \in A$，并设
$\mathcal{F} = \widetilde{M}$，其中 $M$ 为某个 $A$-模。Remark
\ref{perfect-remark-supported-map-c-equations} 中的映射
$c_{f_1, \ldots, f_c}$ 可描述为映射
$$
M/(f_1, \ldots, f_c)M
\longrightarrow
\Coker\left(
\bigoplus M_{f_1 \ldots \hat f_i \ldots f_c} \to M_{f_1 \ldots f_c}
\right)
$$
，它把 $s \in M$ 的类映到余核中 $s/f_1 \ldots f_c$ 的类。
\end{example}

\begin{lemma}
\label{perfect-lemma-supported-map-determinant}
设 $X$、$f_1, \ldots, f_c \in \Gamma(X, \mathcal{O}_X)$ 及
$\mathcal{F}$ 如 Remark \ref{perfect-remark-support-c-equations}。对
$1 \leq i, j \leq c$，设
$a_{ji} \in \Gamma(X, \mathcal{O}_X)$，并令
$g_j = \sum_{i = 1, \ldots, c} a_{ji}f_i$。假设 $g_1, \ldots, g_c$ 在概形
意义下截出 $Z$。若 $\mathcal{F}$ 拟凝聚，则
$$
c_{f_1, \ldots, f_c} = \det(a_{ji}) c_{g_1, \ldots, g_c}
$$
，其中 $c_{f_1, \ldots, f_c}$ 与 $c_{g_1, \ldots, g_c}$ 如 Remark
\ref{perfect-remark-supported-map-c-equations} 所述。
\end{lemma}

\begin{proof}
我们将证明：对任意 $s \in \mathcal{F}(X)$，作为
$\mathcal{H}_Z(\mathcal{F})$ 的整体截面，有
$c_{f_1, \ldots, f_c}(s) =
\det(a_{ij}) c_{g_1, \ldots, g_c}(s)$。这已经足够，因为随后对 $X$ 的任意
开子概形上的截面也得到同样结果。为此，对
$1 \leq i_0 < \ldots < i_p \leq c$ 及
$1 \leq j_0 < \ldots < j_q \leq c$，用
$U_{i_0 \ldots i_p} \subset X$,
$V_{j_0 \ldots j_q} \subset X$, and
$W_{i_0 \ldots i_p, j_0 \ldots j_q} \subset X$
分别表示 $f_{i_0} \ldots f_{i_p}$ 可逆处、
$g_{j_0} \ldots g_{j_q}$ 可逆处以及
$f_{i_0} \ldots f_{i_p}g_{j_0} \ldots g_{j_q}$ 可逆处的开子概形。用
$\mathcal{F}_{i_0 \ldots i_p}$、$\mathcal{F}'_{j_0 \ldots j_q}$、
$\mathcal{F}''_{i_0 \ldots i_p, j_0 \ldots j_q}$ 分别表示 $\mathcal{F}$
在 $U_{i_0 \ldots i_p}$、$V_{j_0 \ldots j_q}$、
$W_{i_0 \ldots i_p, j_0 \ldots j_q}$ 上的限制到 $X$ 的直像。于是得到三个
扩充交错 {\v C}ech 复形
$$
\mathcal{F}^\bullet :
\mathcal{F} \to \bigoplus\nolimits_{i_0} \mathcal{F}_{i_0} \to
\bigoplus\nolimits_{i_0 < i_1} \mathcal{F}_{i_0i_1} \to \ldots
$$
以及
$$
(\mathcal{F}')^\bullet :
\mathcal{F} \to \bigoplus\nolimits_{j_0} \mathcal{F}'_{j_0} \to
\bigoplus\nolimits_{j_0 < j_1} \mathcal{F}'_{j_0j_1} \to \ldots
$$
以及
$$
(\mathcal{F}'')^\bullet :
\mathcal{F} \to
\bigoplus\nolimits_{i_0} \mathcal{F}_{i_0} \oplus
\bigoplus\nolimits_{j_0} \mathcal{F}'_{j_0} \to
\bigoplus\nolimits_{i_0 < i_1} \mathcal{F}_{i_0i_1} \oplus
\bigoplus\nolimits_{i_0, j_0} \mathcal{F}''_{i_0, j_0} \oplus
\bigoplus\nolimits_{j_0 < j_1} \mathcal{F}'_{j_0j_1} \to
\ldots
$$
，其微分就是定义 (\ref{perfect-equation-extended-alternating}) 时所用的微分。存在
复形映射
$$
(\mathcal{F}'')^\bullet \to \mathcal{F}^\bullet
\quad\text{且}\quad
(\mathcal{F}'')^\bullet \to (\mathcal{F}')^\bullet
$$
，它们由各项上的投影映射给出（因而在次数 $0$ 上诱导恒等映射）。注意，由
Lemma \ref{perfect-lemma-extended-alternating-represented}，这些复形中的每一个都
表示 $i_*R\mathcal{H}_Z(\mathcal{F})$，而这些映射表示该对象上的恒等。
因此只需找到元素
$$
\sigma \in H^c((\mathcal{F}'')^\bullet(X))
$$
，使它在这两个映射下分别映到 $c_{f_1, \ldots, f_c}(s)$ 和
$\det(a_{ji})c_{g_1, \ldots, g_c}(s)$。事实上，可以显式给出表示
$\sigma$ 的余圈。具体地，取
$$
\sigma_{1 \ldots c} = \frac{s}{f_1 \ldots f_c} \in \mathcal{F}_{1 \ldots c}(X)
\quad\text{且}\quad
\sigma'_{1 \ldots c} = \frac{\det(a_{ji})s}{g_1 \ldots g_c} \in
\mathcal{F}'_{1 \ldots c}(X)
$$
，并取
$$
\sigma_{i_0 \ldots i_p, j_0 \ldots j_{c - p - 2}} =
\frac{\lambda(i_0 \ldots i_p, j_0 \ldots j_{c - p - 2})s}%
{f_{i_0} \ldots f_{i_p}g_{j_0} \ldots g_{j_{c - p - 2}}}
\in \mathcal{F}''_{i_0 \ldots i_p, j_0 \ldots j_{c - p - 2}}(X)
$$
，其中 $\lambda(i_0 \ldots i_p, j_0 \ldots j_{c - p - 2})$ 是形式表达式
$$
e_{i_0} \wedge \ldots \wedge e_{i_p} \wedge
(a_{j_01} e_1 + \ldots + a_{j_0c}e_c) \wedge \ldots \wedge
(a_{j_{c - p - 2}1} e_1 + \ldots + a_{j_{c - p - 2}c}e_c)
$$
中 $e_1 \wedge \ldots \wedge e_c$ 的系数。为验证 $\sigma$ 是余圈，必须
证明对 $1 \leq i_0 < \ldots < i_p \leq c$ 及
$1 \leq j_0 < \ldots < j_{c - p - 1} \leq c$，有
\begin{align*}
0 & =
\sum\nolimits_{a = 0, \ldots, p} (-1)^a
f_{i_a} \lambda(i_0 \ldots \hat i_a \ldots i_p, j_0 \ldots j_{c - p - 1}) \\
& +
\sum\nolimits_{b = 0, \ldots, c - p - 1} (-1)^{p + b + 1}g_{j_b}
\lambda(i_0 \ldots i_p, j_0 \ldots \hat j_b \ldots j_{c - p - 1})
\end{align*}
最容易的验证方法或许是指出，形式表达式
$$
\xi = e_{i_0} \wedge \ldots \wedge e_{i_p} \wedge
(a_{j_01} e_1 + \ldots + a_{j_0c}e_c) \wedge \ldots \wedge
(a_{j_{c - p - 1}1} e_1 + \ldots + a_{j_{c - p - 1}c}e_c)
$$
为 $0$，因为它是以 $e_1, \ldots, e_c$ 为基的自由模的第 $(c + 1)$ 个
外幂中的元素；并且上面的表达式是 $\xi$ 在把 $e_i \to f_i$ 的 Koszul
微分下的像。略去一些细节。
\end{proof}

\begin{lemma}
\label{perfect-lemma-supported-map-global}
设 $X$ 为概形。设 $Z \to X$ 为有限表示闭浸入，并且其余法层
$\mathcal{C}_{Z/X}$ 是秩为 $c$ 的局部自由层。则有典范映射
$$
c :
\wedge^c(\mathcal{C}_{Z/X})^\vee \otimes_{\mathcal{O}_Z} i^*\mathcal{F}
\longrightarrow
\mathcal{H}_Z^c(\mathcal{F})
$$
它关于拟凝聚模 $\mathcal{F}$ 函子性地成立。
\end{lemma}

\begin{proof}
这由注 \ref{perfect-remark-supported-map-c-equations} 中的构造，以及引理
\ref{perfect-lemma-supported-map-determinant} 所证明的、对理想层生成元选取的
无关性得到。略去一些细节。
\end{proof}

\begin{remark}
\label{perfect-remark-supported-functorial}
设 $g : X' \to X$ 为概形态射，并设
$f_1, \ldots, f_c \in \Gamma(X, \mathcal{O}_X)$。令
$f'_i = g^\sharp(f_i) \in \Gamma(X', \mathcal{O}_{X'})$。分别以
$Z \subset X$ 和 $Z' \subset X'$ 表示由 $f_1, \ldots, f_c$ 和
$f'_1, \ldots, f'_c$ 截出的闭子概形。于是
$Z' = Z \times_X X'$。以 $h : Z' \to Z$ 表示所诱导的概形态射。
设 $\mathcal{F}$ 为 $\mathcal{O}_X$-模，并令
$\mathcal{F}' = g^*\mathcal{F}$。在此情形下，若 $\mathcal{F}$ 拟凝聚，
则图表
$$
\xymatrix{
(i')^{-1}\mathcal{O}_{X'} \otimes_{h^{-1}i^{-1}\mathcal{O}_X}
h^{-1}\mathcal{H}^c_Z(\mathcal{F}) \ar[r] &
\mathcal{H}_{Z'}^c(\mathcal{F}') \\
h^*i^*\mathcal{F} \ar[r] \ar[u]_-{c_{f_1, \ldots, f_c}} &
(i')^*\mathcal{F}' \ar[u]^-{c_{f'_1, \ldots, f'_c}}
}
$$
交换，其中上方水平箭头是在 $c$ 次上同调层上取 Cohomology, Remark
\ref{cohomology-remark-support-functorial} 中的映射。具体而言，分别以
$\mathcal{F}^\bullet$ 和 $(\mathcal{F}')^\bullet$ 表示在注
\ref{perfect-remark-support-c-equations} 中用
$\mathcal{F}, f_1, \ldots, f_c$ 和
$\mathcal{F}', f'_1, \ldots, f'_c$ 构造的扩充交错 {\v C}ech 复形。
注意 $(\mathcal{F}')^\bullet = g^*\mathcal{F}^\bullet$。那么，即使不假设
$\mathcal{F}$ 拟凝聚，图表
$$
\xymatrix{
i'_* L(g|_{Z'})^* R\mathcal{H}_Z(\mathcal{F}) \ar[r] \ar@{=}[d] &
i'_*R\mathcal{H}_{Z'}(Lg^*\mathcal{F}) \ar[d] \\
Lg^*i_*R\mathcal{H}_Z(\mathcal{F}) &
i'_*R\mathcal{H}_{Z'}(\mathcal{F}') \\
Lg^*(\mathcal{F}^\bullet) \ar[u] \ar[r] &
(\mathcal{F}')^\bullet \ar[u]
}
$$
也交换，其中 $g|_{Z'} : (Z', (i')^{-1}\mathcal{O}_{X'}) \to
(Z, i^{-1}\mathcal{O}_X)$ 是所诱导的带环空间态射。这里上方水平箭头
由 Cohomology, Remark \ref{cohomology-remark-support-functorial} 给出，
等号也在那里得到解释。向上的箭头来自注
\ref{perfect-remark-extended-alternating-map-to-support}。下方水平箭头是映射
$Lg^*\mathcal{F}^\bullet
\to g^*\mathcal{F}^\bullet = (\mathcal{F}')^\bullet$，而向下的箭头由
$Lg^*\mathcal{F} \to g^*\mathcal{F} = \mathcal{F}'$ 诱导。此图表交换，
因为沿两条路径绕行图表所得的两个箭头
$Lg^*\mathcal{F}^\bullet \to
i'_*R\mathcal{H}_{Z'}(\mathcal{F}')$，在与
$i'_*R\mathcal{H}_{Z'}(\mathcal{F}') \to \mathcal{F}'$ 复合后，均为
典范映射 $Lg^*\mathcal{F}^\bullet \to \mathcal{F}'$。略去一些细节。
现在考察此图表在 $c$ 次上同调层上的情形，并利用注
\ref{perfect-remark-supported-map-c-equations} 中构造的映射
$i^*\mathcal{F} \to
\Coker(\bigoplus \mathcal{F}_{1 \ldots \hat i \ldots c} \to
\mathcal{F}_{1 \ldots c})$
与拉回相容，即得第一个图表的交换性。
\end{remark}








\section{拟凝聚化子}
\label{perfect-section-coherator}

\noindent
设 $X$ 为概形。所谓 {\it coherator}（拟凝聚化子）是函子
$$
Q_X :
\textit{Mod}(\mathcal{O}_X)
\longrightarrow
\QCoh(\mathcal{O}_X)
$$
它是包含函子
$\QCoh(\mathcal{O}_X) \to \textit{Mod}(\mathcal{O}_X)$ 的右伴随。
它对任意概形 $X$ 都存在；此外，对每个拟凝聚模 $\mathcal{F}$，伴随映射
$Q_X(\mathcal{F}) \to \mathcal{F}$ 都是同构，见 Properties,
Proposition \ref{properties-proposition-coherator}。由于 $Q_X$ 左正合（因为
它是右伴随），可以考虑其右导出延拓
$$
RQ_X :
D(\mathcal{O}_X)
\longrightarrow
D(\QCoh(\mathcal{O}_X)).
$$
由于 $Q_X$ 是包含函子
$\QCoh(\mathcal{O}_X) \to \textit{Mod}(\mathcal{O}_X)$ 的右伴随，
由 Derived Categories, Lemma
\ref{derived-lemma-derived-adjoint-functors} 可知，$RQ_X$ 是典范函子
$D(\QCoh(\mathcal{O}_X)) \to D(\mathcal{O}_X)$ 的右伴随。

\medskip\noindent
本节研究函子 $RQ_X$。在第 \ref{perfect-section-better-coherator} 节中，我们将研究
包含函子 $D_\QCoh(\mathcal{O}_X) \to D(\mathcal{O}_X)$ 的右伴随（当它
存在时）；二者密切相关。

\begin{lemma}
\label{perfect-lemma-affine-pushforward}
设 $f : X \to Y$ 为概形的仿射态射。则 $f_*$ 定义导出函子
$f_* : D(\QCoh(\mathcal{O}_X)) \to D(\QCoh(\mathcal{O}_Y))$，并且图表
$$
\xymatrix{
D(\QCoh(\mathcal{O}_X)) \ar[d]_{f_*} \ar[r] &
D_\QCoh(\mathcal{O}_X) \ar[d]^{Rf_*} \\
D(\QCoh(\mathcal{O}_Y)) \ar[r] &
D_\QCoh(\mathcal{O}_Y)
}
$$
交换。
\end{lemma}

\begin{proof}
函子
$f_* : \QCoh(\mathcal{O}_X) \to \QCoh(\mathcal{O}_Y)$
正合，见 Cohomology of Schemes, Lemma
\ref{coherent-lemma-relative-affine-vanishing}。因此 $f_*$ 定义如下导出函子；
具体只需把 $f_*$ 作用于任意代表复形：
$f_* : D(\QCoh(\mathcal{O}_X)) \to D(\QCoh(\mathcal{O}_Y))$
，见 Derived Categories, Lemma
\ref{derived-lemma-right-derived-exact-functor}。图表由引理
\ref{perfect-lemma-pushforward-affine-morphism} 交换。
\end{proof}

\begin{lemma}
\label{perfect-lemma-flat-pushforward-coherator}
设 $f : X \to Y$ 为概形态射，并假设 $f$ 拟紧、拟分离且平坦。以
$$
\Phi : D(\QCoh(\mathcal{O}_X)) \to D(\QCoh(\mathcal{O}_Y))
$$
表示 $f_* : \QCoh(\mathcal{O}_X) \to \QCoh(\mathcal{O}_Y)$ 的右导出函子，
则有 $RQ_Y \circ Rf_* = \Phi \circ RQ_X$。
\end{lemma}

\begin{proof}
我们将证明 $RQ_Y \circ Rf_*$ 和 $\Phi \circ RQ_X$ 都是同一函子
$D(\QCoh(\mathcal{O}_Y)) \to D(\mathcal{O}_X)$ 的右伴随。

\medskip\noindent
由于 $f$ 拟紧且拟分离，$f_*$ 保持拟凝聚性，见 Schemes, Lemma
\ref{schemes-lemma-push-forward-quasi-coherent}。回忆
$\QCoh(\mathcal{O}_X)$ 是 Grothendieck 阿贝尔范畴（Properties,
Proposition \ref{properties-proposition-coherator}）。因此
$D(\QCoh(\mathcal{O}_X))$ 中任意 $K$ 都可由
$\QCoh(\mathcal{O}_X)$ 中的 K-内射复形 $\mathcal{I}^\bullet$ 表示，见
Injectives, Theorem
\ref{injectives-theorem-K-injective-embedding-grothendieck}。于是可定义
$\Phi(K) = f_*\mathcal{I}^\bullet$。

\medskip\noindent
由于 $f$ 平坦，函子 $f^*$ 正合。因此 $f^*$ 同时定义
$f^* : D(\mathcal{O}_Y) \to D(\mathcal{O}_X)$ 和
$f^* : D(\QCoh(\mathcal{O}_Y)) \to D(\QCoh(\mathcal{O}_X))$。函子
$f^* = Lf^* : D(\mathcal{O}_Y) \to D(\mathcal{O}_X)$ 是
$Rf_* : D(\mathcal{O}_X) \to D(\mathcal{O}_Y)$ 的左伴随，见
Cohomology, Lemma \ref{cohomology-lemma-adjoint}。同理，函子
$f^* : D(\QCoh(\mathcal{O}_Y)) \to D(\QCoh(\mathcal{O}_X))$
是 $\Phi : D(\QCoh(\mathcal{O}_X)) \to D(\QCoh(\mathcal{O}_Y))$ 的左伴随，
这由 Derived Categories, Lemma
\ref{derived-lemma-derived-adjoint-functors} 得到。

\medskip\noindent
设 $A$ 为 $D(\QCoh(\mathcal{O}_Y))$ 的对象，而 $E$ 为
$D(\mathcal{O}_X)$ 的对象。于是
\begin{align*}
\Hom_{D(\QCoh(\mathcal{O}_Y))}(A, RQ_Y(Rf_*E))
& =
\Hom_{D(\mathcal{O}_Y)}(A, Rf_*E) \\
& =
\Hom_{D(\mathcal{O}_X)}(f^*A, E) \\
& =
\Hom_{D(\QCoh(\mathcal{O}_X))}(f^*A, RQ_X(E)) \\
& =
\Hom_{D(\QCoh(\mathcal{O}_Y))}(A, \Phi(RQ_X(E)))
\end{align*}
这便给出所需结论。
\end{proof}

\begin{lemma}
\label{perfect-lemma-affine-coherator}
设 $X = \Spec(A)$ 为仿射概形。则
\begin{enumerate}
\item $Q_X : \textit{Mod}(\mathcal{O}_X) \to \QCoh(\mathcal{O}_X)$
是把 $\mathcal{F}$ 送到与 $A$-模 $\Gamma(X, \mathcal{F})$ 相伴的拟凝聚
$\mathcal{O}_X$-模的函子；
\item $RQ_X : D(\mathcal{O}_X) \to D(\QCoh(\mathcal{O}_X))$
是把 $E$ 送到与 $D(A)$ 的对象 $R\Gamma(X, E)$ 相伴的拟凝聚
$\mathcal{O}_X$-模复形的函子；
\item 函子 $RQ_X$ 限制在 $D_\QCoh(\mathcal{O}_X)$ 上时，给出
(\ref{perfect-equation-compare}) 的拟逆。
\end{enumerate}
\end{lemma}

\begin{proof}
由 Schemes, Lemma \ref{schemes-lemma-compare-constructions}，函子 $Q_X$ 是
$$
\mathcal{F} \mapsto \widetilde{\Gamma(X, \mathcal{F})}
$$
这立即推出 (1) 和 (2)。第三个断言由引理
\ref{perfect-lemma-affine-compare-bounded}（的证明）得到。
\end{proof}

\noindent
现在可以证明函子
$D(\QCoh(\mathcal{O}_X)) \to D_\QCoh(\mathcal{O}_X)$ 何时为等价的一个判据。

\begin{lemma}
\label{perfect-lemma-argument-proves}
设 $X$ 为拟紧拟分离概形。假设对每个仿射开集 $U \subset X$，左正合函子
$j_* : \QCoh(\mathcal{O}_U) \to \QCoh(\mathcal{O}_X)$ 的右导出函子
$$
\Phi : D(\QCoh(\mathcal{O}_U)) \to D(\QCoh(\mathcal{O}_X))
$$
可置于交换图表
$$
\xymatrix{
D(\QCoh(\mathcal{O}_U)) \ar[d]_\Phi \ar[r]_{i_U} &
D_\QCoh(\mathcal{O}_U) \ar[d]^{Rj_*} \\
D(\QCoh(\mathcal{O}_X)) \ar[r]^{i_X} &
D_\QCoh(\mathcal{O}_X)
}
$$
中。则函子 (\ref{perfect-equation-compare})
$$
D(\QCoh(\mathcal{O}_X))
\longrightarrow
D_\QCoh(\mathcal{O}_X)
$$
是等价，且其拟逆由 $RQ_X$ 给出。
\end{lemma}

\begin{proof}
设 $E$ 为 $D_\QCoh(\mathcal{O}_X)$ 的对象，$A$ 为
$D(\QCoh(\mathcal{O}_X))$ 的对象。我们必须证明伴随映射
$$
A \to RQ_X(i_X(A))
\quad\text{且}\quad
i_X(RQ_X(E)) \to E
$$
是同构。考虑命题 $H_n$：只要 $E$ 和 $i_X(A)$ 支持在 $X$ 的某个闭子集
上（定义 \ref{perfect-definition-supported-on}），且该闭子集包含在 $X$ 的 $n$ 个
仿射开集之并中，上述伴随映射就是同构。我们对 $n$ 归纳证明 $H_n$。

\medskip\noindent
初始情形：$n = 0$。此时 $E = 0$，故映射
$i_X(RQ_X(E)) \to E$ 是同构。同样，$i_X(A) = 0$。因此 $i_X(A)$ 的
上同调层为零。由于包含函子
$\QCoh(\mathcal{O}_X) \to \textit{Mod}(\mathcal{O}_X)$ 完全忠实且正合，
可知 $A$ 的上同调对象为零，即 $A = 0$，从而
$A \to RQ_X(i_X(A))$ 也是同构。

\medskip\noindent
归纳步骤。假设 $E$ 和 $i_X(A)$ 支持在 $X$ 的闭子集 $T$ 上，而该闭子集
包含在 $U_1 \cup \ldots \cup U_n$ 中，其中 $U_i \subset X$ 为
仿射开集。令 $U = U_n$。考虑可区别三角
$$
A \to \Phi(A|_U) \to A' \to A[1]
\quad\text{且}\quad
E \to Rj_*(E|_U) \to E' \to E[1]
$$
其中 $\Phi$ 如引理陈述所定义。注意，$E \to Rj_*(E|_U)$ 在
$U = U_n$ 上是拟同构。由假设 $i_X \circ \Phi = Rj_* \circ i_U$，且
$i_X(A)|_U = i_U(A|_U)$，可知 $i_X(A) \to i_X(\Phi(A|_U))$ 在 $U$
上是拟同构。因此 $i_X(A')$ 和 $E'$ 支持在 $X$ 的闭子集
$T \setminus U$ 上，而该闭子集包含在
$U_1 \cup \ldots \cup U_{n - 1}$ 中。由归纳假设，断言对 $A'$ 和 $E'$
成立。根据 Derived Categories, Lemma
\ref{derived-lemma-third-isomorphism-triangle}，只需证明映射
$$
\Phi(A|_U) \to RQ_X(i_X(\Phi(A|_U)))
\quad\text{且}\quad
i_X(RQ_X(Rj_*E|_U)) \to Rj_*(E|_U)
$$
是同构。由假设和引理 \ref{perfect-lemma-flat-pushforward-coherator}（包含态射
$j : U \to X$ 平坦、拟紧且拟分离），有
$$
RQ_X(i_X(\Phi(A|_U))) = RQ_X(Rj_*(i_U(A|_U))) = \Phi(RQ_U(i_U(A|_U)))
$$
and
$$
i_X(RQ_X(Rj_*(E|_U))) = i_X(\Phi(RQ_U(E|_U))) = Rj_*(i_U(RQ_U(E|_U)))
$$
最后，映射
$$
A|_U \to RQ_U(i_U(A|_U))
\quad\text{且}\quad
i_U(RQ_U(E|_U)) \to E|_U
$$
由引理 \ref{perfect-lemma-affine-coherator} 是同构。结论得证。
\end{proof}

\begin{proposition}
\label{perfect-proposition-quasi-compact-affine-diagonal}
设 $X$ 为对角态射仿射的拟紧概形。则函子 (\ref{perfect-equation-compare})
$$
D(\QCoh(\mathcal{O}_X))
\longrightarrow
D_\QCoh(\mathcal{O}_X)
$$
是等价，且其拟逆由 $RQ_X$ 给出。
\end{proposition}

\begin{proof}
设 $U \subset X$ 为仿射开集。由 Morphisms, Lemma
\ref{morphisms-lemma-affine-permanence}，态射 $U \to X$ 是仿射的。因此引理
\ref{perfect-lemma-argument-proves} 的假设由引理 \ref{perfect-lemma-affine-pushforward} 保证，
结论随即得到。
\end{proof}

\begin{lemma}
\label{perfect-lemma-direct-image-coherator}
设 $f : X \to Y$ 为概形态射。假设 $X$ 和 $Y$ 拟紧且具有仿射对角态射。
以
$$
\Phi : D(\QCoh(\mathcal{O}_X)) \to D(\QCoh(\mathcal{O}_Y))
$$
表示 $f_* : \QCoh(\mathcal{O}_X) \to \QCoh(\mathcal{O}_Y)$ 的右导出函子，
则图表
$$
\xymatrix{
D(\QCoh(\mathcal{O}_X)) \ar[d]_\Phi \ar[r] &
D_\QCoh(\mathcal{O}_X) \ar[d]^{Rf_*} \\
D(\QCoh(\mathcal{O}_Y)) \ar[r] &
D_\QCoh(\mathcal{O}_Y)
}
$$
交换。
\end{lemma}

\begin{proof}
注意，由命题 \ref{perfect-proposition-quasi-compact-affine-diagonal}，图表中的水平
箭头均为范畴等价。因此可把这些范畴等同（对其他具有仿射对角态射的拟紧
概形也同样处理）。引理的断言是：对 $D(\QCoh(\mathcal{O}_X))$ 中每个
$K$，典范映射 $\Phi(K) \to Rf_*(K)$ 都是同构。注意，若
$K_1 \to K_2 \to K_3 \to K_1[1]$ 是 $D(\QCoh(\mathcal{O}_X))$ 中的
可区别三角，且断言对其中两个对象成立，则对第三个也成立。

\medskip\noindent
设 $U \subset X$ 为仿射开集。由于 $X$ 的对角态射仿射，包含态射
$j : U \to X$ 也是仿射的（Morphisms, Lemma
\ref{morphisms-lemma-affine-permanence}）。同理，复合
$g = f \circ j : U \to Y$ 仿射。设 $\mathcal{I}^\bullet$ 为
$\QCoh(\mathcal{O}_U)$ 中的 K-内射复形。由于
$j_* : \QCoh(\mathcal{O}_U) \to \QCoh(\mathcal{O}_X)$ 有正合左伴随
$j^* : \QCoh(\mathcal{O}_X) \to \QCoh(\mathcal{O}_U)$，可知
$j_*\mathcal{I}^\bullet$ 是 $\QCoh(\mathcal{O}_X)$ 中的 K-内射复形，见
Derived Categories, Lemma
\ref{derived-lemma-adjoint-preserve-K-injectives}。于是
$$
\Phi(j_*\mathcal{I}^\bullet) =
f_*j_*\mathcal{I}^\bullet =
g_*\mathcal{I}^\bullet
$$
由引理 \ref{perfect-lemma-affine-pushforward}，$j_*\mathcal{I}^\bullet$ 表示
$Rj_*\mathcal{I}^\bullet$，而 $g_*\mathcal{I}^\bullet$ 表示
$Rg_*\mathcal{I}^\bullet$。另一方面，有 $Rf_* \circ Rj_* = Rg_*$。
因此 $f_*j_*\mathcal{I}^\bullet$ 表示 $Rf_*(j_*\mathcal{I}^\bullet)$。
我们断定，引理对任意形如 $j_*\mathcal{G}^\bullet$ 的复形成立，其中
$\mathcal{G}^\bullet$ 是 $U$ 上拟凝聚模的复形。（注意，若
$\mathcal{G}^\bullet \to \mathcal{I}^\bullet$ 是拟同构，则
$j_*\mathcal{G}^\bullet \to j_*\mathcal{I}^\bullet$ 也是拟同构，因为
$j_*$ 在拟凝聚模上是正合函子。）

\medskip\noindent
设 $\mathcal{F}^\bullet$ 为拟凝聚 $\mathcal{O}_X$-模的复形。设
$T \subset X$ 为闭子集，使得每个 $p$ 对应的 $\mathcal{F}^p$ 的支集均
包含在 $T$ 中。令 $n$ 为满足
$T \subset U_1 \cup \ldots \cup U_n$ 的仿射开集 $U_1, \ldots, U_n$
的最少个数，我们对此数作归纳。初始情形 $n = 0$ 显然成立。若
$n \geq 1$，令 $U = U_1$，并以 $j : U \to X$ 表示如上的开浸入。考虑
复形映射 $c : \mathcal{F}^\bullet \to j_*j^*\mathcal{F}^\bullet$，得到
两个复形短正合列：
$$
0 \to \Ker(c) \to \mathcal{F}^\bullet \to \Im(c) \to 0
$$
and
$$
0 \to \Im(c) \to j_*j^*\mathcal{F}^\bullet \to \Coker(c) \to 0
$$
复形 $\Ker(c)$ 和 $\Coker(c)$ 支持在
$T \setminus U \subset U_2 \cup \ldots \cup U_n$ 上，故由归纳假设，
结论对它们成立。由上一段的讨论，结论对
$j_*j^*\mathcal{F}^\bullet$ 成立。考察与这些短正合列相伴的可区别三角，
并使用证明开头的观察，即得结论。
\end{proof}

\begin{remark}[警告]
\label{perfect-remark-warning-coherator}
设 $X$ 为对角态射仿射的拟紧概形。尽管由命题
\ref{perfect-proposition-quasi-compact-affine-diagonal} 已知
$D(\QCoh(\mathcal{O}_X)) = D_\QCoh(\mathcal{O}_X)$，这里仍可能出现反常
现象，也很容易出错。一个陷阱是轻率地认为这个等式意味着导出函子也相同。
例如，假设有拟紧开集 $U \subset X$。可以考虑左正合函子
$\Gamma(U, -)$ 的高阶右导出函子
$$
R^i(\QCoh)\Gamma(U, -) : \QCoh(\mathcal{O}_X) \to \textit{Ab}
$$
这是一个泛 $\delta$-函子；另一方面，对 $X$ 上所有阿贝尔层定义的函子
$H^i(U, -)$ 限制在 $\QCoh(\mathcal{O}_X)$ 上后构成一个 $\delta$-函子。
因而得到典范变换
$$
t^i : R^i(\QCoh)\Gamma(U, -) \to H^i(U, -).
$$
即使 $X = \Spec(A)$ 仿射，这些变换一般也不是同构！事实上，由右导出函子
的构造以及 $\QCoh(\mathcal{O}_X)$ 与 $\text{Mod}_A$ 的等价，若 $I$ 为内射
$A$-模，就有 $R^1(\QCoh)\Gamma(U, \widetilde{I}) = 0$。但 Examples,
Lemma \ref{examples-lemma-nonvanishing} 表明，存在 $A$、$I$ 和 $U$ 使得
$H^1(U, \widetilde{I}) \not = 0$。
\end{remark}




\section{诺特概形的拟凝聚化子}
\label{perfect-section-coherator-Noetherian}

\noindent
对于诺特概形，可以使用下列引理。

\begin{lemma}
\label{perfect-lemma-injective-quasi-coherent-sheaf-Noetherian}
设 $X$ 为诺特概形，$\mathcal{J}$ 为 $\QCoh(\mathcal{O}_X)$ 的内射对象。
则 $\mathcal{J}$ 是 $\mathcal{O}_X$-模的软层。
\end{lemma}

\begin{proof}
设 $U \subset X$ 为开子集，$s \in \mathcal{J}(U)$ 为截面。设
$\mathcal{I} \subset \mathcal{O}_X$ 为定义 $X \setminus U$ 上约化诱导概形
结构的拟凝聚理想层（见 Schemes, Definition
\ref{schemes-definition-reduced-induced-scheme}）。由 Cohomology of Schemes,
Lemma \ref{coherent-lemma-homs-over-open}，截面 $s$ 对应于某个 $n$ 下的映射
$\sigma : \mathcal{I}^n \to \mathcal{J}$。由于 $\mathcal{J}$ 是
$\QCoh(\mathcal{O}_X)$ 的内射对象，可把 $\sigma$ 延拓为映射
$\tilde s : \mathcal{O}_X \to \mathcal{J}$。于是 $\tilde s$ 对应于
$\mathcal{J}$ 的一个整体截面，其在该开集上的限制为 $s$。
\end{proof}

\begin{lemma}
\label{perfect-lemma-Noetherian-pushforward}
设 $f : X \to Y$ 为诺特概形间的态射。则拟凝聚层上的 $f_*$ 有右导出延拓
$\Phi : D(\QCoh(\mathcal{O}_X)) \to D(\QCoh(\mathcal{O}_Y))$
，使图表
$$
\xymatrix{
D(\QCoh(\mathcal{O}_X)) \ar[d]_{\Phi} \ar[r] &
D_\QCoh(\mathcal{O}_X) \ar[d]^{Rf_*} \\
D(\QCoh(\mathcal{O}_Y)) \ar[r] &
D_\QCoh(\mathcal{O}_Y)
}
$$
交换。
\end{lemma}

\begin{proof}
由于 $X$ 和 $Y$ 是诺特概形，该态射拟紧且拟分离（见 Properties, Lemma
\ref{properties-lemma-locally-Noetherian-quasi-separated} 及 Schemes, Remark
\ref{schemes-remark-quasi-compact-and-quasi-separated}）。因此 $f_*$ 保持拟凝聚
性，见 Schemes, Lemma \ref{schemes-lemma-push-forward-quasi-coherent}。接着，
设 $K$ 为 $D(\QCoh(\mathcal{O}_X))$ 的对象。由于
$\QCoh(\mathcal{O}_X)$ 是 Grothendieck 阿贝尔范畴（Properties,
Proposition \ref{properties-proposition-coherator}），可以用 K-内射复形
$\mathcal{I}^\bullet$ 表示 $K$，使每个 $\mathcal{I}^n$ 都是
$\QCoh(\mathcal{O}_X)$ 的内射对象，见 Injectives, Theorem
\ref{injectives-theorem-K-injective-embedding-grothendieck}。因此函子 $\Phi$
由下式定义：
$$
\Phi(K) = f_*\mathcal{I}^\bullet
$$
其中右端视为 $D(\QCoh(\mathcal{O}_Y))$ 的对象。为完成引理的证明，只需证明
典范映射
$$
f_*\mathcal{I}^\bullet \longrightarrow Rf_*\mathcal{I}^\bullet
$$
在 $D(\mathcal{O}_Y)$ 中是同构。由引理 \ref{perfect-lemma-acyclicity-lemma}，只需证明
对所有 $n \in \mathbf{Z}$，$\mathcal{I}^n$ 都是右 $f_*$-无环的。这成立，
因为由引理 \ref{perfect-lemma-injective-quasi-coherent-sheaf-Noetherian}，
$\mathcal{I}^n$ 是软层；而由 Cohomology, Lemma
\ref{cohomology-lemma-flasque-acyclic-pushforward}，软模是右 $f_*$-无环的。
\end{proof}

\begin{proposition}
\label{perfect-proposition-Noetherian}
设 $X$ 为诺特概形。则函子 (\ref{perfect-equation-compare})
$$
D(\QCoh(\mathcal{O}_X))
\longrightarrow
D_\QCoh(\mathcal{O}_X)
$$
是等价，且其拟逆由 $RQ_X$ 给出。
\end{proposition}

\begin{proof}
这由引理 \ref{perfect-lemma-argument-proves} 和引理
\ref{perfect-lemma-Noetherian-pushforward} 得到。
\end{proof}






\section{Koszul 复形}
\label{perfect-section-koszul}

\noindent
设 $A$ 为环，$f_1, \ldots, f_r$ 为 $A$ 中的一列元素。我们已在 More on
Algebra, Definition \ref{more-algebra-definition-koszul-complex} 中定义 Koszul
复形 $K_\bullet(f_1, \ldots, f_r)$。它是位于次数 $r, \ldots, 0$ 的链复形。
令 $K^{-n}(f_1, \ldots, f_r) = K_n(f_1, \ldots, f_r)$ 并采用相同的微分，
即可将其化为上链复形 $K^\bullet(f_1, \ldots, f_r)$。本节余下部分中的所有
复形都是上链复形。

\medskip\noindent
定义复形 $I^\bullet(f_1, \ldots, f_r)$，使得在 $K(A)$ 中有可区别三角
$$
I^\bullet(f_1, \ldots, f_r) \to
A \to
K^\bullet(f_1, \ldots, f_r) \to
I^\bullet(f_1, \ldots, f_r)[1]
$$
换言之，令
$$
I^i(f_1, \ldots, f_r) =
\left\{
\begin{matrix}
K^{i - 1}(f_1, \ldots, f_r) & \text{若 } i \leq 0 \\
0 & \text{否则}
\end{matrix}
\right.
$$
并采用 $K^\bullet(f_1, \ldots, f_r)$ 上微分的负值。可区别三角中的映射是
显然的映射。注意，$I^0(f_1, \ldots, f_r) = A^{\oplus r} \to A$ 在第
$i$ 个因子上由乘以 $f_i$ 给出。因此
$I^\bullet(f_1, \ldots, f_r) \to A$ 分解为
$$
I^\bullet(f_1, \ldots, f_r) \to I \to A
$$
其中 $I = (f_1, \ldots, f_r)$。事实上，有短正合列
$$
0 \to H^{-1}(K^\bullet(f_1, \ldots, f_s)) \to
H^0(I^\bullet(f_1, \ldots, f_s)) \to I \to 0
$$
并且对每个 $i < 0$，有
$H^i(I^\bullet(f_1, \ldots, f_r)) = H^{i - 1}(K^\bullet(f_1, \ldots, f_r))$.
注意，给定 $A$ 中的另一列元素 $g_1, \ldots, g_r$，有典范映射
$$
I^\bullet(f_1g_1, \ldots, f_rg_r) \to I^\bullet(f_1, \ldots, f_r)
\quad\text{且}\quad
K^\bullet(f_1g_1, \ldots, f_rg_r) \to K^\bullet(f_1, \ldots, f_r)
$$
它们与上述映射相容。其中第一个映射在
$I^0(f_1g_1, \ldots, f_rg_r) = A^{\oplus r}$ 的第 $i$ 个直和项上由乘以
$g_i$ 给出。特别地，给定 $f_1, \ldots, f_r$，得到复形的逆向系统
\begin{equation}
\label{perfect-equation-system}
I^\bullet(f_1, \ldots, f_r) \leftarrow
I^\bullet(f_1^2, \ldots, f_r^2) \leftarrow
I^\bullet(f_1^3, \ldots, f_r^3) \leftarrow \ldots
\end{equation}
它将在后文中起重要作用。为便于表述以下引理，我们固定一些记号。

\begin{situation}
\label{perfect-situation-complex}
这里 $A$ 为环，$f_1, \ldots, f_r$ 为 $A$ 中的一列元素。令
$X = \Spec(A)$ 及 $U = D(f_1) \cup \ldots \cup D(f_r) \subset X$。以
$\mathcal{U} : U = \bigcup_{i = 1, \ldots, r} D(f_i)$ 表示 $U$ 的这个给定
开覆盖。
\end{situation}

\noindent
第一个引理说明，上述复形可用于计算 $U$ 上拟凝聚层的上同调。

\begin{lemma}
\label{perfect-lemma-alternating-cech-complex}
在情形 \ref{perfect-situation-complex} 中，设 $M$ 为 $A$-模，并以 $\mathcal{F}$
表示相伴的 $\mathcal{O}_X$-模。则有复形的典范同构
$$
\Psi : \colim_e \Hom_A(I^\bullet(f_1^e, \ldots, f_r^e), M)
\longrightarrow
\check{\mathcal{C}}_{alt}^\bullet(\mathcal{U}, \mathcal{F})
$$
它关于 $M$ 函子性地成立，其中 $\Hom$-复形上的微分是
$I^\bullet(f_1^e, \ldots, f_r^e)$ 上微分的逆步变换。
\end{lemma}

\begin{proof}
回忆，交错 {\v C}ech 复形是通常 {\v C}ech 复形中由交错上链组成的子复形，
见 Cohomology, Section \ref{cohomology-section-alternating-cech}。通常把
$\check{\mathcal{C}}_{alt}^\bullet(\mathcal{U}, \mathcal{F})$ 中的 $p$-上链
视为 $\{1, \ldots, r\}^{p + 1}$ 上的交错函数 $s$，其在
$(i_0, \ldots, i_p)$ 处的值 $s_{i_0\ldots i_p}$ 属于
$M_{f_{i_0}\ldots f_{i_p}} = \mathcal{F}(U_{i_0\ldots i_p})$。另一方面，
$\Hom^\bullet(I^\bullet(f_1^e, \ldots, f_r^e), M)$ 中的 $p$-上链 $t$ 是映射
$t : \wedge^{p + 1}(A^{\oplus r}) \to M$。以 $[i] \in A^{\oplus r}$ 表示
第 $i$ 个基向量，并记
$$
[i_0, \ldots, i_p] = [i_0] \wedge \ldots \wedge [i_p]
\in \wedge^{p + 1}(A^{\oplus r})
$$
对上述 $t$，令
$$
\Psi(t)_{i_0 \ldots i_p} =
(-1)^p \frac{t([i_0, \ldots, i_p])}{f_{i_0}^e\ldots f_{i_p}^e}
$$
显然 $\Psi(t)$ 是交错上链。上述规则与系统的转移映射相容，因为
(\ref{perfect-equation-system}) 的转移映射
$$
I^\bullet(f_1^e, \ldots, f_r^e) \leftarrow
I^\bullet(f_1^{e + 1}, \ldots, f_r^{e + 1}),
$$
把 $[i_0, \ldots, i_p]$ 送到
$f_{i_0}\ldots f_{i_p}[i_0, \ldots, i_p]$。根据 Algebra, Lemma
\ref{algebra-lemma-localization-colimit} 中对局部化
$M_{f_{i_0} \ldots f_{i_p}}$ 的描述，显然规则 $\Psi$ 在余极限的次数 $p$
部分上定义上链模的同构。为完成证明，还需证明此映射与微分相容。对此，
对上述 $t$ 计算：
\begin{align*}
d(\Psi(t))_{i_0 \ldots i_{p + 1}}
& =
\sum\nolimits_{j = 0}^{p + 1} (-1)^j
\Psi(t)_{i_0\ldots \hat i_j \ldots i_{p + 1}} \\
& =
(-1)^p
\sum\nolimits_{j = 0}^{p + 1} (-1)^j t([i_0 \ldots \hat i_j \ldots i_{p + 1}])
(f_{i_0} \ldots \hat f_{i_j} \ldots f_{i_p})^{-e}
\end{align*}
回忆，$I^\bullet(f_1^e, \ldots, f_r^e)$ 上的微分是
$K^\bullet(f_1, \ldots, f_r)$ 上微分的负值。因此
\begin{align*}
\Psi(d(t))_{i_0 \ldots i_{p + 1}} & =
(-1)^{p + 1}
d(t)([i_0, \ldots, i_{p + 1}]) (f_{i_0} \ldots f_{i_{p + 1}})^{-e} \\
& =
(-1)^{p + 1}
t(d([i_0, \ldots, i_{p + 1}])) (f_{i_0} \ldots f_{i_{p + 1}})^{-e} \\
& =
(-1)^{p + 1}
t(-\sum\nolimits_{j = 0}^{p + 1}
(-1)^j f_{i_j}^e [i_0, \ldots, \hat i_j, \ldots i_{p + 1}])
(f_{i_0} \ldots f_{i_{p + 1}})^{-e} \\
& =
-(-1)^{p + 1}
\sum\nolimits_{j = 0}^{p + 1}
(-1)^j t([i_0, \ldots, \hat i_j, \ldots i_{p + 1}])
(f_{i_0} \ldots \hat f_{i_j} \ldots f_{i_p})^{-e}
\end{align*}
两个公式相同，证明完成。
\end{proof}

\noindent
设给定 $A$-模的有限复形 $I^\bullet$ 和 $A$-模复形 $M^\bullet$，则有相应的
Hom 复形 $\Hom^\bullet(I^\bullet, M^\bullet)$。此复形的 $n$ 次项为
$$
\bigoplus\nolimits_{p + q = n} \Hom_A(I^{-q}, M^p)
$$
其微分如 More on Algebra, Section
\ref{more-algebra-section-hom-complexes} 所述。由于复形 $I^\bullet$ 只有有限个
非零项，上述直和是有限直和。由一个复形的 {\v C}ech 复形取相伴总复形的
约定如 Cohomology, Section
\ref{cohomology-section-cech-cohomology-of-complexes} 所述。

\begin{lemma}
\label{perfect-lemma-alternating-cech-complex-complex}
在情形 \ref{perfect-situation-complex} 中，设 $M^\bullet$ 为 $A$-模复形，并以
$\mathcal{F}^\bullet$ 表示相伴的 $\mathcal{O}_X$-模复形。则有复形的典范同构
$$
\colim_e \Hom^\bullet(I^\bullet(f_1^e, \ldots, f_r^e), M^\bullet)
\longrightarrow
\text{Tot}(\check{\mathcal{C}}_{alt}^\bullet(\mathcal{U}, \mathcal{F}^\bullet))
$$
它关于 $M^\bullet$ 函子性地成立。
\end{lemma}

\begin{proof}
考虑 Cohomology, Section
\ref{cohomology-section-cech-cohomology-of-complexes} 中讨论的双复形
$F^{\bullet, \bullet}$，其项为
$F^{p, q} = \mathcal{C}_{alt}^p(\mathcal{U}, \mathcal{F}^q)$。再考虑双复形
$G^{\bullet, \bullet}$，其项为
$G^{p, q} = \colim_e \Hom_A(I^{-p}(f_1^e, \ldots, f_r^e), M^q)$，微分由
函子性给出（不插入符号）。引理 \ref{perfect-lemma-alternating-cech-complex} 的证明中
所构造的映射 $\psi^{p, q} : G^{p, q} \to F^{p, q}$ 是同构，并与微分
$d_1$（由该引理）和 $d_2$（显然）相容。然而，引理中箭头左右两端复形的
微分 $d$ 具有不同符号。具体地，对 $g \in G^{p, q}$，微分为
$$
d(g) =  d_2(g) - (-1)^{p + q} d_1(g) 
$$
（见 More on Algebra, Section \ref{more-algebra-section-hom-complexes}），而对
$f \in F^{p, q}$，微分为
$$
d(f) = d_1(f) + (-1)^p d_2(f)
$$
因此，把 $\psi^{p, q}$ 乘以 $(-1)^{pq + p(p - 1)/2}$ 即可校正符号。
\end{proof}

\begin{lemma}
\label{perfect-lemma-alternating-cech-complex-complex-computes-cohomology}
在情形 \ref{perfect-situation-complex} 中，设 $\mathcal{F}^\bullet$ 为拟凝聚
$\mathcal{O}_X$-模的复形。则在 $D(A)$ 中有典范同构
$$
\text{Tot}(\check{\mathcal{C}}_{alt}^\bullet(\mathcal{U}, \mathcal{F}^\bullet))
\longrightarrow
R\Gamma(U, \mathcal{F}^\bullet)
$$
它关于 $\mathcal{F}^\bullet$ 函子性地成立。
\end{lemma}

\begin{proof}
设 $\mathcal{B}$ 为 $U$ 的仿射开集之集。由于仿射概形上拟凝聚模的高阶
上同调群为零（Cohomology of Schemes, Lemma
\ref{coherent-lemma-quasi-coherent-affine-cohomology-zero}），这是 Cohomology,
Lemma \ref{cohomology-lemma-alternating-cech-complex-complex-ss} 的一个特例。
\end{proof}

\noindent
在情形 \ref{perfect-situation-complex} 中，以 $I_e$ 表示在引理
\ref{perfect-lemma-affine-compare-bounded} 的等价下，与 $A$-模复形
$I^\bullet(f_1^e, \ldots, f_r^e)$ 对应的 $D(\mathcal{O}_X)$ 的对象。映射
(\ref{perfect-equation-system}) 给出系统
$$
I_1 \leftarrow
I_2 \leftarrow
I_3 \leftarrow \ldots
$$
此外，有相容的映射系统 $I_e \to \mathcal{O}_X$，它们限制到 $U$ 上后成为
同构。因此对 $D(\mathcal{O}_X)$ 的每个对象 $E$，都有典范映射
\begin{equation}
\label{perfect-equation-comparison}
\colim_e \Hom_{D(\mathcal{O}_X)}(I_e, E) \longrightarrow H^0(U, E)
\end{equation}
其构造是把映射 $I_e \to E$ 送到它在 $U$ 上的限制，并使用等式
$\Hom_{D(\mathcal{O}_U)}(\mathcal{O}_U, E|_U) = H^0(U, E)$.

\begin{proposition}
\label{perfect-proposition-represent-cohomology-class-on-open}
在情形 \ref{perfect-situation-complex} 中，对 $D_\QCoh(\mathcal{O}_X)$ 的每个对象
$E$，映射 (\ref{perfect-equation-comparison}) 都是同构。
\end{proposition}

\begin{proof}
由引理 \ref{perfect-lemma-affine-compare-bounded}，可假设 $E$ 由拟凝聚层的复形
$\mathcal{F}^\bullet$ 给出。令
$M^\bullet = \Gamma(X, \mathcal{F}^\bullet)$ 为相应的 $A$-模复形。由引理
\ref{perfect-lemma-alternating-cech-complex-complex} 和
\ref{perfect-lemma-alternating-cech-complex-complex-computes-cohomology}，有拟同构
$$
\colim_e \Hom^\bullet(I^\bullet(f_1^e, \ldots, f_r^e), M^\bullet)
\longrightarrow
\text{Tot}(\check{\mathcal{C}}_{alt}^\bullet(\mathcal{U}, \mathcal{F}^\bullet))
\longrightarrow
R\Gamma(U, \mathcal{F}^\bullet)
$$
由 More on Algebra, Lemma \ref{more-algebra-lemma-RHom-out-of-projective} 和
Equation (\ref{more-algebra-equation-h0-RHom})，对复形
$\Hom^\bullet(I^\bullet(f_1^e, \ldots, f_r^e), M^\bullet)$
取 $H^0$ 即计算 $D(A)$ 中的 $\Hom$。因此两端取 $H^0$，得到
$$
\colim_e \Hom_{D(A)}(I^\bullet(f_1^e, \ldots, f_r^e), M^\bullet)
=
H^0(U, E)
$$
由引理 \ref{perfect-lemma-affine-compare-bounded}，
$\Hom_{D(A)}(I^\bullet(f_1^e, \ldots, f_r^e), M^\bullet) =
\Hom_{D(\mathcal{O}_X)}(I_e, E)$，结论得证。
\end{proof}

\noindent
在情形 \ref{perfect-situation-complex} 中，以 $K_e$ 表示在引理
\ref{perfect-lemma-affine-compare-bounded} 的等价下，与 $A$-模复形
$K^\bullet(f_1^e, \ldots, f_r^e)$ 对应的 $D(\mathcal{O}_X)$ 的对象。于是有
可区别三角
$$
I_e \to \mathcal{O}_X \to K_e \to I_e[1]
$$
及系统
$$
K_1 \leftarrow
K_2 \leftarrow
K_3 \leftarrow \ldots
$$
它与系统 $(I_e)$ 相容。此外，还有相容的映射系统
$$
K_e \to H^0(K_e) = \mathcal{O}_X/(f_1^e, \ldots, f_r^e)
$$

\begin{lemma}
\label{perfect-lemma-represent-cohomology-class-on-closed}
在情形 \ref{perfect-situation-complex} 中，设 $E$ 为
$D_\QCoh(\mathcal{O}_X)$ 的对象。假设对
$i = - r + 1, \ldots, 0$ 有 $H^i(E)|_U = 0$。那么给定
$s \in H^0(X, E)$，存在 $e \geq 0$ 及态射 $K_e \to E$，使得 $s$ 属于映射
$H^0(X, K_e) \to H^0(X, E)$ 的像。
\end{lemma}

\begin{proof}
由于 $U$ 被 $r$ 个仿射开集覆盖，对任意拟凝聚模，在 $j \geq r$ 时有
$H^j(U, \mathcal{F}) = 0$（Cohomology of Schemes, Lemma
\ref{coherent-lemma-vanishing-nr-affines}）。由引理
\ref{perfect-lemma-application-nice-K-injective}，$H^0(U, E)$ 等于
$H^0(U, \tau_{\geq -r + 1}E)$。有谱序列
$$
H^j(U, H^i(\tau_{\geq -r + 1}E)) \Rightarrow H^{i + j}(U, \tau_{\geq -N}E)
$$
见 Derived Categories, Lemma \ref{derived-lemma-two-ss-complex-functor}。因此由
对 $E$ 的上同调层所作的消失假设，$H^0(U, E) = 0$，从而 $s|_U = 0$。
把 $s$ 看作 $D(\mathcal{O}_X)$ 中的态射 $\mathcal{O}_X \to E$。由命题
\ref{perfect-proposition-represent-cohomology-class-on-open}，对某个 $e$，复合
$I_e \to \mathcal{O}_X \to E$ 为零。由可区别三角
$I_e \to \mathcal{O}_X \to K_e \to I_e[1]$，得到态射 $K_e \to E$，使
$s$ 是复合 $\mathcal{O}_X \to K_e \to E$。
\end{proof}


\section{伪凝聚复形与完美复形}
\label{perfect-section-spell-out}

\noindent
本节把 Cohomology, Sections \ref{cohomology-section-strictly-perfect},
\ref{cohomology-section-pseudo-coherent},
\ref{cohomology-section-tor}, and
\ref{cohomology-section-perfect} 中定义的一般概念，与 More on Algebra,
Sections
\ref{more-algebra-section-pseudo-coherent},
\ref{more-algebra-section-tor}, and
\ref{more-algebra-section-perfect} 中模复形的相应概念联系起来。

\begin{lemma}
\label{perfect-lemma-pseudo-coherent}
设 $X$ 为概形。若 $E$ 是 $D(\mathcal{O}_X)$ 的 $m$-伪凝聚对象，则对
$i > m$，$H^i(E)$ 是拟凝聚 $\mathcal{O}_X$-模，而 $H^m(E)$ 是某个拟凝聚
$\mathcal{O}_X$-模的商。若 $E$ 伪凝聚，则 $E$ 是
$D_\QCoh(\mathcal{O}_X)$ 的对象。
\end{lemma}

\begin{proof}
在 $X$ 上局部地，存在严格完美复形 $\mathcal{E}^\bullet$，使得对
$i > m$，$H^i(E)$ 同构于 $H^i(\mathcal{E}^\bullet)$，且 $H^m(E)$ 是
$H^m(\mathcal{E}^\bullet)$ 的商。层 $\mathcal{E}^i$ 是有限自由模的直和项，
故拟凝聚。结论得证。
\end{proof}

\begin{lemma}
\label{perfect-lemma-pseudo-coherent-affine}
设 $X = \Spec(A)$ 为仿射概形，$M^\bullet$ 为 $A$-模复形，$E$ 为
$D(\mathcal{O}_X)$ 中相应的对象。那么，$E$ 作为 $D(\mathcal{O}_X)$ 的对象
是 $m$-伪凝聚的（相应地，伪凝聚的），当且仅当 $M^\bullet$ 作为
$A$-模复形是 $m$-伪凝聚的（相应地，伪凝聚的）。
\end{lemma}

\begin{proof}
由定义立即可知，若 $M^\bullet$ 是 $m$-伪凝聚的，则 $E$ 也是。为证明反向，
假设 $E$ 是 $m$-伪凝聚的。由于 $X = \Spec(A)$ 拟紧，且其拓扑有由标准开集
组成的基，可以找到标准开覆盖 $X = D(f_1) \cup \ldots \cup D(f_n)$、
$D(f_i)$ 上的严格完美复形 $\mathcal{E}_i^\bullet$，以及映射
$\alpha_i : \mathcal{E}_i^\bullet \to E|_{U_i}$；这些映射在 $j > m$ 时于
$H^j$ 上诱导同构，并在 $H^m$ 上诱导满射。由 Cohomology, Lemma
\ref{cohomology-lemma-local-actual}，细分开覆盖后，可假设每个 $i$ 对应的
$\alpha_i$ 都由复形映射
$\mathcal{E}_i^\bullet \to \widetilde{M^\bullet}|_{U_i}$
给出。由 Modules, Lemma
\ref{modules-lemma-direct-summand-of-locally-free-is-locally-free}，各项
$\mathcal{E}_i^n$ 都是有限局部自由模。因此再次细分开覆盖后，可假设每个
$\mathcal{E}_i^n$ 都是有限自由 $\mathcal{O}_{U_i}$-模。由定义可知，
$M^\bullet_{f_i}$ 是 $A_{f_i}$-模的 $m$-伪凝聚复形。应用 More on Algebra,
Lemma \ref{more-algebra-lemma-glue-pseudo-coherent} 即得结论。

\medskip\noindent
“伪凝聚”的情形来自如下事实：按定义，$E$ 伪凝聚当且仅当对所有 $m$，
$E$ 都是 $m$-伪凝聚的；而由 More on Algebra, Lemma
\ref{more-algebra-lemma-pseudo-coherent}，$M^\bullet$ 也有相同的性质。
\end{proof}

\begin{lemma}
\label{perfect-lemma-identify-pseudo-coherent-noetherian}
设 $X$ 为诺特概形，$E$ 为 $D_\QCoh(\mathcal{O}_X)$ 的对象。对
$m \in \mathbf{Z}$，以下条件等价：
\begin{enumerate}
\item 对 $i \geq m$，$H^i(E)$ 凝聚，并且对 $i \gg 0$，它为零；
\item $E$ 是 $m$-伪凝聚的。
\end{enumerate}
特别地，$E$ 伪凝聚当且仅当 $E$ 是
$D^-_{\textit{Coh}}(\mathcal{O}_X)$ 的对象。
\end{lemma}

\begin{proof}
由于 $X$ 拟紧，可知在 (1) 和 (2) 两种情形下，对象 $E$ 都有上界。因此问题
关于 $X$ 是局部的，可假设 $X$ 仿射。设 $X = \Spec(A)$，其中 $A$ 为某个
诺特环。此时由引理 \ref{perfect-lemma-affine-compare-bounded}，$E$ 对应于
$A$-模复形 $M^\bullet$。由引理 \ref{perfect-lemma-pseudo-coherent-affine}，$E$ 是
$m$-伪凝聚的，当且仅当 $M^\bullet$ 是 $m$-伪凝聚的。另一方面，$H^i(E)$
凝聚当且仅当 $H^i(M^\bullet)$ 是有限 $A$-模（Properties, Lemma
\ref{properties-lemma-finite-type-module}）。因此结论由 More on Algebra,
Lemma \ref{more-algebra-lemma-Noetherian-pseudo-coherent} 得到。
\end{proof}

\begin{lemma}
\label{perfect-lemma-tor-dimension-affine}
设 $X = \Spec(A)$ 为仿射概形，$M^\bullet$ 为 $A$-模复形，$E$ 为
$D(\mathcal{O}_X)$ 中相应的对象。则
\begin{enumerate}
\item $E$ 的 Tor 振幅在 $[a, b]$ 内，当且仅当 $M^\bullet$ 的 Tor 振幅在
$[a, b]$ 内。
\item $E$ 具有有限 Tor 维数，当且仅当 $M^\bullet$ 具有有限 Tor 维数。
\end{enumerate}
\end{lemma}

\begin{proof}
(2) 显然由 (1) 得到。在 (1) 的证明中，我们将不再特别说明地使用引理
\ref{perfect-lemma-affine-compare-bounded} 的等价 $D(A) = D_\QCoh(X)$。假设
$M^\bullet$ 的 Tor 振幅在 $[a, b]$ 内。则 $K^\bullet$ 在 $D(A)$ 中同构于
由平坦 $A$-模构成的复形 $K^\bullet$，且对 $i \not \in [a, b]$ 有
$K^i = 0$，见 More on Algebra, Lemma
\ref{more-algebra-lemma-tor-amplitude}。于是 $E$ 同构于
$\widetilde{K^\bullet}$。由于每个 $\widetilde{K^i}$ 都是平坦
$\mathcal{O}_X$-模，由 Cohomology, Lemma
\ref{cohomology-lemma-tor-amplitude} 可知 $E$ 的 Tor 振幅在 $[a, b]$ 内。

\medskip\noindent
假设 $E$ 的 Tor 振幅在 $[a, b]$ 内。则 $E$ 有界，从而 $M^\bullet$ 属于
$K^-(A)$。因此可以把 $M^\bullet$ 替换为有上界的 $A$-模复形。甚至可以
选取投射消解，并假设 $M^\bullet$ 是由自由 $A$-模构成的有上界复形。于是
对任意 $A$-模 $N$，有
$$
E \otimes_{\mathcal{O}_X}^\mathbf{L} \widetilde{N}
\cong
\widetilde{M^\bullet} \otimes_{\mathcal{O}_X}^\mathbf{L} \widetilde{N}
\cong
\widetilde{M^\bullet \otimes_A N}
$$
于 $D(\mathcal{O}_X)$ 中。因此左端上同调层的消失性蕴含 $M^\bullet$ 的
Tor 振幅在 $[a, b]$ 内。
\end{proof}

\begin{lemma}
\label{perfect-lemma-tor-dimension-rel-affine}
设 $f : X \to S$ 为对应于环同态 $R \to A$ 的仿射概形间态射。设
$M^\bullet$ 为 $A$-模复形，$E$ 为 $D(\mathcal{O}_X)$ 中相应的对象。则
\begin{enumerate}
\item $E$ 作为 $D(f^{-1}\mathcal{O}_S)$ 的对象，其 Tor 振幅在 $[a, b]$
内，当且仅当 $M^\bullet$ 作为 $D(R)$ 的对象，其 Tor 振幅在 $[a, b]$ 内。
\item $E$ 作为 $D(f^{-1}\mathcal{O}_S)$ 的对象局部具有有限 Tor 维数，
当且仅当 $M^\bullet$ 作为 $D(R)$ 的对象具有有限 Tor 维数。
\end{enumerate}
\end{lemma}

\begin{proof}
考虑位于 $\mathfrak p \subset R$ 之上的素理想 $\mathfrak q \subset A$。
设 $x \in X$ 和 $s = f(x) \in S$ 为相应的点。于是
$(f^{-1}\mathcal{O}_S)_x = \mathcal{O}_{S, s} = R_\mathfrak p$ 且
$E_x = M^\bullet_\mathfrak q$。由此可如下看出等价性。

\medskip\noindent
若 $M^\bullet$ 作为 $R$-模复形的 Tor 振幅在 $[a, b]$ 内，则
$M^\bullet$ 在 $A$ 的任意素理想处局部化后仍有此性质。由 Cohomology,
Lemma \ref{cohomology-lemma-tor-amplitude-stalk} 可知，$E$ 作为
$f^{-1}\mathcal{O}_S$-模层复形的 Tor 振幅在 $[a, b]$ 内。反之，假设 $E$
作为 $D(f^{-1}\mathcal{O}_S)$ 的对象，其 Tor 振幅在 $[a, b]$ 内。再次使用
刚才引用的引理，可知对每个位于 $\mathfrak p \subset R$ 之上的素理想
$\mathfrak q \subset A$，$M^\bullet_\mathfrak q$ 作为
$R_\mathfrak p$-模复形的 Tor 振幅在 $[a, b]$ 内。由 More on Algebra,
Lemma \ref{more-algebra-lemma-tor-amplitude-localization}，$M^\bullet$ 作为
$R$-模复形的 Tor 振幅在 $[a, b]$ 内。这就完成了 (1) 的证明。

\medskip\noindent
由于 $X$ 拟紧，若 $E$ 作为 $f^{-1}\mathcal{O}_S$-模复形局部具有有限 Tor
维数，则事实上存在某些 $a, b$，使 $E$ 作为
$f^{-1}\mathcal{O}_S$-模复形的 Tor 振幅在 $[a, b]$ 内。因此 (2) 由 (1)
得到。
\end{proof}

\begin{lemma}
\label{perfect-lemma-tor-qc-qs}
设 $X$ 为拟分离概形，$E$ 为 $D_\QCoh(\mathcal{O}_X)$ 的对象，并设
$a \leq b$。以下条件等价：
\begin{enumerate}
\item $E$ 的 Tor 振幅在 $[a, b]$ 内；
\item 对 $\QCoh(\mathcal{O}_X)$ 中所有 $\mathcal{F}$ 及
$i \not \in [a, b]$，都有
$H^i(E \otimes_{\mathcal{O}_X}^\mathbf{L} \mathcal{F}) = 0$。
\end{enumerate}
\end{lemma}

\begin{proof}
显然 (1) 蕴含 (2)。假设 (2)。设 $U \subset X$ 为仿射开集。由于 $X$
拟分离，态射 $j : U \to X$ 拟紧且分离，因此 $j_*$ 把拟凝聚模变为拟凝聚
模（Schemes, Lemma \ref{schemes-lemma-push-forward-quasi-coherent}）。故函子
$\QCoh(\mathcal{O}_X) \to \QCoh(\mathcal{O}_U)$ 本质满。由此，条件 (2)
蕴含对所有拟凝聚 $\mathcal{O}_U$-模 $\mathcal{G}$ 及
$i \not \in [a, b]$，
$H^i(E|_U \otimes_{\mathcal{O}_U}^\mathbf{L} \mathcal{G})$
都消失。写成 $U = \Spec(A)$，并设 $M^\bullet$ 为由引理
\ref{perfect-lemma-affine-compare-bounded} 对应于 $E|_U$ 的 $A$-模复形。我们刚刚证明
$M^\bullet \otimes_A^\mathbf{L} N$ 在区间 $[a, b]$ 外的上同调群消失；换言之，
$M^\bullet$ 的 Tor 振幅在 $[a, b]$ 内。由引理
\ref{perfect-lemma-tor-dimension-affine}，$E|_U$ 的 Tor 振幅在 $[a, b]$ 内。引理得证。
\end{proof}

\begin{lemma}
\label{perfect-lemma-perfect-affine}
设 $X = \Spec(A)$ 为仿射概形，$M^\bullet$ 为 $A$-模复形，$E$ 为
$D(\mathcal{O}_X)$ 中相应的对象。则 $E$ 是 $D(\mathcal{O}_X)$ 的完美对象，
当且仅当 $M^\bullet$ 是 $D(A)$ 的完美对象。
\end{lemma}

\begin{proof}
这是引理 \ref{perfect-lemma-pseudo-coherent-affine} 和
\ref{perfect-lemma-tor-dimension-affine}、Cohomology, Lemma
\ref{cohomology-lemma-perfect} 以及 More on Algebra, Lemma
\ref{more-algebra-lemma-perfect} 的逻辑推论。
\end{proof}

\noindent
作为概形上伪凝聚复形之描述的推论，我们可以证明某些内 Hom 是拟凝聚的。

\begin{lemma}
\label{perfect-lemma-quasi-coherence-internal-hom}
设 $X$ 为概形。
\begin{enumerate}
\item 若 $L$ 属于 $D^+_\QCoh(\mathcal{O}_X)$，且 $D(\mathcal{O}_X)$ 中的
$K$ 伪凝聚，则 $R\SheafHom(K, L)$ 属于 $D_\QCoh(\mathcal{O}_X)$，并且
局部有下界。
\item 若 $L$ 属于 $D_\QCoh(\mathcal{O}_X)$，且 $D(\mathcal{O}_X)$ 中的
$K$ 完美，则 $R\SheafHom(K, L)$ 属于 $D_\QCoh(\mathcal{O}_X)$。
\item 若 $X = \Spec(A)$ 仿射且 $K, L \in D(A)$，则
$$
R\SheafHom(\widetilde{K}, \widetilde{L}) = \widetilde{R\Hom_A(K, L)}
$$
在以下两种情形成立：
\begin{enumerate}
\item $K$ 伪凝聚且 $L$ 有下界；
\item $K$ 完美而 $L$ 任意。
\end{enumerate}
\item 若 $X = \Spec(A)$ 且 $K, L$ 属于 $D(A)$，则
$R\SheafHom(\widetilde{K}, \widetilde{L})$ 的第 $n$ 个上同调层是与预层
$$
X \supset D(f) \longmapsto \Ext^n_{A_f}(K \otimes_A A_f, L \otimes_A A_f)
$$
相伴的层，其中 $f \in A$。
\end{enumerate}
\end{lemma}

\begin{proof}
在 $\mathcal{O}_X$ 的导出范畴中构造内 Hom 与局部化交换（见 Cohomology,
Section \ref{cohomology-section-internal-hom}）。因此为证明 (1) 和 (2)，可以
用仿射开集替换 $X$。由引理 \ref{perfect-lemma-affine-compare-bounded}、
\ref{perfect-lemma-pseudo-coherent-affine} 和
\ref{perfect-lemma-perfect-affine}
，为证明 (1) 和 (2)，只需证明 (3)。

\medskip\noindent
把 $K$ 表示为由有限投射 $A$-模构成的有上界（相应地，有限）复形，并把
$L$ 表示为 $A$-模的有下界（相应地，任意）复形，则 (3) 由 Cohomology,
Lemma
\ref{cohomology-lemma-Rhom-complex-of-direct-summands-finite-free}
中对内 Hom 的计算得到。

\medskip\noindent
为证明 (4)，回忆在任意带环空间上，$R\SheafHom(A, B)$ 的第 $n$ 个上同调层
是与预层
$$
U \mapsto \Hom_{D(U)}(A|_U, B|_U[n]) =
\Ext^n_{D(\mathcal{O}_U)}(A|_U, B|_U)
$$
相伴的层，见 Cohomology, Section \ref{cohomology-section-internal-hom}。另一
方面，$\widetilde{K}$ 在主开集 $D(f)$ 上的限制是 $K \otimes_A A_f$ 的像，
对 $L$ 亦然。因此 (4) 由引理 \ref{perfect-lemma-affine-compare-bounded} 的范畴等价
得到。
\end{proof}

\begin{lemma}
\label{perfect-lemma-internal-hom-evaluate-tensor-isomorphism}
设 $X$ 为概形，$K, L, M$ 为 $D_\QCoh(\mathcal{O}_X)$ 的对象。Cohomology,
Lemma \ref{cohomology-lemma-internal-hom-diagonal-better} 中的映射
$$
K \otimes_{\mathcal{O}_X}^\mathbf{L} R\SheafHom(M, L)
\longrightarrow
R\SheafHom(M, K \otimes_{\mathcal{O}_X}^\mathbf{L} L)
$$
在以下情形是同构：
\begin{enumerate}
\item $M$ 完美；
\item $K$ 完美；
\item $M$ 伪凝聚，$L \in D^+(\mathcal{O}_X)$，且 $K$ 具有有限 Tor 维数。
\end{enumerate}
\end{lemma}

\begin{proof}
引理 \ref{perfect-lemma-quasi-coherence-internal-hom} 把情形 (1) 和 (3) 化为仿射
情形，后者已在 More on Algebra, Lemma
\ref{more-algebra-lemma-internal-hom-evaluate-tensor-isomorphism} 中处理。（把
问题转写成代数语言时，还须使用引理 \ref{perfect-lemma-pseudo-coherent-affine}、
\ref{perfect-lemma-perfect-affine} 和 \ref{perfect-lemma-tor-dimension-affine}。）若 $K$ 完美而
不作其他假设，我们并不知道箭头的任何一端属于
$D_\QCoh(\mathcal{O}_X)$；但结论仍成立，因为可以局部工作，把问题化到
$K$ 是有限自由模的有限复形这一显然情形。
\end{proof}





\section{凝聚模的导出范畴}
\label{perfect-section-derived-coherent}

\noindent
设 $X$ 为局部诺特概形。此时，凝聚 $\mathcal{O}_X$-模的范畴
$\textit{Coh}(\mathcal{O}_X) \subset \textit{Mod}(\mathcal{O}_X)$ 是弱 Serre
子范畴，见 Homology, Section \ref{homology-section-serre-subcategories} 及
Cohomology of Schemes, Lemma
\ref{coherent-lemma-coherent-abelian-Noetherian}。以
$$
D_{\textit{Coh}}(\mathcal{O}_X) \subset D(\mathcal{O}_X)
$$
表示上同调层凝聚的复形所成的子范畴，见 Derived Categories, Section
\ref{derived-section-triangulated-sub}。于是得到典范函子
\begin{equation}
\label{perfect-equation-compare-coherent}
D(\textit{Coh}(\mathcal{O}_X))
\longrightarrow
D_{\textit{Coh}}(\mathcal{O}_X)
\end{equation}
见 Derived Categories, Equation (\ref{derived-equation-compare})。

\begin{lemma}
\label{perfect-lemma-coh-to-qcoh}
设 $X$ 为诺特概形。则函子
$$
D^-(\textit{Coh}(\mathcal{O}_X))
\longrightarrow
D^-_{\textit{Coh}(\mathcal{O}_X)}(\QCoh(\mathcal{O}_X))
$$
是等价。
\end{lemma}

\begin{proof}
注意，$\textit{Coh}(\mathcal{O}_X) \subset \QCoh(\mathcal{O}_X)$ 是 Serre
子范畴，见 Homology, Definition \ref{homology-definition-serre-subcategory}、
Lemma \ref{homology-lemma-characterize-serre-subcategory} 及 Cohomology of
Schemes, Lemmas \ref{coherent-lemma-coherent-abelian-Noetherian} 和
\ref{coherent-lemma-coherent-Noetherian-quasi-coherent-sub-quotient}。另一方面，
若 $\mathcal{G} \to \mathcal{F}$ 是从拟凝聚 $\mathcal{O}_X$-模到凝聚
$\mathcal{O}_X$-模的满射，则存在凝聚子模
$\mathcal{G}' \subset \mathcal{G}$ 满射到 $\mathcal{F}$。事实上，由
Properties, Lemma \ref{properties-lemma-quasi-coherent-colimit-finite-type}，
可以把 $\mathcal{G}$ 写成其凝聚子模的滤过并，其中必有一个满足要求。因此
引理由 Derived Categories, Lemma
\ref{derived-lemma-fully-faithful-embedding} 得到。
\end{proof}

\begin{proposition}
\label{perfect-proposition-DCoh}
设 $X$ 为诺特概形。则函子
$$
D^-(\textit{Coh}(\mathcal{O}_X))
\longrightarrow
D^-_{\textit{Coh}}(\mathcal{O}_X)
\quad\text{且}\quad
D^b(\textit{Coh}(\mathcal{O}_X))
\longrightarrow
D^b_{\textit{Coh}}(\mathcal{O}_X)
$$
都是等价。
\end{proposition}

\begin{proof}
考虑交换图表
$$
\xymatrix{
D^-(\textit{Coh}(\mathcal{O}_X)) \ar[r] \ar[d] &
D^-_{\textit{Coh}}(\mathcal{O}_X) \ar[d] \\
D^-(\QCoh(\mathcal{O}_X)) \ar[r] &
D^-_\QCoh(\mathcal{O}_X)
}
$$
由引理 \ref{perfect-lemma-coh-to-qcoh}，左侧竖直箭头完全忠实。由命题
\ref{perfect-proposition-Noetherian}，下方箭头是等价。按构造，右侧竖直箭头完全
忠实。因此上方水平箭头完全忠实。若 $K$ 是
$D^-_{\textit{Coh}}(\mathcal{O}_X)$ 的对象，则由命题
\ref{perfect-proposition-Noetherian} 与之对应的 $D^-(\QCoh(\mathcal{O}_X))$ 的对象
$K'$ 具有凝聚上同调层。因而由引理 \ref{perfect-lemma-coh-to-qcoh}，$K'$ 属于左侧
竖直箭头的本质像，从而上方水平箭头本质满。这完成了有上界情形的证明；
有界情形立即由有上界情形得到。
\end{proof}

\begin{lemma}
\label{perfect-lemma-direct-image-coherent}
设 $S$ 为诺特概形，$f : X \to S$ 为局部有限型概形态射。设 $E$ 为
$D^b_{\textit{Coh}}(\mathcal{O}_X)$ 的对象，使得对所有 $i$，$H^i(E)$ 的
支集在 $S$ 上固有。则 $Rf_*E$ 是
$D^b_{\textit{Coh}}(\mathcal{O}_S)$ 的对象。
\end{lemma}

\begin{proof}
考虑谱序列
$$
R^pf_*H^q(E) \Rightarrow R^{p + q}f_*E
$$
见 Derived Categories, Lemma \ref{derived-lemma-two-ss-complex-functor}。由假设
及 Cohomology of Schemes, Lemma
\ref{coherent-lemma-support-proper-over-base-pushforward}
，层 $R^pf_*H^q(E)$ 凝聚。因此 $R^{p + q}f_*E$ 凝聚，即
$Rf_*E \in D_{\textit{Coh}}(\mathcal{O}_S)$。有下界性显然；有上界性由
Cohomology of Schemes, Lemma
\ref{coherent-lemma-quasi-coherence-higher-direct-images}
或引理 \ref{perfect-lemma-quasi-coherence-direct-image} 得到。
\end{proof}

\begin{lemma}
\label{perfect-lemma-direct-image-coherent-bdd-below}
设 $S$ 为诺特概形，$f : X \to S$ 为局部有限型概形态射。设 $E$ 为
$D^+_{\textit{Coh}}(\mathcal{O}_X)$ 的对象，使得对所有 $i$，$H^i(E)$ 的
支集在 $S$ 上固有。则 $Rf_*E$ 是
$D^+_{\textit{Coh}}(\mathcal{O}_S)$ 的对象。
\end{lemma}

\begin{proof}
证明与引理 \ref{perfect-lemma-direct-image-coherent} 的证明相同。也可以考察正合函子
$Rf_*$ 对可区别三角
$\tau_{\leq a}E \to E \to \tau_{\geq a + 1}E \to \tau_{\leq a}E[1]$.
的作用，由引理 \ref{perfect-lemma-direct-image-coherent} 推出结论。
\end{proof}

\begin{lemma}
\label{perfect-lemma-coherent-internal-hom}
设 $X$ 为局部诺特概形。若 $L$ 属于
$D^+_{\textit{Coh}}(\mathcal{O}_X)$ 且 $K$ 属于
$D^-_{\textit{Coh}}(\mathcal{O}_X)$，则 $R\SheafHom(K, L)$ 属于
$D^+_{\textit{Coh}}(\mathcal{O}_X)$。
\end{lemma}

\begin{proof}
只需在 $X$ 是诺特环 $A$ 的谱时证明。由引理
\ref{perfect-lemma-identify-pseudo-coherent-noetherian}，$K$ 伪凝聚。再由引理
\ref{perfect-lemma-quasi-coherence-internal-hom}，问题转化为如下代数问题：若
$L \in D^+_{\textit{Coh}}(A)$ 且 $K$ 属于
$D^-_{\textit{Coh}}(A)$，则 $R\Hom_A(K, L)$ 属于
$D^+_{\textit{Coh}}(A)$。由于 $L$ 有下界而 $K$ 有上界，存在收敛谱序列
$$
\Ext^p_A(K, H^q(L)) \Rightarrow \text{Ext}^{p + q}_A(K, L)
$$
以及收敛谱序列
$$
\Ext^i_A(H^{-j}(K), H^q(L)) \Rightarrow \text{Ext}^{i + j}_A(K, H^q(L))
$$
见 Injectives, Remarks \ref{injectives-remark-spectral-sequences-ext} 和
\ref{injectives-remark-spectral-sequences-ext-variant}。由 Algebra, Lemma
\ref{algebra-lemma-ext-noetherian}，对有限 $A$-模 $M$、$N$，模
$\Ext^p_A(M, N)$ 有限，故证明完成。
\end{proof}

\begin{lemma}
\label{perfect-lemma-perfect-on-noetherian}
设 $X$ 为诺特概形，$D(\mathcal{O}_X)$ 中的 $E$ 完美。则
\begin{enumerate}
\item $E$ 属于 $D^b_{\textit{Coh}}(\mathcal{O}_X)$；
\item 若 $L$ 属于 $D_{\textit{Coh}}(\mathcal{O}_X)$，则
$E \otimes_{\mathcal{O}_X}^\mathbf{L} L$ 和
$R\SheafHom_{\mathcal{O}_X}(E, L)$ 都属于
$D_{\textit{Coh}}(\mathcal{O}_X)$；
\item 若 $L$ 属于 $D^b_{\textit{Coh}}(\mathcal{O}_X)$，则
$E \otimes_{\mathcal{O}_X}^\mathbf{L} L$ 和
$R\SheafHom_{\mathcal{O}_X}(E, L)$ 都属于
$D^b_{\textit{Coh}}(\mathcal{O}_X)$；
\item 若 $L$ 属于 $D^+_{\textit{Coh}}(\mathcal{O}_X)$，则
$E \otimes_{\mathcal{O}_X}^\mathbf{L} L$ 和
$R\SheafHom_{\mathcal{O}_X}(E, L)$ 都属于
$D^+_{\textit{Coh}}(\mathcal{O}_X)$；
\item 若 $L$ 属于 $D^-_{\textit{Coh}}(\mathcal{O}_X)$，则
$E \otimes_{\mathcal{O}_X}^\mathbf{L} L$ 和
$R\SheafHom_{\mathcal{O}_X}(E, L)$ 都属于
$D^-_{\textit{Coh}}(\mathcal{O}_X)$.
\end{enumerate}
\end{lemma}

\begin{proof}
由于 $X$ 拟紧，每个断言都可在 $X$ 任意开覆盖的各成员上检验。因此可假设
$E$ 由有限自由模的有界复形 $\mathcal{E}^\bullet$ 表示，见 Cohomology,
Lemma \ref{cohomology-lemma-perfect-on-locally-ringed}。此时各断言都很清楚，
因为 $R\SheafHom_{\mathcal{O}_X}(E, L)$ 和
$E \otimes_{\mathcal{O}_X}^\mathbf{L} L$ 都可用 $\mathcal{E}^\bullet$ 在复形
层面计算，见 Cohomology, Lemmas
\ref{cohomology-lemma-Rhom-strictly-perfect} 和
\ref{cohomology-lemma-bounded-flat-K-flat}。略去一些细节。
\end{proof}

\begin{lemma}
\label{perfect-lemma-ext-finite}
设 $A$ 为诺特环，$X$ 为 $A$ 上的固有概形。对属于
$D^+_{\textit{Coh}}(\mathcal{O}_X)$ 的 $L$ 和属于
$D^-_{\textit{Coh}}(\mathcal{O}_X)$ 的 $K$，$A$-模
$\Ext_{\mathcal{O}_X}^n(K, L)$ 是有限的。
\end{lemma}

\begin{proof}
回忆
$$
\Ext_{\mathcal{O}_X}^n(K, L) =
H^n(X, R\SheafHom_{\mathcal{O}_X}(K, L)) =
H^n(\Spec(A), Rf_*R\SheafHom_{\mathcal{O}_X}(K, L))
$$
见 Cohomology, Lemma \ref{cohomology-lemma-section-RHom-over-U} 和 Cohomology,
Section \ref{cohomology-section-Leray}。因此结论由引理
\ref{perfect-lemma-coherent-internal-hom} 和
\ref{perfect-lemma-direct-image-coherent-bdd-below} 得到。
\end{proof}

\begin{lemma}
\label{perfect-lemma-perfect-on-regular}
设 $X$ 为正则概形。则 $D^b_{\textit{Coh}}(\mathcal{O}_X)$ 的每个对象都
完美。若 $X$ 拟紧，即诺特且正则，则反过来，$D(\mathcal{O}_X)$ 的每个完美
对象都属于 $D^b_{\textit{Coh}}(\mathcal{O}_X)$。
\end{lemma}

\begin{proof}
设 $K$ 为 $D^b_{\textit{Coh}}(\mathcal{O}_X)$ 的对象。为检验 $K$ 完美，可以
在 $X$ 上仿射局部地工作（见 Cohomology, Section
\ref{cohomology-section-perfect}）。于是由引理 \ref{perfect-lemma-perfect-affine} 和
More on Algebra, Lemma \ref{more-algebra-lemma-regular-perfect}，$K$ 完美。
反向即引理 \ref{perfect-lemma-perfect-on-noetherian}。
\end{proof}




\section{复形有限性性质的下降}
\label{perfect-section-descent-finiteness}

\noindent
本节是 Descent, Section \ref{descent-section-descent-finiteness} 对概形导出
范畴对象的类似版本。其中最容易的结果或许如下。

\begin{lemma}
\label{perfect-lemma-tor-amplitude-descends}
设 $f : X \to Y$ 为概形（或更一般地，局部带环空间）的满射平坦态射。设
$E \in D(\mathcal{O}_Y)$ 且 $a, b \in \mathbf{Z}$。则 $E$ 的 Tor 振幅在
$[a, b]$ 内，当且仅当 $Lf^*E$ 的 Tor 振幅在 $[a, b]$ 内。
\end{lemma}

\begin{proof}
拉回总保持 Tor 振幅，见 Cohomology, Lemma
\ref{cohomology-lemma-tor-amplitude-pullback}。Tor 振幅在 $[a, b]$ 内这一性质
可在茎上检验，见 Cohomology, Lemma
\ref{cohomology-lemma-tor-amplitude-stalk}。由 Algebra, Lemma
\ref{algebra-lemma-local-flat-ff}，平坦局部环同态忠实平坦。因此结论由 More
on Algebra, Lemma \ref{more-algebra-lemma-flat-descent-tor-amplitude} 得到。
\end{proof}

\begin{lemma}
\label{perfect-lemma-pseudo-coherent-descends-fpqc}
设 $\{f_i : X_i \to X\}$ 为概形的 fpqc 覆盖，
$E \in D_\QCoh(\mathcal{O}_X)$，且 $m \in \mathbf{Z}$。则 $E$ 是
$m$-伪凝聚的，当且仅当每个 $Lf_i^*E$ 都是 $m$-伪凝聚的。
\end{lemma}

\begin{proof}
拉回总保持 $m$-伪凝聚性，见 Cohomology, Lemma
\ref{cohomology-lemma-pseudo-coherent-pullback}。反之，假设对所有 $i$，
$Lf_i^*E$ 都是 $m$-伪凝聚的。设 $U \subset X$ 为仿射开集。只需证明
$E|_U$ 是 $m$-伪凝聚的。由于 $\{f_i : X_i \to X\}$ 是 fpqc 覆盖，可以
找到有限多个仿射开集 $V_j \subset X_{a(j)}$，使得
$f_{a(j)}(V_j) \subset U$ 且 $U = \bigcup f_{a(j)}(V_j)$。令
$V = \coprod V_i$。因此可以用 $U$ 替换 $X$，用 $\{V \to U\}$ 替换
$\{f_i : X_i \to X\}$，并假设 $X$ 仿射且覆盖由仿射概形间的单个满射平坦
态射 $\{f : Y \to X\}$ 给出。此时结论由 More on Algebra, Lemma
\ref{more-algebra-lemma-flat-descent-pseudo-coherent} 以及引理
\ref{perfect-lemma-affine-compare-bounded} 和 \ref{perfect-lemma-pseudo-coherent-affine} 得到。
\end{proof}

\begin{lemma}
\label{perfect-lemma-pseudo-coherent-descends-fppf}
设 $\{f_i : X_i \to X\}$ 为概形的 fppf 覆盖，$E \in D(\mathcal{O}_X)$，
且 $m \in \mathbf{Z}$。则 $E$ 是 $m$-伪凝聚的，当且仅当每个 $Lf_i^*E$
都是 $m$-伪凝聚的。
\end{lemma}

\begin{proof}
拉回总保持 $m$-伪凝聚性，见 Cohomology, Lemma
\ref{cohomology-lemma-pseudo-coherent-pullback}。反之，假设对所有 $i$，
$Lf_i^*E$ 都是 $m$-伪凝聚的。设 $U \subset X$ 为仿射开集。只需证明
$E|_U$ 是 $m$-伪凝聚的。由于 $\{f_i : X_i \to X\}$ 是 fppf 覆盖，可以
找到有限多个仿射开集 $V_j \subset X_{a(j)}$，使得
$f_{a(j)}(V_j) \subset U$ 且 $U = \bigcup f_{a(j)}(V_j)$。令
$V = \coprod V_i$。因此可以用 $U$ 替换 $X$，用 $\{V \to U\}$ 替换
$\{f_i : X_i \to X\}$，并假设 $X$ 仿射且覆盖由单个有限表示满射平坦态射
$\{f : Y \to X\}$ 给出。

\medskip\noindent
由于 $f$ 平坦，导出函子 $Lf^*$ 就由 $f^*$ 给出，且 $f^*$ 正合。因此
$H^i(Lf^*E) = f^*H^i(E)$。由于 $Lf^*E$ 是 $m$-伪凝聚的，有
$Lf^*E \in D^-(\mathcal{O}_Y)$。由于 $f$ 满射且平坦，有
$E \in D^-(\mathcal{O}_X)$。设 $i \in \mathbf{Z}$ 为使 $H^i(E)$ 非零的最大
整数。若 $i < m$，则已经完成。否则，由 Cohomology, Lemma
\ref{cohomology-lemma-finite-cohomology}，$f^*H^i(E)$ 是有限型
$\mathcal{O}_Y$-模。再由 Descent, Lemma
\ref{descent-lemma-finite-type-descends-fppf}，$\mathcal{O}_X$-模 $H^i(E)$
是有限型的。因此用一个有限仿射开覆盖的各成员替换 $X$ 后，可假设存在映射
$$
\alpha : \mathcal{O}_X^{\oplus n}[-i] \longrightarrow E
$$
使 $H^i(\alpha)$ 为满射。设 $C$ 为 $D(\mathcal{O}_X)$ 中 $\alpha$ 的锥。
拉回到 $Y$ 并使用 Cohomology, Lemma
\ref{cohomology-lemma-cone-pseudo-coherent}，可知 $Lf^*C$ 是 $m$-伪凝聚的。
此外，对 $j \geq i$ 有 $H^j(C) = 0$。因此对 $i$ 归纳可知 $C$ 是
$m$-伪凝聚的。再次使用 Cohomology, Lemma
\ref{cohomology-lemma-cone-pseudo-coherent} 即得结论。
\end{proof}

\begin{lemma}
\label{perfect-lemma-perfect-descends-fpqc}
设 $\{f_i : X_i \to X\}$ 为概形的 fpqc 覆盖，$E \in D(\mathcal{O}_X)$。
则 $E$ 完美，当且仅当每个 $Lf_i^*E$ 都完美。
\end{lemma}

\begin{proof}
拉回总保持完美复形，见 Cohomology, Lemma
\ref{cohomology-lemma-perfect-pullback}。反之，假设对所有 $i$，$Lf_i^*E$
都完美。于是由引理 \ref{perfect-lemma-pseudo-coherent} 和 Cohomology, Lemma
\ref{cohomology-lemma-perfect}，每个 $Lf_i^*E$ 的上同调层都拟凝聚。由于态射
$f_i$ 平坦，有 $H^p(Lf_i^*E) = f_i^*H^p(E)$。因此由 Descent,
Proposition \ref{descent-proposition-fpqc-descent-quasi-coherent}，$E$ 的上同调
层拟凝聚。至此，引理由 Cohomology, Lemma
\ref{cohomology-lemma-perfect} 以及引理 \ref{perfect-lemma-tor-amplitude-descends} 和
\ref{perfect-lemma-pseudo-coherent-descends-fpqc} 形式地得到。
\end{proof}

\begin{lemma}
\label{perfect-lemma-closed-push-pseudo-coherent}
设 $i : Z \to X$ 为带环空间态射，使得 $i$ 在底拓扑空间上是闭浸入，且
$i_*\mathcal{O}_Z$ 作为 $\mathcal{O}_X$-模伪凝聚。设
$E \in D(\mathcal{O}_Z)$。则 $E$ 是 $m$-伪凝聚的，当且仅当 $Ri_*E$ 是
$m$-伪凝聚的。
\end{lemma}

\begin{proof}
整个证明中都将使用 $i_*$ 是正合函子、因而 $Ri_* = i_*$ 这一事实，见
Modules, Lemma \ref{modules-lemma-i-star-exact}。

\medskip\noindent
假设 $E$ 是 $m$-伪凝聚的。设 $x \in X$。我们要找 $x$ 的一个邻域，使
$i_*E$ 在其上是 $m$-伪凝聚的。若 $x \not \in Z$，则显然。故可假设
$x \in Z$。我们将使用如下事实：当 $x \in U \subset X$ 且该子集开时，
$U \cap Z$ 构成 $x$ 在 $Z$ 中的一个基本邻域系。缩小 $X$ 后，可假设 $E$
有上界。我们对使 $H^p(E)$ 非零的最大整数 $p$ 作归纳。若 $p < m$，则
无须证明。若 $p \geq m$，则由 Cohomology, Lemma
\ref{cohomology-lemma-finite-cohomology}，$H^p(E)$ 是有限型
$\mathcal{O}_Z$-模。因此缩小 $X$ 后，可选择映射
$\mathcal{O}_Z^{\oplus n}[-p] \to E$，它诱导满射
$\mathcal{O}_Z^{\oplus n} \to H^p(E)$。选择可区别三角
$$
\mathcal{O}_Z^{\oplus n}[-p] \to E \to C \to \mathcal{O}_Z^{\oplus n}[-p + 1]
$$
由 Cohomology, Lemma \ref{cohomology-lemma-cone-pseudo-coherent} 可知，对
$j \geq p$ 有 $H^j(C) = 0$，且 $C$ 是 $m$-伪凝聚的。由归纳可知 $i_*C$
在 $X$ 上是 $m$-伪凝聚的。由于 $i_*\mathcal{O}_Z$ 在 $X$ 上也是
$m$-伪凝聚的，由可区别三角
$$
i_*\mathcal{O}_Z^{\oplus n}[-p] \to i_*E \to i_*C \to
i_*\mathcal{O}_Z^{\oplus n}[-p + 1]
$$
及 Cohomology, Lemma \ref{cohomology-lemma-cone-pseudo-coherent} 可知，
$i_*E$ 是 $m$-伪凝聚的。

\medskip\noindent
反之，假设 $i_*E$ 是 $m$-伪凝聚的。设 $z \in Z$。我们要找 $z$ 的一个
邻域，使 $E$ 在其上是 $m$-伪凝聚的。我们将使用如下事实：当
$z \in U \subset X$ 且该子集开时，$U \cap Z$ 构成 $z$ 在 $Z$ 中的一个基本
邻域系。缩小 $X$ 后，可假设 $i_*E$、从而 $E$ 有上界。我们对使
$H^p(E)$ 非零的最大整数 $p$ 作归纳。若 $p < m$，则无须证明。若
$p \geq m$，则由 Cohomology, Lemma
\ref{cohomology-lemma-finite-cohomology}，
$H^p(i_*E) = i_*H^p(E)$ 是有限型 $\mathcal{O}_X$-模。选择一个表示 $E$ 的
$\mathcal{O}_Z$-模复形 $\mathcal{E}^\bullet$。缩小 $X$ 后，可选择映射
$\alpha : \mathcal{O}_X^{\oplus n}[-p] \to i_*\mathcal{E}^\bullet$，它诱导
满射 $\mathcal{O}_X^{\oplus n} \to i_*H^p(\mathcal{E}^\bullet)$。由伴随得到
映射 $\alpha : \mathcal{O}_Z^{\oplus n}[-p] \to \mathcal{E}^\bullet$，它诱导
满射 $\mathcal{O}_Z^{\oplus n} \to H^p(\mathcal{E}^\bullet)$。选择可区别三角
$$
\mathcal{O}_Z^{\oplus n}[-p] \to E \to C \to \mathcal{O}_Z^{\oplus n}[-p + 1]
$$
可知对 $j \geq p$ 有 $H^j(C) = 0$。由可区别三角
$$
i_*\mathcal{O}_Z^{\oplus n}[-p] \to i_*E \to i_*C \to
i_*\mathcal{O}_Z^{\oplus n}[-p + 1]
$$
、$i_*\mathcal{O}_Z$ 伪凝聚这一事实及 Cohomology, Lemma
\ref{cohomology-lemma-cone-pseudo-coherent}，可知 $i_*C$ 是 $m$-伪凝聚的。
由归纳，$C$ 是 $m$-伪凝聚的。再次应用 Cohomology, Lemma
\ref{cohomology-lemma-cone-pseudo-coherent}，可知 $E$ 是 $m$-伪凝聚的。
\end{proof}

\begin{lemma}
\label{perfect-lemma-finite-push-pseudo-coherent}
设 $f : X \to Y$ 为概形的有限态射，使得 $f_*\mathcal{O}_X$ 作为
$\mathcal{O}_Y$-模伪凝聚\footnote{这意味着 $f$ 是伪凝聚态射，见 More on
Morphisms, Lemma
\ref{more-morphisms-lemma-finite-pseudo-coherent}.}.
设 $E \in D_\QCoh(\mathcal{O}_X)$。则 $E$ 是 $m$-伪凝聚的，当且仅当
$Rf_*E$ 是 $m$-伪凝聚的。
\end{lemma}

\begin{proof}
这是把 More on Algebra, Lemma
\ref{more-algebra-lemma-finite-push-pseudo-coherent} 转写为概形语言。转写时使用
引理 \ref{perfect-lemma-affine-compare-bounded} 和
\ref{perfect-lemma-pseudo-coherent-affine}。
\end{proof}


\section{复形的提升}
\label{perfect-section-lift}

\noindent
设 $U \subset X$ 为带环空间的开子空间，并以 $j : U \to X$ 表示包含态射。
由于 $Rj_*$ 是限制函子的右逆，函子
$D(\mathcal{O}_X) \to D(\mathcal{O}_U)$ 本质满。本节把这一结论推广到具有
拟凝聚上同调层的复形等情形。

\begin{lemma}
\label{perfect-lemma-lift-quasi-coherent}
设 $X$ 为概形，$j : U \to X$ 为拟紧开浸入。函子
$$
D_\QCoh(\mathcal{O}_X) \to D_\QCoh(\mathcal{O}_U)
\quad\text{且}\quad
D^+_\QCoh(\mathcal{O}_X) \to D^+_\QCoh(\mathcal{O}_U)
$$
本质满。若 $X$ 拟紧，则函子
$$
D^-_\QCoh(\mathcal{O}_X) \to D^-_\QCoh(\mathcal{O}_U)
\quad\text{且}\quad
D^b_\QCoh(\mathcal{O}_X) \to D^b_\QCoh(\mathcal{O}_U)
$$
也本质满。
\end{lemma}

\begin{proof}
由引理 \ref{perfect-lemma-quasi-coherence-direct-image}，$Rj_*$ 把
$D_\QCoh(\mathcal{O}_U)$ 映到 $D_\QCoh(\mathcal{O}_X)$，故引理前的论证适用于
第一种情形。显然，$Rj_*$ 把 $D^+_\QCoh(\mathcal{O}_U)$ 映到
$D^+_\QCoh(\mathcal{O}_X)$，这给出关于有下界复形的断言。最后，引理
\ref{perfect-lemma-quasi-coherence-direct-image} 保证当 $X$ 拟紧时，$Rj_*$ 把
$D^-_\QCoh(\mathcal{O}_U)$ 映到 $D^-_\QCoh(\mathcal{O}_X)$。结合二者即得
最后一个断言。
\end{proof}

\begin{lemma}
\label{perfect-lemma-lift-coherent}
设 $X$ 为诺特概形，$j : U \to X$ 为开浸入。函子
$D^b_{\textit{Coh}}(\mathcal{O}_X) \to D^b_{\textit{Coh}}(\mathcal{O}_U)$
本质满。
\end{lemma}

\begin{proof}
设 $K$ 为 $D^b_{\textit{Coh}}(\mathcal{O}_U)$ 的对象。由命题
\ref{perfect-proposition-DCoh}，可用凝聚 $\mathcal{O}_U$-模的有界复形
$\mathcal{F}^\bullet$ 表示 $K$。设对某些 $a \leq b$，当
$i \not \in [a, b]$ 时 $\mathcal{F}^i = 0$。由于 $j$ 拟紧且分离，有界复形
$j_*\mathcal{F}^\bullet$ 的各项是 $X$ 上的拟凝聚模，见 Schemes, Lemma
\ref{schemes-lemma-push-forward-quasi-coherent}。如下归纳选取凝聚子模
$\mathcal{G}^i \subset j_*\mathcal{F}^i$。当 $i = a$ 时，任选一个在 $U$
上的限制为 $\mathcal{F}^a$ 的凝聚子模
$\mathcal{G}^a \subset j_*\mathcal{F}^a$；这由 Properties, Lemma
\ref{properties-lemma-extend} 可行。当 $i > a$ 时，先任选一个在 $U$ 上的
限制为 $\mathcal{F}^i$ 的凝聚子模
$\mathcal{H}^i \subset j_*\mathcal{F}^i$，然后令
$\mathcal{G}^i = \Im(\mathcal{H}^i \oplus \mathcal{G}^{i - 1}
\to j_*\mathcal{F}^i)$。显然，
$\mathcal{G}^\bullet \subset j_*\mathcal{F}^\bullet$ 是凝聚
$\mathcal{O}_X$-模的有界复形，其在 $U$ 上的限制为
$\mathcal{F}^\bullet$，正如所需。
\end{proof}

\begin{lemma}
\label{perfect-lemma-lift-pseudo-coherent}
设 $X$ 为仿射概形，$U \subset X$ 为拟紧开子概形。对
$D(\mathcal{O}_U)$ 中任意伪凝聚对象 $E$，存在有限自由
$\mathcal{O}_X$-模的有上界复形，其在 $U$ 上的限制同构于 $E$。
\end{lemma}

\begin{proof}
由引理 \ref{perfect-lemma-pseudo-coherent}，$E$ 是 $D_\QCoh(\mathcal{O}_U)$ 的对象。
由引理 \ref{perfect-lemma-lift-quasi-coherent}，可假设对
$D_\QCoh(\mathcal{O}_X)$ 的某个对象 $E'$ 有 $E = E'|U$。写成
$X = \Spec(A)$。由引理 \ref{perfect-lemma-affine-compare-bounded}，可以找到
$A$-模复形 $M^\bullet$，其相伴的 $\mathcal{O}_X$-模复形是 $E'$ 的一个
代表。

\medskip\noindent
选择 $f_1, \ldots, f_r \in A$，使得
$U = D(f_1) \cup \ldots \cup D(f_r)$。由引理
\ref{perfect-lemma-pseudo-coherent-affine}，复形 $M^\bullet_{f_j}$ 是
$A_{f_j}$-模的伪凝聚复形。设 $n$ 为整数。假设已有复形映射
$\alpha : F^\bullet \to M^\bullet$，其中 $F^\bullet$ 有上界，对 $i < n$
有 $F^i = 0$，且每个 $F^i$ 都是有限自由 $R$-模，并使得
$$
H^i(\alpha_{f_j}) : H^i(F^\bullet_{f_j}) \to H^i(M^\bullet_{f_j})
$$
在 $i > n$ 时为同构，在 $i = n$ 时为满射。图示如下：
$$
\xymatrix{
& F^n \ar[r] \ar[d]^\alpha & F^{n + 1} \ar[d]^\alpha \ar[r] & \ldots \\
M^{n-1} \ar[r] & M^n \ar[r] & M^{n + 1} \ar[r] & \ldots
}
$$
由于每个 $M^\bullet_{f_j}$ 的上同调在高次数消失，对 $n \gg 0$ 可以找到
这样的映射。我们将对 $n$ 归纳，把它延拓为复形映射
$F^\bullet \to M^\bullet$，使得对所有 $i$，$H^i(\alpha_{f_j})$ 都是同构。
取 $\widetilde{F^\bullet}$ 即得引理。

\medskip\noindent
归纳步骤是在上图中加入 $F^{n - 1}$。设 $C^\bullet$ 为 $\alpha$ 的锥
（Derived Categories, Definition \ref{derived-definition-cone}）。上同调长正合
列说明，对 $i \geq n$ 有 $H^i(C^\bullet_{f_j}) = 0$。由 More on Algebra,
Lemma \ref{more-algebra-lemma-cone-pseudo-coherent}，$C^\bullet_{f_j}$ 是
$(n - 1)$-伪凝聚的。由 More on Algebra, Lemma
\ref{more-algebra-lemma-finite-cohomology}，
$H^{n - 1}(C^\bullet_{f_j})$ 是有限 $A_{f_j}$-模。选择有限自由 $A$-模
$F^{n - 1}$ 和 $A$-模同态 $\beta : F^{n - 1} \to C^{n - 1}$，使得复合
$F^{n - 1} \to C^{n - 1} \to C^n$ 为零，且 $F^{n - 1}_{f_j}$ 满射到
$H^{n - 1}(C^\bullet_{f_j})$。（略去一些细节；提示：清除分母。）由于
$C^{n - 1} = M^{n - 1} \oplus F^n$，可以写成
$\beta = (\alpha^{n - 1}, -d^{n - 1})$。复合
$F^{n - 1} \to C^{n - 1} \to C^n$ 的消失性说明这些映射组成复形态射
$$
\xymatrix{
& F^{n - 1} \ar[d]^{\alpha^{n - 1}} \ar[r]_{d^{n - 1}} &
F^n \ar[r] \ar[d]^\alpha &
F^{n + 1} \ar[d]^\alpha \ar[r] & \ldots \\
\ldots \ar[r] &
M^{n - 1} \ar[r] & M^n \ar[r] & M^{n + 1} \ar[r] & \ldots
}
$$
此外，这些映射定义可区别三角的态射
$$
\xymatrix{
(F^n \to \ldots) \ar[r] \ar[d] &
(F^{n-1} \to \ldots) \ar[r] \ar[d] &
F^{n-1} \ar[r] \ar[d]_\beta &
(F^n \to \ldots)[1] \ar[d] \\
(F^n \to \ldots) \ar[r] &
M^\bullet \ar[r] &
C^\bullet \ar[r] &
(F^n \to \ldots)[1]
}
$$
因此，对 $\beta$ 的选取蕴含复形映射
$(F^{-1} \to \ldots) \to M^\bullet$ 在 $f_j$ 处局部化后的上同调上，于次数
$\geq n$ 诱导同构，并在次数 $n - 1$ 诱导满射。这完成了引理的证明。
\end{proof}

\noindent
以下两个引理或许应放在别处。

\begin{lemma}
\label{perfect-lemma-vanishing-ext}
设 $X$ 为拟紧拟分离概形，$E \in D^b_\QCoh(\mathcal{O}_X)$。存在整数
$n_0 > 0$，使得对每个有限局部自由 $\mathcal{O}_X$-模 $\mathcal{E}$ 和
每个 $n \geq n_0$，都有
$\Ext^n_{D(\mathcal{O}_X)}(\mathcal{E}, E) = 0$
\end{lemma}

\begin{proof}
回忆 $\Ext^n_{D(\mathcal{O}_X)}(\mathcal{E}, E) =
\Hom_{D(\mathcal{O}_X)}(\mathcal{E}, E[n])$。导出范畴中的态射满足
Mayer--Vietoris，见 Cohomology, Lemma
\ref{cohomology-lemma-mayer-vietoris-hom}。因此，若 $X = U \cup V$，且对某个
界 $n_0$，引理的结论对 $E|_U$、$E|_V$ 和 $E|_{U \cap V}$ 成立，则结论对
$E$ 以 $n_0 + 1$ 为界成立。因此只需在 $X$ 仿射时证明，见 Cohomology of
Schemes, Lemma \ref{coherent-lemma-induction-principle}。

\medskip\noindent
假设 $X = \Spec(A)$ 仿射。选择一个 $A$-模复形 $M^\bullet$，使其相伴的
拟凝聚模复形表示 $E$，见引理 \ref{perfect-lemma-affine-compare-bounded}。把
$\mathcal{E} = \widetilde{P}$ 写成某个 $A$-模 $P$ 的相伴层。由于
$\mathcal{E}$ 有限局部自由，$P$ 是有限投射 $A$-模。于是
\begin{align*}
\Hom_{D(\mathcal{O}_X)}(\mathcal{E}, E[n])
& = 
\Hom_{D(A)}(P, M^\bullet[n]) \\
& =
\Hom_{K(A)}(P, M^\bullet[n]) \\
& =
\Hom_A(P, H^n(M^\bullet))
\end{align*}
第一个等号由引理 \ref{perfect-lemma-affine-compare-bounded} 得到，第二个等号由
Derived Categories, Lemma
\ref{derived-lemma-morphisms-from-projective-complex} 得到，最后一个等号是因为
$\Hom_A(P, -)$ 是正合函子。由于 $E$、从而 $M^\bullet$ 有界，对所有充分大的
$n$，结果为零。
\end{proof}

\noindent
以下引理还可以加强（对所有仅在固定范围内具有非零上同调层的 $L$，其消失
性具有一致性）。

\begin{lemma}
\label{perfect-lemma-ext-from-perfect-into-bounded-QCoh}
设 $X$ 为拟紧拟分离概形，$K$ 为 $D(\mathcal{O}_X)$ 的完美对象。则
\begin{enumerate}
\item 存在整数 $a \leq b$，使得对任意 $L \in D_\QCoh(\mathcal{O}_X)$，
若对 $i \in [a, b]$ 有 $H^i(L) = 0$，则
$\Hom_{D(\mathcal{O}_X)}(K, L) = 0$；
\item 若 $L$ 有界，则除有限多个 $n$ 外，
$\Ext^n_{D(\mathcal{O}_X)}(K, L)$ 均为零。
\end{enumerate}
\end{lemma}

\begin{proof}
由于 $\Ext^n_{D(\mathcal{O}_X)}(K, L) =
\Hom_{D(\mathcal{O}_X)}(K, L[n])$，(2) 由 (1) 得到。下面证明 (1)。由于
$K$ 完美，有
$$
\Hom_{D(\mathcal{O}_X)}(K, L) =
H^0(X, K^\vee \otimes_{\mathcal{O}_X}^\mathbf{L} L)
$$
其中 $K^\vee$ 是 $K$ 的“对偶”完美复形，见 Cohomology, Lemma
\ref{cohomology-lemma-dual-perfect-complex}。注意，由引理
\ref{perfect-lemma-quasi-coherence-tensor-product} 和 \ref{perfect-lemma-pseudo-coherent}
（后者用于说明完美复形具有拟凝聚上同调层），
$K^\vee \otimes_{\mathcal{O}_X}^\mathbf{L} L$ 属于 $D_\QCoh(X)$。设
$K^\vee$ 的 Tor 振幅在 $[a, b]$ 内。则谱序列
$$
E_1^{p, q} = H^p(K^\vee \otimes_{\mathcal{O}_X}^\mathbf{L} H^q(L))
\Rightarrow
H^{p + q}(K^\vee \otimes_{\mathcal{O}_X}^\mathbf{L} L)
$$
说明，若对 $q \in [j - b, j - a]$ 有 $H^q(L) = 0$，则
$H^j(K^\vee \otimes_{\mathcal{O}_X}^\mathbf{L} L)$ 为零。令 $N$ 为
Cohomology of Schemes, Lemma
\ref{coherent-lemma-vanishing-nr-affines-quasi-separated} 中的整数 $d$。如果
上同调层
$$
H^{-N}(K^\vee \otimes_{\mathcal{O}_X}^\mathbf{L} L),
\ H^{-N + 1}(K^\vee \otimes_{\mathcal{O}_X}^\mathbf{L} L),
\ \ldots,
\ H^0(K^\vee \otimes_{\mathcal{O}_X}^\mathbf{L} L)
$$
均为零，则 $H^0(X, K^\vee \otimes_{\mathcal{O}_X}^\mathbf{L} L)$ 消失。事实上，
由刚才引用的引理和引理 \ref{perfect-lemma-application-nice-K-injective}，有
$$
H^0(X, K^\vee \otimes_{\mathcal{O}_X}^\mathbf{L} L) =
H^0(X, \tau_{\geq -N}(K^\vee \otimes_{\mathcal{O}_X}^\mathbf{L} L))
$$
并且由上同调层的消失性，它等于
$H^0(X, \tau_{\geq 1}(K^\vee \otimes_{\mathcal{O}_X}^\mathbf{L} L))$
；由 Derived Categories, Lemma \ref{derived-lemma-negative-vanishing}，后者为
零。因此，若对 $i \in [-b - N, -a]$ 有 $H^i(L) = 0$，则
$\Hom_{D(\mathcal{O}_X)}(K, L)$ 为零。
\end{proof}

\begin{lemma}
\label{perfect-lemma-lift-perfect-complex-plus-locally-free}
设 $X$ 为仿射概形，$U \subset X$ 为拟紧开集。对 $D(\mathcal{O}_U)$ 的每个
完美对象 $E$，存在整数 $r$ 和 $U$ 上的有限局部自由层 $\mathcal{F}$，使得
$\mathcal{F}[-r] \oplus E$ 是 $D(\mathcal{O}_X)$ 的某个完美对象的限制。
\end{lemma}

\begin{proof}
设 $X = \Spec(A)$。回忆完美复形是伪凝聚的，见 Cohomology, Lemma
\ref{cohomology-lemma-perfect}。由引理 \ref{perfect-lemma-lift-pseudo-coherent}，可以
找到有限自由 $A$-模的有上界复形 $\mathcal{F}^\bullet$，使 $E$ 在
$D(\mathcal{O}_U)$ 中同构于 $\mathcal{F}^\bullet|_U$。由 Cohomology, Lemma
\ref{cohomology-lemma-perfect} 以及 $U$ 的拟紧性，$E$ 具有有限 Tor 维数；
设 $E$ 的 Tor 振幅在 $[a, b]$ 内。选取 $r < a$，并令
$$
\mathcal{K} = \Ker(\mathcal{F}^{r} \to \mathcal{F}^{r + 1})
= \Im(\mathcal{F}^{r - 1} \to \mathcal{F}^r).
$$
由于 $E$ 的 Tor 振幅在 $[a, b]$ 内，由 Cohomology, Lemma
\ref{cohomology-lemma-last-one-flat}，$\mathcal{F} = \mathcal{K}|_U$ 平坦。因此
$\mathcal{F}$ 平坦且有限表示，从而有限局部自由（Properties, Lemma
\ref{properties-lemma-finite-locally-free}）。于是
$$
\mathcal{F} \to \mathcal{F}^r|_U \to \mathcal{F}^{r + 1}|_U \to \ldots
$$
是 $U$ 上表示 $E$ 的严格完美复形。另一方面，复形
$P = (\mathcal{F}^r \to \mathcal{F}^{r + 1} \to \ldots )$ 是 $X$ 上的完美
复形。使用愚截断，得到可区别三角
$$
P|_U \to E \to \mathcal{F}[-r - 1] \to (P|_U)[1]
$$
若映射 $E \to \mathcal{F}[-r - 1]$ 在 $D(\mathcal{O}_U)$ 中为零，则
$P|_U = \mathcal{F}[-r - 2] \oplus E$，见 Derived Categories, Lemma
\ref{derived-lemma-split}。例如由引理
\ref{perfect-lemma-ext-from-perfect-into-bounded-QCoh}，当 $r \ll 0$ 时这必定成立。
\end{proof}

\begin{lemma}
\label{perfect-lemma-lift-map}
设 $X$ 为仿射概形，$U \subset X$ 为拟紧开集。设 $E, E'$ 为
$D_\QCoh(\mathcal{O}_X)$ 的对象，其中 $E$ 完美。对每个映射
$\alpha : E|_U \to E'|_U$，存在 $X$ 上的复形映射
$$
E \xleftarrow{\beta} E_1 \xrightarrow{\gamma} E'
$$
，其中 $E_1$ 完美，$\beta : E_1 \to E$ 在 $U$ 上限制为同构，且
$\alpha = \gamma|_U \circ \beta|_U^{-1}$。此外，可以假设对 $X$ 上的某个
完美复形 $I$，有 $E_1 = E \otimes_{\mathcal{O}_X}^\mathbf{L} I$。
\end{lemma}

\begin{proof}
写成 $X = \Spec(A)$ 及 $U = D(f_1) \cup \ldots \cup D(f_r)$。选择表示 $E$
的有限投射 $A$-模的有限复形 $M^\bullet$（引理
\ref{perfect-lemma-perfect-affine}），并选择表示 $E'$ 的 $A$-模复形
$(M')^\bullet$（引理 \ref{perfect-lemma-affine-compare-bounded}）。此时复形
$H^\bullet = \Hom_A(M^\bullet, (M')^\bullet)$ 是 $A$-模复形，其相伴的拟凝聚
$\mathcal{O}_X$-模复形表示 $R\SheafHom(E, E')$，见 Cohomology, Lemma
\ref{cohomology-lemma-Rhom-strictly-perfect}。于是 $\alpha$ 确定
$H^0(U, R\SheafHom(E, E'))$ 的元素 $s$，见 Cohomology, Lemma
\ref{cohomology-lemma-section-RHom-over-U}。由命题
\ref{perfect-proposition-represent-cohomology-class-on-open}，存在 $e$ 及对应于 $s$
的映射
$$
\xi : I^\bullet(f_1^e, \ldots, f_r^e) \to \Hom_A(M^\bullet, (M')^\bullet)
$$
。令 $E_1$ 为与如下复形相伴的拟凝聚 $\mathcal{O}_X$-模复形所对应的对象：
$$
\text{Tot}(I^\bullet(f_1^e, \ldots, f_r^e) \otimes_A M^\bullet)
$$
利用典范映射 $I^\bullet(f_1^e, \ldots, f_r^e) \to A$ 得到
$E_1 \to E$；利用 $\xi$ 和 Cohomology, Lemma
\ref{cohomology-lemma-section-RHom-over-U} 得到 $E_1 \to E'$。
\end{proof}

\begin{lemma}
\label{perfect-lemma-lift-perfect-complex-plus-shift}
设 $X$ 为仿射概形，$U \subset X$ 为拟紧开集。对
$D(\mathcal{O}_U)$ 的每个完美对象 $F$，对象 $F \oplus F[1]$ 是
$D(\mathcal{O}_X)$ 的某个完美对象的限制。
\end{lemma}

\begin{proof}
由引理 \ref{perfect-lemma-lift-perfect-complex-plus-locally-free}，可找到
$D(\mathcal{O}_X)$ 的完美对象 $E$，使得对某个有限局部自由
$\mathcal{O}_U$-模 $\mathcal{F}$ 有 $E|_U = \mathcal{F}[r] \oplus F$。由引理
\ref{perfect-lemma-lift-map}，可找到完美复形的态射 $\alpha : E_1 \to E$，使
$(E_1)|_U \cong E|_U$，且 $\alpha|_U$ 是映射
$$
\left(
\begin{matrix}
\text{id}_{\mathcal{F}[r]} & 0 \\
0 & 0
\end{matrix}
\right)
:
\mathcal{F}[r] \oplus F \to \mathcal{F}[r] \oplus F
$$
于是 $\alpha$ 的锥就是一个所求对象。
\end{proof}

\begin{lemma}
\label{perfect-lemma-perfect-into-support-on-T}
设 $X$ 为拟紧拟分离概形，$f \in \Gamma(X, \mathcal{O}_X)$。对
$D_\QCoh(\mathcal{O}_X)$ 中任意态射 $\alpha : E \to E'$，若
\begin{enumerate}
\item $E$ 完美；
\item $E'$ 支持在 $T = V(f)$ 上，
\end{enumerate}
则存在 $n \geq 0$ 使得 $f^n \alpha  = 0$。
\end{lemma}

\begin{proof}
导出范畴中的态射满足 Mayer--Vietoris，见 Cohomology, Lemma
\ref{cohomology-lemma-mayer-vietoris-hom}。因此，若 $X = U \cup V$ 且引理的
结论对 $f|_U$、$f|_V$ 和 $f|_{U \cap V}$ 成立，则结论也对 $f$ 成立。
所以只需在 $X$ 仿射时证明，见 Cohomology of Schemes, Lemma
\ref{coherent-lemma-induction-principle}。

\medskip\noindent
设 $X = \Spec(A)$。则 $f \in A$。以下不再特别说明地使用引理
\ref{perfect-lemma-affine-compare-bounded} 的等价 $D(A) = D_\QCoh(X)$。用有限投射
$A$-模的有限复形 $P^\bullet$ 表示 $E$；这由引理
\ref{perfect-lemma-perfect-affine} 可行。设 $t$ 为使 $P^t$ 非零的最大整数。可区别
三角
$$
P^t[-t] \to P^\bullet \to \sigma_{\leq t - 1}P^\bullet \to P^t[-t + 1]
$$
说明，对复形 $P^\bullet$ 的长度归纳，可化到 $P^\bullet$ 只有一个非零项的
情形。再使用平移函子，可化到 $P^\bullet$ 仅由位于次数 $0$ 的单个有限投射
$A$-模 $P$ 组成的情形。用 $A$-模复形 $M^\bullet$ 表示 $E'$。于是
$\alpha$ 对应于映射 $P \to H^0(M^\bullet)$。由假设 (2)，模
$H^0(M^\bullet)$ 支持在 $V(f)$ 上，故 $H^0(M^\bullet)$ 的每个元素都被 $f$
的某个幂零化。
由于 $P$ 是有限 $A$-模，对某个 $n$，映射
$f^n\alpha : P \to H^0(M^\bullet)$ 为零，正如所需。
\end{proof}

\begin{lemma}
\label{perfect-lemma-lift-perfect-complex-plus-shift-support}
设 $X$ 为仿射概形，$T \subset X$ 为闭子集，使得 $X \setminus T$ 拟紧，并
设 $U \subset X$ 为拟紧开集。对 $D(\mathcal{O}_U)$ 中支持在 $T \cap U$ 上的
每个完美对象 $F$，对象 $F \oplus F[1]$ 是 $D(\mathcal{O}_X)$ 中某个支持在
$T$ 上的完美对象 $E$ 的限制。
\end{lemma}

\begin{proof}
设 $T = V(g_1, \ldots, g_s)$。用 $g_j$ 的某个幂替换它后，可假设乘以
$g_j$ 在 $F$ 上为零，见引理 \ref{perfect-lemma-perfect-into-support-on-T}。按引理
\ref{perfect-lemma-lift-perfect-complex-plus-shift} 选取 $E$。注意，
$g_j : E \to E$ 在 $U$ 上限制为零。选择可区别三角
$$
E \xrightarrow{g_1} E \to C_1 \to E[1]
$$
由 Derived Categories, Lemma \ref{derived-lemma-split}，对象 $C_1$ 在 $U$ 上
限制为 $F \oplus F[1] \oplus F[1] \oplus F[2]$。此外，由 Derived
Categories, Lemma \ref{derived-lemma-third-map-square-zero}，
$g_1 : C_1 \to C_1$ 的平方为零。事实上，图表
$$
\xymatrix{
E \ar[r] \ar[d]_0 & C_1 \ar[d]_{g_1} \ar[r] & E[1] \ar[d]_0 \\
E \ar[r] & C_1 \ar[r] & E[1]
}
$$
交换，因为复合 $E \xrightarrow{g_1} E \to C_1$ 和
$C_1 \to E[1] \xrightarrow{g_1} E[1]$ 均为零。继续进行，令 $C_{i + 1}$
等于映射 $g_i : C_i \to C_i$ 的锥，得到 $X$ 上支持在 $T$ 上的完美复形
$C_s$，它在 $U$ 上的限制为
$$
F \oplus F[1]^{\oplus s} \oplus F[2]^{\oplus {s \choose 2}}
\oplus \ldots \oplus F[s]
$$
按引理 \ref{perfect-lemma-lift-map} 选择完美复形的态射 $\beta : C' \to C_s$ 和
$\gamma : C' \to C_s$，使 $\beta|_U$ 为同构，且
$\gamma|_U \circ \beta|_U^{-1}$ 是态射
$$
F \oplus F[1]^{\oplus s} \oplus F[2]^{\oplus {s \choose 2}}
\oplus \ldots \oplus F[s]
\to
F \oplus F[1]^{\oplus s} \oplus F[2]^{\oplus {s \choose 2}}
\oplus \ldots \oplus F[s]
$$
，它在除 $F$ 外的所有直和项上都是恒等映射，而在该直和项上为零。由引理
\ref{perfect-lemma-lift-map}，对 $X$ 上的某个完美复形 $I$，还有
$C' = C_s \otimes^\mathbf{L} I$。因此 $g_j^2\text{id}_{C_s}$ 为零蕴含在
$C'$ 上也有同样的结论。故 $C'$ 也支持在 $T$ 上。于是
$\text{Cone}(\gamma)$ 就是所求对象。
\end{proof}

\noindent
下列引理的一个特例见 \cite{Neeman-Grothendieck}。

\begin{lemma}
\label{perfect-lemma-lift-map-from-perfect-complex-with-support}
设 $X$ 为拟紧拟分离概形，$U \subset X$ 为拟紧开集，$T \subset X$ 为闭
子集，且 $X \setminus T$ 在 $X$ 中逆紧。设 $E$ 为
$D_\QCoh(\mathcal{O}_X)$ 的对象，$\alpha : P \to E|_U$ 为映射，其中 $P$
是 $D(\mathcal{O}_U)$ 中支持在 $T \cap U$ 上的完美对象。则存在映射
$\beta : R \to E$，其中 $R$ 是 $D(\mathcal{O}_X)$ 中支持在 $T$ 上的完美
对象，使得 $P$ 是 $D(\mathcal{O}_U)$ 中 $R|_U$ 的直和项，并与 $\alpha$
和 $\beta|_U$ 相容。
\end{lemma}

\begin{proof}
由于 $X$ 拟紧，存在整数 $m$ 及 $X$ 的仿射开集 $V_j$，使得
$X = U \cup V_1 \cup \ldots \cup V_m$。对 $m$ 归纳，可假设 $m = 1$。
换言之，可假设 $X = U \cup V$，其中 $V$ 仿射。由引理
\ref{perfect-lemma-lift-perfect-complex-plus-shift-support}，可以在
$D(\mathcal{O}_V)$ 中选择支持在 $T \cap V$ 上的完美对象 $Q$，以及同构
$Q|_{U \cap V} \to (P \oplus P[1])|_{U \cap V}$。由引理
\ref{perfect-lemma-lift-map}，可用 $Q \otimes^\mathbf{L} I$ 替换 $Q$（它仍支持在
$T \cap V$ 上），并假设映射
$$
Q|_{U \cap V} \to (P \oplus P[1])|_{U \cap V}
\longrightarrow P|_{U \cap V}
\longrightarrow
E|_{U \cap V}
$$
可提升为 $Q \to E|_V$。由 Cohomology, Lemma
\ref{cohomology-lemma-glue}，得到 $D(\mathcal{O}_X)$ 的态射 $a : R \to E$，
使得 $a|_U$ 同构于 $P \oplus P[1] \to E|_U$，且 $a|_V$ 同构于
$Q \to E|_V$。于是 $R$ 完美且支持在 $T$ 上，正如所需。
\end{proof}

\begin{remark}
\label{perfect-remark-addendum}
引理 \ref{perfect-lemma-lift-map-from-perfect-complex-with-support} 的证明说明
$$
R|_U = P \oplus P^{\oplus n_1}[1] \oplus \ldots \oplus P^{\oplus n_m}[m]
$$
对某些 $m \geq 0$ 和 $n_j \geq 0$ 成立。因此 $R|_U$ 的最高次上同调层等于
$P$ 的最高次上同调层。对映射
$P^{\oplus n_1}[1] \oplus \ldots \oplus P^{\oplus n_m}[m] \to R|_U$ 重复该
构造、取锥并使用归纳，可以使 $R|_U$ 与 $P$ 在任意给定次数以上的上同调层
相等。
\end{remark}



\section{用完美复形逼近}
\label{perfect-section-approximation}

\noindent
本节讨论 Neeman 与 Lipman 的一个观察：伪凝聚复形可由完美复形“逼近”。

\begin{definition}
\label{perfect-definition-approximation-holds}
设 $X$ 为概形。考虑三元组 $(T, E, m)$，其中
\begin{enumerate}
\item $T \subset X$ 是闭子集；
\item $E$ 是 $D_\QCoh(\mathcal{O}_X)$ 的对象；
\item $m \in \mathbf{Z}$。
\end{enumerate}
若存在 $D(\mathcal{O}_X)$ 中支持在 $T$ 上的完美对象 $P$ 及映射
$\alpha : P \to E$，使它对 $i > m$ 诱导同构 $H^i(P) \to H^i(E)$，并诱导
满射 $H^m(P) \to H^m(E)$，则称{\it 三元组满足逼近性质} $(T, E, m)$。
\end{definition}

\noindent
逼近性质不可能对每个三元组都成立。事实上，若三元组 $(T, E, m)$ 满足逼近
性质，则显然
\begin{enumerate}
\item $E$ 是 $m$-伪凝聚的，见 Cohomology, Definition
\ref{cohomology-definition-pseudo-coherent}；
\item 对 $i \geq m$，上同调层 $H^i(E)$ 支持在 $T$ 上。
\end{enumerate}
此外，完美复形的“支集”是闭子概形，其补集在 $X$ 中逆紧（略去细节）。
因此若不对 $T$ 作此假设，就不能期待逼近性质成立。这部分地解释了下一个
定义中的条件。

\begin{definition}
\label{perfect-definition-approximation}
设 $X$ 为概形。若对任意满足 $X \setminus T$ 在 $X$ 中逆紧的闭子集
$T \subset X$，都存在整数 $r$，使得对定义
\ref{perfect-definition-approximation-holds} 中每个满足下列条件的三元组
$(T, E, m)$：
\begin{enumerate}
\item $E$ 是 $(m - r)$-伪凝聚的；
\item 对 $i \geq m - r$，$H^i(E)$ 支持在 $T$ 上，
\end{enumerate}
逼近性质都成立，则称 $X$ 上{\it 用完美复形逼近的性质成立}。
\end{definition}

\noindent
我们将证明，用完美复形逼近的性质对拟紧拟分离概形成立。第二个条件看来是
我们的证明方法所必需的。仔细分析以下论证，或许可以把第一个条件弱化为
“$E$ 是 $m$-伪凝聚的”。

\begin{lemma}
\label{perfect-lemma-open}
设 $X$ 为概形，$U \subset X$ 为开子概形，$(T, E, m)$ 为定义
\ref{perfect-definition-approximation-holds} 中的三元组。若
\begin{enumerate}
\item $T \subset U$；
\item $(T, E|_U, m)$ 满足逼近性质；
\item 对 $i \geq m$，层 $H^i(E)$ 支持在 $T$ 上，
\end{enumerate}
则 $(T, E, m)$ 满足逼近性质。
\end{lemma}

\begin{proof}
设 $j : U \to X$ 为包含态射。若 $P \to E|_U$ 是 $U$ 上三元组
$(T, E|_U, m)$ 的逼近，则 $j_!P = Rj_*P \to j_!(E|_U) \to E$ 是 $X$ 上
$(T, E, m)$ 的逼近。见 Cohomology, Lemmas
\ref{cohomology-lemma-pushforward-restriction} 和
\ref{cohomology-lemma-pushforward-perfect}。
\end{proof}

\begin{lemma}
\label{perfect-lemma-approximation-affine}
设 $X$ 为仿射概形。对定义 \ref{perfect-definition-approximation-holds} 中每个三元组
$(T, E, m)$，若存在整数 $r \geq 0$ 满足
\begin{enumerate}
\item $E$ 是 $m$-伪凝聚的；
\item 对 $i \geq m - r + 1$，$H^i(E)$ 支持在 $T$ 上；
\item $X \setminus T$ 是 $r$ 个仿射开集之并，
\end{enumerate}
则逼近性质成立。特别地，仿射概形上用完美复形逼近的性质成立。
\end{lemma}

\begin{proof}
设 $X = \Spec(A)$，并写成 $T = V(f_1, \ldots, f_r)$。（当 $r = 0$ 时，
即 $T = X$，结论立即由引理 \ref{perfect-lemma-pseudo-coherent-affine} 和定义得到。）
设 $(T, E, m)$ 为引理中的三元组。令 $t$ 为使 $H^t(E)$ 非零的最大整数，
并对 $t$ 归纳。初始情形为 $t < m$，此时结论显然。现假设 $t \geq m$。
由 Cohomology, Lemma \ref{cohomology-lemma-finite-cohomology}，层 $H^t(E)$
为有限型。由于它拟凝聚，可由有限多个截面生成（Properties, Lemma
\ref{properties-lemma-finite-type-module}）。对每个
$s \in \Gamma(X, H^t(E)) = H^t(X, E)$（见引理
\ref{perfect-lemma-affine-compare-bounded} 的证明），由引理
\ref{perfect-lemma-represent-cohomology-class-on-closed}，可以找到 $e > 0$ 及态射
$K_e[-t] \to E$，使 $s$ 属于
$H^0(K_e) = H^t(K_e[-t]) \to H^t(E)$ 的像。取这些映射的有限直和，得到
映射 $P \to E$，其中 $P$ 是支持在 $T$ 上的完美复形，对 $i > t$ 有
$H^i(P) = 0$，且 $H^t(P) \to E$ 为满射。选择可区别三角
$$
P \to E \to E' \to P[1]
$$
于是 $E'$ 是 $m$-伪凝聚的（Cohomology, Lemma
\ref{cohomology-lemma-cone-pseudo-coherent}），对 $i \geq t$ 有
$H^i(E') = 0$，并且对 $i \geq m - r + 1$，$H^i(E')$ 支持在 $T$ 上。由
归纳，找到 $(T, E', m)$ 的逼近 $P' \to E'$。把复合
$P' \to E' \to P[1]$ 置于可区别三角 $P \to P'' \to P' \to P[1]$ 中，
并使用 TR3，把态射 $P' \to E'$ 和 $P[1] \to P[1]$ 延拓为可区别三角的态射
$$
\xymatrix{
P \ar[r] \ar[d] & P'' \ar[d] \ar[r] & P' \ar[d] \ar[r] & P[1] \ar[d] \\
P \ar[r] &  E \ar[r] & E' \ar[r] & P[1]
}
$$
。于是 $P''$ 是支持在 $T$ 上的完美复形（Cohomology, Lemma
\ref{cohomology-lemma-two-out-of-three-perfect}）。简单的图追逐说明，
$P'' \to E$ 就是所求逼近。
\end{proof}

\begin{lemma}
\label{perfect-lemma-induction-step}
设 $X$ 为概形，$X = U \cup V$ 为开覆盖，其中 $U$ 拟紧，$V$ 仿射，且
$U \cap V$ 拟紧。若 $U$ 上用完美复形逼近的性质成立，则 $X$ 上也成立。
\end{lemma}

\begin{proof}
设 $T \subset X$ 为闭子集，且 $X \setminus T$ 在 $X$ 中逆紧。设 $r_U$ 为
定义 \ref{perfect-definition-approximation} 中适用于二元组 $(U, T \cap U)$ 的整数。
令 $T' = T \setminus U$。注意，$T' \subset V$；并且由对 $T$ 的假设，
$V \setminus T' = (X \setminus T) \cap U \cap V$ 拟紧。设 $r'$ 为覆盖
$V \setminus T'$ 所需的仿射开集数。我们断言
$r = \max(r_U, r')$ 对二元组 $(X, T)$ 适用。

\medskip\noindent
为此，选择三元组 $(T, E, m)$，使 $E$ 是 $(m - r)$-伪凝聚的，并且对
$i \geq m - r$，$H^i(E)$ 支持在 $T$ 上。令 $t$ 为使 $H^t(E)|_U$ 非零的
最大整数。（这样的整数存在，因为 $U$ 拟紧且 $E|_U$ 是
$(m - r)$-伪凝聚的。）我们对 $t$ 归纳证明 $E$ 可被逼近。

\medskip\noindent
初始情形：$t \leq m - r'$。这意味着对 $i \geq m - r'$，$H^i(E)$ 支持在
$T'$ 上。因此引理 \ref{perfect-lemma-approximation-affine} 保证 $V$ 上存在
$(T', E|_V, m)$ 的逼近 $P \to E|_V$。应用引理 \ref{perfect-lemma-open} 可知，
$(T', E, m)$ 可被逼近。这样的逼近也是 $(T, E, m)$ 的逼近。

\medskip\noindent
归纳步骤。选择 $(T \cap U, E|_U, m)$ 的逼近 $P \to E|_U$。这特别给出满射
$H^t(P) \to H^t(E|_U)$。由引理
\ref{perfect-lemma-lift-perfect-complex-plus-shift-support}，可在
$D(\mathcal{O}_V)$ 中选择支持在 $T \cap V$ 上的完美对象 $Q$，以及同构
$Q|_{U \cap V} \to (P \oplus P[1])|_{U \cap V}$。由引理
\ref{perfect-lemma-lift-map}，可用 $Q \otimes^\mathbf{L} I$ 替换 $Q$，并假设映射
$$
Q|_{U \cap V} \to (P \oplus P[1])|_{U \cap V}
\longrightarrow P|_{U \cap V}
\longrightarrow
E|_{U \cap V}
$$
可提升为 $Q \to E|_V$。由 Cohomology, Lemma
\ref{cohomology-lemma-glue}，得到 $D(\mathcal{O}_X)$ 中的态射 $a : R \to E$，
使得 $a|_U$ 同构于 $P \oplus P[1] \to E|_U$，而 $a|_V$ 同构于
$Q \to E|_V$。因此 $R$ 完美且支持在 $T$ 上，并且映射
$H^t(R) \to H^t(E)$ 限制到 $U$ 上后为满射。选择可区别三角
$$
R \to E \to E' \to R[1]
$$
于是 $E'$ 是 $(m - r)$-伪凝聚的（Cohomology, Lemma
\ref{cohomology-lemma-cone-pseudo-coherent}），对 $i \geq t$ 有
$H^i(E')|_U = 0$，并且对 $i \geq m - r$，$H^i(E')$ 支持在 $T$ 上。由
归纳，找到 $(T, E', m)$ 的逼近 $R' \to E'$。把复合
$R' \to E' \to R[1]$ 置于可区别三角 $R \to R'' \to R' \to R[1]$ 中，
并把态射 $R' \to E'$ 和 $R[1] \to R[1]$ 延拓为可区别三角的态射
$$
\xymatrix{
R \ar[r] \ar[d] & R'' \ar[d] \ar[r] & R' \ar[d] \ar[r] & R[1] \ar[d] \\
R \ar[r] &  E \ar[r] & E' \ar[r] & R[1]
}
$$
，这里使用 TR3。于是 $R''$ 是支持在 $T$ 上的完美复形（Cohomology,
Lemma \ref{cohomology-lemma-two-out-of-three-perfect}）。简单的图追逐说明，
$R'' \to E$ 是所求逼近。
\end{proof}

\begin{theorem}
\label{perfect-theorem-approximation}
设 $X$ 为拟紧拟分离概形。则 $X$ 上用完美复形逼近的性质成立。
\end{theorem}

\begin{proof}
这由 Cohomology of Schemes, Lemma
\ref{coherent-lemma-induction-principle} 的归纳原理，以及引理
\ref{perfect-lemma-induction-step} 和 \ref{perfect-lemma-approximation-affine} 得到。
\end{proof}






\section{生成导出范畴}
\label{perfect-section-generating}

\noindent
本节证明，拟紧拟分离概形的导出范畴 $D_\QCoh(\mathcal{O}_X)$ 可由单个完美
对象生成。我们强烈建议读者阅读 Bondal 与 van den Bergh 那篇精彩论文中对此
结果的证明，见 \cite{BvdB}。

\begin{lemma}
\label{perfect-lemma-direct-summand-of-a-restriction}
设 $X$ 为拟紧拟分离概形，$U$ 为拟紧开子概形，$P$ 为
$D(\mathcal{O}_U)$ 的完美对象。则 $P$ 是 $D(\mathcal{O}_X)$ 某个完美对象之
限制的直和项。
\end{lemma}

\begin{proof}
这是引理 \ref{perfect-lemma-lift-map-from-perfect-complex-with-support} 的特例。
\end{proof}

\begin{lemma}
\label{perfect-lemma-orthogonal-koszul-complex}
\begin{reference}
\cite[Proposition 6.1]{Bokstedt-Neeman}
\end{reference}
在情形 \ref{perfect-situation-complex} 中，以 $j : U \to X$ 表示开浸入，并设 $K$
为 $D(\mathcal{O}_X)$ 中对应于 $A$ 上 $f_1, \ldots, f_r$ 的 Koszul 复形的
完美对象。对 $E \in D_\QCoh(\mathcal{O}_X)$，以下条件等价：
\begin{enumerate}
\item $E = Rj_*(E|_U)$；
\item 对所有 $n \in \mathbf{Z}$，
$\Hom_{D(\mathcal{O}_X)}(K[n], E) = 0$。
\end{enumerate}
\end{lemma}

\begin{proof}
选择可区别三角 $E \to Rj_*(E|_U) \to N \to E[1]$。注意
$$
\Hom_{D(\mathcal{O}_X)}(K[n], Rj_*(E|_U)) =
\Hom_{D(\mathcal{O}_U)}(K|_U[n], E) = 0
$$
对所有 $n$ 都成立，因为 $K|_U = 0$。因此只需对 $N$ 证明。换言之，可假设
$E$ 在 $U$ 上的限制为零。注意，有 Koszul 复形的可区别三角
$$
K^\bullet(f_1^{e_1}, \ldots, f_i^{e'_i}, \ldots, f_r^{e_r}) \to
K^\bullet(f_1^{e_1}, \ldots, f_i^{e'_i + e''_i}, \ldots, f_r^{e_r}) \to
K^\bullet(f_1^{e_1}, \ldots, f_i^{e''_i}, \ldots, f_r^{e_r}) \to \ldots
$$
，见 More on Algebra, Lemma \ref{more-algebra-lemma-koszul-mult}。因此，若对
所有 $n \in \mathbf{Z}$ 有 $\Hom_{D(\mathcal{O}_X)}(K[n], E) = 0$，则用引理
\ref{perfect-lemma-represent-cohomology-class-on-closed} 中的 $K_e$ 替换 $K$ 后，同一
结论仍成立。于是本引理由该引理以及如下事实立即得到：$E$ 由 $A$-模复形
$R\Gamma(X, E)$ 确定，见引理 \ref{perfect-lemma-affine-compare-bounded}。
\end{proof}

\begin{theorem}
\label{perfect-theorem-bondal-van-den-Bergh}
设 $X$ 为拟紧拟分离概形。范畴 $D_\QCoh(\mathcal{O}_X)$ 可由单个完美对象
生成。更确切地说，存在 $D(\mathcal{O}_X)$ 的完美对象 $P$，使得对
$E \in D_\QCoh(\mathcal{O}_X)$，以下条件等价：
\begin{enumerate}
\item $E = 0$；
\item 对所有 $n \in \mathbf{Z}$，
$\Hom_{D(\mathcal{O}_X)}(P[n], E) = 0$。
\end{enumerate}
\end{theorem}

\begin{proof}
我们使用 Cohomology of Schemes, Lemma
\ref{coherent-lemma-induction-principle} 的归纳原理证明。

\medskip\noindent
若 $X$ 仿射，则 $\mathcal{O}_X$ 是完美生成元。这由引理
\ref{perfect-lemma-affine-compare-bounded} 得到。

\medskip\noindent
假设 $X = U \cup V$ 为开覆盖，其中 $U$ 拟紧，定理对 $U$ 成立，且 $V$ 为
仿射开集。设 $P$ 为 $D(\mathcal{O}_U)$ 的完美对象，且是
$D_\QCoh(\mathcal{O}_U)$ 的生成元。由引理
\ref{perfect-lemma-direct-summand-of-a-restriction}，可选择 $D(\mathcal{O}_X)$ 的完美
对象 $Q$，使其在 $U$ 上的限制是一个直和，而 $P$ 是其中一个直和项。设
$V = \Spec(A)$。令 $Z = X \setminus U$。这是 $V$ 的闭子集，且
$V \setminus Z$ 拟紧。选择 $f_1, \ldots, f_r \in A$，使
$Z = V(f_1, \ldots, f_r)$。设 $K \in D(\mathcal{O}_V)$ 为对应于 $A$ 上
$f_1, \ldots, f_r$ 的 Koszul 复形的完美对象。注意，由于 $K$ 支持在闭子集
$Z \subset V$ 上，推前 $K' = R(V \to X)_*K$ 是 $D(\mathcal{O}_X)$ 的完美
对象，其在 $V$ 上的限制为 $K$（见 Cohomology, Lemma
\ref{cohomology-lemma-pushforward-perfect}）。我们断言 $Q \oplus K'$ 是
$D_\QCoh(\mathcal{O}_X)$ 的生成元。

\medskip\noindent
设 $E$ 为 $D_\QCoh(\mathcal{O}_X)$ 的对象，使得从 $Q \oplus K'$ 的任意
平移到 $E$ 都没有非零映射。由 Cohomology, Lemma
\ref{cohomology-lemma-pushforward-restriction}，有
$K' =  R(V \to X)_! K$，因而
$$
\Hom_{D(\mathcal{O}_X)}(K'[n], E) = \Hom_{D(\mathcal{O}_V)}(K[n], E|_V)
$$
于是由引理 \ref{perfect-lemma-orthogonal-koszul-complex}，这些群的消失性蕴含
$E|_V$ 同构于 $R(U \cap V \to V)_*E|_{U \cap V}$。这又蕴含
$E = R(U \to X)_*E|_U$（略去一个小细节）。在此情形下，
$$
\Hom_{D(\mathcal{O}_X)}(Q[n], E) = \Hom_{D(\mathcal{O}_U)}(Q|_U[n], E|_U)
$$
其直和项之一是 $\Hom_{D(\mathcal{O}_U)}(P[n], E|_U)$。因此由对 $P$ 的
选取，这些群的消失性蕴含 $E|_U$ 为零，从而 $E$ 为零。
\end{proof}

\noindent
以下结果是定理 \ref{perfect-theorem-bondal-van-den-Bergh} 的加强，且用完全相同的
方法证明。回忆，对概形 $X$ 的闭子集 $T$，以 $D_T(\mathcal{O}_X)$ 表示
$D(\mathcal{O}_X)$ 中由支持在 $T$ 上的对象组成的严格满、饱和三角子范畴
（定义 \ref{perfect-definition-supported-on}）。类似地，以
$D_{\QCoh, T}(\mathcal{O}_X)$ 表示 $D(\mathcal{O}_X)$ 中由上同调层拟凝聚且
支持在 $T$ 上的复形组成的严格满、饱和三角子范畴。

\begin{lemma}
\label{perfect-lemma-generator-with-support}
\begin{reference}
\cite[Theorem 6.8]{Rouquier-dimensions}
\end{reference}
设 $X$ 为拟紧拟分离概形，$T \subset X$ 为闭子集，且 $X \setminus T$ 拟紧。
采用上述记号，范畴 $D_{\QCoh, T}(\mathcal{O}_X)$ 由单个完美对象生成。
\end{lemma}

\begin{proof}
我们使用 Cohomology of Schemes, Lemma
\ref{coherent-lemma-induction-principle} 的归纳原理证明。

\medskip\noindent
假设 $X = \Spec(A)$ 仿射。此时存在
$f_1, \ldots, f_r \in A$，使 $T = V(f_1, \ldots, f_r)$。设 $K$ 为引理
\ref{perfect-lemma-orthogonal-koszul-complex} 中 $f_1, \ldots, f_r$ 上的 Koszul 复形。
于是 $K$ 是上同调支持在 $T$ 上的完美对象，因而是
$D_{\QCoh, T}(\mathcal{O}_X)$ 的完美对象。另一方面，若
$E \in D_{\QCoh, T}(\mathcal{O}_X)$ 且对所有 $n$ 有
$\Hom(K, E[n]) = 0$，则引理 \ref{perfect-lemma-orthogonal-koszul-complex} 告诉我们
$E = Rj_*(E|_{X \setminus T}) = 0$。因此 $K$ 生成
$D_{\QCoh, T}(\mathcal{O}_X)$（按 Derived Categories, Definition
\ref{derived-definition-generators} 中三角范畴生成元的定义）。

\medskip\noindent
假设 $X = U \cup V$ 为开覆盖，其中 $V$ 仿射、$U$ 拟紧，且引理对 $U$
成立。设 $P$ 为 $D(\mathcal{O}_U)$ 中支持在 $T \cap U$ 上的完美对象，且是
$D_{\QCoh, T \cap U}(\mathcal{O}_U)$ 的生成元。由引理
\ref{perfect-lemma-lift-map-from-perfect-complex-with-support}，可选择
$D(\mathcal{O}_X)$ 中支持在 $T$ 上的完美对象 $Q$，使其在 $U$ 上的限制是
一个直和，而 $P$ 是其中一个直和项。写成 $V = \Spec(B)$。令
$Z = X \setminus U$。则 $Z$ 是 $V$ 的闭子集，且 $V \setminus Z$ 拟紧。
由于 $X$ 拟分离，$Z \cap T$ 是 $V$ 的闭子集，且
$W = V \setminus (Z \cap T)$ 拟紧。因此可以选择
$g_1, \ldots, g_s \in B$，使 $Z \cap T = V(g_1, \ldots, g_r)$。设
$K \in D(\mathcal{O}_V)$ 为对应于 $B$ 上 $g_1, \ldots, g_s$ 的 Koszul 复形
的完美对象。注意，由于 $K$ 支持在闭子集 $(Z \cap T) \subset V$ 上，推前
$K' = R(V \to X)_*K$ 是 $D(\mathcal{O}_X)$ 的完美对象，其在 $V$ 上的限制
为 $K$（见 Cohomology, Lemma
\ref{cohomology-lemma-pushforward-perfect}）。我们断言 $Q \oplus K'$ 是
$D_{\QCoh, T}(\mathcal{O}_X)$ 的生成元。

\medskip\noindent
设 $E$ 为 $D_{\QCoh, T}(\mathcal{O}_X)$ 的对象，使得从 $Q \oplus K'$ 的任意
平移到 $E$ 都没有非零映射。由 Cohomology, Lemma
\ref{cohomology-lemma-pushforward-restriction}，有
$K' =  R(V \to X)_! K$，因而
$$
\Hom_{D(\mathcal{O}_X)}(K'[n], E) = \Hom_{D(\mathcal{O}_V)}(K[n], E|_V)
$$
因此由引理 \ref{perfect-lemma-orthogonal-koszul-complex}，有
$E|_V = Rj_*E|_W$，其中 $j : W \to V$ 是包含态射。图示如下：
$$
\xymatrix{
W \ar[r]_j & V & Z \cap T \ar[l] \ar[d] \\
U \cap V \ar[u]^{j'} \ar[ru]_{j''} & & Z \ar[lu]
}
$$
由于 $E$ 支持在 $T$ 上，$E|_W$ 支持在
$T \cap W = T \cap U \cap V$ 上，而后者在 $W$ 中闭。因此
$$
E|_V = Rj_*(E|_W) = Rj_*(Rj'_*(E|_{U \cap V})) = Rj''_*(E|_{U \cap V})
$$
其中第二个等号是 Cohomology, Lemma
\ref{cohomology-lemma-pushforward-restriction} 的 (1)。这蕴含
$E = R(U \to X)_*E|_U$（略去一个小细节）。在此情形下，
$$
\Hom_{D(\mathcal{O}_X)}(Q[n], E) = \Hom_{D(\mathcal{O}_U)}(Q|_U[n], E|_U)
$$
其直和项之一是 $\Hom_{D(\mathcal{O}_U)}(P[n], E|_U)$。因此由对 $P$ 的
选取，这些群的消失性蕴含 $E|_U$ 为零，从而 $E$ 为零。
\end{proof}



\section{一个生成元例子}
\label{perfect-section-example-generator}

\noindent
本节证明，环上射影空间的导出范畴由一个向量丛生成；事实上，它由结构层的
若干平移之直和生成。

\medskip\noindent
下列引理说明，若 $\mathcal{L}$ 丰沛，则
$\bigoplus_{n \geq 0} \mathcal{L}^{\otimes -n}$ 是生成元。

\begin{lemma}
\label{perfect-lemma-nonzero-some-cohomology}
设 $X$ 为概形，$\mathcal{L}$ 为丰沛可逆 $\mathcal{O}_X$-模。若 $K$ 是
$D_\QCoh(\mathcal{O}_X)$ 的非零对象，则对某些 $n \geq 0$ 和
$p \in \mathbf{Z}$，上同调群
$H^p(X, K \otimes_{\mathcal{O}_X}^\mathbf{L} \mathcal{L}^{\otimes n})$
非零。
\end{lemma}

\begin{proof}
回忆，由于 $X$ 具有丰沛可逆层，它拟紧且分离（Properties, Definition
\ref{properties-definition-ample} 和 Lemma
\ref{properties-lemma-affine-s-opens-cover-quasi-separated}）。因此可应用命题
\ref{perfect-proposition-quasi-compact-affine-diagonal}，用拟凝聚模复形
$\mathcal{F}^\bullet$ 表示 $K$。选取任意 $p$，使
$\mathcal{H}^p = \Ker(\mathcal{F}^p \to \mathcal{F}^{p + 1})/
\Im(\mathcal{F}^{p - 1} \to \mathcal{F}^p)$ 非零。选择点 $x \in X$，使茎
$\mathcal{H}^p_x$ 非零。选择 $n \geq 0$ 及
$s \in \Gamma(X, \mathcal{L}^{\otimes n})$，使 $X_s$ 是 $x$ 的仿射开邻域。
选择 $\tau \in \mathcal{H}^p(X_s)$，使它映到茎 $\mathcal{H}^p_x$ 的非零
元素；这是可行的，因为 $\mathcal{H}^p$ 拟凝聚且 $X_s$ 仿射。由于在
$X_s$ 上取截面是拟凝聚模上的正合函子，可以找到截面
$\tau' \in \mathcal{F}^p(X_s)$，它在 $\mathcal{F}^{p + 1}(X_s)$ 中映为零，
并在 $\mathcal{H}^p(X_s)$ 中映到 $\tau$。由 Properties, Lemma
\ref{properties-lemma-invert-s-sections}，存在 $m$，使
$\tau' \otimes s^{\otimes m}$ 是某个截面
$\tau'' \in \Gamma(X, \mathcal{F}^p \otimes \mathcal{L}^{\otimes mn})$ 的像。
再次应用同一引理，得到 $l \geq 0$，使
$\tau'' \otimes s^{\otimes l}$ 在
$\mathcal{F}^{p + 1} \otimes \mathcal{L}^{\otimes (m + l)n}$ 中映为零。于是
$\tau''$ 给出非零类
$H^p(X, K \otimes^\mathbf{L}_{\mathcal{O}_X} \mathcal{L}^{(m + l)n})$
，正如所需。
\end{proof}

\begin{lemma}
\label{perfect-lemma-construct-the-next-one}
设 $A$ 为环，$X = \mathbf{P}^n_A$。对每个 $a \in \mathbf{Z}$，存在 $X$
上向量丛的正合复形
$$
0 \to \mathcal{O}_X(a) \to \ldots
\to \mathcal{O}_X(a + i)^{\oplus {n + 1 \choose i}} \to
\ldots \to \mathcal{O}_X(a + n + 1) \to 0
$$
\end{lemma}

\begin{proof}
回忆 $\mathbf{P}^n_A$ 是 $\text{Proj}(A[X_0, \ldots, X_n])$，见
Constructions, Definition \ref{constructions-definition-projective-space}。
考虑 $S = A[X_0, \ldots, X_n]$ 上 $X_0, \ldots, X_n$ 的 Koszul 复形
$$
K_\bullet = K_\bullet(A[X_0, \ldots, X_n], X_0, \ldots, X_n)
$$
。由于 $X_0, \ldots, X_n$ 显然是多项式环 $S$ 中的正则序列，由 More on
Algebra, Lemma \ref{more-algebra-lemma-regular-koszul-regular}，Koszul 复形
$K_\bullet$ 除次数 $0$ 外均正合，而此次数的上同调为
$S/(X_0, \ldots, X_n)$。若把 $K_i$ 的生成元置于次数 $+i$，则
$K_\bullet$ 成为分次模复形。换言之，有正合复形
$$
0 \to S(-n - 1) \to \ldots \to S(-n - 1 + i)^{\oplus {n \choose i}} \to \ldots
\to S \to S/(X_0, \ldots, X_n) \to 0
$$
应用 Constructions, Lemma \ref{constructions-lemma-proj-sheaves} 的正合函子
$\tilde{\ }$，并利用最后一项属于此函子的核，得到正合复形
$$
0 \to \mathcal{O}_X(-n - 1) \to \ldots
\to \mathcal{O}_X(-n - 1 + i)^{\oplus {n + 1 \choose i}} \to
\ldots \to \mathcal{O}_X \to 0
$$
再用可逆层 $\mathcal{O}_X(n + a + 1)$ 扭曲，即得引理中的正合复形。
\end{proof}

\begin{lemma}
\label{perfect-lemma-generator-P1}
设 $A$ 为环，$X = \mathbf{P}^n_A$。则
$$
E =
\mathcal{O}_X \oplus \mathcal{O}_X(-1) \oplus \ldots \oplus \mathcal{O}_X(-n)
$$
是 $D_\QCoh(X)$ 的一个生成元
（Derived Categories, Definition \ref{derived-definition-generators}）。
\end{lemma}

\begin{proof}
设 $K \in D_\QCoh(\mathcal{O}_X)$。假定对所有 $p \in \mathbf{Z}$ 都有
$\Hom(E, K[p]) = 0$。我们须证明 $K = 0$。
由 Derived Categories, Lemma
\ref{derived-lemma-right-orthogonal}
可知，对所有 $E' \in \langle E \rangle$ 和 $p \in \mathbf{Z}$，
$\Hom(E', K[p])$ 均为零。将 Lemma \ref{perfect-lemma-construct-the-next-one}
用于 $a = -n - 1$，可知
$\mathcal{O}_X(-n - 1) \in \langle E \rangle$，因为它拟同构于一个有限复形，
其各项都是 $E$ 的直和项的有限直和。对 $a = -n - 2$ 重复这一论证，得到
$\mathcal{O}_X(-n - 2) \in \langle E \rangle$。归纳可得，对所有
$m \geq 0$ 都有 $\mathcal{O}_X(-m) \in \langle E \rangle$。由于
$$
\Hom(\mathcal{O}_X(-m), K[p]) =
H^p(X, K \otimes_{\mathcal{O}_X}^\mathbf{L} \mathcal{O}_X(m)) =
H^p(X, K \otimes_{\mathcal{O}_X}^\mathbf{L} \mathcal{O}_X(1)^{\otimes m})
$$
由 Lemma \ref{perfect-lemma-nonzero-some-cohomology} 可知 $K = 0$。（这里还用到了
$\mathcal{O}_X(1)$ 是 $X$ 上的充足可逆层；这由
Properties, Lemma \ref{properties-lemma-open-in-proj-ample} 推出。）
\end{proof}

\begin{remark}
\label{perfect-remark-pullback-generator}
设 $f : X \to Y$ 为拟紧拟分离概形之间的态射，并设
$E \in D_\QCoh(\mathcal{O}_Y)$ 为一个生成元
（见 Theorem \ref{perfect-theorem-bondal-van-den-Bergh}）。则下列条件等价：
\begin{enumerate}
\item 对 $K \in D_\QCoh(\mathcal{O}_X)$，当且仅当 $K = 0$ 时
$Rf_*K = 0$；
\item $Rf_* : D_\QCoh(\mathcal{O}_X) \to D_\QCoh(\mathcal{O}_Y)$
反映同构；
\item $Lf^*E$ 是 $D_\QCoh(\mathcal{O}_X)$ 的一个生成元。
\end{enumerate}
(1) 与 (2) 的等价性是如下事实的形式推论：
$Rf_* : D_\QCoh(\mathcal{O}_X) \to D_\QCoh(\mathcal{O}_Y)$ 是三角范畴的正合函子。
类似地，(1) 与 (3) 的等价性由 $Lf^*$ 是 $Rf_*$ 的左伴随这一事实形式地推出。
当 $f$ 为仿射态射（Lemma \ref{perfect-lemma-affine-morphism}），或 $f$ 为开浸入，
或 $f$ 为这类态射的复合时，
这些条件成立。因此：
\begin{enumerate}
\item 若 $X$ 为拟仿射概形，则 $\mathcal{O}_X$ 是
$D_\QCoh(\mathcal{O}_X)$ 的一个生成元；
\item 若 $X \subset \mathbf{P}^n_A$ 为拟紧局部闭子概形，则由
Lemma \ref{perfect-lemma-generator-P1}，
$\mathcal{O}_X \oplus \mathcal{O}_X(-1) \oplus \ldots \oplus \mathcal{O}_X(-n)$
是 $D_\QCoh(\mathcal{O}_X)$ 的一个生成元。
\end{enumerate}
\end{remark}




\section{紧对象与完美对象}
\label{perfect-section-compact}

\noindent
设 $X$ 为有限维 Noether 概形。由
Cohomology, Proposition \ref{cohomology-proposition-vanishing-Noetherian}
及
Cohomology on Sites, Lemma \ref{sites-cohomology-lemma-when-jshriek-compact}
可知，对所有开集 $U \subset X$，模层 $j_!\mathcal{O}_U$ 都是
$D(\mathcal{O}_X)$ 的紧对象。这些层通常不是拟凝聚的，因而并不给出导出范畴
$D(\mathcal{O}_X)$ 的完美对象。不过，若限制于具有拟凝聚上同调层的复形，
这种情况便不会发生。精确陈述如下。

\begin{proposition}
\label{perfect-proposition-compact-is-perfect}
设 $X$ 为拟紧拟分离概形。$D_\QCoh(\mathcal{O}_X)$ 的一个对象为紧对象，
当且仅当它是完美的。
\end{proposition}

\begin{proof}
若 $K$ 是 $D(\mathcal{O}_X)$ 的完美对象，其对偶为
$K^\vee$（Cohomology, Lemma \ref{cohomology-lemma-dual-perfect-complex}），则
$$
\Hom_{D(\mathcal{O}_X)}(K, M) =
H^0(X, K^\vee \otimes_{\mathcal{O}_X}^\mathbf{L} M)
$$
对 $M$ 函子性地成立。由于 $K^\vee \otimes_{\mathcal{O}_X}^\mathbf{L} -$
与直和可交换，且由 Lemma \ref{perfect-lemma-quasi-coherence-pushforward-direct-sums}，
$H^0(X, -)$ 在 $D_\QCoh(\mathcal{O}_X)$ 上与直和可交换，故 $K$ 是
$D_\QCoh(\mathcal{O}_X)$ 中的紧对象。

\medskip\noindent
反之，设 $K$ 为 $D_\QCoh(\mathcal{O}_X)$ 的紧对象。为证明 $K$ 完美，
只需证明对每个仿射开集 $U \subset X$，$K|_U$ 都完美，见
Cohomology, Lemma \ref{cohomology-lemma-perfect-independent-representative}。
注意 $j : U \to X$ 是拟紧分离态射。因此
$Rj_* : D_\QCoh(\mathcal{O}_U) \to D_\QCoh(\mathcal{O}_X)$
与直和可交换，见 Lemma \ref{perfect-lemma-quasi-coherence-pushforward-direct-sums}。
所以限制到 $U$ 的函子与 $Rj_*$ 的伴随性蕴含 $K|_U$ 是
$D_\QCoh(\mathcal{O}_U)$ 的紧对象。于是归结为 $X$ 仿射的情形。

\medskip\noindent
设 $X = \Spec(A)$ 为仿射概形。由 Lemma \ref{perfect-lemma-affine-compare-bounded}，
问题转化为 $D(A)$ 中的同一问题。对 $D(A)$，结论即
More on Algebra, Proposition \ref{more-algebra-proposition-perfect-is-compact}。
\end{proof}

\begin{remark}
\label{perfect-remark-classical-generator}
设 $X$ 为拟紧拟分离概形，且设 $G$ 为 $D(\mathcal{O}_X)$ 的一个完美对象，
并为 $D_\QCoh(\mathcal{O}_X)$ 的生成元。由
Theorem \ref{perfect-theorem-bondal-van-den-Bergh}，这样的对象至少存在一个。
结合 Lemma \ref{perfect-lemma-quasi-coherence-direct-sums}、
Proposition \ref{perfect-proposition-compact-is-perfect} 与
Derived Categories, Proposition
\ref{derived-proposition-generator-versus-classical-generator}
可知 $G$ 是 $D_{perf}(\mathcal{O}_X)$ 的经典生成元。
\end{remark}

\noindent
下述结果加强了 Proposition \ref{perfect-proposition-compact-is-perfect}。
设 $T \subset X$ 为概形 $X$ 的闭子集。仍以 $D_T(\mathcal{O}_X)$ 表示
$D(\mathcal{O}_X)$ 中由支集于 $T$ 上的对象组成的严格满、饱和三角子范畴
（Definition \ref{perfect-definition-supported-on}）。由于取直和与取上同调层可交换，
$D_T(\mathcal{O}_X)$ 具有直和，且这些直和等于 $D(\mathcal{O}_X)$ 中的直和。

\begin{lemma}
\label{perfect-lemma-compact-is-perfect-with-support}
设 $X$ 为拟紧拟分离概形，$T \subset X$ 为闭子集，且 $X \setminus T$
拟紧。$D_{\QCoh, T}(\mathcal{O}_X)$ 的一个对象为紧对象，当且仅当它作为
$D(\mathcal{O}_X)$ 的对象是完美的。
\end{lemma}

\begin{proof}
由引理前的评注，$D_{\QCoh, T}(\mathcal{O}_X)$ 是具有直和的三角范畴。
由 Proposition \ref{perfect-proposition-compact-is-perfect}，完美对象给出
$D(\mathcal{O}_X)$ 的紧对象，因而更是任一在取直和下保持的子范畴中的紧对象。
反之，我们将用到存在一个生成元 $E \in D_{\QCoh, T}(\mathcal{O}_X)$，
它是 $\mathcal{O}_X$-模的完美复形，见
Lemma \ref{perfect-lemma-generator-with-support}。因此由上所述，$E$ 是紧对象。再由
Derived Categories, Proposition
\ref{derived-proposition-generator-versus-classical-generator}
可知，$E$ 是 $D_{\QCoh, T}(\mathcal{O}_X)$ 的紧对象所成满子范畴的经典生成元。
所以任一紧对象都可由 $E$ 经过有限次如下运算构造：
(a) 取平移，(b) 取有限直和，(c) 取锥，以及 (d) 取直和项。
每种运算都保持完美性，结论随即成立。
\end{proof}

\begin{remark}
\label{perfect-remark-classical-generator-with-support}
设 $X$ 为拟紧拟分离概形，$T \subset X$ 为闭子集，且 $X \setminus T$ 拟紧。
设 $G$ 为 $D_{\QCoh, T}(\mathcal{O}_X)$ 的完美对象，同时也是
$D_{\QCoh, T}(\mathcal{O}_X)$ 的生成元。由
Lemma \ref{perfect-lemma-generator-with-support}，这样的对象至少存在一个。
结合 $D_{\QCoh, T}(\mathcal{O}_X)$ 具有直和这一事实、
Lemma \ref{perfect-lemma-compact-is-perfect-with-support} 以及
Derived Categories, Proposition
\ref{derived-proposition-generator-versus-classical-generator}
可知 $G$ 是 $D_{perf, T}(\mathcal{O}_X)$ 的经典生成元。
\end{remark}

\noindent
下述引理应用了上一个引理证明中的思想。

\begin{lemma}
\label{perfect-lemma-map-from-pseudo-coherent-to-complex-with-support}
设 $X$ 为拟紧拟分离概形，$T \subset X$ 为闭子集，使得
$U = X \setminus T$ 拟紧。设 $\alpha : P \to E$ 为
$D_\QCoh(\mathcal{O}_X)$ 的一个态射，且满足下列条件之一：
\begin{enumerate}
\item $P$ 完美且 $E$ 支集于 $T$ 上；或
\item $P$ 伪凝聚，$E$ 支集于 $T$ 上，且 $E$ 下有界。
\end{enumerate}
则存在一个 $\mathcal{O}_X$-模的完美复形 $I$ 及一个映射
$I \to \mathcal{O}_X[0]$，使得 $I \otimes^\mathbf{L} P \to E$ 为零，
且 $I|_U \to \mathcal{O}_U[0]$ 为同构。
\end{lemma}

\begin{proof}
令 $\mathcal{D} = D_{\QCoh, T}(\mathcal{O}_X)$。在两种情形中，复形
$K = R\SheafHom(P, E)$ 都是 $\mathcal{D}$ 的对象。关于拟凝聚性，见
Lemma \ref{perfect-lemma-quasi-coherence-internal-hom}。由于 $R\SheafHom$ 的构造与
限制到开集可交换，显然 $K$ 支集于 $T$ 上。映射 $\alpha$ 定义了
$H^0(K) = \Hom_{D(\mathcal{O}_X)}(\mathcal{O}_X[0], K)$ 的一个元素。
因此只需对映射 $\alpha : \mathcal{O}_X[0] \to K$ 证明结论。

\medskip\noindent
取完美生成元 $E \in \mathcal{D}$，见
Lemma \ref{perfect-lemma-generator-with-support}。按
Derived Categories, Lemma \ref{derived-lemma-write-as-colimit}，利用生成元 $E$ 写成
$$
K = \text{hocolim} K_n
$$
。由于函子 $\mathcal{D} \to D(\mathcal{O}_X)$ 与直和可交换，
$K = \text{hocolim} K_n$ 在 $D(\mathcal{O}_X)$ 中也成立。由于
$\mathcal{O}_X$ 是 $D(\mathcal{O}_X)$ 的紧对象，存在 $n$ 及一个诱导出
$\alpha$ 的态射 $\alpha_n : \mathcal{O}_X \to K_n$，见
Derived Categories, Lemma \ref{derived-lemma-commutes-with-countable-sums}。
将 Derived Categories, Lemma \ref{derived-lemma-factor-through}
用于环境范畴 $D(\mathcal{O}_X)$ 中的态射 $\mathcal{O}_X[0] \to K_n$，
可知 $\alpha_n$ 分解为 $\mathcal{O}_X[0] \to Q \to K_n$，其中 $Q$ 是
$\langle E \rangle$ 的对象。因此 $Q$ 是支集于 $T$ 上的完美复形。

\medskip\noindent
选取可区别三角
$$
I \to \mathcal{O}_X[0] \to Q \to I[1]
$$
按构造，$I$ 完美，映射 $I \to \mathcal{O}_X[0]$ 限制到 $U$ 后为同构；
又因 $\alpha$ 经 $Q$ 分解，复合 $I \to K$ 为零。这就证明了引理。
\end{proof}






\section{作为模范畴的导出范畴}
\label{perfect-section-derived-is-dga}

\noindent
本节将从前面的结果推出若干结论。在此之前还需要几个引理。

\begin{lemma}
\label{perfect-lemma-tensor-with-QCoh-complex}
设 $X$ 为概形。设 $K^\bullet$ 为一个 $\mathcal{O}_X$-模复形，
其上同调层拟凝聚。令
$(E, d) = \Hom_{\text{Comp}^{dg}(\mathcal{O}_X)}(K^\bullet, K^\bullet)$
为自同态微分分次代数。则
$$
- \otimes_E^\mathbf{L} K^\bullet :
D(E, \text{d}) \longrightarrow D(\mathcal{O}_X)
$$
这一由
Differential Graded Algebra, Lemma
\ref{dga-lemma-tensor-with-complex-derived}
给出的函子之像包含于 $D_\QCoh(\mathcal{O}_X)$。
\end{lemma}

\begin{proof}
设 $P$ 为具有性质 (P) 的微分分次 $E$-模，并设 $F_\bullet$ 为
Differential Graded Algebra, Section \ref{dga-section-P-resolutions}
中所述的 $P$ 上滤过。则
$$
P \otimes_E K^\bullet = \text{hocolim}\ F_iP \otimes_E K^\bullet
$$
每个 $F_iP$ 都有一个有限滤过，其分次片是 $E[k]$ 的直和。结论容易推出。
\end{proof}

\noindent
下述结果取自 \cite{BvdB}。

\begin{theorem}
\label{perfect-theorem-DQCoh-is-Ddga}
设 $X$ 为拟紧拟分离概形。则存在一个微分分次代数 $(E, \text{d})$，
其非零上同调群 $H^i(E)$ 仅有有限多个，并且
$D_\QCoh(\mathcal{O}_X)$ 等价于 $D(E, \text{d})$。
\end{theorem}

\begin{proof}
设 $K^\bullet$ 为一个 $\mathcal{O}$-模的 K-内射复形，它是完美的并生成
$D_\QCoh(\mathcal{O}_X)$。由 Theorem \ref{perfect-theorem-bondal-van-den-Bergh}
及 K-内射解的存在性，这样的复形确实存在。我们将证明，取
$$
(E, \text{d}) = \Hom_{\text{Comp}^{dg}(\mathcal{O}_X)}(K^\bullet, K^\bullet)
$$
时定理成立；其中 $\text{Comp}^{dg}(\mathcal{O}_X)$ 是 $\mathcal{O}$-模复形的
微分分次范畴。参见
Differential Graded Algebra, Section \ref{dga-section-variant-base-change}。
由于 $K^\bullet$ 是 K-内射的，对所有 $n \in \mathbf{Z}$ 都有
\begin{equation}
\label{perfect-equation-E-is-OK}
H^n(E) = \Ext^n_{D(\mathcal{O}_X)}(K^\bullet, K^\bullet)
\end{equation}
。由 Lemma \ref{perfect-lemma-ext-from-perfect-into-bounded-QCoh}，这些 Ext 中只有
有限多个非零。考虑
$$
- \otimes_E^\mathbf{L} K^\bullet :
D(E, \text{d}) \longrightarrow D(\mathcal{O}_X)
$$
这一由
Differential Graded Algebra, Lemma
\ref{dga-lemma-tensor-with-complex-derived} 给出的函子。由于 $K^\bullet$ 完美，
它定义了 $D(\mathcal{O}_X)$ 的一个紧对象，见
Proposition \ref{perfect-proposition-compact-is-perfect}。结合
(\ref{perfect-equation-E-is-OK})，由
Differential Graded Algebra, Lemmas
\ref{dga-lemma-fully-faithful-in-compact-case} 可知上述函子全忠实。由
Differential Graded Algebra, Lemmas
\ref{dga-lemma-tensor-with-complex-hom-adjoint}，它有右伴随
$$
R\Hom(K^\bullet, - ) : D(\mathcal{O}_X) \longrightarrow D(E, \text{d})
$$
；根据伴随函子的一般理论，此函子是左拟逆。另一方面，由
Lemma \ref{perfect-lemma-tensor-with-QCoh-complex} 得到
$$
- \otimes_E^\mathbf{L} K^\bullet :
D(E, \text{d}) \longrightarrow D_\QCoh(\mathcal{O}_X)
$$
，而由于选取 $K^\bullet$ 为 $D_\QCoh(\mathcal{O}_X)$ 的生成元，
该伴随函子限制到 $D_\QCoh(\mathcal{O}_X)$ 后的核为零。
形式论证即给出所需等价，见
Derived Categories, Lemma
\ref{derived-lemma-fully-faithful-adjoint-kernel-zero}.
\end{proof}

\begin{remark}[带支集的变体]
\label{perfect-remark-DQCoh-is-Ddga-with-support}
设 $X$ 为拟紧拟分离概形。设 $T \subset X$ 为闭子集，且 $X \setminus T$
拟紧。Theorem \ref{perfect-theorem-DQCoh-is-Ddga} 对
$D_{\QCoh, T}(\mathcal{O}_X)$ 的类似结论也成立。这由定理证明中完全相同的论证推出，
其中使用 Lemmas \ref{perfect-lemma-generator-with-support}、
\ref{perfect-lemma-compact-is-perfect-with-support} 以及
Lemma \ref{perfect-lemma-tensor-with-QCoh-complex} 的带支集变体。
若以后需要这一结果，我们将在此给出精确陈述和详细证明。
\end{remark}

\begin{remark}[dga 的唯一性]
\label{perfect-remark-independence-choice}
设 $X$ 为环 $R$ 上的拟紧拟分离概形。由
Theorem \ref{perfect-theorem-DQCoh-is-Ddga} 证明中的构造，存在 $R$ 上的微分分次代数
$(A, \text{d})$，使得三角范畴 $D_\QCoh(X)$ 与 $D(A, \text{d})$
$R$-线性等价。可以问：$(A, \text{d})$ 有多唯一？答案（仅仅）略强于说
$(A, \text{d})$ 在导出等价意义下良定义。具体地，设 $(B, \text{d})$
是第二个这样的对。于是
$$
(A, \text{d}) = \Hom_{\text{Comp}^{dg}(\mathcal{O}_X)}(K^\bullet, K^\bullet)
$$
且
$$
(B, \text{d}) = \Hom_{\text{Comp}^{dg}(\mathcal{O}_X)}(L^\bullet, L^\bullet)
$$
，其中 $K^\bullet$ 与 $L^\bullet$ 是某些 $\mathcal{O}_X$-模的 K-内射复形，
分别对应于 $D_\QCoh(\mathcal{O}_X)$ 的完美生成元。令
$$
\Omega = \Hom_{\text{Comp}^{dg}(\mathcal{O}_X)}(K^\bullet, L^\bullet)
\quad
\Omega' = \Hom_{\text{Comp}^{dg}(\mathcal{O}_X)}(L^\bullet, K^\bullet)
$$
则 $\Omega$ 是微分分次 $B^{opp} \otimes_R A$-模，$\Omega'$ 是微分分次
$A^{opp} \otimes_R B$-模。此外，等价
$$
D(A, \text{d}) \to D_\QCoh(\mathcal{O}_X) \to
D(B, \text{d})
$$
由函子 $- \otimes_A^\mathbf{L} \Omega'$ 给出，拟逆也类似。因此正处于
Differential Graded Algebra, Remark \ref{dga-remark-hochschild-cohomology}
的情形。若以后需要此评注，我们将在这里给出精确陈述和详细证明。
\end{remark}






\section{伪凝聚复形的刻画，I}
\label{perfect-section-pseudo-coherent-hocolim}

\noindent
可以用上述方法把伪凝聚对象刻画为完美对象所给逼近的导出同伦极限。

\begin{lemma}
\label{perfect-lemma-pseudo-coherent-hocolim}
设 $X$ 为拟紧拟分离概形，且设 $K \in D(\mathcal{O}_X)$。下列条件等价：
\begin{enumerate}
\item $K$ 是伪凝聚的；
\item $K = \text{hocolim} K_n$，其中 $K_n$ 完美，且对所有 $n$，
$\tau_{\geq -n}K_n \to \tau_{\geq -n}K$ 均为同构。
\end{enumerate}
\end{lemma}

\begin{proof}
在任意环化空间上，(2) $\Rightarrow$ (1) 都成立。事实上，假定 (2) 成立。
回忆导出范畴的完美对象是伪凝聚的，见
Cohomology, Lemma \ref{cohomology-lemma-perfect}。于是由定义，
$\tau_{\geq -n}K_n$ 是 $(-n + 1)$-伪凝聚的，故
$\tau_{\geq -n}K$ 是 $(-n + 1)$-伪凝聚的，进而 $K$ 是
$(-n + 1)$-伪凝聚的。这对所有 $n$ 都成立，所以 $K$ 伪凝聚，见
Cohomology, Definition \ref{cohomology-definition-pseudo-coherent}。

\medskip\noindent
假定 (1)。首先选取由完美复形 $K_1$ 给出的 $(X, K, -2)$ 的一个逼近
$K_1 \to K$，见
Definitions \ref{perfect-definition-approximation-holds} and
\ref{perfect-definition-approximation} and
Theorem \ref{perfect-theorem-approximation}。归纳假设已经有
$$
K_1 \to K_2 \to \ldots \to K_n \to K
$$
，其中 $K_i$ 完美，且对所有 $1 \leq i \leq n$，
$\tau_{\geq -i}K_i \to \tau_{\geq -i}K$ 都是同构。对完美对象 $K_n$，
按 Lemma \ref{perfect-lemma-ext-from-perfect-into-bounded-QCoh} 选取 $a \leq b$。
选取 $(X, K, \min(a - 1, -n - 1))$ 的一个逼近 $K_{n + 1} \to K$。
选取可区别三角
$$
K_{n + 1} \to K \to C \to K_{n + 1}[1]
$$
。于是 $C \in D_\QCoh(\mathcal{O}_X)$ 且当 $i \geq a$ 时
$H^i(C) = 0$。故由 $a, b$ 的选取，
$\Hom_{D(\mathcal{O}_X)}(K_n, C) = 0$。所以复合
$K_n \to K \to C$ 为零。由
Derived Categories, Lemma \ref{derived-lemma-representable-homological}，
$K_n \to K$ 可经 $K_{n + 1}$ 分解，从而完成归纳步骤。

\medskip\noindent
还须证明 $K = \text{hocolim} K_n$。将
Derived Categories, Lemma \ref{derived-lemma-cohomology-of-hocolim}
用于函子
$H^i( - ) : D(\mathcal{O}_X) \to \textit{Mod}(\mathcal{O}_X)$
和我们所选的 $K_n$，即得此结论。
\end{proof}

\begin{lemma}
\label{perfect-lemma-pseudo-coherent-hocolim-with-support}
设 $X$ 为拟紧拟分离概形。设 $T \subset X$ 为闭子集，且 $X \setminus T$
拟紧。设 $K \in D(\mathcal{O}_X)$ 支集于 $T$ 上。下列条件等价：
\begin{enumerate}
\item $K$ 是伪凝聚的；
\item $K = \text{hocolim} K_n$，其中 $K_n$ 完美且支集于 $T$ 上，并且对所有
$n$，$\tau_{\geq -n}K_n \to \tau_{\geq -n}K$ 均为同构。
\end{enumerate}
\end{lemma}

\begin{proof}
本引理的证明与 Lemma \ref{perfect-lemma-pseudo-coherent-hocolim} 的证明完全相同，
唯一差别是在选取逼近时使用三元组 $(T, K, m)$。
\end{proof}









\section{一个等价的例子}
\label{perfect-section-example-equivalence}

\noindent
在 Section \ref{perfect-section-example-generator} 中，我们证明了环 $A$ 上射影空间
$\mathbf{P}^n_A$ 的导出范畴由一个向量丛生成；事实上，该向量丛是结构层若干平移的直和。
本节证明，这确定了 $D_\QCoh(\mathcal{O}_{\mathbf{P}^n_A})$ 与某个
$A$-代数的导出范畴之间的等价。

\medskip\noindent
在陈述结果之前先引入一些记号。设 $A$ 为环，
$X = \mathbf{P}^n_A = \text{Proj}(S)$，其中 $S = A[X_0, \ldots, X_n]$。
由 Lemma \ref{perfect-lemma-generator-P1} 可知
\begin{equation}
\label{perfect-equation-generator-Pn}
P =
\mathcal{O}_X \oplus \mathcal{O}_X(-1) \oplus \ldots \oplus \mathcal{O}_X(-n)
\end{equation}
是 $D_\QCoh(\mathcal{O}_X)$ 的完美生成元。考虑（非交换）$A$-代数
\begin{equation}
\label{perfect-equation-algebra-for-Pn}
R = \Hom_{\mathcal{O}_X}(P, P) =
\left(
\begin{matrix}
S_0 & S_1 & S_2 & \ldots & \ldots \\
0 & S_0 & S_1 & \ldots & \ldots\\
0 & 0 & S_0 & \ldots & \ldots \\
\ldots & \ldots & \ldots & \ldots & \ldots \\
0 & \ldots & \ldots & \ldots & S_0
\end{matrix}
\right)
\end{equation}
，其乘法和加法是显然的。若按通常方式将 $P$ 视为 $\mathcal{O}_X$-模复形
（即 $P$ 位于次数 $0$，其他各次数均为零），则
$$
R = \Hom_{\text{Comp}^{dg}(\mathcal{O}_X)}(P, P)
$$
，其中右端把 $R$ 视为 $A$ 上微分为零的微分分次代数
（即 $R$ 位于次数 $0$，其他各次数均为零）。根据
Differential Graded Algebra, Section \ref{dga-section-variant-base-change}
中的讨论，得到导出函子
$$
- \otimes_R^\mathbf{L} P : D(R) \longrightarrow D(\mathcal{O}_X),
$$
特别见 Differential Graded Algebra, Lemma
\ref{dga-lemma-tensor-with-complex-derived}。由
Lemma \ref{perfect-lemma-tensor-with-QCoh-complex}，此函子的本质像包含于
$D_\QCoh(\mathcal{O}_X)$。

\begin{lemma}
\label{perfect-lemma-Pn-module-category}
\begin{reference}
\cite{Beilinson}
\end{reference}
设 $A$ 为环，$X = \mathbf{P}^n_A = \text{Proj}(S)$，其中
$S = A[X_0, \ldots, X_n]$。取 (\ref{perfect-equation-generator-Pn}) 中的
$P$ 和 (\ref{perfect-equation-algebra-for-Pn}) 中的 $R$，则函子
$$
- \otimes_R^\mathbf{L} P : D(R) \longrightarrow D_\QCoh(\mathcal{O}_X)
$$
是三角范畴的一个 $A$-线性等价，并把 $R$ 送到 $P$。
\end{lemma}

\noindent
换言之，射影空间上拟凝聚模的导出范畴等价于某个（非交换）代数上模的导出范畴。
在 $A$ 上所有射影概形中，射影空间的这一性质似乎相当特殊。

\begin{proof}
为证明该函子全忠实，只需证明当 $i \not = 0$ 时 $\Ext^i_X(P, P)$ 为零，
当 $i = 0$ 时它等于 $R$，见
Differential Graded Algebra, Lemma
\ref{dga-lemma-fully-faithful-in-compact-case}。与
Lemma \ref{perfect-lemma-ext-from-perfect-into-bounded-QCoh} 的证明相同，有
$$
\Ext^i_X(P, P) = H^i(X, P^\wedge \otimes P) =
\bigoplus\nolimits_{0 \leq a, b \leq n} H^i(X, \mathcal{O}_X(a - b))
$$
。由射影空间上同调的计算（Cohomology of Schemes, Lemma
\ref{coherent-lemma-cohomology-projective-space-over-ring}），可知这些
$\Ext$-群除非 $i = 0$ 否则均为零。当 $i = 0$ 时得到 $R$，因为
(\ref{perfect-equation-algebra-for-Pn}) 正是如此定义 $R$ 的。由
Differential Graded Algebra, Lemma
\ref{dga-lemma-tensor-with-complex-hom-adjoint}
，该函子有右伴随 $R\Hom(P, -) : D_\QCoh(\mathcal{O}_X) \to D(R)$。
由 Lemma \ref{perfect-lemma-generator-P1}，$P$ 是 $D_\QCoh(\mathcal{O}_X)$ 的生成元，
故 $R\Hom(P, -)$ 的核为零。因此由 Derived Categories, Lemma
\ref{derived-lemma-fully-faithful-adjoint-kernel-zero}，该函子是三角范畴的等价。
\end{proof}






\section{再论拟凝聚化子}
\label{perfect-section-better-coherator}

\noindent
在 Section \ref{perfect-section-coherator} 中，我们构造并研究了典范函子
$D(\QCoh(\mathcal{O}_X)) \to D(\mathcal{O}_X)$ 的右伴随 $RQ_X$。
它被构造为拟凝聚化子
$Q_X : \textit{Mod}(\mathcal{O}_X) \to \QCoh(\mathcal{O}_X)$
的右导出扩张。本节研究包含函子
$$
D_\QCoh(\mathcal{O}_X) \longrightarrow D(\mathcal{O}_X)
$$
何时具有右伴随。若此右伴随存在，我们将把它记作
\footnote{这可能不是标准记号。不过，我们已经用 $Q_X$ 表示拟凝聚化子，
用 $RQ_X$ 表示其导出扩张。}
$$
DQ_X :
D(\mathcal{O}_X) \longrightarrow D_\QCoh(\mathcal{O}_X)
$$
。事实表明，拟紧拟分离概形具有这样的右伴随。

\begin{lemma}
\label{perfect-lemma-better-coherator}
设 $X$ 为拟紧拟分离概形。包含函子
$D_\QCoh(\mathcal{O}_X) \to D(\mathcal{O}_X)$ 具有右伴随 $DQ_X$。
\end{lemma}

\begin{proof}[第一种证明]
我们使用
Cohomology of Schemes, Lemma \ref{coherent-lemma-induction-principle}
中的归纳原理来证明。若
$D(\QCoh(\mathcal{O}_X)) \to D_\QCoh(\mathcal{O}_X)$ 是等价，则本引理成立，
因为 Section \ref{perfect-section-coherator} 的函子 $RQ_X$ 是函子
$D(\QCoh(\mathcal{O}_X)) \to D(\mathcal{O}_X)$ 的右伴随。特别地，本引理对仿射概形成立，
见 Lemma \ref{perfect-lemma-affine-coherator}。因此只需证明：若
$X = U \cup V$ 是两个拟紧开集之并，且引理对 $U$、$V$ 和 $U \cap V$
成立，则它对 $X$ 成立。

\medskip\noindent
该伴随存在，当且仅当对 $D(\mathcal{O}_X)$ 的每个对象 $K$，都能找到可区别三角
$$
E' \to E \to K \to E'[1]
$$
位于 $D(\mathcal{O}_X)$ 中，使得 $E'$ 属于 $D_\QCoh(\mathcal{O}_X)$，
且对 $D_\QCoh(\mathcal{O}_X)$ 中所有 $M$ 都有 $\Hom(M, K) = 0$。见
Derived Categories, Lemma \ref{derived-lemma-right-adjoint}。考虑
Cohomology, Lemma \ref{cohomology-lemma-exact-sequence-j-star}
在 $D(\mathcal{O}_X)$ 中给出的可区别三角
$$
E \to Rj_{U, *}E|_U \oplus Rj_{V, *}E|_V \to
Rj_{U \cap V, *}E|_{U \cap V} \to E[1]
$$
。由 Derived Categories, Lemma \ref{derived-lemma-prepare-adjoint}，
只需对 $Rj_{U, *}E|_U$、$Rj_{V, *}E|_V$ 和
$Rj_{U \cap V, *}E|_{U \cap V}$ 构造所需的可区别三角。
这把问题化为下一段所讨论的陈述。

\medskip\noindent
设 $j : U \to X$ 是开浸入，其中 $U$ 是本引理成立的拟紧开集。
设 $L$ 为 $D(\mathcal{O}_U)$ 的对象。则存在 $D(\mathcal{O}_X)$ 中的可区别三角
$$
E' \to Rj_*L \to K \to E'[1]
$$
，使得 $E'$ 属于 $D_\QCoh(\mathcal{O}_X)$，且对
$D_\QCoh(\mathcal{O}_X)$ 中所有 $M$ 都有 $\Hom(M, K) = 0$。
为此，选取 $D(\mathcal{O}_U)$ 中的可区别三角
$$
L' \to L \to Q \to L'[1]
$$
，使得 $L'$ 属于 $D_\QCoh(\mathcal{O}_U)$，且对
$D_\QCoh(\mathcal{O}_U)$ 中所有 $N$ 都有 $\Hom(N, Q) = 0$。
这是可能的，因为 Derived Categories, Lemma \ref{derived-lemma-right-adjoint}
中的陈述是充要条件。我们得到 $D(\mathcal{O}_X)$ 中的可区别三角
$$
Rj_*L' \to Rj_*L \to Rj_*Q \to Rj_*L'[1]
$$
。由 Lemma \ref{perfect-lemma-quasi-coherence-direct-image}，
$Rj_*L'$ 属于 $D_\QCoh(\mathcal{O}_X)$。另一方面，若 $M$ 属于
$D_\QCoh(\mathcal{O}_X)$，则
$$
\Hom(M, Rj_*Q) = \Hom(Lj^*M, Q) = 0
$$
，因为由 Lemma \ref{perfect-lemma-quasi-coherence-pullback}，
$Lj^*M$ 属于 $D_\QCoh(\mathcal{O}_U)$。证明完毕。
\end{proof}

\begin{proof}[第二种证明]
由 Derived Categories, Proposition \ref{derived-proposition-brown}，
该伴随存在。其假设均满足：首先，$D_\QCoh(\mathcal{O}_X)$ 具有直和，
且直和与包含函子可交换（Lemma \ref{perfect-lemma-quasi-coherence-direct-sums}）。
另一方面，$D_\QCoh(\mathcal{O}_X)$ 是紧生成的，因为由
Theorem \ref{perfect-theorem-bondal-van-den-Bergh} 它具有完美生成元，
而由 Proposition \ref{perfect-proposition-compact-is-perfect}，完美对象是紧的。
\end{proof}

\begin{lemma}
\label{perfect-lemma-pushforward-better-coherator}
设 $f : X \to Y$ 为概形的拟紧拟分离态射。若对 $X$ 和 $Y$，包含函子
$D_\QCoh \to D$ 的右伴随 $DQ_X$ 与 $DQ_Y$ 存在，则
$$
Rf_* \circ DQ_X = DQ_Y \circ Rf_*
$$
\end{lemma}

\begin{proof}
由 Lemma \ref{perfect-lemma-quasi-coherence-direct-image}，$Rf_*$ 把
$D_\QCoh(\mathcal{O}_X)$ 送入 $D_\QCoh(\mathcal{O}_Y)$，故该陈述有意义。
它之所以成立，是因为 $Lf^*$ 同样把 $D_\QCoh(\mathcal{O}_Y)$ 映入
$D_\QCoh(\mathcal{O}_X)$（Lemma \ref{perfect-lemma-quasi-coherence-pullback}），
因而 $Rf_* \circ DQ_X$ 与 $DQ_Y \circ Rf_*$ 都是
$Lf^* : D_\QCoh(\mathcal{O}_Y) \to D(\mathcal{O}_X)$ 的右伴随。
\end{proof}

\begin{remark}
\label{perfect-remark-explain-consequence}
设 $X$ 为拟紧拟分离概形，并设 $X = U \cup V$，其中 $U$ 与 $V$ 为拟紧开集。
由 Lemma \ref{perfect-lemma-better-coherator}，函子
$DQ_X$、$DQ_U$、$DQ_V$、$DQ_{U \cap V}$ 均存在。此外，对任意
$K \in D(\mathcal{O}_X)$，存在典范可区别三角
$$
DQ_X(K) \to Rj_{U, *}DQ_U(K|_U) \oplus Rj_{V, *}DQ_V(K|_V)
\to Rj_{U \cap V, *}DQ_{U \cap V}(K|_{U \cap V}) \to
$$
。这是将正合函子 $DQ_X$ 作用于
Cohomology, Lemma \ref{cohomology-lemma-exact-sequence-j-star}
的可区别三角，并三次使用 Lemma \ref{perfect-lemma-pushforward-better-coherator}
所得的结果。
\end{remark}

\begin{lemma}
\label{perfect-lemma-boundedness-better-coherator}
设 $X$ 为拟紧拟分离概形。Lemma \ref{perfect-lemma-better-coherator} 的函子
$DQ_X$ 具有如下有界性：存在整数 $N = N(X)$，使得若
$K$ 属于 $D(\mathcal{O}_X)$，并且对 $X$ 中的仿射开集 $U$ 及
$i \not \in [a, b]$ 有 $H^i(U, K) = 0$，则当
$i \not \in [a, b + N]$ 时，上同调层 $H^i(DQ_X(K))$ 为零。
\end{lemma}

\begin{proof}
我们用 Cohomology of Schemes, Lemma
\ref{coherent-lemma-induction-principle} 的归纳原理证明。

\medskip\noindent
若 $X$ 仿射，则取 $N = 0$ 时引理成立，因为此时 $RQ_X = DQ_X$ 是取
$R\Gamma(X, K)$ 所对应的拟凝聚层复形。见
Lemmas \ref{perfect-lemma-affine-compare-bounded} and \ref{perfect-lemma-affine-coherator}。

\medskip\noindent
假设 $U, V$ 是 $X$ 中的拟紧开集，且引理对 $U$、$V$ 及 $U \cap V$ 成立，
相应整数分别记为 $N(U)$、$N(V)$ 与 $N(U \cap V)$。现设 $K$ 属于
$D(\mathcal{O}_X)$，且对所有仿射开集 $W \subset X$ 及所有
$i \not \in [a, b]$ 都有 $H^i(W, K) = 0$。则
$K|_U$、$K|_V$、$K|_{U \cap V}$ 具有同一性质。因此
$RQ_U(K|_U)$、$RQ_V(K|_V)$ 与 $RQ_{U \cap V}(K|_{U \cap V})$
在区间 $[a, b + \max(N(U), N(V), N(U \cap V))$ 外的上同调层均为零。
由 Lemma \ref{perfect-lemma-quasi-coherence-direct-image}，函子
$Rj_{U, *}$、$Rj_{V, *}$、$Rj_{U \cap V, *}$ 在 $D_\QCoh$ 上具有有限上同调维数，
故存在 $N$，使得 $Rj_{U, *}DQ_U(K|_U)$、$Rj_{V, *}DQ_V(K|_V)$ 和
$Rj_{U \cap V, *}DQ_{U \cap V}(K|_{U \cap V})$ 在区间
$[a, b + N]$ 外的上同调层均为零。最后由
Remark \ref{perfect-remark-explain-consequence} 的可区别三角得出结论。
\end{proof}

\begin{example}
\label{perfect-example-inverse-limit-quasi-coherent}
设 $X$ 为拟紧拟分离概形，并设 $(\mathcal{F}_n)$ 为拟凝聚层的逆系统。
由于 $DQ_X$ 是右伴随，它与积可交换，因而也与导出极限可交换。因此
$$
DQ_X(R\lim \mathcal{F}_n) =
(R\lim\text{ 于 }D_\QCoh(\mathcal{O}_X))(\mathcal{F}_n)
$$
，其中第一个 $R\lim$ 在 $D(\mathcal{O}_X)$ 中取。事实上，把它记作
$K = R\lim \mathcal{F}_n$。对任意仿射开集 $U \subset X$，有
$$
H^i(U, K) =
H^i(R\Gamma(U, R\lim \mathcal{F}_n)) =
H^i(R\lim R\Gamma(U, \mathcal{F}_n)) =
H^i(R\lim \Gamma(U, \mathcal{F}_n))
$$
，因为上同调与导出极限可交换，且拟凝聚层 $\mathcal{F}_n$ 在仿射概形上没有高阶上同调。
由阿贝尔群范畴中 $R\lim$ 的计算，除非 $i \in [0, 1]$，否则
$H^i(U, K) = 0$。最后由 Lemma \ref{perfect-lemma-boundedness-better-coherator}
可知，$D_\QCoh(\mathcal{O}_X)$ 中的 $R\lim$，即上面的 $DQ_X(K)$，属于
$D^b_\QCoh(\mathcal{O}_X)$。
\end{example}













\section{上同调与基变换，IV}
\label{perfect-section-cohomology-and-base-change-perfect}

\noindent
本节继续
Cohomology of Schemes, Section
\ref{coherent-section-cohomology-and-base-change-perfect} 中的讨论。
首先，对概形的拟紧拟分离态射以及上同调层拟凝聚的复形，
我们有一个非常一般的投影公式版本。

\begin{lemma}
\label{perfect-lemma-cohomology-base-change}
设 $f : X \to Y$ 为概形的拟紧拟分离态射。对
$D_\QCoh(\mathcal{O}_X)$ 中的 $E$ 和 $D_\QCoh(\mathcal{O}_Y)$ 中的 $K$，
在 Cohomology, Equation
(\ref{cohomology-equation-projection-formula-map}) 中定义的映射
$$
Rf_*(E) \otimes_{\mathcal{O}_Y}^\mathbf{L} K
\longrightarrow
Rf_*(E \otimes_{\mathcal{O}_X}^\mathbf{L} Lf^*K)
$$
是同构。
\end{lemma}

\begin{proof}
为检验该映射是同构，可以在 $Y$ 上局部地工作。因此归结为 $Y$ 仿射的情形。

\medskip\noindent
设 $K = \bigoplus K_i$ 是若干复形 $K_i \in D_\QCoh(\mathcal{O}_Y)$ 的直和。
若陈述对每个 $K_i$ 成立，则也对 $K$ 成立。事实上，函子 $Lf^*$ 与
$\otimes^\mathbf{L}$ 按构造保持直和，而由
Lemma \ref{perfect-lemma-quasi-coherence-pushforward-direct-sums}，$Rf_*$
（对上同调层拟凝聚的复形）与直和可交换。此外，设
$K \to L \to M \to K[1]$ 是 $D_\QCoh(Y)$ 中的可区别三角。
若引理的陈述对 $K, L, M$ 中两个成立，则对第三个也成立
（因为所涉及的函子是三角范畴的正合函子）。

\medskip\noindent
假定 $Y$ 仿射，记 $Y = \Spec(A)$。函子
$\widetilde{\ } : D(A) \to D_\QCoh(\mathcal{O}_Y)$ 是等价
（Lemma \ref{perfect-lemma-affine-compare-bounded}）。令 $T$ 表示 $K \in D(A)$
所满足的如下性质：引理的陈述对 $\widetilde{K}$ 成立。由上述讨论及
More on Algebra, Remark \ref{more-algebra-remark-P-resolution}，
只需证明 $T$ 对 $A[k]$ 成立。由于引理的陈述对结构层的平移显然成立，证明完成。
\end{proof}

\begin{definition}
\label{perfect-definition-tor-independent}
设 $S$ 为概形，$X$、$Y$ 为 $S$ 上概形。若对每个映到同一点 $s \in S$ 的
$x \in X$ 与 $y \in Y$，环 $\mathcal{O}_{X, x}$ 和 $\mathcal{O}_{Y, y}$
在 $\mathcal{O}_{S, s}$ 上 Tor 独立（见 More on Algebra, Definition
\ref{more-algebra-definition-tor-independent}），则称 $X$ 与 $Y$
{\it 在 $S$ 上 Tor 独立}。
\end{definition}

\begin{lemma}
\label{perfect-lemma-tor-independent}
设 $f : X \to S$ 与 $g : Y \to S$ 为概形态射。下列条件等价：
\begin{enumerate}
\item $X$ 与 $Y$ 在 $S$ 上 Tor 独立；
\item 对每组仿射开集 $U \subset X$、$V \subset Y$、$W \subset S$，若
$f(U) \subset W$ 且 $g(V) \subset W$，则环 $\mathcal{O}_X(U)$ 与
$\mathcal{O}_Y(V)$ 在 $\mathcal{O}_S(W)$ 上 Tor 独立；
\item 存在仿射开覆盖 $S = \bigcup W_i$，且对每个 $i$ 存在仿射开覆盖
$f^{-1}(W_i) = \bigcup U_{ij}$ 与 $g^{-1}(W_i) = \bigcup V_{ik}$，
使得对所有 $i, j, k$，环 $\mathcal{O}_X(U_{ij})$ 与
$\mathcal{O}_Y(V_{ik})$ 在 $\mathcal{O}_S(W_i)$ 上 Tor 独立。
\end{enumerate}
\end{lemma}

\begin{proof}
略。提示：使用 More on Algebra, Lemma
\ref{more-algebra-lemma-tor-independent}。
\end{proof}

\begin{lemma}
\label{perfect-lemma-flat-base-change-tor-independent}
设 $X \to S$ 与 $Y \to S$ 为概形态射。设 $S' \to S$ 为概形态射，并记
$X' = X \times_S S'$、$Y' = Y \times_S S'$。若 $X$ 与 $Y$
在 $S$ 上 Tor 独立，且 $S' \to S$ 平坦，则 $X'$ 与 $Y'$ 在 $S'$ 上 Tor 独立。
\end{lemma}

\begin{proof}
略。提示：使用 Lemma \ref{perfect-lemma-tor-independent}，并在仿射开集上使用
More on Algebra, Lemma
\ref{more-algebra-lemma-flat-base-change-tor-independent}。
\end{proof}

\begin{lemma}
\label{perfect-lemma-compare-base-change}
设 $g : S' \to S$ 为概形态射，且设 $f : X \to S$ 拟紧拟分离。
考虑基变换图
$$
\xymatrix{
X' \ar[r]_{g'} \ar[d]_{f'} &
X \ar[d]^f \\
S' \ar[r]^g &
S
}
$$
若 $X$ 与 $S'$ 在 $S$ 上 Tor 独立，则对所有
$E \in D_\QCoh(\mathcal{O}_X)$，典范箭头
$Lg^*Rf_*E \to Rf'_*L(g')^*E$ 均为同构。
\end{lemma}

\begin{proof}
对 $D(\mathcal{O}_X)$ 的任意对象 $E$，由
Cohomology, Remark \ref{cohomology-remark-base-change} 得到典范基变换映射
$Lg^*Rf_*E \to Rf'_*L(g')^*E$。为检验它是同构，可以在 $S'$ 上局部地工作。
因此可假设 $g : S' \to S$ 是仿射概形之间的态射。特别地，$g$ 仿射，
只需证明
$$
Rg_*Lg^*Rf_*E \to Rg_*Rf'_*L(g')^*E = Rf_*(Rg'_* L(g')^* E)
$$
是同构，见 Lemma \ref{perfect-lemma-affine-morphism}（并使用 Lemmas
\ref{perfect-lemma-quasi-coherence-pullback}、
\ref{perfect-lemma-quasi-coherence-tensor-product}, and
\ref{perfect-lemma-quasi-coherence-direct-image}
可知对象 $Rf'_*L(g')^*E$ 与 $Lg^*Rf_*E$ 的上同调层拟凝聚）。注意 $g'$
也仿射（Morphisms, Lemma \ref{morphisms-lemma-base-change-affine}）。
由 Lemma \ref{perfect-lemma-affine-morphism-pull-push}，该映射成为
$$
Rf_*E \otimes_{\mathcal{O}_S}^\mathbf{L} g_*\mathcal{O}_{S'}
\longrightarrow
Rf_*(E \otimes_{\mathcal{O}_X}^\mathbf{L} g'_*\mathcal{O}_{X'})
$$
。注意 $g'_*\mathcal{O}_{X'} = f^*g_*\mathcal{O}_{S'}$（由仿射基变换，见
Cohomology of Schemes, Lemma \ref{coherent-lemma-affine-base-change}）。
因此由 Lemma \ref{perfect-lemma-cohomology-base-change}，只需证明
$Lf^*g_*\mathcal{O}_{S'} = f^*g_*\mathcal{O}_{S'}$。这由 $X$ 与 $S'$
在 $S$ 上 Tor 独立的假设推出。事实上，为检验这一点，可以在 $X$ 上局部地工作，
所以还可假设 $X$ 仿射。记 $X = \Spec(A)$、$S = \Spec(R)$ 且
$S' = \Spec(R')$。我们的假设蕴含 $A$ 与 $R'$ 在 $R$ 上 Tor 独立
（More on Algebra, Lemma \ref{more-algebra-lemma-tor-independent}），即当
$i > 0$ 时 $\text{Tor}_i^R(A, R') = 0$。换言之，
$A \otimes_R^\mathbf{L} R' = A \otimes_R R'$，而这恰好意味着
$Lf^*g_*\mathcal{O}_{S'} = f^*g_*\mathcal{O}_{S'}$
（使用 Lemma \ref{perfect-lemma-quasi-coherence-pullback}）。
\end{proof}

\noindent
下述引理将在对偶化复形一章中使用。

\begin{lemma}
\label{perfect-lemma-affine-morphism-and-hom-out-of-perfect}
考虑拟紧拟分离概形的笛卡尔方块
$$
\xymatrix{
X' \ar[r]_{g'} \ar[d]_{f'} & X \ar[d]^f \\
S' \ar[r]^g & S
}
$$
。假设 $g$ 与 $f$ Tor 独立，且 $S = \Spec(R)$、$S' = \Spec(R')$ 仿射。
对 $M, K \in D(\mathcal{O}_X)$，典范映射
$$
R\Hom_X(M, K) \otimes^\mathbf{L}_R R'
\longrightarrow
R\Hom_{X'}(L(g')^*M, L(g')^*K)
$$
在下列两种情形中是 $D(R')$ 中的同构：
\begin{enumerate}
\item $M \in D(\mathcal{O}_X)$ 完美，且 $K \in D_\QCoh(X)$；或
\item $M \in D(\mathcal{O}_X)$ 伪凝聚，$K \in D_\QCoh^+(X)$，且
$R'$ 在 $R$ 上具有有限 Tor 维数。
\end{enumerate}
\end{lemma}

\begin{proof}
由 Cohomology, Section \ref{cohomology-section-global-RHom}，
在整体 Hom 复形的 $D(\Gamma(X, \mathcal{O}_X))$ 中有典范映射
$R\Hom_X(M, K) \to R\Hom_{X'}(L(g')^*M, L(g')^*K)$
。限制标量后可将其视为 $D(R)$ 中的映射。再利用限制标量与
$- \otimes_R^\mathbf{L} R'$ 的伴随性，便得到引理中所显示的映射。
定义该映射后，只需证明它在阿贝尔群导出范畴中是同构。

\medskip\noindent
由 Lemma \ref{perfect-lemma-affine-morphism-pull-push}，右端等于
$$
R\Hom_X(M, R(g')_*L(g')^*K) =
R\Hom_X(M, K \otimes_{\mathcal{O}_X}^\mathbf{L} g'_*\mathcal{O}_{X'})
$$
。在两种情形中，由 Lemma \ref{perfect-lemma-quasi-coherence-internal-hom}，复形
$R\SheafHom(M, K)$ 都是 $D_\QCoh(\mathcal{O}_X)$ 的对象。有自然映射
$$
R\SheafHom(M, K) \otimes_{\mathcal{O}_X}^\mathbf{L} g'_*\mathcal{O}_{X'}
\longrightarrow
R\SheafHom(M, K \otimes_{\mathcal{O}_X}^\mathbf{L} g'_*\mathcal{O}_{X'})
$$
，由 Lemma \ref{perfect-lemma-internal-hom-evaluate-tensor-isomorphism}，
它在两种情形中都是同构。为说明该引理适用于情形 (2)，注意
$g'_*\mathcal{O}_{X'} = Rg'_*\mathcal{O}_{X'} =
Lf^*g_*\mathcal{O}_{S'}$，其中第二个等式来自
Lemma \ref{perfect-lemma-compare-base-change}。使用
Lemma \ref{perfect-lemma-tor-dimension-affine} 及
Cohomology, Lemma \ref{cohomology-lemma-tor-amplitude-pullback}
可知 $g'_*\mathcal{O}_{X'}$ 具有有限 Tor 维数。因此在两种情形中，
以 $R\SheafHom(M, K)$ 替换 $K$，问题都归结为证明
$$
R\Gamma(X, K) \otimes^\mathbf{L}_A A' \longrightarrow
R\Gamma(X, K \otimes^\mathbf{L}_{\mathcal{O}_X} g'_*\mathcal{O}_{X'})
$$
是同构。注意由 Lemma \ref{perfect-lemma-affine-morphism-pull-push}，左端等于
$R\Gamma(X', L(g')^*K)$。因此结论由 Lemma \ref{perfect-lemma-compare-base-change} 推出。
\end{proof}

\begin{remark}
\label{perfect-remark-multiplication-map}
采用 Lemma \ref{perfect-lemma-affine-morphism-and-hom-out-of-perfect} 的记号。图
$$
\xymatrix{
R\Hom_X(M, Rg'_*L) \otimes_R^\mathbf{L} R' \ar[r] \ar[d]_\mu &
R\Hom_{X'}(L(g')^*M, L(g')^*Rg'_*L) \ar[d]^a \\
R\Hom_X(M, R(g')_*L) \ar@{=}[r] &
R\Hom_{X'}(L(g')^*M, L)
}
$$
可交换，其中上方水平箭头是引理中的映射，$\mu$ 是乘法映射，而 $a$ 来自伴随映射
$L(g')^*Rg'_*L \to L$。乘法映射是任意 $K' \in D(R')$ 的伴随映射
$K' \otimes_R^\mathbf{L} R' \to K'$。
\end{remark}

\begin{lemma}
\label{perfect-lemma-tor-independence-and-tor-amplitude}
考虑概形的笛卡尔方块
$$
\xymatrix{
X' \ar[r]_{g'} \ar[d]_{f'} & X \ar[d]^f \\
S' \ar[r]^g & S
}
$$
。假设 $g$ 与 $f$ Tor 独立。
\begin{enumerate}
\item 若 $E \in D(\mathcal{O}_X)$ 作为 $f^{-1}\mathcal{O}_S$-模复形具有
$[a, b]$ 中的 Tor 振幅，则 $L(g')^*E$ 作为
$f^{-1}\mathcal{O}_{S'}$-模复形具有 $[a, b]$ 中的 Tor 振幅。
\item 若 $\mathcal{G}$ 是在 $S$ 上平坦的 $\mathcal{O}_X$-模，则
$L(g')^*\mathcal{G} = (g')^*\mathcal{G}$。
\end{enumerate}
\end{lemma}

\begin{proof}
可以在茎上计算 Tor 维数，见
Cohomology, Lemma \ref{cohomology-lemma-tor-amplitude-stalk}。
若 $x' \in X'$ 的像为 $x \in X$，则
$$
(L(g')^*E)_{x'} =
E_x \otimes_{\mathcal{O}_{X, x}}^\mathbf{L} \mathcal{O}_{X', x'}
$$
。设 $s' \in S'$ 与 $s \in S$ 分别为 $x'$ 与 $x$ 的像。由于 $X$ 与
$S'$ 在 $S$ 上 Tor 独立，可应用
More on Algebra, Lemma \ref{more-algebra-lemma-base-change-comparison}
，可知所显示公式的右端在 $D(\mathcal{O}_{S', s'})$ 中等于
$E_x \otimes_{\mathcal{O}_{S, s}}^\mathbf{L} \mathcal{O}_{S', s'}$。
于是 (1) 由 More on Algebra, Lemma
\ref{more-algebra-lemma-pull-tor-amplitude} 推出。为证明 (2)，注意
$\mathcal{G}$ 平坦等价于 $\mathcal{G}[0]$ 具有 $[0, 0]$ 中的 Tor 振幅。
应用 (1) 即得结论。
\end{proof}

\begin{lemma}
\label{perfect-lemma-compare-base-change-closed-immersion}
考虑概形的笛卡尔图
$$
\xymatrix{
Z' \ar[r]_{i'} \ar[d]_g & X' \ar[d]^f \\
Z \ar[r]^i & X
}
$$
，其中 $i$ 是闭浸入。若 $Z$ 与 $X'$ 在 $X$ 上 Tor 独立，则作为函子
$D(\mathcal{O}_Z) \to D(\mathcal{O}_{X'})$，有
$Ri'_* \circ Lg^* = Lf^* \circ Ri_*$。
\end{lemma}

\begin{proof}
注意本引理须对所有 $K \in D(\mathcal{O}_Z)$ 成立。函子 $i_*$ 与 $i'_*$
都是正合的，故 $Ri_*$ 与 $Ri'_*$ 可通过把 $i_*$ 与 $i'_*$ 作用于任意代表来计算。
因此，基变换映射
$$
Lf^*(Ri_*(K)) \longrightarrow Ri'_*(Lg^*(K))
$$
在点 $z' \in Z'$ 的茎上由下述映射给出，其中 $z \in Z$ 为其像：
$$
K_z \otimes_{\mathcal{O}_{X, z}}^\mathbf{L} \mathcal{O}_{X', z'}
\longrightarrow
K_z \otimes_{\mathcal{O}_{Z, z}}^\mathbf{L} \mathcal{O}_{Z', z'}
$$
由 More on Algebra, Lemma \ref{more-algebra-lemma-base-change-comparison}
及所假定的 Tor 独立性，该映射是同构。
\end{proof}








\section{K\"unneth 公式，II}
\label{perfect-section-kunneth}

\noindent
关于基底为域的情形，见 Varieties, Section
\ref{varieties-section-kunneth}。考虑概形的笛卡尔图
$$
\xymatrix{
& X \times_S Y \ar[ld]^p \ar[rd]_q \ar[dd]^f \\
X \ar[rd]_a & & Y \ar[ld]^b \\
& S
}
$$
设 $K \in D(\mathcal{O}_X)$ 且 $M \in D(\mathcal{O}_Y)$。存在典范映射
\begin{equation}
\label{perfect-equation-kunneth}
Ra_*K \otimes_{\mathcal{O}_S}^\mathbf{L} Rb_*M
\longrightarrow
Rf_*(Lp^*K \otimes_{\mathcal{O}_{X \times_S Y}}^\mathbf{L} Lq^*M)
\end{equation}
。事实上，可以使用映射
$Ra_*K \to Ra_*Rp_* Lp^*K = Rf_*Lp^*K$ 与
$Rb_*M \to Rb_*Rq_* Lq^*M = Rf_*Lq^*M$，然后使用相对杯积
（Cohomology, Remark \ref{cohomology-remark-cup-product}）。

\medskip\noindent
令 $A = \Gamma(S, \mathcal{O}_S)$。存在整体 K\"unneth 映射
\begin{equation}
\label{perfect-equation-kunneth-global}
R\Gamma(X, K) \otimes_A^\mathbf{L} R\Gamma(Y, M)
\longrightarrow
R\Gamma(X \times_S Y,
Lp^*K \otimes_{\mathcal{O}_{X \times_S Y}}^\mathbf{L} Lq^*M)
\end{equation}
位于 $D(A)$ 中。该映射由拉回映射
$R\Gamma(X, K) \to R\Gamma(X \times_S Y, Lp^*K)$ 与
$R\Gamma(Y, M) \to R\Gamma(X \times_S Y, Lq^*M)$，以及
Cohomology, Section \ref{cohomology-section-cup-product}
中构造的杯积得到。

\begin{lemma}
\label{perfect-lemma-kunneth}
在上述情形中，若 $a$ 与 $b$ 拟紧拟分离，且 $X$ 与 $Y$ 在 $S$ 上 Tor 独立，
则对 $K \in D_\QCoh(\mathcal{O}_X)$ 与 $M \in D_\QCoh(\mathcal{O}_Y)$，
(\ref{perfect-equation-kunneth}) 是同构。若此外 $S = \Spec(A)$ 仿射，则映射
(\ref{perfect-equation-kunneth-global}) 是同构。
\end{lemma}

\begin{proof}[第一种证明]
结论由下列一串同构推出：
\begin{align*}
Rf_*(Lp^*K \otimes_{\mathcal{O}_{X \times_S Y}}^\mathbf{L} Lq^*M)
& =
Ra_*Rp_*(Lp^*K \otimes_{\mathcal{O}_{X \times_S Y}}^\mathbf{L} Lq^*M) \\
& =
Ra_*(K \otimes_{\mathcal{O}_X}^\mathbf{L} Rp_*Lq^*M) \\
& =
Ra_*(K \otimes_{\mathcal{O}_X}^\mathbf{L} La^*Rb_*M) \\
& =
Ra_*K \otimes_{\mathcal{O}_S}^\mathbf{L} Rb_*M
\end{align*}
第一个等式成立是因为 $f = a \circ p$；第二个等式来自
Lemma \ref{perfect-lemma-cohomology-base-change}；第三个等式来自
Lemma \ref{perfect-lemma-compare-base-change}；第四个等式来自
Lemma \ref{perfect-lemma-cohomology-base-change}。略去验证这些同构的复合正是映射
(\ref{perfect-equation-kunneth})。若 $S$ 仿射，则箭头
(\ref{perfect-equation-kunneth-global}) 的源和靶分别是把 $R\Gamma(S, -)$
作用于 (\ref{perfect-equation-kunneth}) 的源和靶所得，故得到最后的陈述；细节从略。
\end{proof}

\begin{proof}[第二种证明]
箭头 (\ref{perfect-equation-kunneth}) 的构造与限制到 $S$ 的开子概形相容，
这由相对杯积的构造立即可见。因此只需在 $S$ 仿射时证明
(\ref{perfect-equation-kunneth}) 是同构。

\medskip\noindent
设 $S = \Spec(A)$ 仿射。由 Leray，
$R\Gamma(S, Rf_*K) = R\Gamma(X, K)$，其他情形也类似。由
Cohomology, Lemma \ref{cohomology-lemma-compose-cup-product}，
对映射 (\ref{perfect-equation-kunneth}) 取 $R\Gamma(S, -)$ 后诱导出映射
(\ref{perfect-equation-kunneth-global})。另一方面，由 Lemmas
\ref{perfect-lemma-quasi-coherence-direct-image} 与
\ref{perfect-lemma-quasi-coherence-tensor-product}
，映射 (\ref{perfect-equation-kunneth}) 的源与靶都属于 $D_\QCoh(\mathcal{O}_S)$。
所以由 Lemma \ref{perfect-lemma-affine-compare-bounded}，只需证明
(\ref{perfect-equation-kunneth-global}) 是同构。

\medskip\noindent
设 $S = \Spec(A)$、$X = \Spec(B)$、$Y = \Spec(C)$ 均仿射。
下文不再特别说明地使用 Lemma \ref{perfect-lemma-affine-compare-bounded}。
此时可选取各项平坦的 $B$-模 K-平坦复形 $K^\bullet$，使得
$K$ 由 $\widetilde{K}^\bullet$ 表示。类似地，可选取各项平坦的
$C$-模 K-平坦复形 $M^\bullet$，使得 $M$ 由 $\widetilde{M}^\bullet$ 表示。
见 More on Algebra, Lemma \ref{more-algebra-lemma-K-flat-resolution}。
于是 $\widetilde{K}^\bullet$ 是 $\mathcal{O}_X$-模的 K-平坦复形，
$\widetilde{M}^\bullet$ 也类似，见 Lemma \ref{perfect-lemma-affine-K-flat}。因此
$La^*K$ 由
$$
a^*\widetilde{K}^\bullet = \widetilde{K^\bullet \otimes_A C}
$$
表示，$Lb^*M$ 也类似。由上述引理及 More on Algebra, Lemma
\ref{more-algebra-lemma-base-change-K-flat}，这又是
$\mathcal{O}_{X \times_S Y}$-模的 K-平坦复形。最终可见，与
$$
\text{Tot}((K^\bullet \otimes_A C) \otimes_{B \otimes_A C}
(B \otimes_A M^\bullet)) =
\text{Tot}(K^\bullet \otimes_A M^\bullet)
$$
相伴的 $\mathcal{O}_{X \times_S Y}$-模复形在 $X \times_S Y$ 的导出范畴中表示
$La^*K \otimes_{\mathcal{O}_{X \times_S Y}}^\mathbf{L} Lb^*M$。
取整体截面得到 $\text{Tot}(K^\bullet \otimes_A M^\bullet)$；当然它也表示
$R\Gamma(X, K) \otimes_A^\mathbf{L} R\Gamma(Y, M)$。该同构由杯积给出这一事实，
可由 Cohomology, Section \ref{cohomology-section-cup-product} 中真正杯积与朴素杯积的关系推出
（特别见 Cohomology, Lemma \ref{cohomology-lemma-cup-compatible-with-naive}
及其后的讨论）。

\medskip\noindent
假设 $S = \Spec(A)$ 与 $Y$ 仿射。我们用
Cohomology of Schemes, Lemma \ref{coherent-lemma-induction-principle}
中的归纳原理证明该陈述。为此，只需证明：若 $X = U \cup V$ 是开覆盖，
且 $U$ 与 $V$ 拟紧，并且 $U$ 与 $Y$ 在 $S$ 上的映射
(\ref{perfect-equation-kunneth-global})
$$
R\Gamma(U, K) \otimes_A^\mathbf{L} R\Gamma(Y, M)
\longrightarrow
R\Gamma(U \times_S Y,
Lp^*K \otimes_{\mathcal{O}_{X \times_S Y}}^\mathbf{L} Lq^*M)
$$
，$V$ 与 $Y$ 在 $S$ 上的映射 (\ref{perfect-equation-kunneth-global})
$$
R\Gamma(V, K) \otimes_A^\mathbf{L} R\Gamma(Y, M)
\longrightarrow
R\Gamma(V \times_S Y,
Lp^*K \otimes_{\mathcal{O}_{X \times_S Y}}^\mathbf{L} Lq^*M)
$$
，以及 $U \cap V$ 与 $Y$ 在 $S$ 上的映射 (\ref{perfect-equation-kunneth-global})
$$
R\Gamma(U \cap V, K) \otimes_A^\mathbf{L} R\Gamma(Y, M)
\longrightarrow
R\Gamma((U \cap V) \times_S Y,
Lp^*K \otimes_{\mathcal{O}_{X \times_S Y}}^\mathbf{L} Lq^*M)
$$
均为同构，则 $X$ 与 $Y$ 在 $S$ 上的映射
(\ref{perfect-equation-kunneth-global}) 也是同构。然而，由 Cohomology, Lemma
\ref{cohomology-lemma-mayer-vietoris-cup}，这些映射组成可区别三角之间的一个映射，
而 (\ref{perfect-equation-kunneth-global}) 是最后一边。因此由
Derived Categories, Lemma \ref{derived-lemma-third-isomorphism-triangle}
得到结论。

\medskip\noindent
设 $S = \Spec(A)$ 仿射。为完成证明，可在 $Y$ 上使用
Cohomology of Schemes, Lemma \ref{coherent-lemma-induction-principle}
的归纳原理。事实上，由上已知当 $Y$ 仿射时该映射是同构。其余论证与上一段完全相同，
只是交换 $X$ 与 $Y$ 的角色。
\end{proof}

\begin{lemma}
\label{perfect-lemma-cohomology-de-rham-base-change}
设 $a : X \to S$ 为概形的拟紧拟分离态射。设 $\mathcal{F}^\bullet$
为 $a^{-1}\mathcal{O}_S$-模的局部有界复形。假设对所有 $n \in \mathbf{Z}$，
层 $\mathcal{F}^n$ 都是平坦 $a^{-1}\mathcal{O}_S$-模，且 $\mathcal{F}^n$ 具有与给定的
$a^{-1}\mathcal{O}_S$-模结构相容的拟凝聚 $\mathcal{O}_X$-模结构
（但复形 $\mathcal{F}^\bullet$ 中的微分不必是 $\mathcal{O}_X$-线性的）。则：
\begin{enumerate}
\item $Ra_*\mathcal{F}^\bullet$ 局部有界；
\item $Ra_*\mathcal{F}^\bullet$ 属于 $D_\QCoh(\mathcal{O}_S)$；
\item $Ra_*\mathcal{F}^\bullet$ 局部具有有限 Tor 维数；
\item $\mathcal{G} \otimes_{\mathcal{O}_S}^\mathbf{L} Ra_*\mathcal{F}^\bullet =
Ra_*(a^{-1}\mathcal{G} \otimes_{a^{-1}\mathcal{O}_S} \mathcal{F}^\bullet)$
对 $\mathcal{G} \in \QCoh(\mathcal{O}_S)$ 成立；
\item $K \otimes_{\mathcal{O}_S}^\mathbf{L} Ra_*\mathcal{F}^\bullet =
Ra_*(a^{-1}K \otimes_{a^{-1}\mathcal{O}_S}^\mathbf{L} \mathcal{F}^\bullet)$
对 $K \in D_\QCoh(\mathcal{O}_S)$ 成立。
\end{enumerate}
\end{lemma}

\begin{proof}
(1)、(2)、(3) 在 $S$ 上是局部的，故可且确实假设 $S$ 仿射。由于 $a$ 拟紧，
$X$ 拟紧。又因 $\mathcal{F}^\bullet$ 局部有界，故 $\mathcal{F}^\bullet$ 有界。

\medskip\noindent
对 (1) 与 (2)，可以使用 Derived Categories, Lemma
\ref{derived-lemma-two-ss-complex-functor} 的第一个谱序列
$R^pa_*\mathcal{F}^q \Rightarrow R^{p + q}a_*\mathcal{F}^\bullet$
。结合 Cohomology of Schemes, Lemma
\ref{coherent-lemma-quasi-coherence-higher-direct-images}
与 Homology, Lemma \ref{homology-lemma-biregular-ss-converges} 即得结论。

\medskip\noindent
用 Cohomology of Schemes, Lemma
\ref{coherent-lemma-induction-principle} 的归纳原理证明 (3)。具体地，对 $X$
的拟紧开集 $U$，考虑 $R(a|_U)_*(\mathcal{F}^\bullet|_U)$ 具有有限 Tor 维数
这一条件。若 $U, V$ 是 $X$ 中的拟紧开集，则由
Cohomology, Lemma \ref{cohomology-lemma-unbounded-relative-mayer-vietoris}
有相对 Mayer--Vietoris 可区别三角
$$
R(a|_{U \cup V})_*\mathcal{F}^\bullet|_{U \cup V} \to
R(a|_U)_*\mathcal{F}^\bullet|_U \oplus
R(a|_V)_*\mathcal{F}^\bullet|_V \to
R(a|_{U \cap V})_*\mathcal{F}^\bullet|_{U \cap V} \to
$$
。由 Tor 振幅在可区别三角中的行为
（见 Cohomology, Lemma \ref{cohomology-lemma-cone-tor-amplitude}），
若已知结论对 $U$、$V$、$U \cap V$ 成立，则也对 $U \cup V$ 成立。
这把问题归结为 $X$ 仿射的情形。此时
$$
Ra_*\mathcal{F}^\bullet = a_*\mathcal{F}^\bullet
$$
，这是由 Leray 无环引理
（Derived Categories, Lemma \ref{derived-lemma-leray-acyclicity}）
及仿射态射下拟凝聚模的高阶直像消失
（Cohomology of Schemes, Lemma \ref{coherent-lemma-relative-affine-vanishing}）
得到的。按假设 $\mathcal{F}^n$ 在 $S$ 上平坦，且 $X$ 仿射，所以对所有 $n$，
模 $a_*\mathcal{F}^n$ 都平坦。因此 $a_*\mathcal{F}^\bullet$ 是平坦
$\mathcal{O}_S$-模的有界复形，从而具有有限 Tor 维数。

\medskip\noindent
证明 (5)。记
$a' : (X, a^{-1}\mathcal{O}_S) \to (S, \mathcal{O}_S)$
为显然的环化空间平坦态射。(5) 断言
$$
K \otimes_{\mathcal{O}_S}^\mathbf{L} Ra'_*\mathcal{F}^\bullet =
Ra'_*(L(a')^*K \otimes_{a^{-1}\mathcal{O}_S}^\mathbf{L}
\mathcal{F}^\bullet)
$$
。因此 Cohomology, Equation
(\ref{cohomology-equation-projection-formula-map}) 给出从左端到右端的函子性映射，
我们要证明它是同构。这个问题在 $S$ 上是局部的，故可且确实假设 $S$ 仿射。
余下证明与 Lemma \ref{perfect-lemma-cohomology-base-change} 的证明{\it 完全}相同，
但须证明函子
$K \mapsto Ra'_*(L(a')^*K \otimes_{a^{-1}\mathcal{O}_S}^\mathbf{L}
\mathcal{F}^\bullet)$ 与直和可交换。此处要用到 $\mathcal{F}^n$
具有拟凝聚 $\mathcal{O}_X$-模结构。事实上，注意函子
$K \mapsto L(a')^*K
\otimes_{a^{-1}\mathcal{O}_S}^\mathbf{L} \mathcal{F}^\bullet$
与任意直和可交换。其次，若 $\mathcal{F}^\bullet$ 仅由单个拟凝聚
$\mathcal{O}_X$-模 $\mathcal{F}^\bullet = \mathcal{F}^n[-n]$ 构成，则由
Cohomology, Lemma \ref{cohomology-lemma-variant-derived-pullback} 有
$L(a')^*G
\otimes_{a^{-1}\mathcal{O}_S}^\mathbf{L} \mathcal{F}^\bullet =
La^*K \otimes_{\mathcal{O}_X}^\mathbf{L} \mathcal{F}^n[-n]$。
故在此情形中，与直和可交换由
Lemma \ref{perfect-lemma-quasi-coherence-pushforward-direct-sums} 推出。一般地，
由于 $S$ 仿射（因而 $X$ 拟紧），且 $\mathcal{F}^\bullet$ 局部有界，可知
$$
\mathcal{F}^\bullet = (\mathcal{F}^a \to \ldots \to \mathcal{F}^b)
$$
有界。对 $b - a$ 归纳并考虑可区别三角
$$
\mathcal{F}^b[-b] \to (\mathcal{F}^a \to \ldots \to \mathcal{F}^b)
\to (\mathcal{F}^a \to \ldots \to \mathcal{F}^{b - 1}) \to
\mathcal{F}^b[-b + 1]
$$
即可完成此部分的证明。略去若干细节。

\medskip\noindent
证明 (4)。令 $a' : (X, a^{-1}\mathcal{O}_S) \to (S, \mathcal{O}_S)$
如上。由于 $\mathcal{F}^\bullet$ 是平坦 $a^{-1}\mathcal{O}_S$-模的局部有界复形，
复形
$a^{-1}\mathcal{G} \otimes_{a^{-1}\mathcal{O}_S} \mathcal{F}^\bullet$
在 $D(a^{-1}\mathcal{O}_S)$ 中表示
$L(a')^*\mathcal{G}
\otimes_{a^{-1}\mathcal{O}_S}^\mathbf{L}
\mathcal{F}^\bullet$。因此 (4) 由 (5) 推出。
\end{proof}

\begin{lemma}
\label{perfect-lemma-K-flat}
设 $f : X \to Y$ 为概形态射，且 $Y = \Spec(A)$ 仿射。设
$\mathcal{U} : X = \bigcup_{i \in I} U_i$ 为有限仿射开覆盖，使得所有有限交
$U_{i_0 \ldots i_p} = U_{i_0} \cap \ldots \cap U_{i_p}$ 都仿射。
设 $\mathcal{F}^\bullet$ 为 $f^{-1}\mathcal{O}_Y$-模的有界复形。假设对所有
$n \in \mathbf{Z}$，层 $\mathcal{F}^n$ 都是平坦 $f^{-1}\mathcal{O}_Y$-模，
且 $\mathcal{F}^n$ 具有与给定的 $p^{-1}\mathcal{O}_Y$-模结构相容的拟凝聚
$\mathcal{O}_X$-模结构（但复形 $\mathcal{F}^\bullet$ 中的微分不必是
$\mathcal{O}_X$-线性的）。则复形
$\text{Tot}(\check{\mathcal{C}}^\bullet(\mathcal{U}, \mathcal{F}^\bullet))$
作为 $A$-模复形是 K-平坦的。
\end{lemma}

\begin{proof}
可以写成
$$
\mathcal{F}^\bullet = (\mathcal{F}^a \to \ldots \to \mathcal{F}^b)
$$
。对 $b - a$ 归纳，考虑可区别三角
$$
\mathcal{F}^b[-b] \to (\mathcal{F}^a \to \ldots \to \mathcal{F}^b)
\to (\mathcal{F}^a \to \ldots \to \mathcal{F}^{b - 1}) \to
\mathcal{F}^b[-b + 1]
$$
并使用
More on Algebra, Lemma \ref{more-algebra-lemma-K-flat-two-out-of-three}
，问题归结为 $\mathcal{F}^\bullet$ 仅由置于次数 $0$ 的单个拟凝聚
$\mathcal{O}_X$-模 $\mathcal{F}$ 构成的情形。此时，$\mathcal{F}$ 与
$\mathcal{U}$ 的 {\v C}ech 复形同伦等价于交错 {\v C}ech 复形，见
Cohomology, Lemma \ref{cohomology-lemma-alternating-usual}。由于
$U_{i_0 \ldots i_p}$ 总是仿射的，$\mathcal{F}(U_{i_0 \ldots i_p})$
在 $A$ 上平坦。因此
$\check{\mathcal{C}}_{alt}^\bullet(\mathcal{U}, \mathcal{F})$
是平坦 $A$-模的有界复形，从而由 More on Algebra, Lemma
\ref{more-algebra-lemma-derived-tor-quasi-isomorphism} 可知它是 K-平坦的。
\end{proof}

\noindent
令 $X, Y, S, a, b, p, q, f$ 如本节引言所述。设 $\mathcal{F}$ 为
$\mathcal{O}_X$-模，$\mathcal{G}$ 为 $\mathcal{O}_Y$-模，并令
$A = \Gamma(S, \mathcal{O}_S)$。考虑 $D(A)$ 中的映射
\begin{equation}
\label{perfect-equation-kunneth-single-sheaves}
R\Gamma(X, \mathcal{F})
\otimes_A^\mathbf{L}
R\Gamma(Y, \mathcal{G})
\longrightarrow
R\Gamma(X \times_S Y,
p^*\mathcal{F}
\otimes_{\mathcal{O}_{X \times_S Y}}
q^*\mathcal{G})
\end{equation}
。该映射由拉回映射
$R\Gamma(X, \mathcal{F}) \to R\Gamma(X \times_S Y, p^*\mathcal{F})$
与 $R\Gamma(Y, \mathcal{G}) \to R\Gamma(X \times_S Y, q^*\mathcal{G})$、
Cohomology, Section \ref{cohomology-section-cup-product} 中构造的杯积，
以及典范映射
$p^*\mathcal{F}
\otimes_{\mathcal{O}_{X \times_S Y}}^\mathbf{L}
q^*\mathcal{G}
\to
p^*\mathcal{F}
\otimes_{\mathcal{O}_{X \times_S Y}}
q^*\mathcal{G}$ 构造。

\begin{lemma}
\label{perfect-lemma-kunneth-single-sheaf}
在上述情形中，若 $S$ 仿射，$\mathcal{F}$ 与 $\mathcal{G}$ 在 $S$ 上平坦且拟凝聚，
并且 $X$ 与 $Y$ 拟紧且具有仿射对角态射，则映射
(\ref{perfect-equation-kunneth-single-sheaves}) 是同构。
\end{lemma}

\begin{proof}
强烈建议读者先阅读 Varieties, Lemma \ref{varieties-lemma-kunneth} 的证明。
选取有限仿射开覆盖 $\mathcal{U} : X = \bigcup_{i \in I} U_i$ 与
$\mathcal{V} : Y = \bigcup_{j \in J} V_j$。这确定了仿射开覆盖
$\mathcal{W} : X \times_S Y = \bigcup_{(i, j) \in I \times J} U_i \times_S V_j$。
注意 $\mathcal{W}$ 是 $\text{pr}_1^{-1}\mathcal{U}$ 与
$\text{pr}_2^{-1}\mathcal{V}$ 的加细。因此由 Cohomology, Section
\ref{cohomology-section-cech-cohomology-of-complexes} 中的讨论，得到映射
$$
\check{\mathcal{C}}^\bullet(\mathcal{U}, \mathcal{F})
\to
\check{\mathcal{C}}^\bullet(\mathcal{W}, p^*\mathcal{F})
\quad\text{且}\quad
\check{\mathcal{C}}^\bullet(\mathcal{V}, \mathcal{G})
\to
\check{\mathcal{C}}^\bullet(\mathcal{W}, q^*\mathcal{G})
$$
；它们在同伦意义下良定义，并与上同调上的拉回映射相容。在
Cohomology, Equation (\ref{cohomology-equation-needs-signs}) 中，
我们构造了复形映射
$$
\text{Tot}(
\check{\mathcal{C}}^\bullet(\mathcal{W}, p^*\mathcal{F})
\otimes_A
\check{\mathcal{C}}^\bullet(\mathcal{W}, q^*\mathcal{G}))
\longrightarrow
\check{\mathcal{C}}^\bullet(\mathcal{W},
p^*\mathcal{F} \otimes_{\mathcal{O}_{X \times_S Y}}
q^*\mathcal{G})
$$
；由 Cohomology, Lemma \ref{cohomology-lemma-diagrams-commute}，
它与上同调上的杯积相容。结合上述映射，得到复形映射
\begin{equation}
\label{perfect-equation-kunneth-on-cech}
\text{Tot}(
\check{\mathcal{C}}^\bullet(\mathcal{U}, \mathcal{F})
\otimes_A
\check{\mathcal{C}}^\bullet(\mathcal{V}, \mathcal{G}))
\to
\check{\mathcal{C}}^\bullet(\mathcal{W},
p^*\mathcal{F}
\otimes_{\mathcal{O}_{X \times_S Y}}
q^*\mathcal{G})
\end{equation}
我们断言，这就是引理陈述中的映射；即此箭头的源和靶分别与
(\ref{perfect-equation-kunneth-single-sheaves}) 的源和靶相同。事实上，由
Cohomology of Schemes, Lemma
\ref{coherent-lemma-quasi-coherent-affine-cohomology-zero}
及
Cohomology, Lemma \ref{cohomology-lemma-cech-complex-complex-computes}
，典范映射
$$
\check{\mathcal{C}}^\bullet(\mathcal{U}, \mathcal{F})
\to
R\Gamma(X, \mathcal{F}),
\quad
\check{\mathcal{C}}^\bullet(\mathcal{V}, \mathcal{G})
\to
R\Gamma(Y, \mathcal{G})
$$
and
$$
\check{\mathcal{C}}^\bullet(\mathcal{W},
p^*\mathcal{F} \otimes_{\mathcal{O}_{X \times_S Y}} q^*\mathcal{G})
\to
R\Gamma(X \times_S Y,
p^*\mathcal{F} \otimes_{\mathcal{O}_{X \times_S Y}}
q^*\mathcal{G})
$$
均为同构。另一方面，复形
$\check{\mathcal{C}}^\bullet(\mathcal{U}, \mathcal{F})$
由 Lemma \ref{perfect-lemma-K-flat} 是 K-平坦的，故
$\text{Tot}(
\check{\mathcal{C}}^\bullet(\mathcal{U}, \mathcal{F})
\otimes_A
\check{\mathcal{C}}^\bullet(\mathcal{V}, \mathcal{G}))$
表示导出张量积
$R\Gamma(X, \mathcal{F}) \otimes_A^\mathbf{L} R\Gamma(Y, \mathcal{G})$
，正如所断言。

\medskip\noindent
还须证明 (\ref{perfect-equation-kunneth-on-cech}) 是拟同构。我们用维数移位来证明。
令 $d(\mathcal{F}) = \max \{d \mid H^d(X, \mathcal{F}) \not = 0\}$。
假设 $d(\mathcal{F}) > 0$，并令 $U = \coprod\nolimits_{i \in I} U_i$。
由于 $I$ 有限，这是仿射概形。记 $j : U \to X$ 为在每个 $U_i$ 上等于包含
$U_i \to X$ 的态射。由于 $X$ 的对角态射仿射，态射 $j$ 仿射，见
Morphisms, Lemma \ref{morphisms-lemma-affine-permanence}。于是
$\mathcal{F}' = j_*j^*\mathcal{F}$ 在 $S$ 上平坦，见
Morphisms, Lemma \ref{morphisms-lemma-pushforward-flat-affine}。结合
Cohomology of Schemes, Lemmas
\ref{coherent-lemma-relative-affine-cohomology} and
\ref{coherent-lemma-quasi-coherent-affine-cohomology-zero}，还可得
$d(\mathcal{F}') = 0$。对所有 $x \in X$，映射
$\mathcal{F}_x  \to \mathcal{F}'_x$ 都是直和项的包含：若 $x \in U_i$，
则 $\mathcal{F}' \to (U_i \to X)_*\mathcal{F}|_{U_i}$ 给出分裂。
故 $\mathcal{F} \to \mathcal{F}'$ 为单射，且
$\mathcal{F}'' = \mathcal{F}'/\mathcal{F}$ 也在 $S$ 上平坦。平坦拟凝聚
$\mathcal{O}_X$-模的短正合列
$0 \to \mathcal{F} \to \mathcal{F}' \to \mathcal{F}'' \to 0$
给出复形的短正合列
$$
0 \to
\text{Tot}(
\check{\mathcal{C}}^\bullet(\mathcal{U}, \mathcal{F})
\otimes_A
\check{\mathcal{C}}^\bullet(\mathcal{V}, \mathcal{G})) \to
\text{Tot}(
\check{\mathcal{C}}^\bullet(\mathcal{U}, \mathcal{F}')
\otimes_A
\check{\mathcal{C}}^\bullet(\mathcal{V}, \mathcal{G})) \to
\text{Tot}(
\check{\mathcal{C}}^\bullet(\mathcal{U}, \mathcal{F}'')
\otimes_A
\check{\mathcal{C}}^\bullet(\mathcal{V}, \mathcal{G})) \to 0
$$
以及复形的短正合列
$$
0 \to
\check{\mathcal{C}}^\bullet(\mathcal{W},
p^*\mathcal{F}
\otimes_{\mathcal{O}_{X \times_S Y}}
q^*\mathcal{G}) \to
\check{\mathcal{C}}^\bullet(\mathcal{W},
p^*\mathcal{F}'
\otimes_{\mathcal{O}_{X \times_S Y}}
q^*\mathcal{G}) \to
\check{\mathcal{C}}^\bullet(\mathcal{W},
p^*\mathcal{F}''
\otimes_{\mathcal{O}_{X \times_S Y}}
q^*\mathcal{G}) \to 0
$$
此外，它们之间的映射 (\ref{perfect-equation-kunneth-on-cech}) 与这些短正合列相容。
所以只需对 $\mathcal{F}'$ 与 $\mathcal{F}''$ 证明
(\ref{perfect-equation-kunneth-on-cech}) 是同构。最后，
$d(\mathcal{F}'') < d(\mathcal{F})$。这样便归结为
$d(\mathcal{F}) = 0$ 的情形。

\medskip\noindent
对 $\mathcal{G}$ 作同样论证，可假设 $\mathcal{F}$ 与 $\mathcal{G}$
都只在次数 $0$ 上具有非零上同调。注意这意味着 $\Gamma(X, \mathcal{F})$
拟同构于位于次数 $\geq 0$ 的 $K$-平坦 $A$-模复形
$\check{\mathcal{C}}^\bullet(\mathcal{U}, \mathcal{F})$。
因此 $\Gamma(X, \mathcal{F})$ 是平坦 $A$-模
（因为与此模的高阶 Tor 可通过同 {\v C}ech 复形作张量来计算）。
设 $V \subset Y$ 为仿射开集。考虑仿射开覆盖
$\mathcal{U}_V : X \times_S V = \bigcup_{i \in I} U_i \times_S V$。
立即有
$$
\check{\mathcal{C}}^\bullet(\mathcal{U}, \mathcal{F})
\otimes_A \mathcal{G}(V) =
\check{\mathcal{C}}^\bullet(\mathcal{U}_V,
p^*\mathcal{F} \otimes_{\mathcal{O}_{X \times Y}}
q^*\mathcal{G})
$$
（复形相等）。由 $\mathcal{G}(V)$ 在 $A$ 上平坦，可知
$\Gamma(X, \mathcal{F}) \otimes_A \mathcal{G}(V) \to
\check{\mathcal{C}}^\bullet(\mathcal{U}, \mathcal{F})
\otimes_A \mathcal{G}(V)$ 是拟同构。由于预层
$V \mapsto \check{\mathcal{C}}^\bullet(\mathcal{U}_V,
p^*\mathcal{F} \otimes_{\mathcal{O}_{X \times Y}}
q^*\mathcal{G})$ 的层化表示
$Rq_*(p^*\mathcal{F} \otimes_{\mathcal{O}_{X \times Y}} q^*\mathcal{G})$
，见 Cohomology of Schemes, Lemma
\ref{coherent-lemma-separated-case-relative-cech}
，所以在 $Y$ 上有
$$
Rq_*(p^*\mathcal{F} \otimes_{\mathcal{O}_{X \times Y}} q^*\mathcal{G})
\cong
\Gamma(X, \mathcal{F}) \otimes_A \mathcal{G}
$$
，其中右端记号表示模
$$
b^*\widetilde{\Gamma(X, \mathcal{F})} \otimes_{\mathcal{O}_Y} \mathcal{G}
$$
。使用 $q$ 的 Leray 谱序列，得到
$$
H^n(X \times_S Y, p^*\mathcal{F} \otimes_{\mathcal{O}_{X \times Y}}
q^*\mathcal{G}) =
H^n(Y,
b^*\widetilde{\Gamma(X, \mathcal{F})} \otimes_{\mathcal{O}_Y} \mathcal{G})
$$
。对态射 $b : Y \to S = \Spec(A)$ 使用
Lemma \ref{perfect-lemma-cohomology-base-change}，并利用 $\Gamma(X, \mathcal{F})$
在 $A$ 上平坦，可知
$H^n(X \times_S Y, p^*\mathcal{F} \otimes_{\mathcal{O}_{X \times Y}}
q^*\mathcal{G})$ 在 $n > 0$ 时为零，在 $n = 0$ 时同构于
$H^0(X, \mathcal{F}) \otimes_A H^0(Y, \mathcal{G})$。当然，这里还用到
$\mathcal{G}$ 只在次数 $0$ 上有上同调。证明至此完成
（只差检验该同构在次数 $0$ 上确实由杯积给出；我们略去此检验）。
\end{proof}

\begin{remark}
\label{perfect-remark-annoying-compatibility}
设 $S = \Spec(A)$ 为仿射概形。设 $a : X \to S$ 与 $b : Y \to S$
为概形态射。设 $\mathcal{F}$、$\mathcal{G}$ 为拟凝聚 $\mathcal{O}_X$-模，
且 $\mathcal{E}$ 为拟凝聚 $\mathcal{O}_Y$-模。设
$\xi \in H^i(X, \mathcal{G})$，其拉回为
$p^*\xi \in H^i(X \times_S Y, p^*\mathcal{G})$。则下图可交换：
$$
\xymatrix{
R\Gamma(X, \mathcal{F})[-i] \otimes_A^\mathbf{L} R\Gamma(Y, \mathcal{E})
\ar[d] \ar[rr]_-{\xi \otimes \text{id}} & &
R\Gamma(X, \mathcal{G} \otimes_{\mathcal{O}_X} \mathcal{F})
\otimes_A^\mathbf{L} R\Gamma(Y, \mathcal{E}) \ar[d] \\
R\Gamma(X \times_S Y, p^*\mathcal{F} \otimes q^*\mathcal{E})[-i]
\ar[rr]^-{p^*\xi} & &
R\Gamma(X \times_S Y,
p^*(\mathcal{G} \otimes_{\mathcal{O}_X} \mathcal{F}) \otimes q^*\mathcal{E})
}
$$
，其中未带下标的张量积均在 $\mathcal{O}_{X \times_S Y}$ 上取。
水平箭头来自 Cohomology, Remark
\ref{cohomology-remark-cup-with-element-map-total-cohomology}
，竖直箭头是 (\ref{perfect-equation-kunneth-global})，因而由先拉回、再在
$X \times_S Y$ 上作杯积给出。该图可交换，是因为整体杯积
（在 $X \times_S Y$ 上对层 $p^*\mathcal{G}$、$p^*\mathcal{F}$ 与
$q^*\mathcal{E}$）具有结合性，见
Cohomology, Lemma \ref{cohomology-lemma-cup-product-associative}。
\end{remark}







\section{K\"unneth 公式，III}
\label{perfect-section-kunneth-complexes}

\noindent
令 $X, Y, S, a, b, p, q, f$ 如 Section \ref{perfect-section-kunneth} 的引言所述。
本节中，给定 $\mathcal{O}_X$-模 $\mathcal{F}$ 与
$\mathcal{O}_Y$-模 $\mathcal{G}$，令
$$
\mathcal{F} \boxtimes \mathcal{G} =
p^*\mathcal{F} \otimes_{\mathcal{O}_{X \times_S Y}} q^*\mathcal{G}
$$
。注意，与后面某节中的情形不同，这里取的是非导出张量积。

\medskip\noindent
在 $X$ 上，设 $\mathcal{F}^\bullet$ 为阿贝尔群层复形，其各项是拟凝聚
$\mathcal{O}_X$-模，且微分
$d^i_\mathcal{F} : \mathcal{F}^i \to \mathcal{F}^{i + 1}$
是 $X/S$ 上有限阶微分算子，见 Morphisms, Section
\ref{morphisms-section-differential-operators}。类似地，在 $Y$ 上，
设 $\mathcal{G}^\bullet$ 为阿贝尔群层复形，其各项是拟凝聚 $\mathcal{O}_Y$-模，
且微分
$d^j_\mathcal{G} : \mathcal{G}^j \to \mathcal{G}^{j + 1}$
是 $Y/S$ 上有限阶微分算子。应用
Morphisms, Lemma \ref{morphisms-lemma-base-change-differential-operators}
中的构造，得到双复形
$$
\xymatrix{
\ldots &
\ldots &
\ldots &
\ldots \\
\ldots \ar[r] &
\mathcal{F}^i \boxtimes \mathcal{G}^{j + 1}
\ar[r]^{d_1^{i, j + 1}} \ar[u] &
\mathcal{F}^{i + 1} \boxtimes \mathcal{G}^{j + 1} \ar[r] \ar[u] &
\ldots \\
\ldots \ar[r] &
\mathcal{F}^i \boxtimes \mathcal{G}^j
\ar[r]^{d_1^{i, j}} \ar[u]^{d_2^{i, j}} &
\mathcal{F}^{i + 1} \boxtimes \mathcal{G}^j \ar[r] \ar[u]_{d_2^{i + 1, j}} &
\ldots \\
\ldots &
\ldots \ar[u] &
\ldots \ar[u] &
\ldots
}
$$
，其各项是拟凝聚模，映射是 $X \times_S Y / S$ 上有限阶微分算子。参见
Morphisms, Remark \ref{morphisms-remark-base-change-differential-operators}
及
Homology, Example \ref{homology-example-double-complex-as-tensor-product-of}.
明确地说，我们令
$$
d_1^{i, j} = d^i_\mathcal{F} \boxtimes 1
\quad\text{且}\quad
d_2^{i, j} = 1 \boxtimes d^j_\mathcal{G}
$$
。下文记号
$$
\text{Tot}(\mathcal{F}^\bullet \boxtimes \mathcal{G}^\bullet)
$$
指与此双复形相伴的全复形。该复形的各项为拟凝聚
$\mathcal{O}_{X \times_S Y}$-模，其微分为 $X \times_S Y / S$ 上有限阶微分算子。

\medskip\noindent
在上述情形中，存在“相对杯积”映射
\begin{equation}
\label{perfect-equation-relative-de-rham-kunneth}
Ra_*(\mathcal{F}^\bullet)
\otimes_{\mathcal{O}_S}^\mathbf{L}
Rb_*(\mathcal{G}^\bullet)
\longrightarrow
Rf_*\left(\text{Tot}(\mathcal{F}^\bullet \boxtimes \mathcal{G}^\bullet)\right)
\end{equation}
。具体地，可组合下列映射来构造它：
\begin{enumerate}
\item $Ra_*(\mathcal{F}^\bullet) \to Rf_*(p^{-1}\mathcal{F}^\bullet)$,
\item $Rb_*(\mathcal{G}^\bullet) \to Rf_*(q^{-1}\mathcal{G}^\bullet)$,
\item
$Rf_*(p^{-1}\mathcal{F}^\bullet)
\otimes_{\mathcal{O}_S}^\mathbf{L}
Rf_*(q^{-1}\mathcal{G}^\bullet)
\to
Rf_*(p^{-1}\mathcal{F}^\bullet \otimes_{f^{-1}\mathcal{O}_S}^\mathbf{L}
q^{-1}\mathcal{G}^\bullet)$,
\item
$p^{-1}\mathcal{F}^\bullet \otimes_{f^{-1}\mathcal{O}_S}^\mathbf{L}
q^{-1}\mathcal{G}^\bullet \to
\text{Tot}(p^{-1}\mathcal{F}^\bullet \otimes_{f^{-1}\mathcal{O}_S}
q^{-1}\mathcal{G}^\bullet)$
\item
$\text{Tot}(p^{-1}\mathcal{F}^\bullet \otimes_{f^{-1}\mathcal{O}_S}
q^{-1}\mathcal{G}^\bullet) \to \text{Tot}(\mathcal{F}^\bullet \boxtimes
\mathcal{G}^\bullet)$.
\end{enumerate}
(1) 与 (2) 是拉回映射；(3) 是相对杯积，见
Cohomology, Remark \ref{cohomology-remark-cup-product}；
(4) 比较导出与非导出张量积；(5) 由底层双复形上的显然映射
$p^{-1}\mathcal{F}^i \otimes_{f^{-1}\mathcal{O}_S} q^{-1}\mathcal{G}^j
\to \mathcal{F}^i \boxtimes \mathcal{G}^j$ 给出。

\medskip\noindent
令 $A = \Gamma(S, \mathcal{O}_S)$。存在“整体杯积”映射
\begin{equation}
\label{perfect-equation-de-rham-kunneth}
R\Gamma(X, \mathcal{F}^\bullet)
\otimes_A^\mathbf{L}
R\Gamma(Y, \mathcal{G}^\bullet)
\longrightarrow
R\Gamma(X \times_S Y,
\text{Tot}(\mathcal{F}^\bullet \boxtimes \mathcal{G}^\bullet))
\end{equation}
位于 $D(A)$ 中。它与上述相对杯积类似，由下列映射构造：
\begin{enumerate}
\item $R\Gamma(X, \mathcal{F}^\bullet) \to
R\Gamma(X \times_S Y, p^{-1}\mathcal{F}^\bullet)$
\item $R\Gamma(Y, \mathcal{G}^\bullet) \to
R\Gamma(X \times_S Y, q^{-1}\mathcal{G}^\bullet)$,
\item $R\Gamma(X \times_S Y, p^{-1}\mathcal{F}^\bullet)
\otimes_A^\mathbf{L} R\Gamma(X \times_S Y, q^{-1}\mathcal{G}^\bullet) \to
R\Gamma(X \times_S Y,
p^{-1}\mathcal{F}^\bullet \otimes_{f^{-1}\mathcal{O}_S}^\mathbf{L}
q^{-1}\mathcal{G}^\bullet)$,
\item
$p^{-1}\mathcal{F}^\bullet \otimes_{f^{-1}\mathcal{O}_S}^\mathbf{L}
q^{-1}\mathcal{G}^\bullet \to
\text{Tot}(p^{-1}\mathcal{F}^\bullet \otimes_{f^{-1}\mathcal{O}_S}
q^{-1}\mathcal{G}^\bullet)$
\item
$\text{Tot}(p^{-1}\mathcal{F}^\bullet \otimes_{f^{-1}\mathcal{O}_S}
q^{-1}\mathcal{G}^\bullet) \to \text{Tot}(\mathcal{F}^\bullet \boxtimes
\mathcal{G}^\bullet)$.
\end{enumerate}
这里 (1) 与 (2) 是拉回映射，(3) 是 Cohomology, Section
\ref{cohomology-section-cup-product} 中构造的杯积。(4) 与 (5) 如上一段所述。

\begin{lemma}
\label{perfect-lemma-kunneth-special}
在上述情形中，若下列假设成立，则杯积 (\ref{perfect-equation-de-rham-kunneth})
是 $D(A)$ 中的同构：
\begin{enumerate}
\item $S = \Spec(A)$ 仿射；
\item $X$ 与 $Y$ 拟紧且具有仿射对角态射；
\item $\mathcal{F}^\bullet$ 有界；
\item $\mathcal{G}^\bullet$ 下有界；
\item $\mathcal{F}^n$ 在 $S$ 上平坦；
\item $\mathcal{G}^m$ 在 $S$ 上平坦。
\end{enumerate}
\end{lemma}

\begin{proof}
我们使用 Morphisms, Remark
\ref{morphisms-remark-base-change-differential-operators} 中引入的记号
$\mathcal{A}_{X/S}$ 与 $\mathcal{A}_{Y/S}$。假设在范畴 $\mathcal{A}_{X/S}$
中有复形映射
$$
\mathcal{F}_1^\bullet \to
\mathcal{F}_2^\bullet \to
\mathcal{F}_3^\bullet \to
\mathcal{F}_1^\bullet[1]
$$
。由杯积的函子性，得到可交换图
$$
\xymatrix{
R\Gamma(X, \mathcal{F}_1^\bullet)
\otimes_A^\mathbf{L}
R\Gamma(Y, \mathcal{G}^\bullet)
\ar[r] \ar[d] &
R\Gamma(X \times_S Y,
\text{Tot}(\mathcal{F}_1^\bullet \boxtimes \mathcal{G}^\bullet)) \ar[d] \\
R\Gamma(X, \mathcal{F}_2^\bullet)
\otimes_A^\mathbf{L}
R\Gamma(Y, \mathcal{G}^\bullet)
\ar[r] \ar[d] &
R\Gamma(X \times_S Y,
\text{Tot}(\mathcal{F}_2^\bullet \boxtimes \mathcal{G}^\bullet)) \ar[d] \\
R\Gamma(X, \mathcal{F}_3^\bullet)
\otimes_A^\mathbf{L}
R\Gamma(Y, \mathcal{G}^\bullet)
\ar[r] \ar[d] &
R\Gamma(X \times_S Y,
\text{Tot}(\mathcal{F}_3^\bullet \boxtimes \mathcal{G}^\bullet)) \ar[d] \\
R\Gamma(X, \mathcal{F}_1^\bullet[1])
\otimes_A^\mathbf{L}
R\Gamma(Y, \mathcal{G}^\bullet)
\ar[r] &
R\Gamma(X \times_S Y,
\text{Tot}(\mathcal{F}_1^\bullet[1] \boxtimes \mathcal{G}^\bullet))
}
$$
若原映射在 $\mathcal{A}_{X/S}$ 的同伦范畴中组成可区别三角，
则此图的各列在 $D(A)$ 中组成可区别三角。

\medskip\noindent
在本引理的情形中，假设当 $n < i$ 时 $\mathcal{F}^n = 0$。可以考虑逐项分裂的
复形短正合列
$$
0 \to \sigma_{\geq i + 1}\mathcal{F}^\bullet \to
\mathcal{F}^\bullet \to \mathcal{F}^i[-i] \to 0
$$
，其中截断如 Homology, Section \ref{homology-section-truncations} 所述。
这在 $\mathcal{A}_{X/S}$ 的同伦范畴中给出可区别三角
$$
\sigma_{\geq i + 1}\mathcal{F}^\bullet \to
\mathcal{F}^\bullet \to
\mathcal{F}^i[-i] \to
(\sigma_{\geq i + 1}\mathcal{F}^\bullet)[1]
$$
，其中最后一个箭头由边界映射 $\mathcal{F}^i \to \mathcal{F}^{i + 1}$ 给出。
由上述讨论，只需对 $\mathcal{F}^i[-i]$ 与
$\sigma_{\geq i + 1}\mathcal{F}^\bullet$ 证明本引理。由于
$\sigma_{\geq i + 1}\mathcal{F}^\bullet$ 的非零项更少，归纳可知，
只要能在 $\mathcal{F}^\bullet$ 仅于单个次数非零时证明本引理，结论就成立。
因此可假设 $\mathcal{F}^\bullet$ 只在一个次数上非零。

\medskip\noindent
假设 $\mathcal{F}^\bullet$ 是如下复形：拟凝聚 $\mathcal{O}_X$-模
$\mathcal{F}$ 在 $S$ 上平坦并位于次数 $0$，其他次数均为零。例如由
Lemma \ref{perfect-lemma-cohomology-de-rham-base-change}，
$R\Gamma(X, \mathcal{F})$ 具有有限 Tor 维数。设其 Tor 振幅在 $[i, j]$ 中。
取 $N \gg 0$，并在 $\mathcal{A}_{Y/S}$ 的同伦范畴中考虑可区别三角
$$
\sigma_{\geq N + 1}\mathcal{G}^\bullet \to
\mathcal{G}^\bullet \to
\sigma_{\leq N}\mathcal{G}^\bullet \to
(\sigma_{\geq N + 1}\mathcal{G}^\bullet)[1]
$$
。注意下列二者
$$
R\Gamma(X, \mathcal{F})
\otimes_A^\mathbf{L}
R\Gamma(Y, \sigma_{\geq N + 1}\mathcal{G}^\bullet)
\quad\text{且}\quad
R\Gamma(X \times_S Y,
\text{Tot}(\mathcal{F} \boxtimes \sigma_{\geq N + 1}\mathcal{G}^\bullet))
$$
在次数 $\leq N + i$ 上的上同调均消失。因此利用上述论证，若要在某个给定次数上
证明陈述，可以假设 $\mathcal{G}^\bullet$ 有界。再重复一次上述论证，还可假设
$\mathcal{G}^\bullet$ 只在一个次数上非零。此情形由
Lemma \ref{perfect-lemma-kunneth-single-sheaf} 处理。
\end{proof}







\section{Ext 的 K\"unneth 公式}
\label{perfect-section-kunneth-Ext}

\noindent
考虑概形的笛卡尔图
$$
\xymatrix{
& X \times_S Y \ar[ld]^p \ar[rd]_q \ar[dd]^f \\
X \ar[rd]_a & & Y \ar[ld]^b \\
& S
}
$$
本节中，对 $K \in D(\mathcal{O}_X)$ 与 $M \in D(\mathcal{O}_Y)$ 定义
$$
K \boxtimes M =
Lp^*K \otimes_{\mathcal{O}_{X \times_S Y}}^\mathbf{L} Lq^*M
$$
。我们断言，对 $K, K' \in D(\mathcal{O}_X)$ 与
$M, M' \in D(\mathcal{O}_Y)$，存在典范映射
\begin{equation}
\label{perfect-equation-kunneth-ext}
Ra_*R\SheafHom(K, K')
\otimes_{\mathcal{O}_S}^\mathbf{L}
Rb_*R\SheafHom(M, M')
\longrightarrow
Rf_*(R\SheafHom(K \boxtimes M, K' \boxtimes M'))
\end{equation}
。事实上，可取下述映射的伴随映射：
$$
\begin{matrix}
Lf^*\left(Ra_*R\SheafHom(K, K')
\otimes_{\mathcal{O}_S}^\mathbf{L}
Rb_* R\SheafHom(M, M')\right) = \\
Lf^* Ra_* R\SheafHom(K, K')
\otimes_{\mathcal{O}_{X \times_S Y}}^\mathbf{L}
Lf^* Rb_* R\SheafHom(M, M') = \\
Lp^* La^* Ra_* R\SheafHom(K, K')
\otimes_{\mathcal{O}_{X \times_S Y}}^\mathbf{L}
Lq^* Lb^* Rb_* R\SheafHom(M, M') \to \\
Lp^* R\SheafHom(K, K')
\otimes_{\mathcal{O}_{X \times_S Y}}^\mathbf{L}
Lq^* R\SheafHom(M, M') \to \\
R\SheafHom(Lp^*K, Lp^*K')
\otimes_{\mathcal{O}_{X \times_S Y}}^\mathbf{L}
R\SheafHom(Lq^*M, Lq^*M') \to \\
R\SheafHom(K \boxtimes M, K' \boxtimes M')
\end{matrix}
$$
这里第一个等式是拉回与张量积的相容性，见
Cohomology, Lemma \ref{cohomology-lemma-pullback-tensor-product}。
第二个等式使用 $f = a \circ p = b \circ q$ 以及拉回的复合，见
Cohomology, Lemma \ref{cohomology-lemma-derived-pullback-composition}。
第一个箭头由伴随映射
$La^* Ra_* \to \text{id}$
与 $Lb^* Rb_* \to \text{id}$ 给出，因为推前与拉回互为伴随，见
Cohomology, Lemma \ref{cohomology-lemma-adjoint}。第二个箭头由
Cohomology, Remark \ref{cohomology-remark-prepare-fancy-base-change} 给出。
第三个亦即最后一个箭头由
Cohomology, Remark \ref{cohomology-remark-tensor-internal-hom} 给出。
一个简单的特殊情形是下述结果。

\begin{lemma}
\label{perfect-lemma-kunneth-Ext}
在上述情形中，假设 $a$ 与 $b$ 拟紧拟分离，且 $X$ 与 $Y$ 在 $S$ 上 Tor 独立。
若 $K$ 完美、$K' \in D_\QCoh(\mathcal{O}_X)$、$M$ 完美且
$M' \in D_\QCoh(\mathcal{O}_Y)$，则 (\ref{perfect-equation-kunneth-ext}) 是同构。
\end{lemma}

\begin{proof}
在此情形中，有 $R\SheafHom(K, K') = K' \otimes^\mathbf{L} K^\vee$、
$R\SheafHom(M, M') = M' \otimes^\mathbf{L} M^\vee$，且
$$
R\SheafHom(K \boxtimes M, K' \boxtimes M') =
(K' \otimes^\mathbf{L} K^\vee) \boxtimes
(M' \otimes^\mathbf{L} M^\vee)
$$
。见 Cohomology, Lemma \ref{cohomology-lemma-dual-perfect-complex}；
这里还用到拉回与张量积保持完美性。因此该情形由
Lemma \ref{perfect-lemma-kunneth} 推出。（略去验证在这些识别下得到的是同一映射。）
\end{proof}








\section{上同调与基变换，V}
\label{perfect-section-cohomology-base-change}

\noindent
在 Section \ref{perfect-section-cohomology-and-base-change-perfect} 中，
我们看到，当态射 Tor 独立时，一个基变换定理成立。即使在仿射情形中，
没有这样的条件也不可能有基变换定理，见 More on Algebra, Section
\ref{more-algebra-section-tor-independence}。本节分析何时可以“一次对一个复形”
得到基变换结果。

\medskip\noindent
为实现这一点，假设有概形的可交换图
$$
\xymatrix{
X' \ar[r]_{g'} \ar[d]_{f'} &
X \ar[d]^f \\
S' \ar[r]^g &
S
}
$$
（通常将假设它是笛卡尔图）。设 $K \in D_\QCoh(\mathcal{O}_X)$，且设
$L(g')^*K \to K'$ 为 $D_\QCoh(\mathcal{O}_{X'})$ 中的映射。
对点 $x' \in X'$，令 $x = g'(x') \in X$、$s' = f'(x') \in S'$ 及
$s = f(x) = g(s')$。于是可考虑映射
$$
K_x \otimes_{\mathcal{O}_{S, s}}^\mathbf{L} \mathcal{O}_{S', s'} \to
K_x \otimes_{\mathcal{O}_{X, x}}^\mathbf{L} \mathcal{O}_{X', x'} \to
K'_{x'}
$$
，其中第一个箭头是 More on Algebra, Equation
(\ref{more-algebra-equation-comparison-map})，第二个来自
$(L(g')^*K)_{x'} =
K_x \otimes_{\mathcal{O}_{X, x}}^\mathbf{L} \mathcal{O}_{X', x'}$
及给定映射 $L(g')^*K \to K'$。对每个 $i \in \mathbf{Z}$，
$H^i(K_x \otimes_{\mathcal{O}_{S, s}}^\mathbf{L} \mathcal{O}_{S', s'})$.
上得到一个
$\mathcal{O}_{X, x} \otimes_{\mathcal{O}_{S, s}} \mathcal{O}_{S', s'}$-模结构。
综合这些数据，得到典范映射
\begin{equation}
\label{perfect-equation-bc}
H^i(K_x \otimes_{\mathcal{O}_{S, s}}^\mathbf{L} \mathcal{O}_{S', s'})
\otimes_{(\mathcal{O}_{X, x} \otimes_{\mathcal{O}_{S, s}} \mathcal{O}_{S', s'})}
\mathcal{O}_{X', x'}
\longrightarrow
H^i(K'_{x'})
\end{equation}
，它是 $\mathcal{O}_{X', x'}$-模映射。

\begin{lemma}
\label{perfect-lemma-single-complex-base-change-condition}
设
$$
\xymatrix{
X' \ar[r]_{g'} \ar[d]_{f'} &
X \ar[d]^f \\
S' \ar[r]^g &
S
}
$$
是概形的笛卡尔图。设 $K \in D_\QCoh(\mathcal{O}_X)$，且设
$L(g')^*K \to K'$ 为 $D_\QCoh(\mathcal{O}_{X'})$ 中的映射。下列条件等价：
\begin{enumerate}
\item 对任意 $x' \in X'$ 与 $i \in \mathbf{Z}$，映射
(\ref{perfect-equation-bc}) 是同构；
\item 对仿射开集 $U \subset X$ 与 $V' \subset S'$，若二者都映入仿射开集
$V \subset S$，且 $U' = V' \times_V U$，则复合
$$
R\Gamma(U, K) \otimes_{\mathcal{O}_S(U)}^\mathbf{L} \mathcal{O}_{S'}(V')
\to
R\Gamma(U, K) \otimes_{\mathcal{O}_X(U)}^\mathbf{L} \mathcal{O}_{X'}(U')
\to
R\Gamma(U', K')
$$
是 $D(\mathcal{O}_{S'}(V'))$ 中的同构；
\item 存在集合 $I$，以及如 (2) 中的四元组 $U_i, V_i', V_i, U_i'$，$i \in I$，
使得 $X' = \bigcup U'_i$。
\end{enumerate}
\end{lemma}

\begin{proof}
(2) 中第二个箭头来自
$$
R\Gamma(U, K) \otimes_{\mathcal{O}_X(U)}^\mathbf{L} \mathcal{O}_{X'}(U') =
R\Gamma(U', L(g')^*K)
$$
这一 Lemma \ref{perfect-lemma-quasi-coherence-pullback} 中的等式及给定箭头
$L(g')^*K \to K'$。(2) 中第一个箭头是 More on Algebra, Equation
(\ref{more-algebra-equation-comparison-map})。显然 (2) 蕴含 (3)。注意 (1)
在 $X'$ 上是局部的。因此只需证明：若 $X$、$S$、$S'$、$X'$ 均仿射，
则 (1) 等价于下述条件：
$$
R\Gamma(X, K) \otimes_{\mathcal{O}_S(S)}^\mathbf{L} \mathcal{O}_{S'}(S')
\to
R\Gamma(X, K) \otimes_{\mathcal{O}_X(X)}^\mathbf{L} \mathcal{O}_{X'}(X')
\to
R\Gamma(X', K')
$$
是 $D(\mathcal{O}_{S'}(S'))$ 中的同构。记
$S = \Spec(R)$、$X = \Spec(A)$、$S' = \Spec(R')$、$X' = \Spec(A')$；
$K$ 对应于 $A$-模复形 $M^\bullet$，$K'$ 对应于 $A'$-模复形
$N^\bullet$。注意 $A' = A \otimes_R R'$。上述条件即复合
$$
M^\bullet \otimes_R^\mathbf{L} R' \to
M^\bullet \otimes_A^\mathbf{L} A' \to
N^\bullet
$$
是 $D(R')$ 中的同构。等价地，对所有 $i \in \mathbf{Z}$，映射
$$
H^i(M^\bullet \otimes_R^\mathbf{L} R') \to
H^i(M^\bullet \otimes_A^\mathbf{L} A') \to
H^i(N^\bullet)
$$
是同构。注意这是 $A \otimes_R R'$-模映射，即 $A'$-模映射。
另一方面，(1) 要求对相容素理想
$\mathfrak q' \subset A'$, $\mathfrak q \subset A$,
$\mathfrak p' \subset R'$, $\mathfrak p \subset R$
，复合
$$
H^i(M^\bullet_\mathfrak q \otimes_{R_\mathfrak p}^\mathbf{L} R'_{\mathfrak p'})
\otimes_{(A_\mathfrak q \otimes_{R_\mathfrak p} R'_{\mathfrak p'})}
A'_{\mathfrak q'} \to
H^i(M^\bullet_{\mathfrak q}
\otimes_{A_\mathfrak q}^\mathbf{L} A'_{\mathfrak q'})
\to H^i(N^\bullet_{\mathfrak q'})
$$
是同构。由于
$$
H^i(M^\bullet_\mathfrak q \otimes_{R_\mathfrak p}^\mathbf{L} R'_{\mathfrak p'})
\otimes_{(A_\mathfrak q \otimes_{R_\mathfrak p} R'_{\mathfrak p'})}
A'_{\mathfrak q'} =
H^i(M^\bullet \otimes_R^\mathbf{L} R') \otimes_{A'} A'_{\mathfrak q'}
$$
是在 $\mathfrak q'$ 处的局部化，由 Algebra, Lemma
\ref{algebra-lemma-characterize-zero-local} 可知这两个条件等价。
\end{proof}

\begin{lemma}
\label{perfect-lemma-single-complex-base-change}
设
$$
\xymatrix{
X' \ar[r]_{g'} \ar[d]_{f'} &
X \ar[d]^f \\
S' \ar[r]^g &
S
}
$$
是概形的笛卡尔图。设 $K \in D_\QCoh(\mathcal{O}_X)$，且设
$L(g')^*K \to K'$ 为 $D_\QCoh(\mathcal{O}_{X'})$ 中的映射。若
\begin{enumerate}
\item Lemma \ref{perfect-lemma-single-complex-base-change-condition} 的等价条件成立；且
\item $f$ 拟紧拟分离，
\end{enumerate}
则复合 $Lg^*Rf_*K \to Rf'_*L(g')^*K \to Rf'_*K'$ 是同构。
\end{lemma}

\begin{proof}
可以使用 Lemma \ref{perfect-lemma-compare-base-change} 证明中的同一方法来证明，
但我们改用归纳原理与相对 Mayer--Vietoris 来证明。

\medskip\noindent
为检验该映射是同构，可以在 $S'$ 上局部地工作。因此可假设
$g : S' \to S$ 是仿射概形之间的态射。特别地，$X$ 是拟紧拟分离概形。
我们使用
Cohomology of Schemes, Lemma \ref{coherent-lemma-induction-principle}
的归纳原理证明：对任意拟紧开集 $U \subset X$，类似构造的映射
$Lg^*R(U \to S)_*K|_U \to R(U' \to S')_*K'|_{U'}$ 是同构；其中
$U' = (g')^{-1}(U)$。

\medskip\noindent
若 $U \subset X$ 为仿射开集，则由假设可知结论成立；见
Lemma \ref{perfect-lemma-single-complex-base-change-condition} 的 (2)，以及
Lemmas \ref{perfect-lemma-affine-compare-bounded} 与
\ref{perfect-lemma-quasi-coherence-pullback} 所给出的代数转译。

\medskip\noindent
归纳步骤。设 $X = U \cup V$ 是开覆盖，$U$、$V$、$U \cap V$ 均拟紧，
且结论对 $U$、$V$ 与 $U \cap V$ 成立。记 $a = f|_U$、$b = f|_V$、
$c = f|_{U \cap V}$。设 $a' : U' \to S'$、$b' : V' \to S'$ 与
$c' : U' \cap V' \to S'$ 分别为 $a$、$b$ 与 $c$ 的基变换。
使用相对 Mayer--Vietoris 的可区别三角
（Cohomology, Lemma \ref{cohomology-lemma-unbounded-relative-mayer-vietoris}），
得到可交换图
$$
\xymatrix{
Lg^*Rf_*K \ar[r] \ar[d] &
Rf'_* K' \ar[d] \\
Lg^*Ra_* K|_U \oplus
Lg^*Rb_* K|_V \ar[r] \ar[d] &
Ra'_* K'|_{U'} \oplus
Rb'_* K'|_{V'} \ar[d] \\
Lg^*Rc_* K|_{U \cap V} \ar[r] \ar[d] &
Rc'_* K'|_{U' \cap V'} \ar[d] \\
Lg^*Rf_* K[1] \ar[r] &
Rf'_* K'[1]
}
$$
由于第二个与第三个水平箭头是同构，第一个也是同构
（Derived Categories, Lemma \ref{derived-lemma-third-isomorphism-triangle}）。
引理得证。
\end{proof}

\begin{lemma}
\label{perfect-lemma-single-complex-base-change-condition-inherited}
设
$$
\xymatrix{
X' \ar[r]_{g'} \ar[d]_{f'} &
X \ar[d]^f \\
S' \ar[r]^g &
S
}
$$
是概形的笛卡尔图。设 $K \in D_\QCoh(\mathcal{O}_X)$，且设
$L(g')^*K \to K'$ 为 $D_\QCoh(\mathcal{O}_{X'})$ 中的映射。若
Lemma \ref{perfect-lemma-single-complex-base-change-condition} 的等价条件成立，则：
\begin{enumerate}
\item 对 $E \in D_\QCoh(\mathcal{O}_X)$，映射
$L(g')^*(E \otimes^\mathbf{L} K) \to L(g')^*E \otimes^\mathbf{L} K'$
满足 Lemma \ref{perfect-lemma-single-complex-base-change-condition} 的等价条件；
\item 若 $D(\mathcal{O}_X)$ 中的 $E$ 完美，则映射
$L(g')^*R\SheafHom(E, K) \to R\SheafHom(L(g')^*E, K')$
满足 Lemma \ref{perfect-lemma-single-complex-base-change-condition} 的等价条件；
\item 若 $K$ 下有界，且 $D(\mathcal{O}_X)$ 中的 $E$ 伪凝聚，则映射
$L(g')^*R\SheafHom(E, K) \to R\SheafHom(L(g')^*E, K')$
满足 Lemma \ref{perfect-lemma-single-complex-base-change-condition} 的等价条件。
\end{enumerate}
\end{lemma}

\begin{proof}
由 Lemmas \ref{perfect-lemma-quasi-coherence-pullback}、
\ref{perfect-lemma-quasi-coherence-tensor-product}、
\ref{perfect-lemma-quasi-coherence-internal-hom} 及 Cohomology, Lemmas
\ref{cohomology-lemma-pseudo-coherent-pullback}、
\ref{cohomology-lemma-perfect-pullback}，所涉及复形的上同调层拟凝聚，
所以该陈述有意义。在情形 (1) 中，利用张量积的传递性计算
(\ref{perfect-equation-bc}) 的源和靶，即可检验映射 (\ref{perfect-equation-bc}) 是同构，见
More on Algebra, Lemma \ref{more-algebra-lemma-triple-tensor-product}。
(2) 与 (3) 中的映射是复合
$$
L(g')^*R\SheafHom(E, K) \to R\SheafHom(L(g')^*E, L(g')^*K)
\to R\SheafHom(L(g')^*E, K')
$$
，其中第一个箭头是 Cohomology, Remark
\ref{cohomology-remark-prepare-fancy-base-change}，第二个箭头来自给定映射
$L(g')^*K \to K'$。为证明映射 (\ref{perfect-equation-bc}) 是同构，在情形 (2) 中
用有限投射 $\mathcal{O}_{X. x}$-模的有界复形表示 $E_x$，在情形 (3) 中
用有限自由模的有上界复形表示它，然后计算箭头的源和靶。略去若干细节。
\end{proof}

\begin{lemma}
\label{perfect-lemma-base-change-tensor}
设 $f : X \to S$ 为概形的拟紧拟分离态射。设
$E \in D_\QCoh(\mathcal{O}_X)$。设 $\mathcal{G}^\bullet$ 为在 $S$ 上平坦的
拟凝聚 $\mathcal{O}_X$-模的有上界复形。则构造
$$
Rf_*(E \otimes^\mathbf{L}_{\mathcal{O}_X} \mathcal{G}^\bullet)
$$
与任意基变换可交换（精确陈述见证明）。
\end{lemma}

\begin{proof}
该陈述意指如下。设 $g : S' \to S$ 为概形态射，并考虑基变换图
$$
\xymatrix{
X' \ar[r]_{g'} \ar[d]_{f'} &
X \ar[d]^f \\
S' \ar[r]^g &
S
}
$$
，换言之，$X' = S' \times_S X$。本引理断言
$$
Lg^*Rf_*(E \otimes^\mathbf{L}_{\mathcal{O}_X} \mathcal{G}^\bullet)
\longrightarrow
Rf'_*\left(
L(g')^*E \otimes^\mathbf{L}_{\mathcal{O}_{X'}} (g')^*\mathcal{G}^\bullet
\right)
$$
是同构。注意右端对 $\mathcal{G}^\bullet$ 并{\bf 未}使用导出拉回。
为证明这一点，应用 Lemmas \ref{perfect-lemma-single-complex-base-change} 与
\ref{perfect-lemma-single-complex-base-change-condition-inherited}，可知只需证明典范映射
$$
L(g')^*\mathcal{G}^\bullet \to (g')^*\mathcal{G}^\bullet
$$
满足 Lemma \ref{perfect-lemma-single-complex-base-change-condition} 的等价条件。
在茎上检验该条件即可；此时它立即由如下事实推出：按我们对复形
$\mathcal{G}^\bullet$ 的假设，
$\mathcal{G}^\bullet_x \otimes_{\mathcal{O}_{S, s}} \mathcal{O}_{S', s'}$
计算导出张量积。
\end{proof}

\begin{lemma}
\label{perfect-lemma-base-change-RHom}
设 $f : X \to S$ 为概形的拟紧拟分离态射。设 $E$ 为
$D(\mathcal{O}_X)$ 的对象，并设 $\mathcal{G}^\bullet$ 为拟凝聚
$\mathcal{O}_X$-模的复形。若
\begin{enumerate}
\item $E$ 完美，$\mathcal{G}^\bullet$ 有上界，且 $\mathcal{G}^n$ 在 $S$ 上平坦；或
\item $E$ 伪凝聚，$\mathcal{G}^\bullet$ 有界，且 $\mathcal{G}^n$ 在 $S$ 上平坦，
\end{enumerate}
则构造
$$
Rf_*R\SheafHom(E, \mathcal{G}^\bullet)
$$
与任意基变换可交换（精确陈述见证明）。
\end{lemma}

\begin{proof}
该陈述意指如下。设 $g : S' \to S$ 为概形态射，并考虑基变换图
$$
\xymatrix{
X' \ar[r]_{g'} \ar[d]_{f'} &
X \ar[d]^f \\
S' \ar[r]^g &
S
}
$$
，换言之，$X' = S' \times_S X$。本引理断言
$$
Lg^*Rf_*R\SheafHom(E, \mathcal{G}^\bullet)
\longrightarrow
R(f')_*R\SheafHom(L(g')^*E, (g')^*\mathcal{G}^\bullet)
$$
是同构。注意右端对 $\mathcal{G}^\bullet$ 并{\bf 未}使用导出拉回。为证明这一点，
应用 Lemmas \ref{perfect-lemma-single-complex-base-change} 与
\ref{perfect-lemma-single-complex-base-change-condition-inherited}，可知只需证明典范映射
$$
L(g')^*\mathcal{G}^\bullet \to (g')^*\mathcal{G}^\bullet
$$
满足 Lemma \ref{perfect-lemma-single-complex-base-change-condition} 的等价条件。
这已在 Lemma \ref{perfect-lemma-base-change-tensor} 的证明中说明。
\end{proof}





\section{构造完美复形}
\label{perfect-section-producing-perfect}

\noindent
下述引理是构造完美复形的主要技术工具。此结果的后续版本将通过 Noether 逼近
归结到这里，见 Section \ref{perfect-section-cohomology-and-base-change-final}。

\begin{lemma}
\label{perfect-lemma-perfect-direct-image}
设 $S$ 为 Noether 概形。设 $f : X \to S$ 为局部有限型概形态射。
设 $E \in D(\mathcal{O}_X)$，并满足：
\begin{enumerate}
\item $E \in D^b_{\textit{Coh}}(\mathcal{O}_X)$；
\item 对所有 $i$，$H^i(E)$ 的支集在 $S$ 上固有；
\item $E$ 作为 $D(f^{-1}\mathcal{O}_S)$ 的对象具有有限 Tor 维数。
\end{enumerate}
则 $Rf_*E$ 是 $D(\mathcal{O}_S)$ 的完美对象。
\end{lemma}

\begin{proof}
由 Lemma \ref{perfect-lemma-direct-image-coherent}，$Rf_*E$ 是
$D^b_{\textit{Coh}}(\mathcal{O}_S)$ 的对象。因此 $Rf_*E$ 伪凝聚
（Lemma \ref{perfect-lemma-identify-pseudo-coherent-noetherian}）。所以只需证明
$Rf_*E$ 具有有限 Tor 维数，见 Cohomology, Lemma
\ref{cohomology-lemma-perfect}。由 Lemma \ref{perfect-lemma-tor-qc-qs}，
只需检验对 $S$ 上所有拟凝聚层 $\mathcal{F}$，
$Rf_*(E) \otimes_{\mathcal{O}_S}^\mathbf{L} \mathcal{F}$
都具有一致有界的上同调。上有界性显然，因为 $Rf_*(E)$ 上有界。
令 $T \subset X$ 为所有 $i$ 的 $H^i(E)$ 的支集之并。由假设 (1) 与 (2)，
$T$ 在 $S$ 上固有，见 Cohomology of Schemes, Lemma
\ref{coherent-lemma-union-closed-proper-over-base}。特别地，存在包含 $T$ 的拟紧开集
$X' \subset X$。令 $f' = f|_{X'}$；由于 $E$ 限制到 $X \setminus T$ 上为零，
有 $Rf_*(E) = Rf'_*(E|_{X'})$。所以可用 $X'$ 替换 $X$ 并假设 $f$ 拟紧。
此外，由 Morphisms, Lemma
\ref{morphisms-lemma-finite-type-Noetherian-quasi-separated}，$f$ 拟分离。现在
$$
Rf_*(E) \otimes_{\mathcal{O}_S}^\mathbf{L} \mathcal{F} =
Rf_*\left(E \otimes_{\mathcal{O}_X}^\mathbf{L} Lf^*\mathcal{F}\right) =
Rf_*\left(E \otimes_{f^{-1}\mathcal{O}_S}^\mathbf{L} f^{-1}\mathcal{F}\right)
$$
，这是由 Lemma \ref{perfect-lemma-cohomology-base-change} 及
Cohomology, Lemma \ref{cohomology-lemma-variant-derived-pullback} 得到的。
由假设 (3)，复形
$E \otimes_{f^{-1}\mathcal{O}_S}^\mathbf{L} f^{-1}\mathcal{F}$
的上同调层位于某个给定有限区间内，记为 $[a, b]$。于是对它取 $Rf_*$ 后，
上同调位于区间 $[a, \infty)$ 内，结论得证。
\end{proof}

\begin{lemma}
\label{perfect-lemma-tensor-perfect}
设 $S$ 为 Noether 概形。设 $f : X \to S$ 为局部有限型概形态射。
设 $E \in D(\mathcal{O}_X)$ 完美。设 $\mathcal{G}^\bullet$ 为凝聚
$\mathcal{O}_X$-模的有界复形，它在 $S$ 上平坦且支集在 $S$ 上固有。
则 $K = Rf_*(E \otimes_{\mathcal{O}_X}^\mathbf{L} \mathcal{G}^\bullet)$
是 $D(\mathcal{O}_S)$ 的完美对象。
\end{lemma}

\begin{proof}
由 Lemma \ref{perfect-lemma-perfect-direct-image}，对象 $K$ 完美。我们检验该引理适用：
局部地，$E$ 同构于有限自由 $\mathcal{O}_X$-模的有限复形。因此局部地，
$E \otimes^\mathbf{L}_{\mathcal{O}_X} \mathcal{G}^\bullet$ 同构于一个有限复形，
其各项形如
$$
\bigoplus\nolimits_{i = a, \ldots, b} (\mathcal{G}^i)^{\oplus r_i}
$$
，其中 $a, b, r_a, \ldots, r_b$ 是某些整数。这立即蕴含上同调层
$H^i(E \otimes^\mathbf{L}_{\mathcal{O}_X} \mathcal{G})$ 凝聚。
Tor 维数的假设也由 $\mathcal{G}^i$ 在 $f^{-1}\mathcal{O}_S$ 上平坦推出。
\end{proof}

\begin{lemma}
\label{perfect-lemma-ext-perfect}
设 $S$ 为 Noether 概形。设 $f : X \to S$ 为局部有限型概形态射。
设 $E \in D(\mathcal{O}_X)$ 完美。设 $\mathcal{G}^\bullet$ 为凝聚
$\mathcal{O}_X$-模的有界复形，它在 $S$ 上平坦且支集在 $S$ 上固有。
则 $K = Rf_*R\SheafHom(E, \mathcal{G}^\bullet)$ 是
$D(\mathcal{O}_S)$ 的完美对象。
\end{lemma}

\begin{proof}
由于 $E$ 是完美复形，存在对偶完美复形 $E^\vee$，见
Cohomology, Lemma \ref{cohomology-lemma-dual-perfect-complex}。注意
$R\SheafHom(E, \mathcal{G}^\bullet) =
E^\vee \otimes^\mathbf{L}_{\mathcal{O}_X} \mathcal{G}^\bullet$。
因此 $K$ 的完美性由 Lemma \ref{perfect-lemma-tensor-perfect} 推出。
\end{proof}

\noindent
我们将在 Lemma \ref{perfect-lemma-flat-proper-perfect-direct-image-general}
中把下述引理推广到一般基底上的平坦固有态射，并在
More on Morphisms, Lemma
\ref{more-morphisms-lemma-perfect-proper-perfect-direct-image}
中推广到完美固有态射。

\begin{lemma}
\label{perfect-lemma-flat-proper-perfect-direct-image}
设 $S$ 为 Noether 概形。设 $f : X \to S$ 为平坦固有概形态射。
设 $E \in D(\mathcal{O}_X)$ 完美。则 $Rf_*E$ 是
$D(\mathcal{O}_S)$ 的完美对象。
\end{lemma}

\begin{proof}
我们断言 Lemma \ref{perfect-lemma-perfect-direct-image} 适用。条件 (1) 与 (2) 是立即的。
条件 (3) 在 $X$ 上是局部的。因此可假设 $X$ 与 $S$ 仿射，且 $E$
由 $\mathcal{O}_X$-模的严格完美复形表示。由于 $\mathcal{O}_X$ 作为
$f^{-1}\mathcal{O}_S$-模层是平坦的，条件 (3) 成立。
\end{proof}






\section{Ext 的投影公式}
\label{perfect-section-ext}

\noindent
Lemma \ref{perfect-lemma-compute-ext}（或文献中的类似结果）有时用于验证 Quot 函子、
Hilbert 概形及其他模问题的某个 Artin 判据。假设 $f : X \to S$
是固有、平坦、有限表示的概形态射，且 $E \in D(\mathcal{O}_X)$ 完美。
此时该引理断言
$$
\Ext^i_X(E, f^*\mathcal{F}) =
\Ext^i_S((Rf_*E^\vee)^\vee, \mathcal{F})
$$
对 $S$ 上拟凝聚的 $\mathcal{F}$ 成立。这样写使它看起来像 Ext 的投影公式；
事实上，该结果确实相当容易地由 Lemma \ref{perfect-lemma-cohomology-base-change} 推出。

\begin{lemma}
\label{perfect-lemma-compute-tensor-perfect}
采用 Lemma \ref{perfect-lemma-tensor-perfect} 的假设与记号。则对 $S$ 上拟凝聚的
$\mathcal{F}$，存在与边界映射相容的函子性同构（见证明）
$$
H^i(S, K \otimes^\mathbf{L}_{\mathcal{O}_S} \mathcal{F})
\longrightarrow
H^i(X, E \otimes_{\mathcal{O}_X}^\mathbf{L}
(\mathcal{G}^\bullet \otimes_{\mathcal{O}_X} f^*\mathcal{F}))
$$
。
\end{lemma}

\begin{proof}
有
$$
\mathcal{G}^\bullet \otimes_{\mathcal{O}_X}^\mathbf{L} Lf^*\mathcal{F} =
\mathcal{G}^\bullet \otimes_{f^{-1}\mathcal{O}_S}^\mathbf{L} f^{-1}\mathcal{F} =
\mathcal{G}^\bullet \otimes_{f^{-1}\mathcal{O}_S} f^{-1}\mathcal{F} =
\mathcal{G}^\bullet \otimes_{\mathcal{O}_X} f^*\mathcal{F}
$$
；第一个等式由 Cohomology, Lemma
\ref{cohomology-lemma-variant-derived-pullback} 给出；第二个等式是因为
$\mathcal{G}^n$ 是平坦 $f^{-1}\mathcal{O}_S$-模；第三个等式来自拉回的定义。
因此得到
\begin{align*}
H^i(X, E \otimes^\mathbf{L}_{\mathcal{O}_X}
(\mathcal{G}^\bullet \otimes_{\mathcal{O}_X} f^*\mathcal{F}))
& =
H^i(X, E \otimes^\mathbf{L}_{\mathcal{O}_X} \mathcal{G}^\bullet
\otimes_{\mathcal{O}_X}^\mathbf{L} Lf^*\mathcal{F}) \\
& =
H^i(S,
Rf_*(E \otimes^\mathbf{L}_{\mathcal{O}_X} \mathcal{G}^\bullet
\otimes^\mathbf{L}_{\mathcal{O}_X} Lf^*\mathcal{F})) \\
& =
H^i(S, Rf_*(E \otimes^\mathbf{L}_{\mathcal{O}_X} \mathcal{G}^\bullet)
\otimes^\mathbf{L}_{\mathcal{O}_S} \mathcal{F}) \\
& =
H^i(S, K \otimes^\mathbf{L}_{\mathcal{O}_S} \mathcal{F}) 
\end{align*}
。第一个等式由上文给出；第二个等式由 Leray 给出
（Cohomology, Lemma \ref{cohomology-lemma-before-Leray}）；第三个等式由
Lemma \ref{perfect-lemma-cohomology-base-change} 给出。关于边界映射的陈述意指如下：
给定拟凝聚 $\mathcal{O}_S$-模的短正合列
$0 \to \mathcal{F}_1 \to \mathcal{F}_2 \to \mathcal{F}_3 \to 0$，
这些同构组成可交换图
$$
\xymatrix{
H^i(S, K \otimes^\mathbf{L}_{\mathcal{O}_S} \mathcal{F}_3)
\ar[r] \ar[d]_\delta &
H^i(X, E \otimes^\mathbf{L}_{\mathcal{O}_X}
(\mathcal{G}^\bullet \otimes_{\mathcal{O}_X} f^*\mathcal{F}_3))
\ar[d]^\delta \\
H^{i + 1}(S, K \otimes^\mathbf{L}_{\mathcal{O}_S} \mathcal{F}_1)
\ar[r] &
H^{i + 1}(X, E \otimes^\mathbf{L}_{\mathcal{O}_X}
(\mathcal{G}^\bullet \otimes_{\mathcal{O}_X} f^*\mathcal{F}_1))
}
$$
，其中边界映射来自可区别三角
$$
K \otimes^\mathbf{L}_{\mathcal{O}_S} \mathcal{F}_1 \to
K \otimes^\mathbf{L}_{\mathcal{O}_S} \mathcal{F}_2 \to
K \otimes^\mathbf{L}_{\mathcal{O}_S} \mathcal{F}_3 \to
K \otimes^\mathbf{L}_{\mathcal{O}_S} \mathcal{F}_1[1]
$$
以及与 $D(\mathcal{O}_X)$ 中下述短正合列相伴的可区别三角：
$$
0 \to
\mathcal{G}^\bullet \otimes_{\mathcal{O}_X} f^*\mathcal{F}_1 \to
\mathcal{G}^\bullet \otimes_{\mathcal{O}_X} f^*\mathcal{F}_2 \to
\mathcal{G}^\bullet \otimes_{\mathcal{O}_X} f^*\mathcal{F}_3 \to 0
$$
它是 $\mathcal{O}_X$-模复形的短正合列。该列正合是因为 $\mathcal{G}^n$
在 $S$ 上平坦。略去验证所显示图的可交换性。
\end{proof}

\begin{lemma}
\label{perfect-lemma-compute-ext-perfect}
采用 Lemma \ref{perfect-lemma-ext-perfect} 的假设与记号。则对 $S$ 上拟凝聚的
$\mathcal{F}$，存在与边界映射相容的函子性同构（见证明）
$$
H^i(S, K \otimes^\mathbf{L}_{\mathcal{O}_S} \mathcal{F})
\longrightarrow
\Ext^i_{\mathcal{O}_X}(E,
\mathcal{G}^\bullet \otimes_{\mathcal{O}_X} f^*\mathcal{F})
$$
。
\end{lemma}

\begin{proof}
如 Lemma \ref{perfect-lemma-ext-perfect} 的证明，令 $E^\vee$ 为对偶完美复形，并回忆
$K = Rf_*(E^\vee \otimes_{\mathcal{O}_X}^\mathbf{L} \mathcal{G}^\bullet)$.
又有
$$
\Ext^i_{\mathcal{O}_X}(E,
\mathcal{G}^\bullet \otimes_{\mathcal{O}_X} f^*\mathcal{F})
=
H^i(X, E^\vee \otimes^\mathbf{L}_{\mathcal{O}_X}
(\mathcal{G}^\bullet \otimes_{\mathcal{O}_X} f^*\mathcal{F}))
$$
，这是由 $E^\vee$ 的构造得到的，所以对 $E^\vee$ 与 $\mathcal{G}^\bullet$
应用 Lemma \ref{perfect-lemma-compute-tensor-perfect} 即得这些同构的存在性。
关于边界映射的陈述意指如下：给定短正合列
$0 \to \mathcal{F}_1 \to \mathcal{F}_2 \to \mathcal{F}_3 \to 0$，
这些同构组成可交换图
$$
\xymatrix{
H^i(S, K \otimes^\mathbf{L}_{\mathcal{O}_S} \mathcal{F}_3)
\ar[r] \ar[d]_\delta &
\Ext^i_{\mathcal{O}_X}(E,
\mathcal{G}^\bullet \otimes_{\mathcal{O}_X} f^*\mathcal{F}_3) \ar[d]^\delta \\
H^{i + 1}(S, K \otimes^\mathbf{L}_{\mathcal{O}_S} \mathcal{F}_1)
\ar[r] &
\Ext^{i + 1}_{\mathcal{O}_X}(E,
\mathcal{G}^\bullet \otimes_{\mathcal{O}_X} f^*\mathcal{F}_1)
}
$$
，其中边界映射来自可区别三角
$$
K \otimes^\mathbf{L}_{\mathcal{O}_S} \mathcal{F}_1 \to
K \otimes^\mathbf{L}_{\mathcal{O}_S} \mathcal{F}_2 \to
K \otimes^\mathbf{L}_{\mathcal{O}_S} \mathcal{F}_3 \to
K \otimes^\mathbf{L}_{\mathcal{O}_S} \mathcal{F}_1[1]
$$
以及与 $D(\mathcal{O}_X)$ 中下述短正合列相伴的可区别三角：
$$
0 \to
\mathcal{G}^\bullet \otimes_{\mathcal{O}_X} f^*\mathcal{F}_1 \to
\mathcal{G}^\bullet \otimes_{\mathcal{O}_X} f^*\mathcal{F}_2 \to
\mathcal{G}^\bullet \otimes_{\mathcal{O}_X} f^*\mathcal{F}_3 \to 0
$$
。这是复形的短正合列。该列正合是因为 $\mathcal{G}$ 在 $S$ 上平坦。
略去验证所显示图的可交换性。
\end{proof}

\begin{lemma}
\label{perfect-lemma-compute-ext}
设 $f : X \to S$ 为概形态射，$E \in D(\mathcal{O}_X)$，且
$\mathcal{G}^\bullet$ 为 $\mathcal{O}_X$-模复形。假设：
\begin{enumerate}
\item $S$ 是 Noether 概形；
\item $f$ 局部有限型；
\item $E \in D^-_{\textit{Coh}}(\mathcal{O}_X)$；
\item $\mathcal{G}^\bullet$ 是凝聚 $\mathcal{O}_X$-模的有界复形，
它在 $S$ 上平坦且支集在 $S$ 上固有。
\end{enumerate}
则下列两个陈述成立：
\begin{enumerate}
\item[(A)] 对每个 $m \in \mathbf{Z}$，存在 $D(\mathcal{O}_S)$ 的完美对象 $K$
及函子性映射
$$
\alpha^i_\mathcal{F} :
\Ext^i_{\mathcal{O}_X}(E,
\mathcal{G}^\bullet \otimes_{\mathcal{O}_X} f^*\mathcal{F})
\longrightarrow
H^i(S, K \otimes^\mathbf{L}_{\mathcal{O}_S} \mathcal{F})
$$
，对 $S$ 上拟凝聚的 $\mathcal{F}$ 与边界映射相容（见证明），
并且当 $i \leq m$ 时 $\alpha^i_\mathcal{F}$ 是同构；
\item[(B)] 存在伪凝聚对象 $L \in D(\mathcal{O}_S)$ 及函子性同构
$$
\Ext^i_{\mathcal{O}_S}(L, \mathcal{F}) \longrightarrow
\Ext^i_{\mathcal{O}_X}(E,
\mathcal{G}^\bullet \otimes_{\mathcal{O}_X} f^*\mathcal{F})
$$
，对 $S$ 上拟凝聚的 $\mathcal{F}$ 与边界映射相容。
\end{enumerate}
\end{lemma}

\begin{proof}
先证明 (A)。
设 $\mathcal{G}^i$ 仅在 $i \in [a, b]$ 时非零。
我们可以把 $X$ 替换为
$\mathcal{G}^i$ 的支集之并的一个拟紧开邻域。
因此可以假设 $X$ 是诺特的。
此时 $X$ 与 $f$ 都是拟紧且拟分离的。
选取由完美复形 $P$ 给出的
$(X, E, -m - 1 + a)$
的一个逼近 $P \to E$
（其存在性由定理 \ref{perfect-theorem-approximation} 保证）。
于是诱导映射
$$
\Ext^i_{\mathcal{O}_X}(E,
\mathcal{G}^\bullet \otimes_{\mathcal{O}_X} f^*\mathcal{F})
\longrightarrow
\Ext^i_{\mathcal{O}_X}(P,
\mathcal{G}^\bullet \otimes_{\mathcal{O}_X} f^*\mathcal{F})
$$
在 $i \leq m$ 时是同构。事实上，此映射的核（相应地，余核）
是下列群的一个商（相应地，子模）：
$$
\Ext^i_{\mathcal{O}_X}(C,
\mathcal{G}^\bullet \otimes_{\mathcal{O}_X} f^*\mathcal{F})
\quad\text{分别}\quad
\Ext^{i + 1}_{\mathcal{O}_X}(C,
\mathcal{G}^\bullet \otimes_{\mathcal{O}_X} f^*\mathcal{F})
$$
其中 $C$ 是 $P \to E$ 的锥。由于 $C$ 在次数
$\geq -m - 1 + a$ 的上同调层均为零，导出范畴，
引理 \ref{derived-lemma-negative-exts} 表明，当
$i \leq m + 1$ 时这些 $\Ext$ 群为零。
这就把问题归结到 $E$ 是完美复形的情形，
而这正是引理 \ref{perfect-lemma-compute-ext-perfect}。
关于边界映射的断言已在
引理 \ref{perfect-lemma-compute-ext-perfect} 的证明中说明。

\medskip\noindent
证明 (B)。与 (A) 的证明相同，可以假设 $X$ 是诺特的。
由引理 \ref{perfect-lemma-identify-pseudo-coherent-noetherian}，
$E$ 是伪凝聚的。
根据引理 \ref{perfect-lemma-pseudo-coherent-hocolim}，可写成
$E = \text{hocolim} E_n$，其中 $E_n$ 完美，并且 $E_n \to E$
在截断 $\tau_{\geq -n}$ 上诱导同构。令 $E_n^\vee$
为其对偶完美复形
（上同调，引理 \ref{cohomology-lemma-dual-perfect-complex}）。
由此得到完美对象的逆系统
$\ldots \to E_3^\vee \to E_2^\vee \to E_1^\vee$。
它又给出逆系统
$$
\ldots \to K_3 \to K_2 \to K_1\quad\text{其中}\quad
K_n = Rf_*(E_n^\vee \otimes_{\mathcal{O}_X}^\mathbf{L} \mathcal{G}^\bullet)
$$
其各项在 $S$ 上都是完美的，见引理 \ref{perfect-lemma-tensor-perfect}。
由引理 \ref{perfect-lemma-compute-ext-perfect} 及其证明，
再结合上一段取 $P = E_n$ 时的论证，
对 $S$ 上任意拟凝聚层 $\mathcal{F}$，都有函子性的典范映射
$$
\xymatrix{
& \Ext^i_{\mathcal{O}_X}(E,
\mathcal{G}^\bullet \otimes_{\mathcal{O}_X} f^*\mathcal{F})
\ar[ld] \ar[rd] \\
H^i(S, K_{n + 1} \otimes_{\mathcal{O}_S}^\mathbf{L} \mathcal{F})
\ar[rr] & &
H^i(S, K_n \otimes_{\mathcal{O}_S}^\mathbf{L} \mathcal{F})
}
$$
它们在 $i \leq n + a$ 时是同构。
令 $L_n = K_n^\vee$ 为对偶完美复形。
于是 $L_1 \to L_2 \to L_3 \to \ldots$
是 $D(\mathcal{O}_S)$ 中的完美对象系统，
并且对 $S$ 上任意拟凝聚层 $\mathcal{F}$，映射
$$
\Ext^i_{\mathcal{O}_S}(L_{n + 1}, \mathcal{F})
\longrightarrow
\Ext^i_{\mathcal{O}_S}(L_n, \mathcal{F})
$$
在 $i \leq n + a - 1$ 时均为同构。这说明
$L_n \to L_{n + 1}$ 在截断
$\tau_{\geq -n - a + 2}$ 上诱导同构（提示：取
$L_n \to L_{n + 1}$ 的锥，并考察它最后一个非零上同调层）。
因此 $L = \text{hocolim} L_n$ 是伪凝聚的，见
引理 \ref{perfect-lemma-pseudo-coherent-hocolim}。同伦余极限的映射性质给出
$\Ext^i_{\mathcal{O}_S}(L, \mathcal{F}) =
\Ext^i_{\mathcal{O}_S}(L_n, \mathcal{F})$
在 $i \leq n + a - 3$ 时成立，证明完毕。
\end{proof}

\begin{remark}
\label{perfect-remark-base-change-of-L}
引理 \ref{perfect-lemma-compute-ext} (B) 中的伪凝聚复形 $L$
由该情形典范地确定。例如，按 (B) 构造 $L$
与基变换相容。换言之，给定笛卡尔图
$$
\xymatrix{
X' \ar[r]_{g'} \ar[d]_{f'} &
X \ar[d]^f \\
S' \ar[r]^g &
S
}
$$
以及 $S'$ 上的拟凝聚层 $\mathcal{F}'$，存在函子性的典范同构
$$
\Ext^i_{\mathcal{O}_{S'}}(Lg^*L, \mathcal{F}') \longrightarrow
\Ext^i_{\mathcal{O}_X}(L(g')^*E,
(g')^*\mathcal{G}^\bullet \otimes_{\mathcal{O}_{X'}} (f')^*\mathcal{F}')
$$
注意，在右侧我们对 $\mathcal{G}^\bullet$ {\bf 不}使用导出拉回。
若以后需要这种形式，我们将在此表述一个精确结果并给出详细证明。
\end{remark}






\section{极限与导出范畴}
\label{perfect-section-limits}

\noindent
本节汇集若干关于一个概形的导出范畴的结果，
其中该概形是某个概形逆系统的极限。
更确切地说，我们将在下述设定中工作。

\begin{situation}
\label{perfect-situation-descent}
设 $S = \lim_{i \in I} S_i$ 是一个概形有向系统的极限，
其转移态射 $f_{i'i} : S_{i'} \to S_i$ 均为仿射态射。
假设对所有 $i \in I$，$S_i$ 都是拟紧且拟分离的。
以 $f_i : S \to S_i$ 表示投影，并固定一个元素 $0 \in I$。
\end{situation}

\begin{lemma}
\label{perfect-lemma-descend-homomorphisms}
在情形 \ref{perfect-situation-descent} 中，令 $E_0$ 与 $K_0$ 为
$D(\mathcal{O}_{S_0})$ 的对象。
当 $i \geq 0$ 时置
$E_i = Lf_{i0}^*E_0$、$K_i = Lf_{i0}^*K_0$，
并置 $E = Lf_0^*E_0$、$K = Lf_0^*K_0$。则映射
$$
\colim_{i \geq 0} \Hom_{D(\mathcal{O}_{S_i})}(E_i, K_i)
\longrightarrow
\Hom_{D(\mathcal{O}_S)}(E, K)
$$
在下列任一情形下是同构：
\begin{enumerate}
\item $E_0$ 是完美的，且 $K_0 \in D_\QCoh(\mathcal{O}_{S_0})$；或者
\item $E_0$ 是伪凝聚的，且
$K_0 \in D_\QCoh(\mathcal{O}_{S_0})$ 具有有限 Tor 维数。
\end{enumerate}
\end{lemma}

\begin{proof}
对每个开集 $U_0 \subset S_0$，考虑条件 $P$：
典范映射
$$
\colim_{i \geq 0} \Hom_{D(\mathcal{O}_{U_i})}(E_i|_{U_i}, K_i|_{U_i})
\longrightarrow
\Hom_{D(\mathcal{O}_U)}(E|_U, K|_U)
$$
是同构，其中 $U = f_0^{-1}(U_0)$ 且 $U_i = f_{i0}^{-1}(U_0)$。
我们将利用概形的上同调，引理
\ref{coherent-lemma-induction-principle} 的归纳原理，
证明 $P$ 对所有拟紧开集 $U_0$ 成立。
该引理的条件 (2) 立即由导出范畴中的 Hom 的
Mayer--Vietoris 性质得到，见上同调，引理
\ref{cohomology-lemma-mayer-vietoris-hom}。
所以只需在 $S_0$ 仿射时证明本引理。

\medskip\noindent
假设 $S_0$ 仿射。记 $S_0 = \Spec(A_0)$、$S_i = \Spec(A_i)$，
以及 $S = \Spec(A)$。下文将不再说明地使用
引理 \ref{perfect-lemma-affine-compare-bounded}。

\medskip\noindent
在情形 (1) 中，对象 $E_0^\bullet$ 对应于有限生成投射
$A_0$-模的一个有限复形，见引理 \ref{perfect-lemma-perfect-affine}。
可以用 $A_0$-模的一个 K-平坦复形 $K_0^\bullet$
表示对象 $K_0$。在此情形中，需要证明
$$
\colim_{i \geq 0} \Hom_{D(A_i)}(E_0^\bullet \otimes_{A_0} A_i,
K_0^\bullet \otimes_{A_0} A_i)
\longrightarrow
\Hom_{D(A)}(E_0^\bullet \otimes_{A_0} A, K_0^\bullet \otimes_{A_0} A)
$$
由于 $E_0^\bullet$ 是投射模的上有界复形，
可将其改写为
$$
\colim_{i \geq 0} \Hom_{K(A_0)}(E_0^\bullet,
K_0^\bullet \otimes_{A_0} A_i)
\longrightarrow
\Hom_{K(A_0)}(E_0^\bullet, K_0^\bullet \otimes_{A_0} A)
$$
由于非零模 $E_0^n$ 只有有限多个，且它们都是有限表示模，
此映射是同构。

\medskip\noindent
在情形 (2) 中，对象 $E_0$ 对应于有限自由 $A_0$-模的一个
上有界复形 $E_0^\bullet$，见引理
\ref{perfect-lemma-pseudo-coherent-affine}。
可以用平坦 $A_0$-模的一个有限复形 $K_0^\bullet$
表示 $K_0$，见引理 \ref{perfect-lemma-tor-dimension-affine}
与代数续篇，引理 \ref{more-algebra-lemma-tor-amplitude}。
特别地，$K_0^\bullet$ 是 K-平坦的，因而可以像前面一样
归结到映射
$$
\colim_{i \geq 0} \Hom_{K(A_0)}(E_0^\bullet,
K_0^\bullet \otimes_{A_0} A_i)
\longrightarrow
\Hom_{K(A_0)}(E_0^\bullet, K_0^\bullet \otimes_{A_0} A)
$$
显然，此映射是同构（因为 $K_0^\bullet$ 有界，
其中只涉及有限多个项）。
\end{proof}

\begin{lemma}
\label{perfect-lemma-descend-perfect}
在情形 \ref{perfect-situation-descent} 中，$D(\mathcal{O}_S)$
的完美对象范畴是各 $D(\mathcal{O}_{S_i})$
的完美对象范畴的余极限。
\end{lemma}

\begin{proof}
对每个开集 $U_0 \subset S_0$，考虑条件 $P$：
函子
$$
\colim_{i \geq 0} D_{perf}(\mathcal{O}_{U_i})
\longrightarrow
D_{perf}(\mathcal{O}_U)
$$
是等价，其中 ${}_{perf}$ 表示完美对象构成的全子范畴，
且 $U = f_0^{-1}(U_0)$、$U_i = f_{i0}^{-1}(U_0)$。
我们将利用概形的上同调，引理
\ref{coherent-lemma-induction-principle} 的归纳原理，
证明 $P$ 对所有拟紧开集 $U_0$ 成立。
首先，由引理 \ref{perfect-lemma-descend-homomorphisms}，
已经知道该函子是全忠实的。因此只需证明本质满性。

\medskip\noindent
先验证归纳原理的条件 (2)。设
$S_0 = U_0 \cup V_0$，且 $P$ 对
$U_0$、$V_0$ 与 $U_0 \cap V_0$ 成立。令 $E$ 为
$D(\mathcal{O}_S)$ 的完美对象。可以找到 $i \geq 0$，
以及 $U_i$ 上的完美对象 $E_{U, i}$ 和 $V_i$ 上的完美对象
$E_{V, i}$，使它们到 $U$ 与 $V$ 的拉回分别同构于
$E|_U$ 与 $E|_V$。记
$$
a : E_{U, i} \to (Rf_{i, *}E)|_{U_i}
\quad\text{且}\quad
b : E_{V, i} \to (Rf_{i, *}E)|_{V_i}
$$
为同构 $Lf_i^*E_{U, i} \to E|_U$
与 $Lf_i^*E_{V, i} \to E|_V$ 的伴随映射。
由全忠实性，增大 $i$ 后，可以找到同构
$c : E_{U, i}|_{U_i \cap V_i} \to E_{V, i}|_{U_i \cap V_i}$
，其拉回为复合识别
$$
Lf_i^*E_{U, i}|_{U \cap V} \to E|_{U \cap V} \to Lf_i^*E_{V, i}|_{U \cap V}.
$$
应用上同调，引理 \ref{cohomology-lemma-glue}，
得到 $S_i$ 上的对象 $E_i$ 与映射 $d : E_i \to Rf_{i, *}E$，
其在 $U_i$ 与 $V_i$ 上分别限制为 $a$ 与 $b$。
于是显然 $E_i$ 是完美的，并且
$d$ 是某个同构 $Lf_i^*E_i \to E$ 的伴随映射。

\medskip\noindent
最后验证归纳原理的条件 (1)，即验证
$S_0$ 仿射时本引理成立。
记 $S_0 = \Spec(A_0)$、$S_i = \Spec(A_i)$，以及
$S = \Spec(A)$。利用引理 \ref{perfect-lemma-affine-compare-bounded}
与 \ref{perfect-lemma-perfect-affine}，只需证明
$$
D_{perf}(A) = \colim D_{perf}(A_i)
$$
这由以下事实立即可得：环上的完美复形由有限生成投射模
（因而是有限表示模）的有限复形给出。细节见代数续篇，
引理 \ref{more-algebra-lemma-colimit-perfect-complexes}。
\end{proof}





\section{上同调与基变换 VI}
\label{perfect-section-cohomology-and-base-change-final}

\noindent
这是关于上同调与基变换的最后一节，继续各节
\ref{perfect-section-cohomology-and-base-change-perfect},
\ref{perfect-section-cohomology-base-change} 与
\ref{perfect-section-producing-perfect} 的讨论。
评注 \ref{perfect-remark-explain-perfect-direct-image}
给出了一个容易理解的特殊情形。

\begin{lemma}
\label{perfect-lemma-base-change-tensor-perfect}
设 $f : X \to S$ 是有限表示态射，
$E \in D(\mathcal{O}_X)$ 是完美对象。设 $\mathcal{G}^\bullet$
是有限表示 $\mathcal{O}_X$-模的有界复形，
它在 $S$ 上平坦且支集在 $S$ 上固有。则
$$
K = Rf_*(E \otimes_{\mathcal{O}_X}^\mathbf{L} \mathcal{G}^\bullet)
$$
是 $D(\mathcal{O}_S)$ 的完美对象，且其构造与任意基变换相容。
\end{lemma}

\begin{proof}
关于基变换的断言就是引理 \ref{perfect-lemma-base-change-tensor}。
所以只需证明 $K$ 是完美对象。当 $S$ 是诺特概形时，
这由引理 \ref{perfect-lemma-tensor-perfect} 得到。
我们将利用诺特逼近归结到这种情形。
建议读者跳过本证明的余下部分。

\medskip\noindent
问题在 $S$ 上是局部的，所以可以假设 $S$ 仿射。
记 $S = \Spec(R)$。把 $R = \colim R_i$
写成诺特环 $R_i$ 的滤过余极限。由极限，引理
\ref{limits-lemma-descend-finite-presentation}
，存在某个 $i$ 及一个在 $R_i$ 上有限表示的概形 $X_i$，
其到 $R$ 的基变换为 $X$。由极限，引理
\ref{limits-lemma-descend-modules-finite-presentation}，
增大 $i$ 后，可以假设存在有限表示 $\mathcal{O}_{X_i}$-模的
一个有界复形 $\mathcal{G}_i^\bullet$，
其到 $X$ 的拉回为 $\mathcal{G}^\bullet$。
再次增大 $i$，可以假设 $\mathcal{G}_i^n$ 在 $R_i$ 上平坦，
见极限，引理 \ref{limits-lemma-descend-module-flat-finite-presentation}。
再增大 $i$，可以假设 $\mathcal{G}_i^n$ 的支集在 $R_i$ 上固有，
见极限，引理 \ref{limits-lemma-eventually-proper-support}
与概形的上同调，引理
\ref{coherent-lemma-module-support-proper-over-base}。
最后，由引理 \ref{perfect-lemma-descend-perfect}，增大 $i$ 后，
可以假设存在 $D(\mathcal{O}_{X_i})$ 的完美对象 $E_i$，
其到 $X$ 的拉回为 $E$。把引理 \ref{perfect-lemma-tensor-perfect}
应用于 $X_i \to \Spec(R_i)$、$E_i$、$\mathcal{G}_i^\bullet$，
并利用已经证明的基变换性质，即得结论。
\end{proof}

\begin{remark}
\label{perfect-remark-explain-perfect-direct-image}
设 $R$ 是环，$X$ 是在 $R$ 上有限表示的概形。
设 $\mathcal{G}$ 是有限表示 $\mathcal{O}_X$-模，
它在 $R$ 上平坦且支集在 $R$ 上固有。由
引理 \ref{perfect-lemma-base-change-tensor-perfect}，
存在有限生成投射 $R$-模的有限复形 $M^\bullet$，使得
$$
R\Gamma(X_{R'}, \mathcal{G}_{R'}) = M^\bullet \otimes_R R'
$$
对 $R$-代数 $R'$ 函子性地成立。
\end{remark}

\begin{lemma}
\label{perfect-lemma-base-change-tensor-pseudo-coherent}
设 $f : X \to S$ 是有限表示态射，
$E \in D(\mathcal{O}_X)$ 是伪凝聚对象。
设 $\mathcal{G}^\bullet$ 是有限表示 $\mathcal{O}_X$-模的上有界复形，
它在 $S$ 上平坦且支集在 $S$ 上固有。则
$$
K = Rf_*(E \otimes_{\mathcal{O}_X}^\mathbf{L} \mathcal{G}^\bullet)
$$
是 $D(\mathcal{O}_S)$ 的伪凝聚对象，且其构造与任意基变换相容。
\end{lemma}

\begin{proof}
关于基变换的断言就是引理 \ref{perfect-lemma-base-change-tensor}。
所以只需证明 $K$ 是伪凝聚对象。
这将由完美复形逼近及引理
\ref{perfect-lemma-base-change-tensor-perfect} 得到。
建议读者跳过本证明的余下部分。

\medskip\noindent
问题在 $S$ 上是局部的，所以可以假设 $S$ 仿射。
于是 $X$ 拟紧且拟分离。此外，存在整数 $N$，使全直像
$Rf_* : D_\QCoh(\mathcal{O}_X) \to D_\QCoh(\mathcal{O}_S)$
具有上同调维数 $N$，见引理
\ref{perfect-lemma-quasi-coherence-direct-image} 中的说明。
选取整数 $b$，使得当 $i > b$ 时 $\mathcal{G}^i = 0$。
只需证明对每个 $m$，$K$ 都是 $m$-伪凝聚的。
选取由完美复形 $P$ 给出的 $(X, E, m - N - 1 - b)$
的一个逼近 $P \to E$。其存在性由
定理 \ref{perfect-theorem-approximation} 保证。
在 $D_\QCoh(\mathcal{O}_X)$ 中选取可区别三角
$$
P \to E \to C \to P[1]
$$
。$C$ 的上同调层在次数 $\geq m - N - 1 - b$ 均为零。
因此 $C \otimes^\mathbf{L} \mathcal{G}^\bullet$
的上同调层在次数 $\geq m - N - 1$ 均为零。于是
$Rf_*(C \otimes^\mathbf{L} \mathcal{G}^\bullet)$
的上同调层在次数 $\geq m - 1$ 均为零。因此
$$
Rf_*(P \otimes^\mathbf{L} \mathcal{G}^\bullet) \to
Rf_*(E \otimes^\mathbf{L} \mathcal{G}^\bullet)
$$
在次数 $\geq m$ 的上同调层上是同构。
其次，设当 $i > a$ 时 $H^i(P) = 0$。则
$
P \otimes^\mathbf{L} \sigma_{\geq m - N - 1 - a}\mathcal{G}^\bullet
\longrightarrow
P \otimes^\mathbf{L} \mathcal{G}^\bullet
$
在次数 $\geq m - N - 1$ 的上同调层上是同构。
因而再次得到
$$
Rf_*(P \otimes^\mathbf{L} \sigma_{\geq m - N - 1 - a}\mathcal{G}^\bullet) \to
Rf_*(P \otimes^\mathbf{L} \mathcal{G}^\bullet)
$$
在次数 $\geq m$ 的上同调层上是同构。
由引理 \ref{perfect-lemma-base-change-tensor-perfect}，其源是完美复形。
因此 $K$ 如所需是 $m$-伪凝聚的。
\end{proof}

\begin{lemma}
\label{perfect-lemma-flat-proper-perfect-direct-image-general}
设 $S$ 是概形，$f : X \to S$ 是有限表示的固有态射。
\begin{enumerate}
\item 若 $E \in D(\mathcal{O}_X)$ 完美且 $f$ 平坦，则
$Rf_*E$ 是 $D(\mathcal{O}_S)$ 的完美对象，且其构造与任意基变换相容。
\item 若 $\mathcal{G}$ 是有限表示的 $\mathcal{O}_X$-模且在 $S$ 上平坦，
则 $Rf_*\mathcal{G}$ 是 $D(\mathcal{O}_S)$ 的完美对象，
且其构造与任意基变换相容。
\end{enumerate}
\end{lemma}

\begin{proof}
这是引理 \ref{perfect-lemma-base-change-tensor-perfect} 的特殊情形：
(1) 取 $\mathcal{G}^\bullet$ 为置于次数 $0$ 的 $\mathcal{O}_X$；
(2) 取 $E = \mathcal{O}_X$，并令 $\mathcal{G}^\bullet$
为置于次数 $0$ 的 $\mathcal{G}$。
\end{proof}

\begin{lemma}
\label{perfect-lemma-flat-proper-pseudo-coherent-direct-image-general}
设 $S$ 是概形，$f : X \to S$ 是有限表示的平坦固有态射。
设 $E \in D(\mathcal{O}_X)$ 是伪凝聚的。则 $Rf_*E$
是 $D(\mathcal{O}_S)$ 的伪凝聚对象，且其构造与任意基变换相容。
\end{lemma}

\noindent
更一般地，若 $f : X \to S$ 固有，且 $X$ 上的 $E$
相对于 $S$ 伪凝聚（态射续篇，定义
\ref{more-morphisms-definition-relative-pseudo-coherence}),
则 $Rf_*E$ 是伪凝聚的
（但在这种一般性下，其构造不与基变换相容）。
见 \cite{Kiehl}。

\begin{proof}
这是引理 \ref{perfect-lemma-base-change-tensor-pseudo-coherent} 中
取 $\mathcal{G}^\bullet$ 为置于次数 $0$ 的 $\mathcal{O}_X$
所得的特殊情形。
\end{proof}

\begin{lemma}
\label{perfect-lemma-pullback-and-limits}
设 $R$ 是环，$X$ 是概形，并设
$f : X \to \Spec(R)$ 是有限表示的平坦固有态射。
设 $(M_n)$ 为转移映射均满射的 $R$-模逆系统。
则典范映射
$$
\mathcal{O}_X \otimes_R (\lim M_n)
\longrightarrow
\lim \mathcal{O}_X \otimes_R M_n
$$
诱导从其源到对其目标应用 $DQ_X$ 后所得对象的同构。
\end{lemma}

\begin{proof}
该断言意指，对 $D_\QCoh(\mathcal{O}_X)$ 的任意对象 $E$，
诱导映射
$$
\Hom(E, \mathcal{O}_X \otimes_R (\lim M_n))
\longrightarrow
\Hom(E, \lim \mathcal{O}_X \otimes_R M_n)
$$
是同构。由于 $D_\QCoh(\mathcal{O}_X)$ 有完美生成元
（定理 \ref{perfect-theorem-bondal-van-den-Bergh}），
只需对完美的 $E$ 验证此事。
由引理 \ref{perfect-lemma-Rlim-quasi-coherent}，
$\lim \mathcal{O}_X \otimes_R M_n = R\lim \mathcal{O}_X \otimes_R M_n$.
正合函子
$R\Hom_X(E, -) : D_\QCoh(\mathcal{O}_X) \to D(R)$
（见上同调，第 \ref{cohomology-section-global-RHom} 节）
与积相容，因而也与导出极限相容，所以
$$
R\Hom_X(E, \lim \mathcal{O}_X \otimes_R M_n) =
R\lim R\Hom_X(E, \mathcal{O}_X \otimes_R M_n)
$$
令 $E^\vee$ 为对偶完美复形，见上同调，引理
\ref{cohomology-lemma-dual-perfect-complex}。有
$$
R\Hom_X(E, \mathcal{O}_X \otimes_R M_n) =
R\Gamma(X, E^\vee \otimes_{\mathcal{O}_X}^\mathbf{L} Lf^*M_n) =
R\Gamma(X, E^\vee) \otimes_R^\mathbf{L} M_n
$$
，这里使用了引理 \ref{perfect-lemma-cohomology-base-change}。
由引理 \ref{perfect-lemma-flat-proper-perfect-direct-image-general}，
$R\Gamma(X, E^\vee)$ 是 $R$-模的完美复形。
特别地，它是伪凝聚复形，因而由代数续篇，引理
\ref{more-algebra-lemma-pseudo-coherent-tensor-limit} 得到
$$
R\lim R\Gamma(X, E^\vee) \otimes_R^\mathbf{L} M_n =
R\Gamma(X, E^\vee) \otimes_R^\mathbf{L} \lim M_n
$$
，正如所需。
\end{proof}

\begin{lemma}
\label{perfect-lemma-base-change-RHom-perfect}
设 $f : X \to S$ 是有限表示态射，
$E \in D(\mathcal{O}_X)$ 是完美对象。设 $\mathcal{G}^\bullet$
是有限表示 $\mathcal{O}_X$-模的有界复形，
它在 $S$ 上平坦且支集在 $S$ 上固有。则
$$
K = Rf_*R\SheafHom(E, \mathcal{G}^\bullet)
$$
是 $D(\mathcal{O}_S)$ 的完美对象，且其构造与任意基变换相容。
\end{lemma}

\begin{proof}
关于基变换的断言就是引理 \ref{perfect-lemma-base-change-RHom}。
所以只需证明 $K$ 是完美对象。当 $S$ 是诺特概形时，
这由引理 \ref{perfect-lemma-ext-perfect} 得到。
我们将利用诺特逼近归结到这种情形。
建议读者跳过本证明的余下部分。

\medskip\noindent
问题在 $S$ 上是局部的，所以可以假设 $S$ 仿射。
记 $S = \Spec(R)$。把 $R = \colim R_i$
写成诺特环 $R_i$ 的滤过余极限。由极限，引理
\ref{limits-lemma-descend-finite-presentation}
，存在某个 $i$ 及一个在 $R_i$ 上有限表示的概形 $X_i$，
其到 $R$ 的基变换为 $X$。由极限，引理
\ref{limits-lemma-descend-modules-finite-presentation}，
增大 $i$ 后，可以假设存在有限表示 $\mathcal{O}_{X_i}$-模的
一个有界复形 $\mathcal{G}_i^\bullet$，
其到 $X$ 的拉回为 $\mathcal{G}^\bullet$。
再次增大 $i$，可以假设 $\mathcal{G}_i^n$ 在 $R_i$ 上平坦，
见极限，引理 \ref{limits-lemma-descend-module-flat-finite-presentation}。
再增大 $i$，可以假设 $\mathcal{G}_i^n$ 的支集在 $R_i$ 上固有，
见极限，引理 \ref{limits-lemma-eventually-proper-support}
与概形的上同调，引理
\ref{coherent-lemma-module-support-proper-over-base}。
最后，由引理 \ref{perfect-lemma-descend-perfect}，增大 $i$ 后，
可以假设存在 $D(\mathcal{O}_{X_i})$ 的完美对象 $E_i$，
其到 $X$ 的拉回为 $E$。把引理 \ref{perfect-lemma-ext-perfect}
应用于 $X_i \to \Spec(R_i)$、$E_i$、$\mathcal{G}_i^\bullet$，
并利用已经证明的基变换性质，即得结论。
\end{proof}









\section{完美复形}
\label{perfect-section-perfect-complexes}

\noindent
先讨论完美复形 Betti 数的跳跃轨迹。
给定概形 $X$ 上的复形 $E$ 与 $X$ 的点 $x$，
我们常以 $E \otimes_{\mathcal{O}_X}^\mathbf{L} \kappa(x)$
代替更准确的写法 $Li_x^*E$，其中
$i_x : x \to X$ 是典范态射。

\begin{lemma}
\label{perfect-lemma-jump-loci}
设 $X$ 是概形，$E \in D(\mathcal{O}_X)$ 伪凝聚
（例如它是完美的）。对任意 $i \in \mathbf{Z}$，考虑函数
$$
\beta_i : X \longrightarrow \{0, 1, 2, \ldots\},\quad
x \longmapsto
\dim_{\kappa(x)}
H^i(E \otimes_{\mathcal{O}_X}^\mathbf{L} \kappa(x))
$$
则有：
\begin{enumerate}
\item $\beta_i$ 的构造与任意基变换相容；
\item 函数 $\beta_i$ 上半连续；
\item $\beta_i$ 的等值集在 $X$ 中局部可构造。
\end{enumerate}
\end{lemma}

\begin{proof}
考虑概形态射 $f : Y \to X$ 与点 $y \in Y$。
令 $x$ 为 $y$ 的像，并考虑交换图
$$
\xymatrix{
y \ar[r]_j \ar[d]_g & Y \ar[d]^f \\
x \ar[r]^i & X
}
$$
于是 $Lg^* \circ Li^* = Lj^* \circ Lf^*$。
这说明，与伪凝聚复形 $Lf^*E$ 相伴的函数 $\beta'_i$
是函数 $\beta_i$ 的拉回，即
$\beta'_i = \beta_i \circ f$。这正是 (1) 的含义。

\medskip\noindent
固定 $i$ 并取 $x \in X$。只需在 $x$ 的某个开邻域中
证明 (2) 与 (3)，故可假设 $X$ 仿射。
于是可用有限自由模的一个上有界复形 $\mathcal{F}^\bullet$
表示 $E$（引理 \ref{perfect-lemma-lift-pseudo-coherent}）。
此时 $P = \sigma_{\geq i - 1}\mathcal{F}^\bullet$ 是完美对象，
且 $P \to E$ 对所有 $x' \in X$ 都诱导同构
$$
H^i(P \otimes_{\mathcal{O}_X}^\mathbf{L} \kappa(x')) \to
H^i(E \otimes_{\mathcal{O}_X}^\mathbf{L} \kappa(x'))
$$
。因此可以假设 $E$ 完美。在这种情形中，由代数续篇，引理
\ref{more-algebra-lemma-lift-perfect-from-residue-field}
，存在 $x$ 的仿射开邻域 $U$ 以及 $a \leq b$，
使 $E|_U$ 可由复形
$$
\ldots \to 0 \to \mathcal{O}_U^{\oplus \beta_a(x)}
\to \mathcal{O}_U^{\oplus \beta_{a + 1}(x)} \to
\ldots \to
\mathcal{O}_U^{\oplus \beta_{b - 1}(x)} \to
\mathcal{O}_U^{\oplus \beta_b(x)} \to 0
\to \ldots
$$
表示。（这里还使用了把问题化为代数问题的较早结果，例如
引理 \ref{perfect-lemma-affine-compare-bounded} 与
\ref{perfect-lemma-perfect-affine}。）
立即可见，对所有 $x' \in U$，
$\beta_i(x') \leq \beta_i(x)$。这证明了 $\beta_i$ 上半连续。

\medskip\noindent
为证明 (3)，可以假设 $X$ 仿射，并且
$E$ 由有限自由 $\mathcal{O}_X$-模的一个复形给出
（例如沿用上一段的论证，或使用上同调，引理
\ref{cohomology-lemma-perfect-on-locally-ringed}）。
因此需要证明：给定复形
$$
\mathcal{O}_X^{\oplus a} \to
\mathcal{O}_X^{\oplus b} \to
\mathcal{O}_X^{\oplus c}
$$
，把点 $x \in X$ 映到
$\kappa_x^{\oplus a} \to \kappa_x^{\oplus b} \to \kappa_x^{\oplus c}$
在中间项处上同调维数的函数具有可构造等值集。
令 $A \in \text{Mat}(a \times b, \Gamma(X, \mathcal{O}_X))$
为第一个箭头的矩阵。$A$ 在
$\text{Mat}(a \times b, \kappa(x))$ 中的像秩等于 $r$，
当且仅当 $A$ 的所有 $(r + 1) \times (r + 1)$ 子式在 $x$ 处为零，
并且 $A$ 的某个 $r \times r$ 子式在 $x$ 处非零。
所以秩等于 $r$ 的点集是可构造的局部闭集。
对第二个箭头作同样论证，再把这些结果合在一起，即得所需结论。
\end{proof}

\begin{lemma}
\label{perfect-lemma-chi-locally-constant}
设 $X$ 是概形，$E \in D(\mathcal{O}_X)$ 完美。函数
$$
\chi_E : X \longrightarrow \mathbf{Z},\quad
x \longmapsto \sum (-1)^i
\dim_{\kappa(x)} H^i(E \otimes_{\mathcal{O}_X}^\mathbf{L} \kappa(x))
$$
在 $X$ 上局部常值。
\end{lemma}

\begin{proof}
由上同调，引理
\ref{cohomology-lemma-perfect-on-locally-ringed}
可知，在 $X$ 上局部地，可以用有限自由
$\mathcal{O}_X$-模的有限复形 $\mathcal{E}^\bullet$ 表示 $E$。
在这样的开集上，函数 $\chi_E$ 为常值，其值是
$\sum (-1)^i \text{rank}(\mathcal{E}^i)$.
\end{proof}

\begin{lemma}
\label{perfect-lemma-open-where-cohomology-in-degree-i-rank-r}
设 $X$ 是概形，$E \in D(\mathcal{O}_X)$ 完美。
给定 $i, r \in \mathbf{Z}$，存在开子概形 $U \subset X$，
由下列性质刻画：
\begin{enumerate}
\item $E|_U \cong H^i(E|_U)[-i]$，且 $H^i(E|_U)$
是秩为 $r$ 的局部自由 $\mathcal{O}_U$-模；
\item 态射 $f : Y \to X$ 经过 $U$ 分解，当且仅当
$Lf^*E$ 同构于置于次数 $i$ 的一个秩为 $r$ 的局部自由模。
\end{enumerate}
\end{lemma}

\begin{proof}
对 $j \in \mathbf{Z}$，令
$\beta_j : X \to \{0, 1, 2, \ldots\}$
为引理 \ref{perfect-lemma-jump-loci} 中的函数。则集合
$$
W = \{x \in X \mid \beta_j(x) \leq 0\text{ 对所有 }j \not = i\}
$$
在 $X$ 中是开集，且其构造与到任意 $X$-概形 $Y$ 的拉回相容。
这是该引理的推论，因为先验地，在任一点的某个邻域中
只有有限多个 $\beta_j$ 非零。因此可以用 $W$ 替换 $X$，
并假设对所有 $x \in X$ 与所有 $j \not = i$，
$\beta_j(x) = 0$。在这种情形中，$H^i(E)$ 是有限局部自由模，
且 $E \cong H^i(E)[-i]$；例如见代数续篇，引理
\ref{more-algebra-lemma-lift-perfect-from-residue-field}。
于是 $X$ 是若干开子概形的不交并，并且 $H^i(E)$
在每个这样的开子概形上具有固定的秩，结论得证。
\end{proof}

\begin{lemma}
\label{perfect-lemma-locally-closed-where-H0-locally-free}
设 $X$ 是概形。设 $E \in D(\mathcal{O}_X)$ 完美，
并且对某些 $a, b \in \mathbf{Z}$，其 Tor 振幅包含于 $[a, b]$。
令 $r \geq 0$。则存在局部闭子概形 $j : Z \to X$，
由下列性质刻画：
\begin{enumerate}
\item $H^a(Lj^*E)$ 是秩为 $r$ 的局部自由 $\mathcal{O}_Z$-模；
\item 态射 $f : Y \to X$ 经过 $Z$ 分解，当且仅当
对所有态射 $g : Y' \to Y$，$\mathcal{O}_{Y'}$-模
$H^a(L(f \circ g)^*E)$ 都是秩为 $r$ 的局部自由模。
\end{enumerate}
此外，$j : Z \to X$ 是有限表示的，并且：
\begin{enumerate}
\item[(3)] 若 $f : Y \to X$ 分解为 $Y \xrightarrow{g} Z \to X$，则
$H^a(Lf^*E) = g^*H^a(Lj^*E)$,
\item[(4)] 若对所有 $x \in X$ 都有 $\beta_a(x) \leq r$，
则 $j$ 是闭浸入，并且对给定的 $f : Y \to X$，下列条件等价：
\begin{enumerate}
\item $f : Y \to X$ 经过 $Z$ 分解；
\item $H^a(Lf^*E)$ 是秩为 $r$ 的局部自由 $\mathcal{O}_Y$-模；
\end{enumerate}
若 $r = 1$，它们还等价于
\begin{enumerate}
\item[(c)] $\mathcal{O}_Y \to \SheafHom_{\mathcal{O}_Y}(H^a(Lf^*E), H^a(Lf^*E))$
是单射。
\end{enumerate}
\end{enumerate}
\end{lemma}

\begin{proof}
首先，令 $U \subset X$ 为局部可构造开子概形，
在其上引理 \ref{perfect-lemma-jump-loci} 的函数
$\beta_a$ 取值 $\leq r$。令 $f : Y \to X$ 如 (2) 所述。
则对任意 $y \in Y$，都有 $\beta_a(Lf^*E) = r$，
故由引理 \ref{perfect-lemma-jump-loci}，$y$ 映入 $U$。
所以 (2) 中的 $f$ 经过 $U$ 分解。
因此可以用 $U$ 替换 $X$，并假设对所有 $x \in X$，
$\beta_a(x) \in \{0, 1, \ldots, r\}$。
我们将证明，在这种情形中，存在由有限型拟凝聚理想截出的
闭子概形 $Z \subset X$，它由 (4) 中 (a)、(b)
（以及当 $r = 1$ 时的 (4)(c)）的等价性刻画，并且 (3) 成立。
这将完成证明，因为它尤其说明 (2) 中的态射都经过 $Z$ 分解。

\medskip\noindent
若 $x \in X$ 且 $\beta_a(x) < r$，则存在 $x$ 的开邻域，
在其上 $\beta_a < r$（引理 \ref{perfect-lemma-jump-loci}）。
由此可见，至少在集合论意义下，$Z$ 是闭子集。

\medskip\noindent
为赋予它概形论结构，考虑满足 $\beta_a(x) = r$ 的点
$x \in X$，并置 $\beta = \beta_{a + 1}(x)$。
由代数续篇，引理
\ref{more-algebra-lemma-lift-perfect-from-residue-field}
，存在 $x$ 的仿射开邻域 $U$，使得 $K|_U$ 可由复形
$$
\ldots \to 0 \to \mathcal{O}_U^{\oplus r}
\xrightarrow{(f_{ij})} \mathcal{O}_U^{\oplus \beta} \to
\ldots \to
\mathcal{O}_U^{\oplus \beta_{b - 1}(x)} \to
\mathcal{O}_U^{\oplus \beta_b(x)} \to 0
\to \ldots
$$
表示。（这里还使用了把问题化为代数问题的较早结果，例如
引理 \ref{perfect-lemma-affine-compare-bounded} 与
\ref{perfect-lemma-perfect-affine}。）现在，若 $g : Y \to U$ 是概形态射，
并且对某一对 $i, j$，$g^\sharp(f_{ij})$ 非零，则
$H^a(Lg^*E)$ 不是秩为 $r$ 的局部自由 $\mathcal{O}_Y$-模。
见代数续篇，引理 \ref{more-algebra-lemma-coker-injective-free}。
若对所有 $i, j$ 都有 $g^\sharp(f_{ij}) = 0$，
则 $H^a(Lg^*E)$ 显然是局部自由 $\mathcal{O}_Y$-模。
因此在 $U$ 上，由所有 $f_{ij}$ 截出的闭子概形满足 (3)，
并且 (4)(a) 与 (b) 等价。
$Z$ 的刻画说明这些局部构造的部分可以粘合（略去细节）。
最后，若 $r = 1$，则 (4)(c) 与 (4)(b) 等价，
因为此时局部地，$H^a(Lg^*E) \subset \mathcal{O}_Y$
是由各元素 $g^\sharp(f_{ij})$ 生成的理想的零化子。
\end{proof}





\section{应用}
\label{perfect-section-applications}

\noindent
这里大多是上同调与基变换的应用。将来可以把这些结果
推广到引理 \ref{perfect-lemma-base-change-tensor-perfect} 所讨论的情形。

\begin{lemma}
\label{perfect-lemma-jump-loci-geometric}
设 $f : X \to S$ 是有限表示的固有态射。
设 $\mathcal{F}$ 是有限表示的 $\mathcal{O}_X$-模，
并且在 $S$ 上平坦。对固定的 $i \in \mathbf{Z}$，考虑函数
$$
\beta_i : S \to \{0, 1, 2, \ldots\},\quad
s \longmapsto \dim_{\kappa(s)} H^i(X_s, \mathcal{F}_s)
$$
则有：
\begin{enumerate}
\item $\beta_i$ 的构造与任意基变换相容；
\item 函数 $\beta_i$ 上半连续；
\item $\beta_i$ 的等值集在 $S$ 中局部可构造。
\end{enumerate}
\end{lemma}

\begin{proof}
由上同调与基变换（更确切地说，由引理
\ref{perfect-lemma-flat-proper-perfect-direct-image-general}），
对象 $K = Rf_*\mathcal{F}$ 是 $S$ 的导出范畴中的完美对象，
且其构造与任意基变换相容。特别地，有
$$
H^i(X_s, \mathcal{F}_s) = H^i(K \otimes_{\mathcal{O}_S}^\mathbf{L} \kappa(s))
$$
因此本引理由引理 \ref{perfect-lemma-jump-loci} 得到。
\end{proof}

\begin{lemma}
\label{perfect-lemma-chi-locally-constant-geometric}
设 $f : X \to S$ 是有限表示的固有态射。
设 $\mathcal{F}$ 是有限表示的 $\mathcal{O}_X$-模，
并且在 $S$ 上平坦。函数
$$
s \longmapsto \chi(X_s, \mathcal{F}_s)
$$
在 $S$ 上局部常值，且此函数的构造与基变换相容。
\end{lemma}

\begin{proof}
由上同调与基变换（更确切地说，由引理
\ref{perfect-lemma-flat-proper-perfect-direct-image-general}），
对象 $K = Rf_*\mathcal{F}$ 是 $S$ 的导出范畴中的完美对象，
且其构造与任意基变换相容。因此需要证明映射
$$
s \longmapsto \sum (-1)^i \dim_{\kappa(s)}
H^i(K \otimes^\mathbf{L}_{\mathcal{O}_S} \kappa(s))
$$
在 $S$ 上局部常值。这正是引理 \ref{perfect-lemma-chi-locally-constant}。
\end{proof}

\begin{lemma}
\label{perfect-lemma-open-where-cohomology-in-degree-i-rank-r-geometric}
设 $f : X \to S$ 是有限表示的固有态射。
设 $\mathcal{F}$ 是有限表示的 $\mathcal{O}_X$-模，
并且在 $S$ 上平坦。固定 $i, r \in \mathbf{Z}$。
则存在具有下述性质的开子概形 $U \subset S$：
态射 $T \to S$ 经过 $U$ 分解，当且仅当
$Rf_{T, *}\mathcal{F}_T$ 同构于置于次数 $i$
的秩为 $r$ 的有限局部自由模。
\end{lemma}

\begin{proof}
由上同调与基变换（更确切地说，由引理
\ref{perfect-lemma-flat-proper-perfect-direct-image-general}），
对象 $K = Rf_*\mathcal{F}$ 是 $S$ 的导出范畴中的完美对象，
且其构造与任意基变换相容。因此本引理立即由
引理 \ref{perfect-lemma-open-where-cohomology-in-degree-i-rank-r} 得到。
\end{proof}

\begin{lemma}
\label{perfect-lemma-vanishing-implies-locally-free}
设 $f : X \to S$ 是有限表示态射。
设 $\mathcal{F}$ 是有限表示的 $\mathcal{O}_X$-模，
它在 $S$ 上平坦且支集在 $S$ 上固有。若对 $i > 0$，
$R^if_*\mathcal{F} = 0$，则 $f_*\mathcal{F}$ 局部自由，
且其构造与任意基变换相容（说明见证明）。
\end{lemma}

\begin{proof}
由引理 \ref{perfect-lemma-base-change-tensor-perfect}，
$D(\mathcal{O}_S)$ 的对象 $E = Rf_*\mathcal{F}$ 完美，
且其构造与任意基变换相容；这意味着，对任意笛卡尔图
$$
\xymatrix{
X' \ar[r]_{g'} \ar[d]_{f'} &
X \ar[d]^f \\
S' \ar[r]^g &
S
}
$$
，都有 $Rf'_*(g')^*\mathcal{F} = Lg^*E$。
由于次数 $< 0$ 绝不会出现上同调，可知对某个 $b$，
$E$（局部地）具有包含于 $[0, b]$ 的 Tor 振幅。
若对 $i > 0$，$H^i(E) = R^if_*\mathcal{F} = 0$，
则 $E$ 的 Tor 振幅包含于 $[0, 0]$。所以
$E = H^0(E)[0]$。由代数续篇，引理
\ref{more-algebra-lemma-perfect}
（以及代数，引理 \ref{algebra-lemma-finite-projective}
中对有限投射模的刻画），
$H^0(E) = f_*\mathcal{F}$ 是有限局部自由的。
与基变换相容意味着，对上述图有
$g^*f_*\mathcal{F} = f'_*(g')^*\mathcal{F}$，
而这由已经确立的 $E$ 与基变换相容性得到。
\end{proof}

\begin{lemma}
\label{perfect-lemma-proper-flat-h0}
设 $f : X \to S$ 是概形态射。假设
\begin{enumerate}
\item $f$ 固有、平坦且有限表示；
\item 对所有 $s \in S$，都有
$\kappa(s) = H^0(X_s, \mathcal{O}_{X_s})$。
\end{enumerate}
则有：
\begin{enumerate}
\item[(a)] $f_*\mathcal{O}_X = \mathcal{O}_S$，且在任意基变换后仍成立；
\item[(b)] 在 $S$ 上局部地，有
$$
Rf_*\mathcal{O}_X = \mathcal{O}_S \oplus P
$$
作为 $D(\mathcal{O}_S)$ 中的等式，
其中 $P$ 完美且 Tor 振幅包含于 $[1, \infty)$。
\end{enumerate}
\end{lemma}

\begin{proof}
由上同调与基变换
（引理 \ref{perfect-lemma-flat-proper-perfect-direct-image-general}），
复形 $E = Rf_*\mathcal{O}_X$ 完美，
且其构造与任意基变换相容。这首先说明 $E$
的 Tor 振幅包含于 $[0, \infty)$。其次，它说明对 $s \in S$ 有
$H^0(E \otimes^\mathbf{L} \kappa(s)) =
H^0(X_s, \mathcal{O}_{X_s}) = \kappa(s)$.
所以映射 $\mathcal{O}_S \to Rf_*\mathcal{O}_X = E$
诱导同构
$\mathcal{O}_S \otimes \kappa(s) \to H^0(E \otimes^\mathbf{L} \kappa(s))$.
因此
$H^0(E) \otimes \kappa(s) \to H^0(E \otimes^\mathbf{L} \kappa(s))$
是满射。应用代数续篇，引理
\ref{more-algebra-lemma-better-cut-complex-in-two} 可知，
把 $S$ 替换为 $s$ 的一个仿射开邻域后，有分解
$E = H^0(E) \oplus \tau_{\geq 1}E$，
其中 $\tau_{\geq 1}E$ 完美且 Tor 振幅包含于 $[1, \infty)$。
由于 $E$ 的 Tor 振幅包含于 $[0, \infty)$，
$H^0(E)$ 是平坦 $\mathcal{O}_S$-模。
因此 $H^0(E)$ 是平坦且完美的 $\mathcal{O}_S$-模，
从而有限局部自由，见代数续篇，引理
\ref{more-algebra-lemma-perfect}
（以及代数，引理 \ref{algebra-lemma-finite-projective}
所述的有限投射模为有限局部自由模这一事实）。
于是映射 $\mathcal{O}_S \to H^0(E)$ 是同构，
因为这可在与剩余域作张量后验证
（代数，引理 \ref{algebra-lemma-cokernel-flat}）。
\end{proof}

\begin{lemma}
\label{perfect-lemma-proper-flat-geom-red-connected}
设 $f : X \to S$ 是概形态射。假设
\begin{enumerate}
\item $f$ 固有、平坦且有限表示；
\item $f$ 的几何纤维既约且连通。
\end{enumerate}
则 $f_*\mathcal{O}_X = \mathcal{O}_S$，且在任意基变换后仍成立。
\end{lemma}

\begin{proof}
由引理 \ref{perfect-lemma-proper-flat-h0}，
只需证明对所有 $s \in S$ 都有
$\kappa(s) = H^0(X_s, \mathcal{O}_{X_s})$。
这由簇，引理
\ref{varieties-lemma-proper-geometrically-reduced-global-sections}
以及 $X_s$ 几何连通且几何既约这一事实得到。
\end{proof}

\begin{lemma}
\label{perfect-lemma-proper-idempotent-on-fibre}
设 $f : X \to S$ 是概形的固有态射。令 $s \in S$，
并令 $e \in H^0(X_s, \mathcal{O}_{X_s})$ 为幂等元。
则 $e$ 属于映射
$(f_*\mathcal{O}_X)_s \to H^0(X_s, \mathcal{O}_{X_s})$
的像。
\end{lemma}

\begin{proof}
令 $X_s = T_1 \amalg T_2$ 为与 $e$ 对应的不交并分解，
其中 $T_1$ 与 $T_2$ 非空，且都是 $X_s$ 中的开闭子集；
也就是说，$e$ 在 $T_1$ 上恒等于 $1$，
在 $T_2$ 上恒等于 $0$。

\medskip\noindent
假设 $S$ 是诺特概形。我们将使用概形的上同调，引理
\ref{coherent-lemma-formal-functions-stalk}
所给形式的形式函数定理。它告诉我们
$$
(f_*\mathcal{O}_X)_s^\wedge = \lim_n H^0(X_n, \mathcal{O}_{X_n})
$$
，其中 $X_n$ 是 $X_s$ 的第 $n$ 个无穷小邻域。
由于 $X_n$ 的底拓扑空间等于 $X_s$ 的底拓扑空间，
对所有 $n $ 都得到概形的不交并分解
$X_n = T_{1, n} \amalg T_{2, n}$，
其中当 $i = 1, 2$ 时，$T_{i, n}$ 的底拓扑空间是 $T_i$。
这意味着 $H^0(X_n, \mathcal{O}_{X_n})$
含有非平凡幂等元 $e_n$，即在 $T_{1, n}$ 上恒等于 $1$、
在 $T_{2, n}$ 上恒等于 $0$ 的函数。显然，
$e_{n + 1}$ 在 $X_n$ 上限制为 $e_n$。
所以 $e_\infty = \lim e_n$ 是该极限的非平凡幂等元。
于是 $e_\infty$ 是 $(f_*\mathcal{O}_X)_s$ 的完备化中的元素，
并映到 $H^0(X_s, \mathcal{O}_{X_s})$ 中的 $e$。
由于映射
$(f_*\mathcal{O}_X)_s^\wedge \to H^0(X_s, \mathcal{O}_{X_s})$
经过
$(f_*\mathcal{O}_X)^\wedge_s / \mathfrak m_s (f_*\mathcal{O}_X)_s^\wedge =
(f_*\mathcal{O}_X)_s / \mathfrak m_s (f_*\mathcal{O}_X)_s$
分解（代数，引理 \ref{algebra-lemma-hathat-finitely-generated}），
可知 $e$ 属于映射
$(f_*\mathcal{O}_X)_s \to H^0(X_s, \mathcal{O}_{X_s})$
的像，正如所需。

\medskip\noindent
一般情形：利用极限论证把一般情形归结到诺特情形。
建议读者跳过这段证明。
可以把 $S$ 替换为 $s$ 的一个仿射开邻域。
所以可以且确实假设 $S$ 仿射。由极限，引理
\ref{limits-lemma-proper-limit-of-proper-finite-presentation-noetherian}
，可以写成
$(f : X \to S) = \lim (f_i : X_i \to S_i)$，
其中 $f_i$ 固有且 $S_i$ 诺特。以 $s_i \in S_i$ 表示 $s$ 的像。
则 $s = \lim s_i$，见极限，引理
\ref{limits-lemma-inverse-limit-irreducibles}。
又因为极限与极限可交换（范畴，引理
\ref{categories-lemma-colimits-commute}），有
$X_s = X \times_S s = \lim X_i \times_{S_i} s_i = \lim X_{i, s_i}$。
因此由极限，引理 \ref{limits-lemma-descend-section}，
$e$ 是某个幂等元
$e_i \in H^0(X_{i, s_i}, \mathcal{O}_{X_{i, s_i}})$
的像。由诺特情形，茎
$(f_{i, *}\mathcal{O}_{X_i})_{s_i}$ 中存在元素 $\tilde e_i$
映到 $e_i$。拉回 $\tilde e_i$，得到
$(f_*\mathcal{O}_X)_s$ 中的元素 $\tilde e$，
它映到 $e$，证明完毕。
\end{proof}

\begin{lemma}
\label{perfect-lemma-proper-flat-geom-red}
设 $f : X \to S$ 是概形态射，$s \in S$。假设
\begin{enumerate}
\item $f$ 固有、平坦且有限表示；
\item 纤维 $X_s$ 几何既约。
\end{enumerate}
则把 $S$ 替换为 $s$ 的某个开邻域后，存在直和分解
$Rf_*\mathcal{O}_X = f_*\mathcal{O}_X \oplus P$
作为 $D(\mathcal{O}_S)$ 中的等式，其中 $f_*\mathcal{O}_X$
是有限平展 $\mathcal{O}_S$-代数，而
$P$ 是 Tor 振幅包含于 $[1, \infty)$ 的完美复形。
\end{lemma}

\begin{proof}
本引理的证明与引理 \ref{perfect-lemma-proper-flat-h0} 的证明相似，
建议读者先阅读后者。
由上同调与基变换
（引理 \ref{perfect-lemma-flat-proper-perfect-direct-image-general}），
复形 $E = Rf_*\mathcal{O}_X$ 完美，
且其构造与任意基变换相容。这首先说明
$E$ 的 Tor 振幅包含于 $[0, \infty)$。

\medskip\noindent
我们声称，把 $S$ 替换为 $s$ 的某个开邻域后，
可以在 $D(\mathcal{O}_S)$ 中找到直和分解
$E = H^0(E) \oplus \tau_{\geq 1}E$，
其中 $\tau_{\geq 1}E$ 的 Tor 振幅包含于 $[1, \infty)$。
暂且假设该断言成立，并假设已经作了这一替换，
因而具有上述直和分解。由于 $E$ 的 Tor 振幅包含于
$[0, \infty)$，$H^0(E)$ 是平坦 $\mathcal{O}_S$-模。
所以 $H^0(E)$ 是平坦且完美的 $\mathcal{O}_S$-模，
从而有限局部自由，见代数续篇，引理
\ref{more-algebra-lemma-perfect}
（以及代数，引理 \ref{algebra-lemma-finite-projective}
所述的有限投射模为有限局部自由模这一事实）。
当然，$H^0(E) = f_*\mathcal{O}_X$ 是 $\mathcal{O}_S$-代数。
由上同调与基变换，得到
$H^0(E) \otimes \kappa(s) = H^0(X_s, \mathcal{O}_{X_s})$.
由簇，引理
\ref{varieties-lemma-proper-geometrically-reduced-global-sections}
与 $X_s$ 几何既约这一假设可知，
$\kappa(s) \to H^0(E) \otimes \kappa(s)$ 是有限平展的。
对有限局部自由态射
$\underline{\Spec}_S(H^0(E)) \to S$
应用态射，引理
\ref{morphisms-lemma-set-points-where-fibres-etale}
可知，缩小 $S$ 后，$\mathcal{O}_S$-代数
$H^0(E)$ 是有限平展的。

\medskip\noindent
还需证明该断言。为此，只需证明映射
$$
(f_*\mathcal{O}_X)_s
\longrightarrow
H^0(X_s, \mathcal{O}_{X_s}) = H^0(E \otimes^\mathbf{L} \kappa(s))
$$
是满射，见代数续篇，引理
\ref{more-algebra-lemma-better-cut-complex-in-two}。
选取平坦局部环同态 $\mathcal{O}_{S, s} \to A$，
使 $A$ 的剩余域 $k$ 代数闭，见代数，引理
\ref{algebra-lemma-flat-local-given-residue-field}。
由平坦基变换（概形的上同调，引理
\ref{coherent-lemma-flat-base-change-cohomology})
，得到 $H^0(X_A, \mathcal{O}_{X_A}) =
(f_*\mathcal{O}_X)_s \otimes_{\mathcal{O}_{S, s}} A$
与 $H^0(X_k, \mathcal{O}_{X_k}) =
H^0(X_s, \mathcal{O}_{X_s}) \otimes_{\kappa(s)} k$.
所以只需证明
$H^0(X_A, \mathcal{O}_{X_A}) \to H^0(X_k, \mathcal{O}_{X_k})$
是满射。由于 $X_k$ 是 $k$ 上的既约固有概形，
且 $k$ 代数闭，由上面已经使用过的簇，引理
\ref{varieties-lemma-proper-geometrically-reduced-global-sections}，
$H^0(X_k, \mathcal{O}_{X_k})$ 是有限多个 $k$ 的乘积。
又由引理 \ref{perfect-lemma-proper-idempotent-on-fibre}，
这个 $k$-代数的幂等元都属于映射
$H^0(X_A, \mathcal{O}_{X_A}) \to H^0(X_k, \mathcal{O}_{X_k})$
的像，故结论成立。
\end{proof}





\section{其他应用}
\label{perfect-section-other-applications}

\noindent
本节陈述并证明若干可由上面建立的理论推出的结果。

\begin{lemma}
\label{perfect-lemma-cohomology-over-coherent-ring}
设 $R$ 是凝聚环，$X$ 是在 $R$ 上有限表示的概形。
设 $\mathcal{G}$ 是有限表示的 $\mathcal{O}_X$-模，
它在 $R$ 上平坦且支集在 $R$ 上固有。则
$H^i(X, \mathcal{G})$ 是凝聚 $R$-模。
\end{lemma}

\begin{proof}
结合引理 \ref{perfect-lemma-base-change-tensor-perfect}
与代数续篇，引理 \ref{more-algebra-lemma-coherent-pseudo-coherent}
及 \ref{more-algebra-lemma-perfect}。
\end{proof}

\begin{lemma}
\label{perfect-lemma-countable-cohomology}
设 $X$ 是拟紧拟分离概形。设 $K$ 为
$D_\QCoh(\mathcal{O}_X)$ 的对象，并且上同调层 $H^i(K)$
在仿射开集上的截面集均为可数集。则对任意拟紧开集
$U \subset X$ 以及 $D(\mathcal{O}_X)$ 中任意完美对象 $E$，集合
$$
H^i(U, K \otimes^\mathbf{L} E),\quad \Ext^i(E|_U, K|_U)
$$
都是可数的。
\end{lemma}

\begin{proof}
利用上同调，引理 \ref{cohomology-lemma-dual-perfect-complex}，
只需对群 $H^i(U, K \otimes^\mathbf{L} E)$ 证明结论。
我们将使用归纳原理证明本引理，见概形的上同调，引理
\ref{coherent-lemma-induction-principle}。

\medskip\noindent
先证明 $U = \Spec(A)$ 仿射时结论成立。可以用
$A$-模复形 $K^\bullet$ 表示 $K$，并用有限生成投射
$A$-模的有限复形 $P^\bullet$ 表示 $E$。见引理
\ref{perfect-lemma-affine-compare-bounded}、
\ref{perfect-lemma-perfect-affine}
以及我们对完美 $A$-模复形的定义
（代数续篇，定义 \ref{more-algebra-definition-perfect}）。
于是 $(E \otimes^\mathbf{L} K)|_U$ 由双复形
$P^\bullet \otimes_A K^\bullet$
相伴的全复形表示
（引理 \ref{perfect-lemma-quasi-coherence-tensor-product}）。
对复形 $P^\bullet$ 的长度作归纳
（或使用适当的谱序列），可知只需证明
对每个 $a$，$H^i(P^a \otimes_A K^\bullet)$ 都可数。
由于对某个 $n$，$P^a$ 是 $A^{\oplus n}$ 的直和因子，
这由上同调群 $H^i(K^\bullet)$ 可数这一假设得到。

\medskip\noindent
为完成证明，只需说明：若 $U = V \cup W$，
且结论对 $V$、$W$ 与 $V \cap W$ 成立，则它也对 $U$ 成立。
这是 Mayer--Vietoris 序列的直接推论，见上同调，引理
\ref{cohomology-lemma-unbounded-mayer-vietoris}。
\end{proof}

\begin{lemma}
\label{perfect-lemma-countable}
设 $X$ 是拟紧拟分离概形，并且 $\mathcal{O}_X$
在仿射开集上的截面集均为可数集。
设 $K$ 为 $D_\QCoh(\mathcal{O}_X)$ 的对象。下列条件等价：
\begin{enumerate}
\item $K = \text{hocolim} E_n$，其中 $E_n$ 是
$D(\mathcal{O}_X)$ 的完美对象；
\item 上同调层 $H^i(K)$ 在仿射开集上的截面集均为可数集。
\end{enumerate}
\end{lemma}

\begin{proof}
若 (1) 成立，则 (2) 成立，因为取同伦余极限与取上同调层相容
（由导出范畴，引理 \ref{derived-lemma-cohomology-of-hocolim}），
并且完美复形局部同构于有限自由 $\mathcal{O}_X$-模的有限复形，
故由对 $X$ 的假设满足 (2)。

\medskip\noindent
假设 (2)。选取表示 $K$ 的 K-内射复形 $\mathcal{K}^\bullet$。
选取 $D_\QCoh(\mathcal{O}_X)$ 的完美生成元 $E$，
并用 K-内射复形 $\mathcal{I}^\bullet$ 表示它。
根据定理 \ref{perfect-theorem-DQCoh-is-Ddga} 及其证明，
存在三角范畴的等价
$F : D_\QCoh(\mathcal{O}_X) \to D(A, \text{d})$，
其中 $(A, \text{d})$ 是微分分次代数
$$
(A, \text{d}) =
\Hom_{\text{Comp}^{dg}(\mathcal{O}_X)}
(\mathcal{I}^\bullet, \mathcal{I}^\bullet)
$$
，而该等价把 $K$ 送到微分分次模
$$
M = \Hom_{\text{Comp}^{dg}(\mathcal{O}_X)}
(\mathcal{I}^\bullet, \mathcal{K}^\bullet)
$$
注意，$H^i(A) = \Ext^i(E, E)$，并且
$H^i(M) = \Ext^i(E, K)$。
此外，由于 $F$ 是等价，它及其拟逆都与同伦余极限相容。
因此只需把 $M$ 写成 $D(A, \text{d})$ 的紧对象的同伦余极限。
由微分分次代数，引理 \ref{dga-lemma-countable}，
只需证明对每个 $i$，$\Ext^i(E, E)$ 与
$\Ext^i(E, K)$ 都可数。这由引理
\ref{perfect-lemma-countable-cohomology} 得到。
\end{proof}

\begin{lemma}
\label{perfect-lemma-computing-sections-as-colim}
设 $A$ 是环，$X$ 是在 $A$ 上有限表示的概形。
设 $f : U \to X$ 是有限表示的平坦态射。则：
\begin{enumerate}
\item 存在 $D(\mathcal{O}_X)$ 的完美对象 $L_n$ 的一个逆系统，使得
$$
R\Gamma(U, Lf^*K) = \text{hocolim}\ R\Hom_X(L_n, K)
$$
作为 $D(A)$ 中的等式，对 $D_\QCoh(\mathcal{O}_X)$ 中的 $K$
函子性地成立；
\item 存在 $D(\mathcal{O}_X)$ 的完美对象 $E_n$ 的一个系统，使得
$$
R\Gamma(U, Lf^*K) = \text{hocolim}\ R\Gamma(X, E_n \otimes^\mathbf{L} K)
$$
作为 $D(A)$ 中的等式，对 $D_\QCoh(\mathcal{O}_X)$ 中的 $K$
函子性地成立。
\end{enumerate}
\end{lemma}

\begin{proof}
由引理 \ref{perfect-lemma-cohomology-base-change}，有
$$
R\Gamma(U, Lf^*K) = R\Gamma(X, Rf_*\mathcal{O}_U \otimes^\mathbf{L} K)
$$
，并且对 $K$ 函子性地成立。注意，$R\Gamma(X, -)$
与同伦余极限相容，因为由引理
\ref{perfect-lemma-quasi-coherence-pushforward-direct-sums}，
它与直和相容。类似地，$- \otimes^\mathbf{L} K$
与导出余极限相容，因为 $- \otimes^\mathbf{L} K$ 与直和相容
（$D(\mathcal{O}_X)$ 中的直和由表示复形的直和给出）。
所以为证明 (2)，只需对 $D(\mathcal{O}_X)$ 的某个完美对象系统
$E_n$ 写出 $Rf_*\mathcal{O}_U = \text{hocolim} E_n$。
一旦做到这一点，置 $L_n = E_n^\vee$ 即得到 (1)，见上同调，
引理 \ref{cohomology-lemma-dual-perfect-complex}。

\medskip\noindent
写成 $A = \colim A_i$，其中 $A_i$ 在 $\mathbf{Z}$ 上有限型。
由极限，引理 \ref{limits-lemma-descend-finite-presentation}，
可以找到某个 $i$ 以及有限表示态射
$U_i \to X_i \to \Spec(A_i)$，其到 $\Spec(A)$ 的基变换恢复
$U \to X \to \Spec(A)$。
增大 $i$ 后，可以假设 $f_i : U_i \to X_i$ 平坦，见极限，
引理 \ref{limits-lemma-descend-flat-finite-presentation}。
由引理 \ref{perfect-lemma-compare-base-change}，
$Rf_{i, *}\mathcal{O}_{U_i}$ 沿 $g : X \to X_i$ 的导出拉回
等于 $Rf_*\mathcal{O}_U$。由于 $Lg^*$ 与导出余极限相容，
只需对 $f_i$ 证明所需结论。因此可以假设
$U$ 与 $X$ 都在 $\mathbf{Z}$ 上有限型。

\medskip\noindent
假设 $f : U \to X$ 是 $\mathbf{Z}$ 上有限型概形之间的态射。
为完成证明，我们将说明 $Rf_*\mathcal{O}_U$
是完美复形的同伦余极限。为此应用引理 \ref{perfect-lemma-countable}。
所以只需证明 $R^if_*\mathcal{O}_U$
在仿射开集上的截面集均为可数集。
这由把引理 \ref{perfect-lemma-countable-cohomology} 应用于结构层得到。
\end{proof}




\section{伪凝聚复形的刻画 II}
\label{perfect-section-pseudo-coherent}

\noindent
本节接续第 \ref{perfect-section-pseudo-coherent-hocolim} 节。
我们讨论用上同调刻画伪凝聚复形的方法。
更多这类结果见态射续篇，第
\ref{more-morphisms-section-characterize-pseudo-coherent} 节。

\begin{lemma}
\label{perfect-lemma-pseudo-coherent-over-algebra}
设 $A$ 是环，$R$ 是一个（可以非交换的）$A$-代数，
并且作为 $A$-模有限自由。则 $D(R)$ 中任何在 $D(A)$
里伪凝聚的对象 $M$，都可由有限自由（右）$R$-模的上有界复形表示。
\end{lemma}

\begin{proof}
选取表示 $M$ 的右 $R$-模复形 $M^\bullet$。
由于 $M$ 伪凝聚，对充分大的 $i$，有 $H^i(M) = 0$。
令 $m$ 为使 $H^m(M)$ 非零的最小指标。
由代数续篇，引理 \ref{more-algebra-lemma-finite-cohomology}，
$H^m(M)$ 是有限 $A$-模。
因此可以选取有限自由 $R$-模 $F^m$ 与映射
$F^m \to M^m$，使复合 $F^m \to M^m \to M^{m + 1}$ 为零，
并且 $F^m \to H^m(M)$ 满射。图示如下：
$$
\xymatrix{
& F^m \ar[d]^\alpha \ar[r] & 0 \ar[d] \ar[r] & \ldots \\
M^{m - 1} \ar[r] & M^m \ar[r] & M^{m + 1} \ar[r] & \ldots
}
$$
我们将对 $n \leq m$ 作降序归纳，构造：
当 $i \geq n$ 时的有限自由 $R$-模 $F^i$；
当 $i \geq n$ 时的微分 $d^i : F^i \to F^{i + 1}$；
以及与微分相容的映射 $\alpha : F^i \to K^i$，使得
(1) 当 $i > n$ 时 $H^i(\alpha)$ 是同构，而当 $i = n$ 时满射；
(2) 当 $i > m$ 时 $F^i = 0$。图示为
$$
\xymatrix{
& F^n \ar[r] \ar[d]^\alpha & F^{n + 1} \ar[d]^\alpha \ar[r] & \ldots \ar[r]
& F^i \ar[d]^\alpha \ar[r] & 0 \ar[d] \ar[r] & \ldots \\
M^{n - 1} \ar[r] & M^n \ar[r] & M^{n + 1} \ar[r] & \ldots \ar[r] &
M^i \ar[r] & M^{i + 1} \ar[r] & \ldots
}
$$
基本情形 $n = m$ 已在上面完成。
归纳步骤：令 $C^\bullet$ 为 $\alpha$ 的锥
（导出范畴，定义 \ref{derived-definition-cone}）。
上同调长正合列说明，当 $i \geq n$ 时，
$H^i(C^\bullet) = 0$。注意，因为 $R$ 作为 $A$-模有限自由，
$F^\bullet$ 作为 $A$-模复形是伪凝聚的。
所以由代数续篇，引理
\ref{more-algebra-lemma-cone-pseudo-coherent}，
$C^\bullet$ 作为 $A$-模复形是 $(n - 1)$-伪凝聚的。
由代数续篇，引理 \ref{more-algebra-lemma-finite-cohomology}，
$H^{n - 1}(C^\bullet)$ 是有限 $A$-模。
选取有限自由 $R$-模 $F^{n - 1}$ 与映射
$\beta : F^{n - 1} \to C^{n - 1}$，使复合
$F^{n - 1} \to C^{n - 1} \to C^n$ 为零，
并且 $F^{n - 1}$ 满射到 $H^{n - 1}(C^\bullet)$。
由于 $C^{n - 1} = M^{n - 1} \oplus F^n$，
可写成 $\beta = (\alpha^{n - 1}, -d^{n - 1})$。
复合 $F^{n - 1} \to C^{n - 1} \to C^n$ 为零说明
这些映射组成复形态射
$$
\xymatrix{
& F^{n - 1} \ar[d]^{\alpha^{n - 1}} \ar[r]_{d^{n - 1}} &
F^n \ar[r] \ar[d]^\alpha &
F^{n + 1} \ar[d]^\alpha \ar[r] & \ldots \\
\ldots \ar[r] &
M^{n - 1} \ar[r] & M^n \ar[r] & M^{n + 1} \ar[r] & \ldots
}
$$
此外，这些映射定义可区别三角之间的态射
$$
\xymatrix{
(F^n \to \ldots) \ar[r] \ar[d] &
(F^{n - 1} \to \ldots) \ar[r] \ar[d] &
F^{n - 1} \ar[r] \ar[d]_\beta &
(F^n \to \ldots)[1] \ar[d] \\
(F^n \to \ldots) \ar[r] &
M^\bullet \ar[r] &
C^\bullet \ar[r] &
(F^n \to \ldots)[1]
}
$$
因此，对 $\beta$ 的选择说明复形映射
$(F^{n - 1} \to \ldots) \to M^\bullet$
在次数 $\geq n$ 的上同调上诱导同构，并在次数 $n - 1$ 上诱导满射。
这完成了本引理的证明。
\end{proof}

\begin{lemma}
\label{perfect-lemma-pseudo-coherent-on-projective-space}
设 $A$ 是环，$n \geq 0$，并设
$K \in D_\QCoh(\mathcal{O}_{\mathbf{P}^n_A})$。
下列条件等价：
\begin{enumerate}
\item $K$ 伪凝聚；
\item 对 $D(\mathcal{O}_{\mathbf{P}^n_A})$ 的每个伪凝聚对象 $E$，
$R\Gamma(\mathbf{P}^n_A, E \otimes^\mathbf{L} K)$
都是 $D(A)$ 的伪凝聚对象；
\item 对 $D(\mathcal{O}_{\mathbf{P}^n_A})$ 的每个完美对象 $E$，
$R\Gamma(\mathbf{P}^n_A, E \otimes^\mathbf{L} K)$
都是 $D(A)$ 的伪凝聚对象；
\item 对 $D(\mathcal{O}_{\mathbf{P}^n_A})$ 的每个完美对象 $E$，
$R\Hom_{\mathbf{P}^n_A}(E, K)$ 都是 $D(A)$ 的伪凝聚对象；
\item $R\Gamma(\mathbf{P}^n_A,
K \otimes^\mathbf{L} \mathcal{O}_{\mathbf{P}^n_A}(d))$
对 $d = 0, 1, \ldots, n$ 都是 $D(A)$ 的伪凝聚对象。
\end{enumerate}
\end{lemma}

\begin{proof}
回忆
$$
R\Hom_{\mathbf{P}^n_A}(E, K) =
R\Gamma(\mathbf{P}^n_A, R\SheafHom_{\mathcal{O}_{\mathbf{P}^n_A}}(E, K))
$$
按定义成立，见上同调，第 \ref{cohomology-section-global-RHom} 节。
因此由上同调，引理 \ref{cohomology-lemma-dual-perfect-complex}，
(4) 与 (3) 等价。

\medskip\noindent
由于每个完美复形都是伪凝聚的，显然 (2) 蕴含 (3)。

\medskip\noindent
假设 (1) 成立。则对每个伪凝聚的 $E$，
$E \otimes^\mathbf{L} K$ 伪凝聚，见上同调，引理
\ref{cohomology-lemma-tensor-pseudo-coherent}。
由引理 \ref{perfect-lemma-flat-proper-pseudo-coherent-direct-image-general}，
这种伪凝聚复形的直像仍伪凝聚，故 (2) 成立。

\medskip\noindent
(3) 蕴含 (5)，因为可以对 $d = 0, 1, \ldots, n$
取 $E = \mathcal{O}_{\mathbf{P}^n_A}(d)$。

\medskip\noindent
为完成证明，还需说明 (5) 蕴含 (1)。
令 $P$ 如 (\ref{perfect-equation-generator-Pn})，并令
$R$ 如 (\ref{perfect-equation-algebra-for-Pn})。
由引理 \ref{perfect-lemma-Pn-module-category}，有等价
$$
- \otimes^\mathbf{L}_R P :
D(R) \longrightarrow D_\QCoh(\mathcal{O}_{\mathbf{P}^n_A})
$$
令 $M \in D(R)$ 为满足 $M \otimes^\mathbf{L} P = K$ 的对象。
由微分分次代数，引理
\ref{dga-lemma-upgrade-tensor-with-complex-derived}
，有同构
$$
R\Hom(R, M) = R\Hom_{\mathbf{P}^n_A}(P, K)
$$
作为 $D(A)$ 中的同构。与上面一样论证，得到
$$
R\Hom_{\mathbf{P}^n_A}(P, K) = R\Gamma(\mathbf{P}^n_A,
R\SheafHom_{\mathcal{O}_{\mathbf{P}^n_A}}(E, K)) =
R\Gamma(\mathbf{P}^n_A,
P^\vee \otimes^\mathbf{L}_{\mathcal{O}_{\mathbf{P}^n_A}} K).
$$
利用 $P^\vee$ 是各
$\mathcal{O}_{\mathbf{P}^n_A}(d)$（$d = 0, 1, \ldots, n$）
的直和以及 (5)，可知 $R\Hom(R, M)$ 作为 $A$-模复形伪凝聚。
当然，在 $D(A)$ 中 $M = R\Hom(R, M)$。
所以 $M$ 作为 $A$-模复形伪凝聚。
由引理 \ref{perfect-lemma-pseudo-coherent-over-algebra}，
可以用有限自由 $R$-模的上有界复形 $F^\bullet$ 表示 $M$。于是
$F^\bullet = \bigcup_{p \geq 0} \sigma_{\geq p}F^\bullet$
是一个滤过，它说明 $F^\bullet$ 是具有性质 (P)
的微分分次 $R$-模，见微分分次代数，第
\ref{dga-section-P-resolutions} 节。
因此 $K = M \otimes^\mathbf{L}_R P$ 由
$F^\bullet \otimes_R P$ 表示
（这由导出张量函子的构造得到；例如见微分分次代数，
引理 \ref{dga-lemma-tensor-with-complex-derived} 的证明）。
由于 $F^\bullet \otimes_R P$ 是上有界复形，
其各项均为若干个 $P$ 的直和，本引理得证。
\end{proof}

\begin{lemma}
\label{perfect-lemma-perfect-enough}
设 $A$ 是环，$X$ 是拟紧拟分离的 $A$-概形。
设 $K \in D^-_\QCoh(\mathcal{O}_X)$。
若对 $D(\mathcal{O}_X)$ 中每个完美的 $E$，
$R\Gamma(X, E \otimes^\mathbf{L} K)$ 在 $D(A)$ 中伪凝聚，
则对 $D(\mathcal{O}_X)$ 中每个伪凝聚的 $E$，
$R\Gamma(X, E \otimes^\mathbf{L} K)$ 在 $D(A)$ 中伪凝聚。
\end{lemma}

\begin{proof}
存在整数 $N$，使得
$R\Gamma(X, -) : D_\QCoh(\mathcal{O}_X) \to D(A)$
具有上同调维数 $N$，见引理
\ref{perfect-lemma-quasi-coherence-direct-image} 中的说明。
取 $b \in \mathbf{Z}$，使得当 $i > b$ 时 $H^i(K) = 0$。
令 $E$ 为 $X$ 上的伪凝聚对象。
只需证明对每个 $m$，
$R\Gamma(X, E \otimes^\mathbf{L} K)$ 都是 $m$-伪凝聚的。
选取由完美复形 $P$ 给出的 $(X, E, m - N - 1 - b)$
的一个逼近 $P \to E$。其存在性由
定理 \ref{perfect-theorem-approximation} 保证。
在 $D_\QCoh(\mathcal{O}_X)$ 中选取可区别三角
$$
P \to E \to C \to P[1]
$$
。$C$ 的上同调层在次数 $\geq m - N - 1 - b$ 均为零。
因此 $C \otimes^\mathbf{L} K$ 的上同调层
在次数 $\geq m - N - 1$ 均为零。
所以 $R\Gamma(X, C \otimes^\mathbf{L} K)$ 的上同调
在次数 $\geq m - 1$ 均为零。于是
$$
R\Gamma(X, P \otimes^\mathbf{L} K) \to R\Gamma(X, E \otimes^\mathbf{L} K)
$$
在次数 $\geq m$ 的上同调上是同构。
由假设，其源伪凝聚。
因此 $R\Gamma(X, E \otimes^\mathbf{L} K)$
如所需是 $m$-伪凝聚的。
\end{proof}









\section{相对完美对象}
\label{perfect-section-relatively-perfect}

\noindent
本节引入 \cite{lieblich-complexes} 中的一个概念。

\begin{definition}
\label{perfect-definition-relatively-perfect}
设 $f : X \to S$ 是平坦且局部有限表示的概形态射。
若 $D(\mathcal{O}_X)$ 的对象 $E$ 伪凝聚
（上同调，定义 \ref{cohomology-definition-pseudo-coherent}），
并且 $E$ 作为
$D(f^{-1}\mathcal{O}_S)$
的对象局部具有有限 Tor 维数
（上同调，定义 \ref{cohomology-definition-tor-amplitude}），
则称 $E$ {\it 相对于 $S$ 完美}，或称其为
{\it $S$-完美的}。
\end{definition}

\noindent
相关讨论见评注 \ref{perfect-remark-discuss-rel-perfect}。

\begin{example}
\label{perfect-example-relatively-perfect-field}
设 $k$ 是域，$X$ 是在 $k$ 上有限表示的概形
（特别地，$X$ 拟紧）。则 $D(\mathcal{O}_X)$ 的对象 $E$
是 $k$-完美的，当且仅当它有界且伪凝聚（由定义）；
也就是说，当且仅当它属于 $D^b_{\textit{Coh}}(X)$
（由引理 \ref{perfect-lemma-identify-pseudo-coherent-noetherian}）。
所以，相对完美{\bf 不}意味着“在纤维上完美”。
\end{example}

\noindent
相应的代数概念见代数续篇，第
\ref{more-algebra-section-relatively-perfect} 节。
概形上的概念与代数概念可如下联系。

\begin{lemma}
\label{perfect-lemma-affine-locally-rel-perfect}
设 $f : X \to S$ 是平坦且局部有限表示的概形态射。
设 $E$ 为 $D_\QCoh(\mathcal{O}_X)$ 的对象。下列条件等价：
\begin{enumerate}
\item $E$ 是 $S$-完美的；
\item 对任意映入仿射开集 $V \subset S$ 的仿射开集
$U \subset X$，复形 $R\Gamma(U, E)$ 是
$\mathcal{O}_S(V)$-完美的；
\item 存在仿射开覆盖 $S = \bigcup V_i$，并且对每个 $i$
存在仿射开覆盖 $f^{-1}(V_i) = \bigcup U_{ij}$，
使复形 $R\Gamma(U_{ij}, E)$ 是 $\mathcal{O}_S(V_i)$-完美的。
\end{enumerate}
\end{lemma}

\begin{proof}
伪凝聚性是局部性质，“局部具有有限 Tor 维数”也是局部性质。
因此本引理立即归结为下述断言：
若 $X$ 与 $S$ 仿射，则 $E$ 是 $S$-完美的，
当且仅当 $K = R\Gamma(X, E)$ 是 $\mathcal{O}_S(S)$-完美的。
记 $X = \Spec(A)$、$S = \Spec(R)$，并设 $E$ 对应于
$K \in D(A)$；也就是说，$K = R\Gamma(X, E)$，见引理
\ref{perfect-lemma-affine-compare-bounded}。

\medskip\noindent
注意，$K$ 是 $R$-完美的，当且仅当 $K$
作为 $R$-模复形伪凝聚且具有有限 Tor 维数
（代数续篇，定义 \ref{more-algebra-definition-relatively-perfect}）。
由引理 \ref{perfect-lemma-pseudo-coherent-affine}，
$E$ 伪凝聚当且仅当 $K$ 伪凝聚。
由引理 \ref{perfect-lemma-tor-dimension-rel-affine}，
$E$ 在 $f^{-1}\mathcal{O}_S$ 上具有有限 Tor 维数，
当且仅当 $K$ 作为 $R$-模复形具有有限 Tor 维数。
\end{proof}

\begin{lemma}
\label{perfect-lemma-triangulated}
设 $f : X \to S$ 是平坦且局部有限表示的概形态射。
$D(\mathcal{O}_X)$ 中由 $S$-完美对象组成的全子范畴
是饱和的\footnote{导出范畴，定义
\ref{derived-definition-saturated}。}三角子范畴。
\end{lemma}

\begin{proof}
这由上同调，诸引理
\ref{cohomology-lemma-cone-pseudo-coherent},
\ref{cohomology-lemma-summands-pseudo-coherent},
\ref{cohomology-lemma-cone-tor-amplitude} 与
\ref{cohomology-lemma-summands-tor-amplitude} 得到。
\end{proof}

\begin{lemma}
\label{perfect-lemma-perfect-relatively-perfect}
设 $f : X \to S$ 是平坦且局部有限表示的概形态射。
$D(\mathcal{O}_X)$ 的完美对象是 $S$-完美的。
若 $K, M \in D(\mathcal{O}_X)$，且 $K$ 完美而 $M$ 是
$S$-完美的，则
$K \otimes_{\mathcal{O}_X}^\mathbf{L} M$ 是 $S$-完美的。
\end{lemma}

\begin{proof}
第一种证明：利用引理 \ref{perfect-lemma-affine-locally-rel-perfect}
归结到仿射情形，再应用代数续篇，引理
\ref{more-algebra-lemma-perfect-relatively-perfect}。
\end{proof}

\begin{lemma}
\label{perfect-lemma-base-change-relatively-perfect}
设 $f : X \to S$ 是平坦且局部有限表示的概形态射。
设 $g : S' \to S$ 是概形态射。置 $X' = S' \times_S X$，
并以 $g' : X' \to X$ 表示投影。
若 $K \in D(\mathcal{O}_X)$ 是 $S$-完美的，则
$L(g')^*K$ 是 $S'$-完美的。
\end{lemma}

\begin{proof}
第一种证明：利用引理 \ref{perfect-lemma-affine-locally-rel-perfect}
归结到仿射情形，再应用代数续篇，引理
\ref{more-algebra-lemma-base-change-relatively-perfect}。

\medskip\noindent
第二种证明：由上同调，引理
\ref{cohomology-lemma-pseudo-coherent-pullback}，
$L(g')^*K$ 伪凝聚；有限 Tor 维数性质则由
引理 \ref{perfect-lemma-tor-independence-and-tor-amplitude} 得到。
\end{proof}

\begin{situation}
\label{perfect-situation-relative-descent}
设 $S = \lim_{i \in I} S_i$ 是概形有向系统的极限，
其转移态射 $g_{i'i} : S_{i'} \to S_i$ 均为仿射态射。
假设对所有 $i \in I$，$S_i$ 都拟紧且拟分离。
以 $g_i : S \to S_i$ 表示投影。固定元素 $0 \in I$
以及有限表示的平坦态射 $X_0 \to S_0$。
置 $X_i = S_i \times_{S_0} X_0$ 与 $X = S \times_{S_0} X_0$，
并以 $f_{i'i} : X_{i'} \to X_i$ 表示转移态射，
以 $f_i : X \to X_i$ 表示投影。
\end{situation}

\begin{lemma}
\label{perfect-lemma-relative-descend-homomorphisms}
在情形 \ref{perfect-situation-relative-descent} 中，
设 $K_0$ 与 $L_0$ 为 $D(\mathcal{O}_{X_0})$ 的对象。
当 $i \geq 0$ 时置
$K_i = Lf_{i0}^*K_0$、$L_i = Lf_{i0}^*L_0$，
并置 $K = Lf_0^*K_0$、$L = Lf_0^*L_0$。则映射
$$
\colim_{i \geq 0} \Hom_{D(\mathcal{O}_{X_i})}(K_i, L_i)
\longrightarrow
\Hom_{D(\mathcal{O}_X)}(K, L)
$$
在下列条件下是同构：$K_0$ 伪凝聚，并且
$L_0 \in D_\QCoh(\mathcal{O}_{X_0})$ 作为
$D((X_0 \to S_0)^{-1}\mathcal{O}_{S_0})$
的对象（局部）具有有限 Tor 维数。
\end{lemma}

\begin{proof}
对每个拟紧开集 $U_0 \subset X_0$，考虑条件 $P$：
$$
\colim_{i \geq 0} \Hom_{D(\mathcal{O}_{U_i})}(K_i|_{U_i}, L_i|_{U_i})
\longrightarrow
\Hom_{D(\mathcal{O}_U)}(K|_U, L|_U)
$$
是同构，其中 $U = f_0^{-1}(U_0)$ 且 $U_i = f_{i0}^{-1}(U_0)$。
若 $P$ 对 $U_0$、$V_0$ 与 $U_0 \cap V_0$ 成立，
则由导出范畴中的 Hom 的 Mayer--Vietoris 性质，
它也对 $U_0 \cup V_0$ 成立；见上同调，引理
\ref{cohomology-lemma-mayer-vietoris-hom}。

\medskip\noindent
以 $\pi_0 : X_0 \to S_0$ 表示给定态射。
先考虑 $U_0 = \pi_0^{-1}(W_0)$，其中
$W_0 \subset S_0$ 是拟紧开集。把概形的上同调，引理
\ref{coherent-lemma-induction-principle} 的归纳原理
应用于 $S_0$ 的拟紧开集，并结合上面的观察，
可知只需对 $U_0 = \pi_0^{-1}(W_0)$ 且 $W_0$ 仿射证明 $P$。
换言之，已经归结到 $S_0$ 仿射的情形。
接着再次应用归纳原理，这一次作用于 $X_0$
的所有拟紧拟分离开集，从而进一步归结到 $X_0$ 也仿射的情形。

\medskip\noindent
若 $X_0$ 与 $S_0$ 仿射，结论由代数续篇，引理
\ref{more-algebra-lemma-colimit-relatively-perfect} 得到。
事实上，由引理 \ref{perfect-lemma-pseudo-coherent} 与
\ref{perfect-lemma-affine-compare-bounded}
，该断言转化为模的导出范畴中的 Hom 计算。
引理 \ref{perfect-lemma-pseudo-coherent-affine} 说明
与 $K_0$ 对应的模复形伪凝聚，而引理
\ref{perfect-lemma-tor-dimension-rel-affine} 说明
与 $L_0$ 对应的模复形在 $\mathcal{O}_{S_0}(S_0)$
上具有有限 Tor 维数。因此代数续篇，引理
\ref{more-algebra-lemma-colimit-relatively-perfect}
的假设均满足，结论得证。
\end{proof}

\begin{lemma}
\label{perfect-lemma-descend-relatively-perfect}
在情形 \ref{perfect-situation-relative-descent} 中，
$D(\mathcal{O}_X)$ 的 $S$-完美对象范畴，是各
$D(\mathcal{O}_{X_i})$ 的 $S_i$-完美对象范畴的余极限。
\end{lemma}

\begin{proof}
对每个拟紧开集 $U_0 \subset X_0$，考虑条件 $P$：
函子
$$
\colim_{i \geq 0} D_{S_i\text{-完美}}(\mathcal{O}_{U_i})
\longrightarrow
D_{S\text{-完美}}(\mathcal{O}_U)
$$
是等价，其中 $U = f_0^{-1}(U_0)$ 且 $U_i = f_{i0}^{-1}(U_0)$。
由引理 \ref{perfect-lemma-relative-descend-homomorphisms}，
已经知道该函子全忠实。因此只需证明本质满性。

\medskip\noindent
设 $P$ 对 $X_0$ 的拟紧开集 $U_0$、$V_0$ 成立。
我们声称 $P$ 对 $U_0 \cup V_0$ 也成立。使用记号
$U_i = f_{i0}^{-1}U_0$、$U = f_0^{-1}U_0$、
$V_i = f_{i0}^{-1}V_0$ 与 $V = f_0^{-1}V_0$；
并且略微滥用记号，以 $f_i$ 统一表示态射
$U \to U_i$、$V \to V_i$、
$U \cap V \to U_i \cap V_i$ 与
$U \cup V \to U_i \cup V_i$。
设 $E$ 是 $D(\mathcal{O}_{U \cup V})$ 的 $S$-完美对象。
目标是证明 $E$ 属于该函子的本质像。
由假设，可以找到 $i \geq 0$、$U_i$ 上的 $S_i$-完美对象
$E_{U, i}$、$V_i$ 上的 $S_i$-完美对象 $E_{V, i}$，
以及同构 $Lf_i^*E_{U, i} \to E|_U$ 与
$Lf_i^*E_{V, i} \to E|_V$。令
$$
a : E_{U, i} \to (Rf_{i, *}E)|_{U_i}
\quad\text{且}\quad
b : E_{V, i} \to (Rf_{i, *}E)|_{V_i}
$$
为上述同构 $Lf_i^*E_{U, i} \to E|_U$
与 $Lf_i^*E_{V, i} \to E|_V$ 的伴随映射。
由全忠实性，增大 $i$ 后，可以找到同构
$c : E_{U, i}|_{U_i \cap V_i} \to E_{V, i}|_{U_i \cap V_i}$
，其拉回为复合识别
$$
Lf_i^*E_{U, i}|_{U \cap V} \to E|_{U \cap V} \to Lf_i^*E_{V, i}|_{U \cap V}.
$$
应用上同调，引理 \ref{cohomology-lemma-glue}，
得到 $U_i \cup V_i$ 上的对象 $E_i$
以及映射 $d : E_i \to Rf_{i, *}E$，
其在 $U_i$ 与 $V_i$ 上分别限制为 $a$ 与 $b$。
显然 $E_i$ 是 $S_i$-完美的
（因为相对完美性是局部性质），且
$d$ 是某个同构 $Lf_i^*E_i \to E$ 的伴随映射。

\medskip\noindent
采用引理 \ref{perfect-lemma-relative-descend-homomorphisms}
证明中完全相同的论证，并使用归纳原理
（概形的上同调，引理 \ref{coherent-lemma-induction-principle}），
可归结到 $X_0$ 与 $S_0$ 都仿射的情形。
（先处理 $S_0$ 的开集以归结到 $S_0$ 仿射，
再处理 $X_0$ 的开集以归结到 $X_0$ 仿射。）
在仿射情形中，结论由代数续篇，引理
\ref{more-algebra-lemma-colimit-relatively-perfect} 得到。
到代数问题的转化由引理
\ref{perfect-lemma-affine-locally-rel-perfect} 完成。
\end{proof}

\begin{lemma}
\label{perfect-lemma-derived-pushforward-rel-perfect}
设 $f : X \to S$ 是有限表示的平坦固有概形态射。
设 $E \in D(\mathcal{O}_X)$ 是 $S$-完美的。
则 $Rf_*E$ 是 $D(\mathcal{O}_S)$ 的完美对象，
且其构造与任意基变换相容。
\end{lemma}

\begin{proof}
关于基变换的断言就是引理 \ref{perfect-lemma-compare-base-change}。
所以只需证明 $Rf_*E$ 是完美对象。
我们将用极限论证归结到 $S$ 是诺特仿射概形的情形。

\medskip\noindent
问题在 $S$ 上是局部的，所以可以假设 $S$ 仿射。
记 $S = \Spec(R)$。把 $R = \colim R_i$
写成诺特环 $R_i$ 的滤过余极限。由极限，引理
\ref{limits-lemma-descend-finite-presentation}
，存在某个 $i$ 及一个在 $R_i$ 上有限表示的概形 $X_i$，
其到 $R$ 的基变换为 $X$。由极限，诸引理
\ref{limits-lemma-eventually-proper} 与
\ref{limits-lemma-descend-flat-finite-presentation}
，可以假设 $X_i$ 在 $R_i$ 上固有且平坦。
由引理 \ref{perfect-lemma-descend-relatively-perfect}，
可以假设存在 $D(\mathcal{O}_{X_i})$ 的 $R_i$-完美对象 $E_i$，
其到 $X$ 的拉回为 $E$。
把引理 \ref{perfect-lemma-perfect-direct-image}
应用于 $X_i \to \Spec(R_i)$ 与 $E_i$，
并利用已经证明的基变换性质，即得结论。
\end{proof}

\begin{lemma}
\label{perfect-lemma-compute-ext-rel-perfect}
设 $f : X \to S$ 是概形态射，并设
$E, K \in D(\mathcal{O}_X)$。假设
\begin{enumerate}
\item $S$ 拟紧且拟分离；
\item $f$ 固有、平坦且有限表示；
\item $E$ 是 $S$-完美的；
\item $K$ 伪凝聚。
\end{enumerate}
则存在伪凝聚对象 $L \in D(\mathcal{O}_S)$，使得
$$
Rf_*R\SheafHom(K, E) = R\SheafHom(L, \mathcal{O}_S)
$$
，且在任意基变换后仍成立：给定
$$
\vcenter{
\xymatrix{
X' \ar[r]_{g'} \ar[d]_{f'} &
X \ar[d]^f \\
S' \ar[r]^g &
S
}
}
\quad\quad
\begin{matrix}
\text{为 Cartesian，此时有 } \\
Rf'_*R\SheafHom(L(g')^*K, L(g')^*E) \\
= R\SheafHom(Lg^*L, \mathcal{O}_{S'})
\end{matrix}
$$
\end{lemma}

\begin{proof}
由于 $S$ 拟紧且拟分离，$X$ 也具有这些性质。
由引理 \ref{perfect-lemma-pseudo-coherent-hocolim}，可写成
$K = \text{hocolim} K_n$，其中 $K_n$ 完美，并且 $K_n \to K$
在截断 $\tau_{\geq -n}$ 上诱导同构。令 $K_n^\vee$
为对偶完美复形
（上同调，引理 \ref{cohomology-lemma-dual-perfect-complex}）。
由此得到完美对象的逆系统
$\ldots \to K_3^\vee \to K_2^\vee \to K_1^\vee$。
由引理 \ref{perfect-lemma-perfect-relatively-perfect}，
$K_n^\vee \otimes_{\mathcal{O}_X} E$ 是 $S$-完美的。
所以可以把引理 \ref{perfect-lemma-derived-pushforward-rel-perfect}
应用于 $K_n^\vee \otimes_{\mathcal{O}_X} E$，从而得到
$S$ 上完美复形的逆系统
$$
\ldots \to M_3 \to M_2 \to M_1
$$
，其中
$$
M_n = Rf_*(K_n^\vee \otimes_{\mathcal{O}_X}^\mathbf{L} E) =
Rf_*R\SheafHom(K_n, E)
$$
此外，这些复形的构造与任意基变换相容，即 $Lg^*M_n =
Rf'_*((L(g')^*K_n)^\vee \otimes_{\mathcal{O}_{X'}}^\mathbf{L} L(g')^*E) =
Rf'_*R\SheafHom(L(g')^*K_n, L(g')^*E)$.

\medskip\noindent
由于 $K_n \to K$ 在 $\tau_{\geq -n}$ 上诱导同构，
$K_n \to K_{n + 1}$ 也在 $\tau_{\geq -n}$ 上诱导同构。
又因为 $K_n^\vee = R\SheafHom(K_n, \mathcal{O}_X)$，
$K_{n + 1}^\vee \to K_n^\vee$
在 $\tau_{\leq n}$ 上诱导同构。
设 $E$ 作为 $f^{-1}\mathcal{O}_Y$-模复形
具有包含于 $[a, b]$ 的 Tor 振幅。
由引理 \ref{perfect-lemma-tor-independence-and-tor-amplitude}，
任意基变换后仍然如此。因此
$K_{n + 1}^\vee \otimes_{\mathcal{O}_X} E \to
K_n^\vee \otimes_{\mathcal{O}_X} E$
在 $\tau_{\leq n + a}$ 上诱导同构，
且在任意基变换后仍然如此。
应用右导出函子 $Rf_*$，可知映射
$M_{n + 1} \to M_n$ 在 $\tau_{\leq n + a}$
上诱导同构，且任意基变换后仍然如此。
选取可区别三角
$$
M_{n + 1} \to M_n \to C_n \to M_{n + 1}[1]
$$
取 $S'$ 为某点 $s \in S$ 的剩余域的谱，并作拉回，可见
$C_n \otimes_{\mathcal{O}_S}^\mathbf{L} \kappa(s)$
仅在次数 $\geq n + a$ 可能有非零上同调。由代数续篇，引理
\ref{more-algebra-lemma-lift-perfect-from-residue-field}
，对某个整数 $m_n$，完美复形 $C_n$
的 Tor 振幅包含于 $[n + a, m_n]$。
特别地，对偶完美复形 $C_n^\vee$
的 Tor 振幅包含于 $[-m_n, -n - a]$。

\medskip\noindent
令 $L_n = M_n^\vee$ 为对偶完美复形。上一段的讨论给出：
$L_n \to L_{n + 1}$ 在 $\tau_{\geq -n - a}$ 上诱导同构。
所以 $L = \text{hocolim} L_n$ 伪凝聚，见引理
\ref{perfect-lemma-pseudo-coherent-hocolim}。
由于
$$
R\SheafHom(K, E) = R\SheafHom(\text{hocolim} K_n, E) =
R\lim R\SheafHom(K_n, E) = R\lim K_n^\vee \otimes_{\mathcal{O}_X} E
$$
（上同调，引理 \ref{cohomology-lemma-colim-and-lim-of-duals}），
并且 $R\lim$ 与 $Rf_*$ 相容，得到
$$
Rf_*R\SheafHom(K, E) = R\lim M_n = R\lim R\SheafHom(L_n, \mathcal{O}_S) =
R\SheafHom(L, \mathcal{O}_S)
$$
这证明了 $S$ 上的公式。由于 $M_n$ 的构造与基变换相容，
任意基变换后该公式仍成立。
\end{proof}

\begin{remark}
\label{perfect-remark-compare-L}
读者或许已注意到引理 \ref{perfect-lemma-compute-ext-rel-perfect}
与引理 \ref{perfect-lemma-compute-ext} 之间的相似性。
事实上，引理 \ref{perfect-lemma-compute-ext-rel-perfect}
中的伪凝聚复形 $L$，可刻画为 $S$ 上唯一的伪凝聚复形，
使得存在函子性同构
$$
\Ext^i_{\mathcal{O}_S}(L, \mathcal{F}) \longrightarrow
\Ext^i_{\mathcal{O}_X}(K,
E \otimes_{\mathcal{O}_X}^\mathbf{L} Lf^*\mathcal{F})
$$
；当 $\mathcal{F}$ 取遍 $\QCoh(\mathcal{O}_S)$ 时，
这些同构与边界映射相容。若以后需要这一形式，
我们将在此表述精确结果并给出详细证明。
\end{remark}

\begin{lemma}
\label{perfect-lemma-bounded-on-fibres}
设 $f : X \to S$ 是平坦且局部有限表示的概形态射。
设 $E$ 为 $D(\mathcal{O}_X)$ 的伪凝聚对象。下列条件等价：
\begin{enumerate}
\item $E$ 是 $S$-完美的；
\item $E$ 局部下有界，并且对每个点 $s \in S$，
$D(\mathcal{O}_{X_s})$ 的对象 $L(X_s \to X)^*E$
局部下有界。
\end{enumerate}
\end{lemma}

\begin{proof}
由于一切都是局部的，立即归结到 $X$ 与 $S$ 都仿射的情形，
见引理 \ref{perfect-lemma-affine-locally-rel-perfect}。
设 $X \to S$ 对应于 $\Spec(A) \to \Spec(R)$，
而 $E$ 对应于 $D(A)$ 中的 $K$。若 $s$ 对应于素理想
$\mathfrak p \subset R$，则由于 $R \to A$ 平坦，
$L(X_s \to X)^*E$ 对应于
$K \otimes_R^\mathbf{L} \kappa(\mathfrak p)$；
例如见引理 \ref{perfect-lemma-compare-base-change}。
因此本引理由相应的代数结果得到，见代数续篇，引理
\ref{more-algebra-lemma-bounded-on-fibres-relatively-perfect}。
\end{proof}

\begin{remark}
\label{perfect-remark-discuss-rel-perfect}
在我们的定义适用时，定义 \ref{perfect-definition-relatively-perfect}
给出的相对完美复形概念，与 \cite{lieblich-complexes}
中的概念等价\footnote{为说明这一点，使用引理
\ref{perfect-lemma-affine-locally-rel-perfect} 与代数续篇，引理
\ref{more-algebra-lemma-structure-relatively-perfect}。}。
接着，假设只要求 $f : X \to S$ 局部有限型
（不必平坦，也不必局部有限表示）。
上述论文中的定义是：$E \in D(\mathcal{O}_X)$ 相对完美，如果
\begin{enumerate}
\item[(A)] 在 $X$ 上局部地，对象 $E$ 与一个有限复形拟同构，
该复形由在 $S$ 上平坦的有限表示 $\mathcal{O}_X$-模组成。
\end{enumerate}
另一方面，我们的定义 \ref{perfect-definition-relatively-perfect}
的自然推广是：
\begin{enumerate}
\item[(B)] $E$ 相对于 $S$ 伪凝聚
（态射续篇，定义
\ref{more-morphisms-definition-relative-pseudo-coherence})
，并且 $E$ 作为
$D(f^{-1}\mathcal{O}_S)$
的对象局部具有有限 Tor 维数
（上同调，定义 \ref{cohomology-definition-tor-amplitude}）。
\end{enumerate}
条件 (B) 的优点是，它显然定义了 $D(\mathcal{O}_X)$
的一个三角子范畴；而我们怀疑条件 (A) 并非如此。
条件 (A) 的优点是更易于操作，尤其是在涉及极限时。
\end{remark}









\section{消解性质}
\label{perfect-section-resolution-property}

\noindent
论文 \cite{totaro_resolution} 讨论了这一概念；
\cite{Gross-thesis}、\cite{Gross-surface} 与 \cite{Gross-stack}
继续了相关讨论。目前尚不知道域上的固有概形是否总具有消解性质，
还是这一断言为假。若你知道该问题的答案，请发送电子邮件至
\href{mailto:stacks.project@gmail.com}{stacks.project@gmail.com}.

\medskip\noindent
可以作出下述定义，尽管对一般概形考虑它几乎没有意义。

\begin{definition}
\label{perfect-definition-resolution-property}
设 $X$ 是概形。若每个有限型拟凝聚 $\mathcal{O}_X$-模
都是某个有限局部自由 $\mathcal{O}_X$-模的商，
则称 $X$ 具有{\it 消解性质}。
\end{definition}

\noindent
若 $X$ 是拟紧拟分离概形，则只需验证每个有限表示
$\mathcal{O}_X$-模（自动拟凝聚）都是有限局部自由
$\mathcal{O}_X$-模的商，见性质，引理
\ref{properties-lemma-finite-directed-colimit-surjective-maps}.
若 $X$ 是诺特概形，则有限型拟凝聚模恰好就是凝聚
$\mathcal{O}_X$-模，见概形的上同调，引理
\ref{coherent-lemma-coherent-Noetherian}。

\begin{lemma}
\label{perfect-lemma-resolution-property-ample}
设 $X$ 是概形。若 $X$ 有充足可逆 $\mathcal{O}_X$-模，
则 $X$ 具有消解性质。
\end{lemma}

\begin{proof}
这是性质，命题 \ref{properties-proposition-characterize-ample}
的直接推论。
\end{proof}

\begin{lemma}
\label{perfect-lemma-resolution-property-ample-relative}
设 $f : X \to Y$ 是概形态射。假设
\begin{enumerate}
\item $Y$ 拟紧、拟分离并具有消解性质；
\item $X$ 上存在一个 $f$-充足可逆模。
\end{enumerate}
则 $X$ 具有消解性质。
\end{lemma}

\begin{proof}
令 $\mathcal{F}$ 为有限型拟凝聚 $\mathcal{O}_X$-模，
并令 $\mathcal{L}$ 为 $f$-充足可逆模。
选取仿射开覆盖 $Y = V_1 \cup \ldots \cup V_m$。
置 $U_j = f^{-1}(V_j)$。由性质，命题
\ref{properties-proposition-characterize-ample}
，对每个 $j$ 都存在有限多个映射
$s_{j, i} : \mathcal{L}^{\otimes n_{j, i}}|_{U_j} \to \mathcal{F}|_{U_j}$
，它们联合满射。考虑拟凝聚 $\mathcal{O}_Y$-模
$$
\mathcal{H}_n =
f_*(\mathcal{F} \otimes_{\mathcal{O}_X} \mathcal{L}^{\otimes n})
$$
可以把 $s_{j, i}$ 视为层 $\mathcal{H}_{-n_{j, i}}$
在 $V_j$ 上的截面。假设能找到有限局部自由
$\mathcal{O}_Y$-模 $\mathcal{E}_{i, j}$ 以及映射
$\mathcal{E}_{i, j} \to \mathcal{H}_{-n_{j, i}}$，
使 $s_{j, i}$ 属于其像。则相应映射
$$
f^*\mathcal{E}_{i, j}
\otimes_{\mathcal{O}_X}
\mathcal{L}^{\otimes n_{i, j}} \longrightarrow \mathcal{F}
$$
将联合满射，从而本引理得证。由性质，引理
\ref{properties-lemma-quasi-coherent-colimit-finite-type}，
对每个 $i, j$ 都可以找到有限型拟凝聚子模
$\mathcal{H}'_{i, j} \subset  \mathcal{H}_{-n_{j, i}}$
，它包含 $V_j$ 上的截面 $s_{i, j}$。
所以 $Y$ 的消解性质给出满射
$\mathcal{E}_{i, j} \to \mathcal{H}'_{j, i}$，结论成立。
\end{proof}

\begin{lemma}
\label{perfect-lemma-resolution-property-goes-up-affine}
设 $f : X \to Y$ 是仿射或拟仿射概形态射，
且 $Y$ 拟紧拟分离。若 $Y$ 具有消解性质，则 $X$ 也具有。
\end{lemma}

\begin{proof}
由态射，引理 \ref{morphisms-lemma-quasi-affine-O-ample}，
这是引理 \ref{perfect-lemma-resolution-property-ample-relative} 的特殊情形。
\end{proof}

\noindent
下面是一种可证明消解性质下降的情形。

\begin{lemma}
\label{perfect-lemma-resolution-property-goes-down-finite-flat}
设 $f : X \to Y$ 是满射的有限局部自由概形态射。
若 $X$ 具有消解性质，则 $Y$ 也具有。
\end{lemma}

\begin{proof}
该条件意味着 $f$ 仿射，并且 $f_*\mathcal{O}_X$
是正秩有限局部自由 $\mathcal{O}_Y$-模。
令 $\mathcal{G}$ 为有限型拟凝聚 $\mathcal{O}_Y$-模。
由假设，对某个有限局部自由 $\mathcal{O}_X$-模 $\mathcal{E}$，
存在满射 $\mathcal{E} \to f^*\mathcal{G}$。
由于 $f_*$ 在拟凝聚模上正合
（概形的上同调，引理 \ref{coherent-lemma-relative-affine-vanishing}），
得到满射
$$
f_*\mathcal{E}
\longrightarrow
f_*f^*\mathcal{G} = \mathcal{G} \otimes_{\mathcal{O}_Y} f_*\mathcal{O}_X
$$
取对偶，得到满射
$$
f_*\mathcal{E}
\otimes_{\mathcal{O}_Y}
\SheafHom_{\mathcal{O}_Y}(f_*\mathcal{O}_X, \mathcal{O}_Y)
\longrightarrow
\mathcal{G}
$$
由于 $f_*\mathcal{E}$ 有限局部自由\footnote{具体地，
若 $A \to B$ 是有限局部自由环映射，且 $N$
是有限局部自由 $B$-模，则 $N$ 是有限局部自由 $A$-模。
首先，$N$ 在 $B$ 上有限局部自由蕴含 $N$ 作为 $B$-模
平坦且有限表示，见代数，引理
\ref{algebra-lemma-finite-projective}。其次，由代数，引理
\ref{algebra-lemma-finite-finitely-presented-extension}，
$N$ 是有限表示 $A$-模；由代数，引理
\ref{algebra-lemma-composition-flat}，它是平坦 $A$-模。
最后对 $A$ 应用代数，引理
\ref{algebra-lemma-finite-projective} 即得结论。}，结论成立。
\end{proof}

\begin{lemma}
\label{perfect-lemma-resolution-property-finite-number}
设 $X$ 是概形。假设给定：
\begin{enumerate}
\item 有限仿射开覆盖 $X = U_1 \cup \ldots \cup U_m$；
\item 有限型拟凝聚理想 $\mathcal{I}_j$，满足
$V(\mathcal{I}_j) = X \setminus U_j$。
\end{enumerate}
则 $X$ 具有消解性质，当且仅当对
$j = 1, \ldots, m$，$\mathcal{I}_j$
都是有限局部自由 $\mathcal{O}_X$-模的商。
\end{lemma}

\begin{proof}
一个方向显然。反过来，设
$\mathcal{E}_j \to \mathcal{I}_j$ 是满射，
且 $\mathcal{E}_j$ 有限局部自由。下一段将归结到诺特情形；
建议读者跳过。

\medskip\noindent
首先注意，$U_j \cap U_{j'}$ 拟紧，因为它是仿射概形
$U_j$ 上有限型拟凝聚理想
$\mathcal{I}_{j'}|{U_j}$ 的零点概形之补，见性质，引理
\ref{properties-lemma-quasi-coherent-finite-type-ideals}。
所以 $X$ 拟紧且拟分离，见概形，引理
\ref{schemes-lemma-characterize-quasi-separated}。
由极限，命题 \ref{limits-proposition-approximate}，
可把 $X = \lim X_i$ 写成诺特概形有向系统的极限，
其转移态射均仿射。对每个 $j$，可以找到某个 $i$
以及有限局部自由 $\mathcal{O}_{X_i}$-模
$\mathcal{E}_{i, j}$，其拉回为 $\mathcal{E}_j$，见极限，
引理 \ref{limits-lemma-descend-invertible-modules}。
增大 $i$ 后，可以假设复合
$\mathcal{E}_j \to \mathcal{I}_j \to \mathcal{O}_X$
是某映射 $\mathcal{E}_{i, j} \to \mathcal{O}_{X_i}$ 的拉回，
见极限，引理 \ref{limits-lemma-descend-modules-finite-presentation}。
以 $\mathcal{I}_{i, j} \subset \mathcal{O}_{X_i}$
表示该映射的像；这是诺特概形 $X_i$ 上的拟凝聚理想层，
其到 $X$ 的拉回为 $\mathcal{I}_j$。
以 $U_{i, j} \subset X_i$ 表示相应补开集；
由极限，引理 \ref{limits-lemma-limit-affine}，
可以假设对所有 $i, j$，这些开集都仿射。
若能对 $X_i$ 上的开集 $U_{i, j}$ 与理想层
$\mathcal{I}_{i, j}$ 证明本引理，则由于 $X$ 在 $X_i$ 上仿射，
引理 \ref{perfect-lemma-resolution-property-goes-up-affine}
说明它具有消解性质。这样便归结到诺特概形的情形。

\medskip\noindent
假设 $X$ 诺特。对每个凝聚模 $\mathcal{F}$，
可以选取有限列截面 $s_{jk} \in \mathcal{F}(U_j)$，
$k = 1, \ldots, e_j$，它们生成 $\mathcal{F}$ 到 $U_j$ 的限制。
由概形的上同调，引理 \ref{coherent-lemma-homs-over-open}，
对某个 $n_{jk} \geq 1$，可以把 $s_{jk}$ 延拓为映射
$s'_{jk} : \mathcal{I}_i^{n_{jk}} \to \mathcal{F}$。
于是考虑复合
$$
\mathcal{E}_j^{\otimes n_{jk}} \to \mathcal{I}_j^{n_{jk}} \to \mathcal{F}
$$
即可得结论。
\end{proof}

\begin{lemma}
\label{perfect-lemma-resolution-property-ample-family}
设 $X$ 是概形。若 $X$ 有可逆模的充足族
（态射，定义
\ref{morphisms-definition-family-ample-invertible-modules}),
则 $X$ 具有消解性质。
\end{lemma}

\begin{proof}
由于 $X$ 拟紧，存在 $n$ 以及各对
$(\mathcal{L}_i, s_i)$，$i = 1, \ldots, n$，
其中 $\mathcal{L}_i$ 是可逆 $\mathcal{O}_X$-模，
$s_i \in \Gamma(X, \mathcal{L}_i)$ 是截面；
并且 $s_i$ 非零的点集 $U_i \subset X$ 仿射，
且 $X = U_1 \cup \ldots \cup U_n$。
令 $\mathcal{I}_i \subset \mathcal{O}_X$ 为
$s_i : \mathcal{L}_i^{\otimes -1} \to \mathcal{O}_X$ 的像。
应用引理 \ref{perfect-lemma-resolution-property-finite-number}，
可知 $X$ 具有消解性质。
\end{proof}

\begin{lemma}
\label{perfect-lemma-regular-resolution-property}
设 $X$ 是具有仿射对角的拟紧正则概形。
则 $X$ 具有消解性质。
\end{lemma}

\begin{proof}
结合除子，引理 \ref{divisors-lemma-regular-ample-family}
与上面的引理 \ref{perfect-lemma-resolution-property-ample-family}。
\end{proof}

\begin{lemma}
\label{perfect-lemma-resolution-property-descends}
设 $X = \lim X_i$ 是拟紧拟分离概形有向系统的极限，
其转移态射均仿射。则 $X$ 具有消解性质，
当且仅当对某个 $i$，$X_i$ 具有消解性质。
\end{lemma}

\begin{proof}
若 $X_i$ 具有消解性质，则由引理
\ref{perfect-lemma-resolution-property-goes-up-affine}，$X$ 也具有。
反过来，假设 $X$ 具有消解性质。选取 $i \in I$，
并选取有限仿射开覆盖
$X_i = U_{i, 1} \cup \ldots \cup U_{i, m}$.
对每个 $j$，选取有限型拟凝聚理想层
$\mathcal{I}_{i, j} \subset \mathcal{O}_{X_i}$，使
$X_i \setminus V(\mathcal{I}_{i, j}) = U_{i, j}$，见性质，引理
\ref{properties-lemma-quasi-coherent-finite-type-ideals}。
以 $U_j \subset X$ 表示 $U_{i, j}$ 的逆像，
并以 $\mathcal{I}_j \subset \mathcal{O}_X$
表示 $\mathcal{I}_{i, j}$ 的拉回。
由于 $X$ 具有消解性质，可以选取有限局部自由
$\mathcal{O}_X$-模 $\mathcal{E}_j$
以及满射 $\mathcal{E}_j \to \mathcal{I}_j$。
由极限，诸引理 \ref{limits-lemma-descend-invertible-modules} 与
\ref{limits-lemma-descend-modules-finite-presentation}
，增大 $i$ 后，可以找到有限局部自由
$\mathcal{O}_{X_i}$-模 $\mathcal{E}_{i, j}$
与映射 $\mathcal{E}_{i, j} \to \mathcal{O}_{X_i}$，
它们到 $X$ 的基变换恢复各复合
$\mathcal{E}_j \to \mathcal{I}_j \to \mathcal{O}_X$, $j = 1, \ldots, m$.
有限表示 $\mathcal{O}_{X_i}$-模
$\Coker(\mathcal{E}_{i, j} \to \mathcal{O}_{X_i})$
与 $\mathcal{O}_{X_i}/\mathcal{I}_{i, j}$
到 $X$ 的拉回，作为 $\mathcal{O}_X$ 的商是一致的。
所以由极限，引理
\ref{limits-lemma-descend-modules-finite-presentation}
，可以假设这两个模相同；换言之，
$\mathcal{E}_{i, j} \to \mathcal{O}_{X_i}$ 的像等于
$\mathcal{I}_{i, j}$。于是由引理
\ref{perfect-lemma-resolution-property-finite-number}，
$X_i$ 具有消解性质。
\end{proof}

\begin{lemma}
\label{perfect-lemma-resolution-property-affine-diagonal}
\begin{reference}
\cite[Proposition 1.3]{totaro_resolution} 的特殊情形。
\end{reference}
设 $X$ 是具有消解性质的拟紧拟分离概形。
则 $X$ 具有仿射对角。
\end{lemma}

\begin{proof}
结合极限，命题 \ref{limits-proposition-approximate}
与引理 \ref{perfect-lemma-resolution-property-descends}，
可归结到 $X$ 诺特的情形（略去一个小细节）。
假设 $X$ 诺特。回忆，$X \times X$ 被仿射开集
$U \times V$ 覆盖，其中 $U$、$V$ 为 $X$ 的仿射开集；
见概形，第 \ref{schemes-section-fibre-products} 节。
所以为证明对角 $\Delta : X \to X \times X$ 仿射，
只需证明对 $X$ 的所有仿射开集 $U$、$V$，
$U \cap V = \Delta^{-1}(U \times V)$ 仿射；见态射，引理
\ref{morphisms-lemma-characterize-affine}。
特别地，只需证明：若 $U$ 是 $X$ 的仿射开集，
则包含态射 $j : U \to X$ 仿射。
由概形的上同调，引理
\ref{coherent-lemma-criterion-affine-morphism}，
只需证明对任意拟凝聚 $\mathcal{O}_U$-模 $\mathcal{G}$，
$R^1j_*\mathcal{G} = 0$。
由命题 \ref{perfect-proposition-Noetherian}
（这里使用了已归结到诺特情形这一点），
可以用拟凝聚 $\mathcal{O}_X$-模复形
$\mathcal{H}^\bullet$ 表示 $Rj_*\mathcal{G}$。
为得矛盾，假设
$$
H^1(\mathcal{H}^\bullet) =
\Ker(\mathcal{H}^1 \to \mathcal{H}^2)/\Im(\mathcal{H}^0 \to \mathcal{H}^1)
$$
非零。于是可以找到凝聚 $\mathcal{O}_X$-模
$\mathcal{F}$ 与映射
$$
\mathcal{F} \longrightarrow
\Ker(\mathcal{H}^1 \to \mathcal{H}^2)
$$
，使其与到 $H^1(\mathcal{H}^\bullet)$ 的投影之复合非零。
事实上，由性质，引理
\ref{properties-lemma-quasi-coherent-colimit-finite-type}，
可以把 $\Ker(\mathcal{H}^1 \to \mathcal{H}^2)$
写成其凝聚子模的滤过并，而其中某一个会满足要求。
接着，利用 $X$ 的消解性质，选取有限局部自由
$\mathcal{O}_X$-模 $\mathcal{E}$
以及满射 $\mathcal{E} \to \mathcal{F}$。
这在导出范畴中产生映射
$$
\mathcal{E}[-1] \longrightarrow Rj_*\mathcal{G}
$$
，它在上同调层上非零，因而在 $D(\mathcal{O}_X)$ 中非零。
由伴随，这等同于映射
$$
j^*\mathcal{E}[-1] \to \mathcal{G}
$$
在 $D(\mathcal{O}_U)$ 中非零。由于 $\mathcal{E}$ 有限局部自由，
这又等同于下列群中的非零元素：
$$
H^1(U, j^*\mathcal{E}^\vee \otimes_{\mathcal{O}_U} \mathcal{G})
$$
其中
$\mathcal{E}^\vee = \SheafHom_{\mathcal{O}_X}(\mathcal{E}, \mathcal{O}_X)$
是对偶有限局部自由模。但由概形的上同调，引理
\ref{coherent-lemma-quasi-coherent-affine-cohomology-zero}
，该群为零，这就是所需矛盾。
（若对使用对偶的步骤有疑问，请见更一般的上同调，引理
\ref{cohomology-lemma-dual-perfect-complex}。）
\end{proof}






\section{消解性质与完美复形}
\label{perfect-section-resolution-property-perfect}

\noindent
本节讨论具有消解性质的概形上，完美复形与严格完美复形之间的关系。

\begin{lemma}
\label{perfect-lemma-construct-strictly-perfect}
设 $X$ 是具有消解性质的拟紧拟分离概形。
设 $\mathcal{F}^\bullet$ 是拟凝聚 $\mathcal{O}_X$-模的下有界复形，
它表示 $D(\mathcal{O}_X)$ 的一个完美对象。
则存在有限局部自由 $\mathcal{O}_X$-模的有界复形
$\mathcal{E}^\bullet$，以及拟同构
$\mathcal{E}^\bullet \to \mathcal{F}^\bullet$。
\end{lemma}

\begin{proof}
取整数 $a, b \in \mathbf{Z}$，使得 $\mathcal{F}^\bullet$
的 Tor 振幅包含于 $[a, b]$，并且当 $n < a$ 时
$\mathcal{F}^n = 0$。这种整数对的存在性由上同调，引理
\ref{cohomology-lemma-perfect} 与 $X$ 拟紧这一事实得到。
若 $b < a$，则 $\mathcal{F}^\bullet$ 在导出范畴中为零，
本引理成立。我们将对 $b - a \geq 0$ 作归纳，证明存在复形
$\mathcal{E}^a \to \ldots \to \mathcal{E}^b$
，其中 $\mathcal{E}^i$ 有限局部自由，并且存在拟同构
$\mathcal{E}^\bullet \to \mathcal{F}^\bullet$。

\medskip\noindent
基本情形为 $b - a = 0$。此时
$H^b(\mathcal{F}^\bullet) = H^a(\mathcal{F}^\bullet) =
\Ker(\mathcal{F}^a \to \mathcal{F}^{a + 1})$
有限局部自由。事实上，它是 Tor 维数为 $0$
的有限表示 $\mathcal{O}_X$-模，因而有限局部自由。
见上同调，诸引理 \ref{cohomology-lemma-perfect}、
\ref{cohomology-lemma-finite-cohomology} 以及性质，引理
\ref{properties-lemma-finite-locally-free}。
所以可以取 $\mathcal{E}^\bullet$ 为置于次数 $b$ 的
$H^b(\mathcal{F}^\bullet)$。证明的余下部分处理归纳步骤。

\medskip\noindent
假设 $b > a$。注意
$$
H^b(\mathcal{F}^\bullet) =
\Ker(\mathcal{F}^b \to \mathcal{F}^{b + 1})/
\Im(\mathcal{F}^{b - 1} \to \mathcal{F}^b)
$$
是有限型拟凝聚 $\mathcal{O}_X$-模，见上同调，诸引理
\ref{cohomology-lemma-perfect} 与
\ref{cohomology-lemma-finite-cohomology}。于是可以找到
有限型拟凝聚 $\mathcal{O}_X$-模 $\mathcal{F}$ 以及映射
$$
\mathcal{F} \longrightarrow
\Ker(\mathcal{F}^b \to \mathcal{F}^{b + 1})
$$
，使它与到 $H^b(\mathcal{F}^\bullet)$ 的投影之复合满射。
事实上，由性质，引理
\ref{properties-lemma-quasi-coherent-colimit-finite-type}，
可以把 $\Ker(\mathcal{F}^b \to \mathcal{F}^{b + 1})$
写成其有限型拟凝聚子模的滤过并，而其中某一个会满足要求。
接着，利用 $X$ 的消解性质，选取有限局部自由
$\mathcal{O}_X$-模 $\mathcal{E}^b$
以及满射 $\mathcal{E}^b \to \mathcal{F}$。考虑复形映射
$$
\alpha : \mathcal{E}^b[-b] \to \mathcal{F}^\bullet
$$
及其锥 $C(\alpha)^\bullet$，见导出范畴，定义
\ref{derived-definition-cone}。
注意，$C(\alpha)^\bullet$ 仅在次数 $\geq a$ 时可能非零；
由上同调，引理 \ref{cohomology-lemma-cone-tor-amplitude}，
其 Tor 振幅包含于 $[a, b]$；并且由构造，
$H^b(C(\alpha)^\bullet) = 0$。
所以实际上 $C(\alpha)^\bullet$ 的 Tor 振幅包含于
$[a, b - 1]$。因此可对 $C(\alpha)^\bullet$
应用归纳假设，得到具有归纳假设中所述性质的复形映射
$$
(\mathcal{E}^a \to \ldots \to \mathcal{E}^{b - 1})
\longrightarrow
C(\alpha)^\bullet
$$
。展开锥的定义，得到交换图
$$
\xymatrix{
\ldots \ar[r] &
\mathcal{E}^{b - 2} \ar[r] \ar[d] &
\mathcal{E}^{b - 1} \ar[r] \ar[d] &
0 \ar[r] \ar[d] &
\ldots \\
\ldots \ar[r] &
\mathcal{F}^{b - 2} \ar[r] &
\mathcal{F}^{b - 1} \oplus \mathcal{E}^b \ar[r] &
\mathcal{F}^b \ar[r] &
\ldots
}
$$
显然，由此得到复形映射
$(\mathcal{E}^a \to \ldots \to \mathcal{E}^b) \to \mathcal{F}^\bullet$.
略去验证此映射是拟同构。
\end{proof}

\begin{lemma}
\label{perfect-lemma-resolution-property-perfect-complex}
设 $X$ 是具有消解性质的拟紧拟分离概形。
则 $D(\mathcal{O}_X)$ 的每个完美对象，
都可由有限局部自由 $\mathcal{O}_X$-模的有界复形表示。
\end{lemma}

\begin{proof}
令 $E$ 为 $D(\mathcal{O}_X)$ 的完美对象。
由引理 \ref{perfect-lemma-resolution-property-affine-diagonal}，
$X$ 具有仿射对角。所以由命题
\ref{perfect-proposition-quasi-compact-affine-diagonal}，
可用拟凝聚 $\mathcal{O}_X$-模复形
$\mathcal{F}^\bullet$ 表示 $E$。
注意，因为 $X$ 拟紧，$E$ 属于 $D^b(\mathcal{O}_X)$。
所以当 $n \ll 0$ 时，
$\tau_{\geq n}\mathcal{F}^\bullet$
是表示 $E$ 的拟凝聚 $\mathcal{O}_X$-模下有界复形。
于是应用引理 \ref{perfect-lemma-construct-strictly-perfect} 即得结论。
\end{proof}

\begin{lemma}
\label{perfect-lemma-resolution-property-map-perfect-complex}
设 $X$ 是具有消解性质的拟紧拟分离概形。
设 $\mathcal{E}^\bullet$ 与 $\mathcal{F}^\bullet$
为有限局部自由 $\mathcal{O}_X$-模的有限复形。
则任意
$\alpha \in \Hom_{D(\mathcal{O}_X)}(\mathcal{E}^\bullet, \mathcal{F}^\bullet)$
都可由图
$$
\mathcal{E}^\bullet \leftarrow \mathcal{G}^\bullet \to \mathcal{F}^\bullet
$$
表示，其中 $\mathcal{G}^\bullet$
是有限局部自由 $\mathcal{O}_X$-模的有界复形，
且 $\mathcal{G}^\bullet \to \mathcal{E}^\bullet$ 是拟同构。
\end{lemma}

\begin{proof}
由引理 \ref{perfect-lemma-resolution-property-affine-diagonal}，
$X$ 具有仿射对角。所以由命题
\ref{perfect-proposition-quasi-compact-affine-diagonal}，
可用图
$$
\mathcal{E}^\bullet \leftarrow \mathcal{H}^\bullet \to \mathcal{F}^\bullet
$$
表示 $\alpha$，其中 $\mathcal{H}^\bullet$
是拟凝聚 $\mathcal{O}_X$-模复形，且
$\mathcal{H}^\bullet \to \mathcal{E}^\bullet$
是拟同构。当 $n \ll 0$ 时，映射
$\mathcal{H}^\bullet \to \mathcal{E}^\bullet$ 与
$\mathcal{H}^\bullet \to \mathcal{F}^\bullet$
都经过拟同构
$\mathcal{H}^\bullet \to \tau_{\geq n}\mathcal{H}^\bullet$
分解，这仅仅因为 $\mathcal{E}^\bullet$ 与
$\mathcal{F}^\bullet$ 都是有界复形。
因此可以用 $\tau_{\geq n}\mathcal{H}^\bullet$
替换 $\mathcal{H}^\bullet$，并假设 $\mathcal{H}^\bullet$ 下有界。
随后应用引理 \ref{perfect-lemma-construct-strictly-perfect} 即得结论。
\end{proof}

\begin{lemma}
\label{perfect-lemma-resolution-property-homotopy-map-perfect-complex}
设 $X$ 是具有消解性质的拟紧拟分离概形。
设 $\mathcal{E}^\bullet$ 与 $\mathcal{F}^\bullet$
为有限局部自由 $\mathcal{O}_X$-模的有限复形。
设 $\alpha^\bullet, \beta^\bullet :\mathcal{E}^\bullet \to \mathcal{F}^\bullet$
为两个复形映射，它们在 $D(\mathcal{O}_X)$ 中定义同一映射。
则存在拟同构
$\gamma^\bullet : \mathcal{G}^\bullet \to \mathcal{E}^\bullet$
，其中 $\mathcal{G}^\bullet$
是有限局部自由 $\mathcal{O}_X$-模的有界复形，
并且 $\alpha^\bullet \circ \gamma^\bullet$ 与
$\beta^\bullet \circ \gamma^\bullet$ 是同伦的复形映射。
\end{lemma}

\begin{proof}
由引理 \ref{perfect-lemma-resolution-property-affine-diagonal}，
$X$ 具有仿射对角。所以由命题
\ref{perfect-proposition-quasi-compact-affine-diagonal}
（以及导出范畴的定义），存在拟同构
$\gamma^\bullet : \mathcal{G}^\bullet \to \mathcal{E}^\bullet$
，其中 $\mathcal{G}^\bullet$ 是拟凝聚
$\mathcal{O}_X$-模复形，并且
$\alpha^\bullet \circ \gamma^\bullet$ 与
$\beta^\bullet \circ \gamma^\bullet$ 是同伦的复形映射。
选取证明此事的同伦
$h^i : \mathcal{G}^i \to \mathcal{F}^{i - 1}$。
选取 $n \ll 0$。由于 $\mathcal{E}^\bullet$ 下有界，
映射 $\gamma^\bullet$ 典范地经过商映射
$\mathcal{G}^\bullet \to \tau_{\geq n}\mathcal{G}^\bullet$
分解。出于完全相同的原因，各映射 $h^i$ 经过满射
$\mathcal{G}^i \to (\tau_{\geq n}\mathcal{G})^i$.
分解。因此可以用 $\tau_{\geq n}\mathcal{G}^\bullet$
替换 $\mathcal{G}^\bullet$。
随后应用引理 \ref{perfect-lemma-construct-strictly-perfect} 即得结论。
\end{proof}

\begin{proposition}
\label{perfect-proposition-perfect-resolution-property}
设 $X$ 是具有消解性质的拟紧拟分离概形。记
\begin{enumerate}
\item $\mathcal{A}$ 为有限局部自由 $\mathcal{O}_X$-模的加性范畴；
\item $K^b(\mathcal{A})$ 为 $\mathcal{A}$
中有界复形的同伦范畴，见导出范畴，第
\ref{derived-section-homotopy} 节；
\item $D_{perf}(\mathcal{O}_X)$ 为 $D(\mathcal{O}_X)$
中由完美对象组成的严格全、饱和三角子范畴。
\end{enumerate}
在此记号下，显然的函子
$$
K^b(\mathcal{A}) \longrightarrow D_{perf}(\mathcal{O}_X)
$$
是三角范畴的正合函子，并经过三角范畴的等价
$S^{-1}K^b(\mathcal{A}) \to D_{perf}(\mathcal{O}_X)$ 分解，
其中 $S$ 是 $K^b(\mathcal{A})$ 中由拟同构组成的饱和乘法系。
\end{proposition}

\begin{proof}
若能读懂本命题的陈述，请跳过第一段。关于其中使用的一些定义，见
导出范畴，定义
\ref{derived-definition-triangulated-subcategory}
（三角子范畴），导出范畴，定义
\ref{derived-definition-saturated}（饱和三角子范畴），
导出范畴，定义 \ref{derived-definition-localization}
（与三角结构相容的乘法系），以及范畴，定义
\ref{categories-definition-saturated-multiplicative-system}
（饱和乘法系）。
由上同调，诸引理
\ref{cohomology-lemma-two-out-of-three-perfect} 与
\ref{cohomology-lemma-summands-perfect}，
$D_{perf}(\mathcal{O}_X)$ 是 $D(\mathcal{O}_X)$
的饱和三角子范畴。还要注意，$K^b(\mathcal{A})$
是三角范畴，见导出范畴，引理
\ref{derived-lemma-bounded-triangulated-subcategories}.

\medskip\noindent
显然，该函子把可区别三角送到可区别三角，即它是正合的。
于是由导出范畴，引理
\ref{derived-lemma-triangle-functor-localize}，
$S$ 是与 $K^b(\mathcal{A})$ 上三角结构相容的饱和乘法系。
所以由导出范畴，命题
\ref{derived-proposition-construct-localization}.
，局部化 $S^{-1}K^b(\mathcal{A})$ 存在且为三角范畴。
由导出范畴，引理
\ref{derived-lemma-universal-property-localization}.
，得到正合分解
$S^{-1}K^b(\mathcal{A}) \to D_{perf}(\mathcal{O}_X)$。
由诸引理 \ref{perfect-lemma-resolution-property-perfect-complex}、
\ref{perfect-lemma-resolution-property-map-perfect-complex} 与
\ref{perfect-lemma-resolution-property-homotopy-map-perfect-complex}
，该函子是等价。最后，函子
$S^{-1}K^b(\mathcal{A}) \to D_{perf}(\mathcal{O}_X)$
由导出范畴，引理 \ref{derived-lemma-exact-equivalence}
是三角范畴的等价（即可区别三角相互对应）。
\end{proof}






\section{K-群}
\label{perfect-section-K-groups}

\noindent
简要讨论与概形相伴的各类范畴的 $K_0$。
先前材料见代数，第 \ref{algebra-section-K-groups} 节，
同调，第 \ref{homology-section-K-groups} 节，
导出范畴，第 \ref{derived-section-K-groups} 节，以及
代数续篇，引理 \ref{more-algebra-lemma-perfect-to-K-group-universal}。

\medskip\noindent
类似于代数，第 \ref{algebra-section-K-groups} 节，
对任意诺特概形 $X$，我们将定义两个 $K$-群
$K'_0(X)$ 与 $K_0(X)$。前者使用凝聚 $\mathcal{O}_X$-模，
后者使用有限局部自由 $\mathcal{O}_X$-模。

\begin{lemma}
\label{perfect-lemma-Noetherian-Kprime}
设 $X$ 是诺特概形。则
$$
K_0(\textit{Coh}(\mathcal{O}_X)) =
K_0(D^b(\textit{Coh}(\mathcal{O}_X)) =
K_0(D^b_{\textit{Coh}}(\mathcal{O}_X))
$$
\end{lemma}

\begin{proof}
第一个等式由导出范畴，引理
\ref{derived-lemma-K-bounded-derived} 给出。
由导出范畴，引理 \ref{derived-lemma-DBA-map-K}，有
$K_0(\textit{Coh}(\mathcal{O}_X)) =
K_0(D^b_{\textit{Coh}}(\mathcal{O}_X))$
。这证明了本引理。
（也可以利用命题 \ref{perfect-proposition-DCoh} 给出的
$D^b(\textit{Coh}(\mathcal{O}_X)) = D^b_{\textit{Coh}}(\mathcal{O}_X)$
来得到第二个等式。）
\end{proof}

\noindent
定义如下。

\begin{definition}
\label{perfect-definition-K-group}
设 $X$ 是概形。
\begin{enumerate}
\item 以 $K_0(X)$ 表示 {\it $X$ 的 Grothendieck 群}。
它是 $D(\mathcal{O}_X)$ 中由完美对象组成的严格全、饱和三角子范畴
$D_{perf}(\mathcal{O}_X)$ 的第零 K-群。公式为
$$
K_0(X) = K_0(D_{perf}(\mathcal{O}_X))
$$
\item 若 $X$ 局部诺特，则以 $K'_0(X)$ 表示
{\it $X$ 上凝聚层的 Grothendieck 群}。
它是凝聚 $\mathcal{O}_X$-模的阿贝尔范畴的第零 $K$-群。公式为
$$
K'_0(X) = K_0(\textit{Coh}(\mathcal{O}_X))
$$
\end{enumerate}
\end{definition}

\noindent
我们将证明，若 $X$ 具有消解性质
（例如，若 $X$ 有充足可逆模），则我们对 $K_0(X)$ 的定义
与常用的有限局部自由模定义一致。见引理 \ref{perfect-lemma-K-is-old-K}。

\begin{lemma}
\label{perfect-lemma-K-agrees-affine}
设 $X = \Spec(R)$ 是仿射概形。则 $K_0(X) = K_0(R)$；
若 $R$ 诺特，则 $K'_0(X) = K'_0(R)$。
\end{lemma}

\begin{proof}
回忆，$K'_0(R)$ 与 $K_0(R)$ 已在代数，
第 \ref{algebra-section-K-groups} 节中定义。

\medskip\noindent
由代数续篇，引理
\ref{more-algebra-lemma-perfect-to-K-group-universal}，
有 $K_0(R) = K_0(D_{perf}(R))$。
由诸引理 \ref{perfect-lemma-perfect-affine} 与
\ref{perfect-lemma-affine-compare-bounded}，
有 $D_{perf}(R) = D_{perf}(\mathcal{O}_X)$。
这证明了等式 $K_0(R) = K_0(X)$。

\medskip\noindent
等式 $K'_0(R) = K'_0(X)$ 成立，因为由概形的上同调，引理
\ref{coherent-lemma-coherent-Noetherian}，
$\textit{Coh}(\mathcal{O}_X)$ 等价于有限 $R$-模的范畴。
此外，由定义显然可知，$K'_0(R)$ 正是有限 $R$-模范畴的第零 K-群。
\end{proof}

\noindent
设 $X$ 是诺特概形。此时 $K'_0(X)$ 与 $K_0(X)$ 都有定义，
并且存在典范映射
$$
K_0(X) = K_0(D_{perf}(\mathcal{O}_X))
\longrightarrow
K_0(D^b_{\textit{Coh}}(\mathcal{O}_X)) = K'_0(X)
$$
事实上，完美复形属于 $D^b_{\textit{Coh}}(\mathcal{O}_X)$
（由引理 \ref{perfect-lemma-identify-pseudo-coherent-noetherian}），
包含函子
$D_{perf}(\mathcal{O}_X) \to D^b_{\textit{Coh}}(\mathcal{O}_X)$
在第零 $K$-群上诱导映射
（导出范畴，引理 \ref{derived-lemma-map-K}），
而右侧的等式由引理 \ref{perfect-lemma-Noetherian-Kprime} 给出。

\begin{lemma}
\label{perfect-lemma-Kprime-K}
设 $X$ 是诺特正则概形。则映射
$K_0(X) \to K'_0(X)$ 是同构。
\end{lemma}

\begin{proof}
这立即由引理 \ref{perfect-lemma-perfect-on-regular}
以及上面对映射 $K_0(X) \to K'_0(X)$ 的构造得到。
\end{proof}

\noindent
设 $X$ 是概形。以 $\textit{Vect}(X)$ 表示有限局部自由
$\mathcal{O}_X$-模的范畴。尽管 $\textit{Vect}(X)$
一般不是阿贝尔范畴，但 $\textit{Vect}(X)$ 中短正合列的含义是清楚的。
以 $K_0(\textit{Vect}(X))$ 表示具有下列性质的唯一阿贝尔群
\footnote{这里正确的一般性应当是对任意正合范畴定义 $K_0$，
见内射对象，评注 \ref{injectives-remark-embed-exact-category}。}：
\begin{enumerate}
\item 对每个有限局部自由 $\mathcal{O}_X$-模 $\mathcal{E}$，
给定 $K_0(\textit{Vect}(X))$ 中的元素 $[\mathcal{E}]$；
\item 对有限局部自由 $\mathcal{O}_X$-模的每个短正合列
$0 \to \mathcal{E}' \to \mathcal{E} \to \mathcal{E}'' \to 0$
，在 $K_0(\textit{Vect}(X))$ 中有关系
$[\mathcal{E}] = [\mathcal{E}'] + [\mathcal{E}'']$；
\item 群 $K_0(\textit{Vect}(X))$ 由各元素 $[\mathcal{E}]$ 生成；
\item $K_0(\textit{Vect}(X))$ 中生成元
$[\mathcal{E}]$ 之间的所有关系，都是上述正合列所给关系的
$\mathbf{Z}$-线性组合。
\end{enumerate}
略去 $K_0(\textit{Vect}(X))$ 的详细构造。存在自然映射
$$
K_0(\textit{Vect}(X)) \longrightarrow K_0(X)
$$
具体地，对有限局部自由 $\mathcal{O}_X$-模 $\mathcal{E}$，
以 $\mathcal{E}[0]$ 表示 $X$ 上的完美复形，
它在次数 $0$ 为 $\mathcal{E}$，在其他次数为零。
给定有限局部自由 $\mathcal{O}_X$-模的短正合列
$0 \to \mathcal{E} \to \mathcal{E}' \to \mathcal{E}'' \to 0$
，得到可区别三角
$\mathcal{E}[0] \to \mathcal{E}'[0] \to \mathcal{E}''[0] \to \mathcal{E}[1]$，
见导出范畴，第 \ref{derived-section-canonical-delta-functor} 节。
这说明，把 $[\mathcal{E}]$ 送到 $[\mathcal{E}[0]]$，
就得到映射
$K_0(\textit{Vect}(X)) \to K_0(D_{perf}(\mathcal{O}_X)) = K_0(X)$
；对这一可怕记号深表歉意。

\begin{lemma}
\label{perfect-lemma-K-is-old-K}
设 $X$ 是具有消解性质的拟紧拟分离概形。
则映射 $K_0(\textit{Vect}(X)) \to K_0(X)$ 是同构。
\end{lemma}

\begin{proof}
本引理将直接由诸引理
\ref{perfect-lemma-resolution-property-perfect-complex}、
\ref{perfect-lemma-resolution-property-map-perfect-complex} 与
\ref{perfect-lemma-resolution-property-homotopy-map-perfect-complex}
得到；下文将不再说明地使用这些引理的结果。
构造逆映射
$$
c : K_0(X) = K_0(D_{perf}(\mathcal{O}_X)) \longrightarrow
K_0(\textit{Vect}(X))
$$
具体地，$D_{perf}(\mathcal{O}_X)$ 的任意对象，
都可由有限局部自由 $\mathcal{O}_X$-模的有界复形
$\mathcal{E}^\bullet$ 表示。置
$$
c([\mathcal{E}^\bullet]) = \sum (-1)^i[\mathcal{E}^i]
$$
当然，还需证明这一定义良好。暂且把 $c$
看成定义在有限局部自由 $\mathcal{O}_X$-模有界复形上的映射。

\medskip\noindent
假设 $\mathcal{E}^\bullet \to \mathcal{F}^\bullet$
是有限局部自由 $\mathcal{O}_X$-模有界复形之间的满射。
令 $\mathcal{K}^\bullet$ 为其核。于是得到
$\mathcal{O}_X$-模的短正合列
$$
0 \to \mathcal{K}^n \to \mathcal{E}^n \to \mathcal{F}^n \to 0
$$
。由于 $\mathcal{F}^n$ 有限局部自由，这些正合列局部分裂。
所以 $\mathcal{K}^\bullet$ 也是有限局部自由
$\mathcal{O}_X$-模的有界复形，并且
$c(\mathcal{E}^\bullet) = c(\mathcal{K}^\bullet) + c(\mathcal{F}^\bullet)$
在 $K_0(\textit{Vect}(X))$ 中成立。

\medskip\noindent
设给定有限局部自由 $\mathcal{O}_X$-模的无环有界复形
$\mathcal{E}^\bullet$。设当 $n \not \in [a, b]$ 时
$\mathcal{E}^n = 0$。可以把 $\mathcal{E}^\bullet$
分解成短正合列
$$
\begin{matrix}
0 \to \mathcal{E}^a \to \mathcal{E}^{a + 1} \to \mathcal{F}^{a + 1} \to 0,\\
0 \to \mathcal{F}^{a + 1} \to \mathcal{E}^{a + 2} \to
\mathcal{F}^{a + 3} \to 0, \\
\ldots \\
0 \to \mathcal{F}^{b - 3} \to \mathcal{E}^{b - 2} \to
\mathcal{F}^{b - 2} \to 0, \\
0 \to \mathcal{F}^{b - 2} \to \mathcal{E}^{b - 1} \to \mathcal{E}^b \to 0
\end{matrix}
$$
作降序归纳可知
$\mathcal{F}^{b - 2}, \ldots, \mathcal{F}^{a + 1}$
都是有限局部自由 $\mathcal{O}_X$-模，并且
$$
c(\mathcal{E}^\bullet) = \sum (-1)[\mathcal{E}^n] =
\sum (-1)^n([\mathcal{F}^{n - 1}] + [\mathcal{F}^n]) = 0
$$
所以我们的构造在无环复形上取零。

\medskip\noindent
由前两段的结果可知，$c$ 定义良好。具体地，设有限局部自由
$\mathcal{O}_X$-模的有界复形
$\mathcal{E}^\bullet$ 与 $\mathcal{F}^\bullet$
表示 $D(\mathcal{O}_X)$ 中的同一对象。则可找到拟同构
$a : \mathcal{G}^\bullet \to \mathcal{E}^\bullet$ 与
$b : \mathcal{G}^\bullet \to \mathcal{F}^\bullet$，
其中 $\mathcal{G}^\bullet$ 是有限局部自由
$\mathcal{O}_X$-模的有界复形。得到复形的短正合列
$$
0 \to \mathcal{E}^\bullet \to C(a)^\bullet \to \mathcal{G}^\bullet[1] \to 0
$$
，见导出范畴，定义 \ref{derived-definition-cone}。
由于 $a$ 是拟同构，其锥 $C(a)^\bullet$ 无环
（例如这可由导出范畴，第
\ref{derived-section-canonical-delta-functor} 节中的讨论推出）。因此
$$
0 = c(C(f)^\bullet) = c(\mathcal{E}^\bullet) + c(\mathcal{G}^\bullet[1]) =
c(\mathcal{E}^\bullet) - c(\mathcal{G}^\bullet)
$$
，正如所需。对 $b$ 使用同一论证可知
$0 = c(\mathcal{F}^\bullet) - c(\mathcal{G}^\bullet)$.
所以 $c(\mathcal{E}^\bullet) = c(\mathcal{F}^\bullet)$，
而 $c$ 定义良好。

\medskip\noindent
类似地，对有限局部自由 $\mathcal{O}_X$-模有界复形之间的映射
$\mathcal{E}^\bullet \to \mathcal{F}^\bullet$ 的锥作同样论证，
可知若 $X \to Y \to Z$ 是 $D_{perf}(\mathcal{O}_X)$
中的可区别三角，则 $c(Y) = c(X) + c(Z)$。略去细节。
于是得到所需的阿贝尔群同态
$c : K_0(X) \to K_0(\textit{Vect}(X))$。

\medskip\noindent
显然，复合
$K_0(\textit{Vect}(X)) \to K_0(X) \to K_0(\textit{Vect}(X))$
是恒等映射。另一方面，设 $\mathcal{E}^\bullet$
是有限局部自由 $\mathcal{O}_X$-模的有界复形。
朴素截断的可区别三角
（见同调，第 \ref{homology-section-truncations} 节）
$$
\sigma_{\geq n}\mathcal{E}^\bullet \to
\sigma_{\geq n - 1}\mathcal{E}^\bullet \to
\mathcal{E}^{n - 1}[-n + 1] \to
(\sigma_{\geq n}\mathcal{E}^\bullet)[1]
$$
的存在性与归纳法表明
$$
[\mathcal{E}^\bullet] = \sum (-1)^i[\mathcal{E}^i[0]]
$$
在 $K_0(X) = K_0(D_{perf}(\mathcal{O}_X))$ 中成立；
仍请原谅这里的记号。因此映射
$K_0(\textit{Vect}(X)) \to K_0(D_{perf}(\mathcal{O}_X)) = K_0(X)$
是满射，证明完毕。
\end{proof}

\begin{remark}
\label{perfect-remark-K-ring}
设 $X$ 是概形。K 群 $K_0(X)$ 典范地是交换环。
具体地，利用导出张量积
$$
\otimes = \otimes^\mathbf{L}_{\mathcal{O}_X} :
D_{perf}(\mathcal{O}_X) \times D_{perf}(\mathcal{O}_X)
\longrightarrow
D_{perf}(\mathcal{O}_X)
$$
以及导出范畴，引理 \ref{derived-lemma-bilinear-map-K}，
得到双线性乘法。由于 $K \otimes L \cong L \otimes K$，
这个乘法是交换的。由于
$(K \otimes L) \otimes M = K \otimes (L \otimes M)$
，这个乘法是结合的。最后，$K_0(X)$ 的单位元是
$1 = [\mathcal{O}_X]$。

\medskip\noindent
若 $\textit{Vect}(X)$ 与 $K_0(\textit{Vect}(X))$ 如上，
则显然 $K_0(\textit{Vect}(X))$ 也有环结构：若
$\mathcal{E}$ 与 $\mathcal{F}$ 是有限局部自由
$\mathcal{O}_X$-模，则置
$$
[\mathcal{E}] \cdot [\mathcal{F}] =
[\mathcal{E} \otimes_{\mathcal{O}_X} \mathcal{F}]
$$
读者容易验证，这确实定义了双线性、交换且结合的乘法。
略去细节。上面构造的映射
$$
K_0(\textit{Vect}(X)) \longrightarrow K_0(X)
$$
对于这些定义是环映射。

\medskip\noindent
现假设 $X$ 是 Noether 概形。导出张量积还给出映射
$$
\otimes = \otimes^\mathbf{L}_{\mathcal{O}_X} :
D_{perf}(\mathcal{O}_X) \times D^b_{\textit{Coh}}(\mathcal{O}_X)
\longrightarrow
D^b_{\textit{Coh}}(\mathcal{O}_X)
$$
再次利用导出范畴，引理 \ref{derived-lemma-bilinear-map-K}，
得到双线性乘法 $K_0(X) \times K'_0(X) \to K'_0(X)$；
这是因为由引理 \ref{perfect-lemma-Noetherian-Kprime}，
$K'_0(X) = K_0(D^b_{\textit{Coh}}(\mathcal{O}_X))$。
读者容易证明，这赋予 $K'_0(X)$ 以环 $K_0(X)$ 上模的结构。
\end{remark}

\begin{remark}
\label{perfect-remark-pushforward-K}
设 $f : X \to Y$ 是局部 Noether 概形之间的固有态射。
存在映射
$$
f_* : K'_0(X) \longrightarrow K'_0(Y)
$$
，它把 $[\mathcal{F}]$ 送到
$$
[\bigoplus\nolimits_{i \geq 0} R^{2i}f_*\mathcal{F}] -
[\bigoplus\nolimits_{i \geq 0} R^{2i + 1}f_*\mathcal{F}]
$$
这是定义良好的，因为层 $R^if_*\mathcal{F}$ 是凝聚的
（概形上同调，命题
\ref{coherent-proposition-proper-pushforward-coherent}），
因为局部只有有限多个这样的层非零，也因为 $X$
上凝聚层的短正合列会给出 $Y$ 上 $R^if_*$ 的长正合列。
若 $Y$ 拟紧（这是实践中最常使用的情形），则上式可改写为
$$
f_*[\mathcal{F}] = \sum (-1)^i[R^if_*\mathcal{F}] = [Rf_*\mathcal{F}]
$$
，这里使用了引理 \ref{perfect-lemma-Noetherian-Kprime} 中的等式
$K'_0(Y) = K_0(D^b_{\textit{Coh}}(Y))$。
\end{remark}

\begin{lemma}
\label{perfect-lemma-projection-formula}
设 $f : X \to Y$ 是局部 Noether 概形之间的固有态射。
则对 $\alpha \in K'_0(X)$ 与 $\beta \in K_0(Y)$，有
$f_*(\alpha \cdot f^*\beta) = f_*\alpha \cdot \beta$。
\end{lemma}

\begin{proof}
这由引理 \ref{perfect-lemma-cohomology-base-change}、评注
\ref{perfect-remark-pushforward-K} 中的讨论，以及评注
\ref{perfect-remark-K-ring} 中乘法
$K'_0(X) \times K_0(X) \to K'_0(X)$ 的定义得出。
\end{proof}

\begin{remark}
\label{perfect-remark-perf-Z}
设 $X$ 是概形，$Z \subset X$ 是闭子概形。考虑严格满、饱和的三角子范畴
$$
D_{Z, perf}(\mathcal{O}_X) \subset D(\mathcal{O}_X)
$$
，其对象是这样的 $\mathcal{O}_X$-模完美复形：
它们的同调层在集合论意义下支集于 $Z$。
这个三角范畴的第零 $K$ 群
$K_0(D_{Z, perf}(\mathcal{O}_X))$
有时记作 $K_Z(X)$ 或 $K_{0, Z}(X)$。
完全如评注 \ref{perfect-remark-K-ring} 中那样利用导出张量积，
可知 $K_0(D_{Z, perf}(\mathcal{O}_X))$
具有结合且交换的乘法，但一般
$K_0(D_{Z, perf}(\mathcal{O}_X))$ 没有单位元。
\end{remark}













\section{复形的行列式}
\label{perfect-section-det-two-terms}

\noindent
本节续接代数续篇，第
\ref{more-algebra-section-determinants-complexes} 节。
对任意赋环空间 $(X, \mathcal{O}_X)$，存在函子
$$
\det :
\left\{
\begin{matrix}
\text{完美复形范畴} \\
\text{态射为同构}
\end{matrix}
\right\}
\longrightarrow
\left\{
\begin{matrix}
\text{可逆模范畴} \\
\text{态射为同构}
\end{matrix}
\right\}
$$
此外，给定滤过导出范畴 $DF(\mathcal{O}_X)$ 的对象 $(L, F)$，
若其滤过有限且各分次部分都是完美复形，则存在典范同构
$\det(\text{gr}L) \to \det(L)$。原始论述见
\cite{determinant}。我们稍后会加入这一材料（插入将来的引用）。

\medskip\noindent
目前先给出如下特设构造：$X$ 是概形，而我们考虑
$D(\mathcal{O}_X)$ 中 Tor 振幅位于 $[-1, 0]$ 的完美对象 $L$。

\begin{lemma}
\label{perfect-lemma-determinant-two-term-complexes}
设 $X$ 是概形。存在函子
$$
\det :
\left\{
\begin{matrix}
\text{Tor 振幅位于 }[-1, 0]\text{ 的} \\
\text{完美复形范畴} \\
\text{态射为同构}
\end{matrix}
\right\}
\longrightarrow
\left\{
\begin{matrix}
\text{可逆模范畴} \\
\text{态射为同构}
\end{matrix}
\right\}
$$
此外，给定 $D(\mathcal{O}_X)$ 中秩为 $0$、Tor 振幅位于
$[-1, 0]$ 的完美对象 $L$，存在典范元素
$\delta(L) \in \Gamma(X, \det(L))$，使得对
$D(\mathcal{O}_X)$ 中任意同构 $a : L \to K$，都有
$\det(a)(\delta(L)) = \delta(K)$。
而且此构造在仿射局部由代数续篇，第
\ref{more-algebra-section-determinants-complexes} 节中的构造给出。
\end{lemma}

\begin{proof}
设 $L$ 是左侧范畴的对象。若 $\Spec(A) = U \subset X$
是仿射开集，则 $L|_U$ 对应于 Tor 振幅位于 $[-1, 0]$
的 $A$-模完美复形 $L^\bullet$，见诸引理
\ref{perfect-lemma-affine-compare-bounded}、
\ref{perfect-lemma-tor-dimension-affine} 与
\ref{perfect-lemma-perfect-affine}。
于是可以考虑代数续篇，引理
\ref{more-algebra-lemma-determinant-two-term-complexes}
中构造的可逆 $A$-模 $\det(L^\bullet)$。
若 $\Spec(B) = V \subset U$ 是包含于 $U$ 的另一仿射开集，则
$\det(L^\bullet) \otimes_A B = \det(L^\bullet \otimes_A B)$，
所以此构造与限制映射相容
（见引理 \ref{perfect-lemma-quasi-coherence-pullback}，并注意 $A \to B$ 平坦）。
因此可将这些可逆模黏合，得到 $X$ 上的可逆模 $\det(L)$。
函子性与典范截面完全以同样方式构造。略去细节。
\end{proof}

\begin{remark}
\label{perfect-remark-functorial-det}
引理 \ref{perfect-lemma-determinant-two-term-complexes} 的构造与拉回相容。
更确切地，给定概形态射 $f : X \to Y$ 以及
$D(\mathcal{O}_Y)$ 中 Tor 振幅位于 $[-1, 0]$ 的完美对象 $K$，
则 $Lf^*K$ 是 $D(\mathcal{O}_X)$ 中 Tor 振幅位于
$[-1, 0]$ 的完美对象 $K$，并且有典范等同
$$
f^*\det(K) \longrightarrow \det(Lf^*K)
$$
此外，若 $K$ 的秩为 $0$，则在此映射下，
$\delta(K)$ 拉回为 $\delta(Lf^*K)$。
这由行列式的仿射局部构造显然可见。
\end{remark}






\section{判定有界性}
\label{perfect-section-detecting-boundedness}

\noindent
本节证明，第 \ref{perfect-section-generating} 节所构造的拟紧拟分离概形
之 $D_\QCoh$ 的紧生成元具有一项特殊性质。
由于该节与本节十分相似，建议先阅读该节。

\begin{lemma}
\label{perfect-lemma-orthogonal-koszul-first-variant}
在环境 \ref{perfect-situation-complex} 中，将开浸入记作
$j : U \to X$，并令 $K$ 为 $D(\mathcal{O}_X)$ 中对应于
$A$ 上由 $f_1, \ldots, f_r$ 所给 Koszul 复形的完美对象。
设 $E \in D_\QCoh(\mathcal{O}_X)$ 且 $a \in \mathbf{Z}$。
考虑下列条件：
\begin{enumerate}
\item 典范映射 $\tau_{\geq a}E \to \tau_{\geq a} Rj_*(E|_U)$
是同构。
\item 对所有 $n \geq a$，有
$\Hom_{D(\mathcal{O}_X)}(K[-n], E) = 0$。
\end{enumerate}
则 (2) 蕴含 (1)，而 (1) 蕴含把 $a$ 替换为 $a + 1$ 后的 (2)。
\end{lemma}

\begin{proof}
选取可区别三角 $N \to E \to Rj_*(E|_U) \to N[1]$。
于是 (1) 蕴含 $\tau_{\geq a + 1} N = 0$，而
$\tau_{\geq a}N = 0$ 蕴含 (1)。注意
$$
\Hom_{D(\mathcal{O}_X)}(K[-n], Rj_*(E|_U)) =
\Hom_{D(\mathcal{O}_U)}(K|_U[-n], E) = 0
$$
对所有 $n$ 都成立，因为 $K|_U = 0$。所以 (2) 等价于：
对所有 $n \geq a$，有
$\Hom_{D(\mathcal{O}_X)}(K[-n], N) = 0$。
注意，存在 Koszul 复形的可区别三角
$$
K^\bullet(f_1^{e_1}, \ldots, f_i^{e'_i}, \ldots, f_r^{e_r}) \to
K^\bullet(f_1^{e_1}, \ldots, f_i^{e'_i + e''_i}, \ldots, f_r^{e_r}) \to
K^\bullet(f_1^{e_1}, \ldots, f_i^{e''_i}, \ldots, f_r^{e_r}) \to \ldots
$$
，见代数续篇，引理 \ref{more-algebra-lemma-koszul-mult}。
因此，对所有 $n \geq a$ 有
$\Hom_{D(\mathcal{O}_X)}(K[-n], N) = 0$，等价于：
对所有 $n \geq a$ 及所有 $e \geq 1$，有
$\Hom_{D(\mathcal{O}_X)}(K_e[-n], N) = 0$，
其中 $K_e$ 如引理
\ref{perfect-lemma-represent-cohomology-class-on-closed} 所述。
由于 $N|_U = 0$，该引理又表明，这等价于当 $n \geq a$ 时
$H^n(X, N) = 0$。故 (2) 等价于 $\tau_{\geq a}N = 0$，
因为 $N$ 由 $A$-模复形 $R\Gamma(X, N)$ 决定，见引理
\ref{perfect-lemma-affine-compare-bounded}。引理得证。
\end{proof}

\begin{lemma}
\label{perfect-lemma-orthogonal-koszul-second-variant}
在环境 \ref{perfect-situation-complex} 中，将开浸入记作
$j : U \to X$，并令 $K$ 为 $D(\mathcal{O}_X)$ 中对应于
$A$ 上由 $f_1, \ldots, f_r$ 所给 Koszul 复形的完美对象。
设 $E \in D_\QCoh(\mathcal{O}_X)$ 且 $a \in \mathbf{Z}$。
考虑下列条件：
\begin{enumerate}
\item 典范映射 $\tau_{\leq a}E \to \tau_{\leq a} Rj_*(E|_U)$
是同构；并且
\item $\Hom_{D(\mathcal{O}_X)}(K[-n], E) = 0$ for all $n \leq a$.
\end{enumerate}
则 (2) 蕴含 (1)，而 (1) 蕴含把 $a$ 替换为 $a - 1$ 后的 (2)。
\end{lemma}

\begin{proof}
选取可区别三角 $E \to Rj_*(E|_U) \to N \to E[1]$。
于是 (1) 蕴含 $\tau_{\leq a - 1}N = 0$，而
$\tau_{\leq a}N = 0$ 蕴含 (1)。注意
$$
\Hom_{D(\mathcal{O}_X)}(K[-n], Rj_*(E|_U)) =
\Hom_{D(\mathcal{O}_U)}(K|_U[-n], E) = 0
$$
对所有 $n$ 都成立，因为 $K|_U = 0$。所以 (2) 等价于：
对所有 $n \leq a$，有
$\Hom_{D(\mathcal{O}_X)}(K[-n], N) = 0$。
注意，存在 Koszul 复形的可区别三角
$$
K^\bullet(f_1^{e_1}, \ldots, f_i^{e'_i}, \ldots, f_r^{e_r}) \to
K^\bullet(f_1^{e_1}, \ldots, f_i^{e'_i + e''_i}, \ldots, f_r^{e_r}) \to
K^\bullet(f_1^{e_1}, \ldots, f_i^{e''_i}, \ldots, f_r^{e_r}) \to \ldots
$$
，见代数续篇，引理 \ref{more-algebra-lemma-koszul-mult}。
因此，对所有 $n \leq a$ 有
$\Hom_{D(\mathcal{O}_X)}(K[-n], N) = 0$，等价于：
对所有 $n \leq a$ 及所有 $e \geq 1$，有
$\Hom_{D(\mathcal{O}_X)}(K_e[-n], N) = 0$，
其中 $K_e$ 如引理
\ref{perfect-lemma-represent-cohomology-class-on-closed} 所述。
由于 $N|_U = 0$，该引理又表明，这等价于当 $n \leq a$ 时
$H^n(X, N) = 0$。故 (2) 等价于 $\tau_{\leq a}N = 0$，
因为 $N$ 由 $A$-模复形 $R\Gamma(X, N)$ 决定，见引理
\ref{perfect-lemma-affine-compare-bounded}。引理得证。
\end{proof}

\begin{lemma}
\label{perfect-lemma-bounded-truncation}
设 $X$ 是拟紧拟分离概形。
设 $P \in D_{perf}(\mathcal{O}_X)$、
$E \in D_{\QCoh}(\mathcal{O}_X)$ 且 $a \in \mathbf{Z}$。
下列条件等价：
\begin{enumerate}
\item $\Hom_{D(\mathcal{O}_X)}(P[-i], E) = 0$ for $i \gg 0$, and
\item $\Hom_{D(\mathcal{O}_X)}(P[-i], \tau_{\geq a} E) = 0$ for $i \gg 0$.
\end{enumerate}
\end{lemma}

\begin{proof}
利用三角 $ \tau_{< a} E \to E \to \tau_{\geq a} E \to$，
可知只需证明
$$
\Hom_{D(\mathcal{O}_X)}(P[-i], \tau_{< a} E) =
\Hom_{D(\mathcal{O}_X)}(P, (\tau_{< a} E)[i]) = 0 
$$
当 $i \gg 0$ 时成立。由于 $P$ 完美，这由引理
\ref{perfect-lemma-ext-from-perfect-into-bounded-QCoh} 得出。
\end{proof}

\begin{lemma}
\label{perfect-lemma-bounded-below-truncation}
设 $X$ 是拟紧拟分离概形。设
$P \in D_{perf}(\mathcal{O}_X)$、
$E \in D_{\QCoh}(\mathcal{O}_X)$ 且 $a \in \mathbf{Z}$。
下列条件等价：
\begin{enumerate}
\item $\Hom_{D(\mathcal{O}_X)}(P[-i], E) = 0$ for $i \ll 0$, and
\item $\Hom_{D(\mathcal{O}_X)}(P[-i], \tau_{\leq a} E) = 0$ for $i \ll 0$.
\end{enumerate}
\end{lemma}

\begin{proof}
利用三角 $ \tau_{\leq a} E \to E \to \tau_{> a} E \to$，
可知只需证明
$$
\Hom_{D(\mathcal{O}_X)}(P[-i], \tau_{> a} E) =
\Hom_{D(\mathcal{O}_X)}(P, (\tau_{> a} E)[i]) = 0
$$
当 $i \ll 0$ 时成立。由于 $P$ 完美，这由引理
\ref{perfect-lemma-ext-from-perfect-into-bounded-QCoh} 得出。
\end{proof}

\begin{proposition}
\label{perfect-proposition-detecting-bounded-above}
设 $X$ 是拟紧拟分离概形。设
$G \in D_{perf}(\mathcal{O}_X)$ 是生成
$D_\QCoh (\mathcal{O}_X)$ 的完美复形，并设
$E \in D_\QCoh (\mathcal{O}_X)$。下列条件等价：
\begin{enumerate}
\item $E \in D^-_\QCoh (\mathcal{O}_X)$,
\item $\Hom_{D(\mathcal{O}_X)}(G[-i], E) = 0$ for $i \gg 0$,
\item $\Ext^i_X(G, E) = 0$ for $i \gg 0$,
\item $R\Hom_X(G, E)$ 属于 $D^-(\mathbf{Z})$；
\item $H^i(X, G^\vee \otimes_{\mathcal{O}_X}^\mathbf{L} E) = 0$
，其中 $i \gg 0$；
\item $R\Gamma(X, G^\vee \otimes_{\mathcal{O}_X}^\mathbf{L} E)$
属于 $D^-(\mathbf{Z})$；
\item 对 $D(\mathcal{O}_X)$ 的每个完美对象 $P$，
\begin{enumerate}
\item 把 $G$ 替换为 $P$ 后，断言 (2)、(3)、(4) 成立；
\item $H^i(X, P \otimes_{\mathcal{O}_X}^\mathbf{L} E) = 0$ for $i \gg 0$,
\item $R\Gamma(X, P \otimes_{\mathcal{O}_X}^\mathbf{L} E)$
属于 $D^-(\mathbf{Z})$。
\end{enumerate}
\end{enumerate}
\end{proposition}

\begin{proof}
假设 (1) 成立。由于
$\Hom_{D(\mathcal{O}_X)}(G[-i], E) = \Hom_{D(\mathcal{O}_X)}(G, E[i])$
，由引理 \ref{perfect-lemma-ext-from-perfect-into-bounded-QCoh}，
当 $i \gg 0$ 时该群为零。这证明了 (1) 蕴含 (2)。

\medskip\noindent
由同调，第 \ref{cohomology-section-global-RHom} 节中的讨论，
(2)、(3)、(4) 等价。按定义
$H^i(X, -) = H^i(R\Gamma(X, -))$，故 (5) 与 (6) 等价。
由同调，引理 \ref{cohomology-lemma-dual-perfect-complex}
以及 $(G[-i])^\vee = G^\vee[i]$ 这一事实，
彼此等价的条件 (2)、(3)、(4) 又等价于彼此等价的条件 (5)、(6)。

\medskip\noindent
显然 (7) 蕴含 (2)。反过来，证明彼此等价的条件
(2) -- (6) 蕴含 (7)。由评注 \ref{perfect-remark-classical-generator}，
$G$ 是 $D_{perf}(\mathcal{O}_X)$ 的经典生成元。
对 $P \in D_{perf}(\mathcal{O}_X)$，令 $T(P)$ 表示断言：
$R\Hom_X(P, E)$ 属于 $D^-(\mathbf{Z})$。
显然，$T$ 对直和遗传，对可区别三角满足三取二性质，
对直和项遗传，并在平移下保持。因此由导出范畴，评注
\ref{derived-remark-check-on-generator}，(4) 蕴含 $T$
在整个 $D_{perf}(\mathcal{O}_X)$ 上成立。
同一论证适用于其余所有性质；不过对性质 (7)(b) 与 (7)(c)，
还要使用 $P \mapsto P^\vee$ 是 $D_{perf}(\mathcal{O}_X)$
的自等价。略去这一小细节。

\medskip\noindent
我们将利用概形上同调，引理
\ref{coherent-lemma-induction-principle} 的归纳原理，
证明彼此等价的条件 (2) -- (7) 蕴含 (1)。

\medskip\noindent
首先在 $X$ 仿射时证明 (2) -- (7) $\Rightarrow$ (1)。
置 $P = \mathcal{O}_X[0]$。由 (7)，当 $i \gg 0$ 时
$H^i (X, E) = 0$。由于 $E$ 由 $R\Gamma (X, E)$ 决定，
(1) 随即成立，见引理 \ref{perfect-lemma-affine-compare-bounded}。

\medskip\noindent
现设 $X = U \cup V$，其中 $U$ 是 $X$ 的拟紧开集，
$V$ 是仿射开集；并假设对概形 $U$、$V$ 与 $U \cap V$，
已知蕴含关系 (2) -- (7) $\Rightarrow$ (1)。
设 $E \in D_\QCoh(\mathcal{O}_X)$ 满足 (2) -- (7)。
由引理 \ref{perfect-lemma-direct-summand-of-a-restriction} 与定理
\ref{perfect-theorem-bondal-van-den-Bergh}，存在 $X$ 上完美复形 $Q$，
使得 $Q|_U$ 生成 $D_\QCoh (\mathcal{O}_U)$。
取 $f_1, \dots , f_r \in \Gamma (V, \mathcal{O}_V)$，
使得作为 $V$ 的子集有
$V \setminus U = V(f_1, \dots , f_r)$。
令 $K \in D_{perf}(\mathcal{O}_V)$ 为对应于
$f_1, \dots , f_r$ 上 Koszul 复形的对象，并令
$K' \in D_{perf}(\mathcal{O}_X)$ 为
\begin{equation}
\label{perfect-equation-detecting-bounded-above}
K' = R (V \to X)_* K = R (V \to X)_! K,
\end{equation}
，见同调，诸引理 \ref{cohomology-lemma-pushforward-restriction}
与 \ref{cohomology-lemma-pushforward-perfect}。
这是 $X$ 上支集于闭集 $X \setminus U \subset V$
的完美复形，并且在 $V$ 上同构于 $K$。由假设，
$R\Hom_{\mathcal{O}_X}(Q, E)$ 与
$R\Hom_{\mathcal{O}_X}(K', E)$ 都上有界。

\medskip\noindent
由 (\ref{perfect-equation-detecting-bounded-above}) 中对 $K'$ 的第二种描述，有
$$
\Hom_{D(\mathcal{O}_V)}(K[-i], E|_V) = \Hom_{D(\mathcal{O}_X)}(K'[-i], E) = 0
$$
，其中 $i \gg 0$。因此可对 $E|_V$ 应用引理
\ref{perfect-lemma-orthogonal-koszul-first-variant}，得到整数 $a$，使得
$\tau_{\geq a}(E|_V) = \tau_{\geq a} R (U \cap V \to V)_* (E|_{U \cap V})$.
于是 $\tau_{\geq a} E = \tau_{\geq a} R (U \to X)_* (E |_U)$
（限制到 $U$ 与 $V$ 后检查典范映射是同构）。
因而两次使用引理 \ref{perfect-lemma-bounded-truncation} 可知
$$
\Hom_{D(\mathcal{O}_U)}(Q|_U [-i], E|_U) =
\Hom_{D(\mathcal{O}_X)}(Q[-i], R (U \to X)_* (E|_U)) = 0
$$
，其中 $i \gg 0$。由于本命题对 $U$ 及生成元 $Q|_U$ 成立，
有 $E|_U \in D^-_\QCoh(\mathcal{O}_U)$。
又因函子 $R (U \to X)_*$ 保持 $D^-_\QCoh$
（由引理 \ref{perfect-lemma-quasi-coherence-direct-image}），得到
$\tau_{\geq a}E \in D^-_\QCoh(\mathcal{O}_X)$。
因此 $E \in D^-_\QCoh (\mathcal{O}_X)$。
\end{proof}

\begin{proposition}
\label{perfect-proposition-detecting-bounded-below}
设 $X$ 是拟紧拟分离概形。设
$G \in D_{perf}(\mathcal{O}_X)$
是生成 $D_\QCoh (\mathcal{O}_X)$ 的完美复形，并设
$E \in D_\QCoh (\mathcal{O}_X)$。下列条件等价：
\begin{enumerate}
\item $E \in D^+_\QCoh (\mathcal{O}_X)$,
\item $\Hom_{D(\mathcal{O}_X)}(G[-i], E) = 0$ for $i \ll 0$,
\item $\Ext^i_X(G, E) = 0$ for $i \ll 0$,
\item $R\Hom_X(G, E)$ 属于 $D^+(\mathbf{Z})$；
\item $H^i(X, G^\vee \otimes_{\mathcal{O}_X}^\mathbf{L} E) = 0$
，其中 $i \ll 0$；
\item $R\Gamma(X, G^\vee \otimes_{\mathcal{O}_X}^\mathbf{L} E)$
属于 $D^+(\mathbf{Z})$；
\item 对 $D(\mathcal{O}_X)$ 的每个完美对象 $P$，
\begin{enumerate}
\item 把 $G$ 替换为 $P$ 后，断言 (2)、(3)、(4) 成立；
\item $H^i(X, P \otimes_{\mathcal{O}_X}^\mathbf{L} E) = 0$ for $i \ll 0$,
\item $R\Gamma(X, P \otimes_{\mathcal{O}_X}^\mathbf{L} E)$
属于 $D^+(\mathbf{Z})$。
\end{enumerate}
\end{enumerate}
\end{proposition}

\begin{proof}
假设 (1) 成立。由于
$\Hom_{D(\mathcal{O}_X)}(G[-i], E) = \Hom_{D(\mathcal{O}_X)}(G, E[i])$
，由引理 \ref{perfect-lemma-ext-from-perfect-into-bounded-QCoh}，
当 $i \ll 0$ 时该群为零。这证明了 (1) 蕴含 (2)。

\medskip\noindent
由同调，第 \ref{cohomology-section-global-RHom} 节中的讨论，
(2)、(3)、(4) 等价。按定义
$H^i(X, -) = H^i(R\Gamma(X, -))$，故 (5) 与 (6) 等价。
由同调，引理 \ref{cohomology-lemma-dual-perfect-complex}
以及 $(G[-i])^\vee = G^\vee[i]$ 这一事实，
彼此等价的条件 (2)、(3)、(4) 又等价于彼此等价的条件 (5)、(6)。

\medskip\noindent
显然 (7) 蕴含 (2)。反过来，证明彼此等价的条件
(2) -- (6) 蕴含 (7)。由评注 \ref{perfect-remark-classical-generator}，
$G$ 是 $D_{perf}(\mathcal{O}_X)$ 的经典生成元。
对 $P \in D_{perf}(\mathcal{O}_X)$，令 $T(P)$ 表示断言：
$R\Hom_X(P, E)$ 属于 $D^+(\mathbf{Z})$。
显然，$T$ 对直和遗传，对可区别三角满足三取二性质，
对直和项遗传，并在平移下保持。因此由导出范畴，评注
\ref{derived-remark-check-on-generator}，(4) 蕴含 $T$
在整个 $D_{perf}(\mathcal{O}_X)$ 上成立。
同一论证适用于其余所有性质；不过对性质 (7)(b) 与 (7)(c)，
还要使用 $P \mapsto P^\vee$ 是 $D_{perf}(\mathcal{O}_X)$
的自等价。略去这一小细节。

\medskip\noindent
我们将利用概形上同调，引理
\ref{coherent-lemma-induction-principle} 的归纳原理，
证明彼此等价的条件 (2) -- (7) 蕴含 (1)。

\medskip\noindent
首先在 $X$ 仿射时证明 (2) -- (7) $\Rightarrow$ (1)。
取 $P = \mathcal{O}_X[0]$。由 (7)，当 $i \ll 0$ 时
$H^i (X, E) = 0$。由于 $E$ 由 $R\Gamma (X, E)$ 决定，
(1) 随即成立，见引理 \ref{perfect-lemma-affine-compare-bounded}。

\medskip\noindent
现设 $X = U \cup V$，其中 $U$ 是 $X$ 的拟紧开集，
$V$ 是仿射开集；并假设对概形 $U$、$V$ 与 $U \cap V$，
已知蕴含关系 (2) -- (7) $\Rightarrow$ (1)。
设 $E \in D_\QCoh(\mathcal{O}_X)$ 满足 (2) -- (7)。
由引理 \ref{perfect-lemma-direct-summand-of-a-restriction} 与定理
\ref{perfect-theorem-bondal-van-den-Bergh}，存在 $X$ 上完美复形 $Q$，
使得 $Q|_U$ 生成 $D_\QCoh (\mathcal{O}_U)$。
取 $f_1, \dots , f_r \in \Gamma (V, \mathcal{O}_V)$，
使得作为 $V$ 的子集有
$V \setminus U = V(f_1, \dots , f_r)$。
令 $K \in D_{perf}(\mathcal{O}_V)$ 为对应于
$f_1, \dots , f_r$ 上 Koszul 复形的对象，并令
$K' \in D_{perf}(\mathcal{O}_X)$ 为
\begin{equation}
\label{perfect-equation-detecting-bounded-below}
K' = R (V \to X)_* K = R (V \to X)_! K,
\end{equation}
，见同调，诸引理 \ref{cohomology-lemma-pushforward-restriction}
与 \ref{cohomology-lemma-pushforward-perfect}。
这是 $X$ 上支集于闭集 $X \setminus U \subset V$
的完美复形，并且在 $V$ 上同构于 $K$。由假设，
$R\Hom_{\mathcal{O}_X}(Q, E)$ 与
$R\Hom_{\mathcal{O}_X}(K', E)$ 都下有界。

\medskip\noindent
由 (\ref{perfect-equation-detecting-bounded-below}) 中对 $K'$ 的第二种描述，有
$$
\Hom_{D(\mathcal{O}_V)}(K[-i], E|_V) = \Hom_{D(\mathcal{O}_X)}(K'[-i], E) = 0
$$
，其中 $i \ll 0$。因此可对 $E|_V$ 应用引理
\ref{perfect-lemma-orthogonal-koszul-second-variant}，得到整数 $a$，使得
$\tau_{\leq a}(E|_V) = \tau_{\leq a} R(U \cap V \to V)_*(E|_{U \cap V})$.
于是 $\tau_{\leq a} E = \tau_{\leq a} R(U \to X)_*(E|_U)$
（限制到 $U$ 与 $V$ 后检查典范映射是同构）。
因而两次使用引理 \ref{perfect-lemma-bounded-below-truncation} 可知
$$
\Hom_{D(\mathcal{O}_U)}(Q|_U [-i], E|_U) =
\Hom_{D(\mathcal{O}_X)}(Q[-i], R (U \to X)_* (E|_U)) = 0
$$
，其中 $i \ll 0$。由于本命题对 $U$ 及生成元 $Q|_U$ 成立，
有 $E|_U \in D^+_\QCoh (\mathcal{O}_U)$。
又因函子 $R(U \to X)_*$ 保持下有界对象
（见同调，第 \ref{cohomology-section-derived-functors} 节），得到
$\tau_{\leq a} E \in D^+_\QCoh(\mathcal{O}_X)$。
因此 $E \in D^+_\QCoh (\mathcal{O}_X)$。
\end{proof}










\section{导出范畴中的拟凝聚对象}
\label{perfect-section-QC}

\noindent
设 $X$ 是概形。回忆 $X_{affine, Zar}$ 表示 $X$
的仿射开集所成的范畴，其拓扑由标准 Zariski 覆盖给出，
见拓扑，定义 \ref{topologies-definition-big-small-Zariski}。
提醒读者，$X_{affine, Zar}$ 的拓扑斯就是 $X$ 的小
Zariski 拓扑斯，见拓扑，引理
\ref{topologies-lemma-alternative-zariski}。
位点 $X_{affine, Zar}$ 带有结构层 $\mathcal{O}$，
并且存在赋环拓扑斯的等价
$$
(\Sh(X_{affine, Zar}), \mathcal{O})
\longrightarrow
(\Sh(X_{Zar}), \mathcal{O})
$$
见下降，等式 (\ref{descent-equation-alternative-small-ringed})
及第 \ref{descent-section-alternative-quasi-coherent} 节中围绕该式的讨论；
那里使用了略有不同的记号。

\medskip\noindent
本节以 $X_{affine}$ 表示 $X_{affine, Zar}$
的底层范畴，并赋予它混沌拓扑，即使层与预层相同的拓扑。
特别地，结构层 $\mathcal{O}$ 也成为 $X_{affine}$ 上的层。
由此得到赋环位点态射
$$
\epsilon :
(X_{affine, Zar}, \mathcal{O})
\longrightarrow
(X_{affine}, \mathcal{O})
$$
，如位点上同调，第 \ref{sites-cohomology-section-compare} 节所述。
本节将把 $D_\QCoh(\mathcal{O}_X)$ 与位点上同调，第
\ref{sites-cohomology-section-modules-cohomology} 节中引入的范畴
$\mathit{QC}(X_{affine}, \mathcal{O})$ 等同起来。

\begin{lemma}
\label{perfect-lemma-DQCoh-alternative-small}
在上述情形中，下列三角范畴之间存在典范正合等价：
\begin{enumerate}
\item $D_\QCoh(\mathcal{O}_X)$,
\item $D_\QCoh(X_{Zar}, \mathcal{O})$,
\item $D_\QCoh(X_{affine, Zar}, \mathcal{O})$,
\item $D_\QCoh(X_{affine}, \mathcal{O}_X)$, and
\item $\mathit{QC}(X_{affine}, \mathcal{O})$.
\end{enumerate}
\end{lemma}

\begin{proof}
若 $U \subset V \subset X$ 是仿射开集，则环映射
$\mathcal{O}(V) \to \mathcal{O}(U)$ 平坦。
因此 (4) 与 (5) 之间的等价是位点上同调，引理
\ref{sites-cohomology-lemma-cartesian-flat-quasi-coherent}
的特例（该证明也阐明了陈述的含义）。

\medskip\noindent
由下降，评注 \ref{descent-remark-Zariski-site-space}，
赋环位点 $(X_{Zar}, \mathcal{O})$ 与赋环空间
$(X, \mathcal{O}_X)$ 具有相同的模范畴。
由下降，命题 \ref{descent-proposition-equivalence-quasi-coherent}，
在此等价下拟凝聚模彼此对应。
所以得到 (1) 与 (2) 中三角范畴之间的典范正合等价。

\medskip\noindent
引理之前的讨论表明，存在赋环拓扑斯的等价
$(\Sh(X_{affine, Zar}), \mathcal{O}) \to (\Sh(X_{Zar}), \mathcal{O})$，
从而存在模的阿贝尔范畴之间的等价。
由于拟凝聚模的概念是内在的
（位点上的模，引理 \ref{sites-modules-lemma-special-locally-free}），
这个等价保持拟凝聚模子范畴。
因此得到 (2) 与 (3) 中三角范畴之间的典范正合等价。

\medskip\noindent
为得到 (3) 与 (4) 中三角范畴之间的正合等价，
我们把位点上同调，引理
\ref{sites-cohomology-lemma-compare-topologies-derived-adequate-modules}
应用于上述态射
$\epsilon : (X_{affine, Zar}, \mathcal{O}) \to
(X_{affine}, \mathcal{O})$。取
$\mathcal{B} = \Ob(X_{affine})$，并令
$\mathcal{A} \subset \textit{PMod}(X_{affine}, \mathcal{O})$
为这样的预层 $\mathcal{F}$ 所成的满子范畴：
$\mathcal{F}(V) \otimes_{\mathcal{O}(V)} \mathcal{O}(U)
\to \mathcal{F}(U)$ 是同构。注意，由下降，引理
\ref{descent-lemma-quasi-coherent-alternative-small}，
$\mathcal{A}$ 的对象恰好是 Zariski 拓扑下这样的层：
它们是 $(X_{affine, Zar}, \mathcal{O})$ 上的拟凝聚模。
另一方面，由位点上的模，引理
\ref{sites-modules-lemma-chaotic-quasi-coherent}，
$\mathcal{A}$ 的对象也恰好是 $(X_{affine}, \mathcal{O})$
上，即混沌拓扑下的拟凝聚模。因此，只要证明位点上同调，引理
\ref{sites-cohomology-lemma-compare-topologies-derived-adequate-modules}
可以应用，就确实可利用 $\epsilon^*$ 与 $R\epsilon_*$
得到 (3) 与 (4) 中范畴之间的典范等价。

\medskip\noindent
需要验证四个条件：
\begin{enumerate}
\item $\mathcal{A}$ 的每个对象都是 Zariski 拓扑下的层。
上面已经看到这一点。
\item $\mathcal{A}$ 是
$\textit{Mod}(X_{affine, Zar}, \mathcal{O})$ 的弱 Serre 子范畴。
上面已经看到
$\mathcal{A} = \QCoh(X_{affine, Zar}, \mathcal{O})$，
并且在等价
$\textit{Mod}(X_{affine, Zar}, \mathcal{O}) =
\textit{Mod}(X, \mathcal{O}_X)$
下，这些对象对应于 $X$ 上的拟凝聚模。因此结论由概形，第
\ref{schemes-section-quasi-coherent} 节中的讨论得出。
\item $X_{affine}$ 的每个对象在混沌拓扑下都有一个覆盖，
其成员属于 $\mathcal{B}$。这是因为 $\mathcal{B}$ 包含所有对象。
\item 对 $X_{affine}$ 的每个对象 $U$ 以及
$\mathcal{A}$ 中的每个 $\mathcal{F}$，当 $p > 0$ 时有
$H^p_{Zar}(U, \mathcal{F}) = 0$。
这由仿射概形上拟凝聚模的同调消灭性得出，见概形上同调，引理
\ref{coherent-lemma-quasi-coherent-affine-cohomology-zero}。
\end{enumerate}
证明完毕。
\end{proof}

\begin{remark}
\label{perfect-remark-QC-compare}
设 $S$ 是概形。稍后还将证明
$\mathit{QC}((\textit{Aff}/S), \mathcal{O})$
典范等价于 $D_\QCoh(\mathcal{O}_S)$。
见叠上的层，命题 \ref{stacks-sheaves-proposition-QC-compare}。
\end{remark}
